% !TeX program = pdflatex
% ======================================================================
%  CLASSIC McELIECE — APUNTES
%  Documento didáctico sobre el cifrado de clave pública Classic McEliece
%  Compila con: latexmk -pdf ClassicMcEliece-Apuntes.tex
%  (requiere también makeindex para el índice alfabético)
% ======================================================================
\documentclass[11pt,a4paper,oneside]{report}

% ---------- codificación y lengua ----------
\usepackage[T1]{fontenc}
\usepackage[utf8]{inputenc}
\usepackage[spanish,es-nodecimaldot,es-noquoting,es-noshorthands]{babel}
\usepackage{lmodern}
\usepackage{microtype}
\usepackage{csquotes}

% ---------- matemáticas ----------
\usepackage{amsmath,mathtools,amssymb,amsthm}
\usepackage{dsfont}
\usepackage{centernot}
\usepackage{mathdots}

% ---------- utilidades ----------
\usepackage[a4paper,top=2.7cm,bottom=2.7cm,left=2.5cm,right=2.5cm,headheight=24pt,headsep=14pt]{geometry}
\usepackage{array,booktabs,multirow,longtable,tabularx,tablefootnote}
\usepackage{enumitem}
\usepackage[dvipsnames]{xcolor}
\usepackage{graphicx}
\usepackage{adjustbox}
\usepackage{caption}
\usepackage{subcaption}
\usepackage{siunitx}
\usepackage{fontawesome5}
\usepackage{listings}
\usepackage[ruled,vlined,linesnumbered,spanish]{algorithm2e}

% ---------- gráficos ----------
\usepackage{tikz}
\usetikzlibrary{arrows.meta,calc,positioning,shapes.geometric,shapes.misc,fit,backgrounds,patterns,decorations.pathreplacing,decorations.markings,matrix,chains,quotes}
\usepackage{tikz-cd}
\usepackage{pgfplots}
\pgfplotsset{compat=1.15}
\usepgfplotslibrary{groupplots}
\pgfplotsset{
  every axis/.append style={font=\small},
  every tick label/.append style={font=\footnotesize},
}

% ---------- cajas ----------
\usepackage[most]{tcolorbox}
\tcbuselibrary{theorems,breakable,skins}

% ---------- estructura del documento ----------
\usepackage{titlesec}
\usepackage{titletoc}
\usepackage{tocloft}
\usepackage{fancyhdr}
\usepackage{imakeidx}

% ======================================================================
%  PALETA DE COLORES
% ======================================================================
\definecolor{mcblue}{HTML}{1E4E79}
\definecolor{mcblue2}{HTML}{2E75B6}
\definecolor{mcsky}{HTML}{DCE9F7}
\definecolor{mcorange}{HTML}{C55A11}
\definecolor{mcorange2}{HTML}{FBE5D6}
\definecolor{mcgreen}{HTML}{2E7D32}
\definecolor{mcgreen2}{HTML}{E2F0E3}
\definecolor{mcred}{HTML}{B3261E}
\definecolor{mcred2}{HTML}{FCE8E6}
\definecolor{mcviolet}{HTML}{6A1B9A}
\definecolor{mcviolet2}{HTML}{F3E5F5}
\definecolor{mcteal}{HTML}{00695C}
\definecolor{mcteal2}{HTML}{E0F2F1}
\definecolor{mcgray}{HTML}{455A64}
\definecolor{mclight}{HTML}{F4F7FA}
\definecolor{mcgold}{HTML}{8C6D1F}
\definecolor{mcgold2}{HTML}{FBF3DA}
\definecolor{mcsepia}{HTML}{7B4A12}
\definecolor{mcsepia2}{HTML}{F7EFE3}

% ======================================================================
%  ESTILOS DE PÁGINA Y TÍTULOS
% ======================================================================
\pagestyle{fancy}
\fancyhf{}
\fancyhead[L]{\small\color{mcblue}\nouppercase{\leftmark}}
\fancyhead[R]{\small\color{mcgray}\rightmark}
\fancyfoot[C]{\small\color{mcgray}\thepage}
\renewcommand{\headrulewidth}{0.6pt}
\renewcommand{\headrule}{\hbox to\headwidth{\color{mcblue}\leaders\hrule height \headrulewidth\hfill}}
\fancypagestyle{plain}{\fancyhf{}\fancyfoot[C]{\small\color{mcgray}\thepage}\renewcommand{\headrulewidth}{0pt}}
% marcas cortas para evitar solapes en la cabecera
\renewcommand{\chaptermark}[1]{\markboth{Capítulo \thechapter}{}}
\renewcommand{\sectionmark}[1]{\markright{\thesection.\ #1}}

\titleformat{\chapter}[display]
  {\normalfont\bfseries\color{mcblue}}
  {\filleft\Large\color{mcblue2}CAPÍTULO \thechapter}
  {6pt}
  {\titlerule[1.2pt]\vspace{6pt}\Huge\filright}
  [\vspace{4pt}{\titlerule[1.2pt]}]
\titlespacing*{\chapter}{0pt}{-20pt}{28pt}

\titleformat{\section}
  {\normalfont\Large\bfseries\color{mcblue}}
  {\thesection}{0.6em}{}[{\color{mcblue2}\quad\raisebox{-2pt}{\rule{0.4\linewidth}{0.8pt}}}]
\titleformat{\subsection}{\normalfont\large\bfseries\color{mcblue2}}{\thesubsection}{0.5em}{}
\titleformat{\subsubsection}{\normalfont\normalsize\bfseries\color{mcgray}}{\thesubsubsection}{0.5em}{}

% partes con banda de color
\titleformat{\part}[display]
  {\normalfont\bfseries\centering}
  {\begin{tikzpicture}
      \node[fill=mcblue, text=white, rounded corners=4pt, inner sep=18pt,
            minimum width=0.86\textwidth, align=center, font=\Huge\bfseries] (p)
        {PARTE \thepart};
   \end{tikzpicture}}
  {16pt}
  {\color{mcblue}\LARGE}
\titlespacing*{\part}{0pt}{40pt}{40pt}

% TOC
\renewcommand{\cftchapfont}{\bfseries\color{mcblue}}
\renewcommand{\cftsecfont}{\color{mcgray}}
\renewcommand{\cftchapleader}{\color{mcblue2}\cftdotfill{\cftdotsep}}
\renewcommand{\cftchappagefont}{\bfseries\color{mcblue}}
\setlength{\cftbeforechapskip}{4pt}

% ======================================================================
%  HIPERVÍNCULOS
% ======================================================================
\usepackage[colorlinks=true,linkcolor=mcblue,citecolor=mcteal,urlcolor=mcorange,bookmarksnumbered=true]{hyperref}
\usepackage{cleveref}

% nombres en español para que \cref funcione con las cajas-teorema
\crefname{definicion}{definición}{definiciones}
\Crefname{definicion}{Definición}{Definiciones}
\crefname{teorema}{teorema}{teoremas}
\Crefname{teorema}{Teorema}{Teoremas}
\crefname{lema}{lema}{lemas}
\Crefname{lema}{Lema}{Lemas}
\crefname{proposicion}{proposición}{proposiciones}
\Crefname{proposicion}{Proposición}{Proposiciones}
\crefname{corolario}{corolario}{corolarios}
\Crefname{corolario}{Corolario}{Corolarios}
\crefname{ejemplo}{ejemplo}{ejemplos}
\Crefname{ejemplo}{Ejemplo}{Ejemplos}
\crefname{observacion}{observación}{observaciones}
\Crefname{observacion}{Observación}{Observaciones}

% ======================================================================
%  TEOREMAS Y CAJAS DIDÁCTICAS
% ======================================================================
\tcbset{
  thmbox/.style={
    enhanced, breakable,
    boxrule=0pt, leftrule=3pt,
    arc=2pt, outer arc=2pt,
    left=5pt, right=5pt, top=4pt, bottom=4pt,
    fonttitle=\bfseries,
    attach boxed title to top left={yshift=-2pt, xshift=6pt},
    boxed title style={boxrule=0pt, arc=1.5pt, left=4pt, right=4pt, top=1pt, bottom=1pt},
  }
}
\newtcbtheorem[number within=chapter]{definicion}{Definición}%
  {thmbox, colback=mcsky!45, colframe=mcblue, colbacktitle=mcblue, coltitle=white}{def}
\newtcbtheorem[number within=chapter,use counter from=definicion]{teorema}{Teorema}%
  {thmbox, colback=mcgreen2!80, colframe=mcgreen, colbacktitle=mcgreen, coltitle=white}{thm}
\newtcbtheorem[number within=chapter,use counter from=definicion]{lema}{Lema}%
  {thmbox, colback=mcgreen2!80, colframe=mcgreen, colbacktitle=mcgreen, coltitle=white}{lem}
\newtcbtheorem[number within=chapter,use counter from=definicion]{proposicion}{Proposición}%
  {thmbox, colback=mcgreen2!80, colframe=mcgreen, colbacktitle=mcgreen, coltitle=white}{prop}
\newtcbtheorem[number within=chapter,use counter from=definicion]{corolario}{Corolario}%
  {thmbox, colback=mcgreen2!80, colframe=mcgreen, colbacktitle=mcgreen, coltitle=white}{cor}
\newtcbtheorem[number within=chapter,use counter from=definicion]{ejemplo}{Ejemplo}%
  {thmbox, colback=mcorange2!70, colframe=mcorange, colbacktitle=mcorange, coltitle=white}{ej}
\newtcbtheorem[number within=chapter,use counter from=definicion]{observacion}{Observación}%
  {thmbox, colback=mcgray!8, colframe=mcgray, colbacktitle=mcgray, coltitle=white}{obs}

\newtcolorbox{roadmap}[1][En este capítulo]{
  enhanced, breakable, colback=mcblue!5, colframe=mcblue2, boxrule=0.6pt,
  left=6pt, right=6pt, top=5pt, bottom=5pt, arc=3pt,
  fonttitle=\bfseries\color{mcblue},
  title={\faMapSigns\ #1}, coltitle=mcblue,
  attach boxed title to top left={yshift=-2pt, xshift=6pt},
  boxed title style={boxrule=0pt, colback=mcblue!12, arc=2pt},
}
\newtcolorbox{idea}[1][Idea clave]{
  enhanced, breakable, colback=mcgold2!80, colframe=mcgold, boxrule=0.6pt,
  left=6pt, right=6pt, top=5pt, bottom=5pt, arc=3pt,
  title={\faLightbulb\ #1}, fonttitle=\bfseries, coltitle=mcgold,
  attach boxed title to top left={yshift=-2pt, xshift=6pt},
  boxed title style={boxrule=0pt, colback=mcgold2, arc=2pt},
}
\newtcolorbox{rigor}[1][Rigor: cuidado con lo que se afirma]{
  enhanced, breakable, colback=mcred2!70, colframe=mcred, boxrule=0.6pt,
  left=6pt, right=6pt, top=5pt, bottom=5pt, arc=3pt,
  title={\faBalanceScale\ #1}, fonttitle=\bfseries, coltitle=mcred,
  attach boxed title to top left={yshift=-2pt, xshift=6pt},
  boxed title style={boxrule=0pt, colback=mcred2, arc=2pt},
}
\newtcolorbox{practica}[1][Conexión con la práctica]{
  enhanced, breakable, colback=mcteal2!70, colframe=mcteal, boxrule=0.6pt,
  left=6pt, right=6pt, top=5pt, bottom=5pt, arc=3pt,
  title={\faCogs\ #1}, fonttitle=\bfseries, coltitle=mcteal,
  attach boxed title to top left={yshift=-2pt, xshift=6pt},
  boxed title style={boxrule=0pt, colback=mcteal2, arc=2pt},
}
\newtcolorbox{historia}[1][Sabías que\ldots]{
  enhanced, breakable, colback=mcsepia2!80, colframe=mcsepia, boxrule=0.6pt,
  left=6pt, right=6pt, top=5pt, bottom=5pt, arc=3pt,
  title={\faBookOpen\ #1}, fonttitle=\bfseries, coltitle=mcsepia,
  attach boxed title to top left={yshift=-2pt, xshift=6pt},
  boxed title style={boxrule=0pt, colback=mcsepia2, arc=2pt},
}
\newtcolorbox{trampa}[1][Trampa común]{
  enhanced, breakable, colback=mcviolet2!70, colframe=mcviolet, boxrule=0.6pt,
  left=6pt, right=6pt, top=5pt, bottom=5pt, arc=3pt,
  title={\faExclamationTriangle\ #1}, fonttitle=\bfseries, coltitle=mcviolet,
  attach boxed title to top left={yshift=-2pt, xshift=6pt},
  boxed title style={boxrule=0pt, colback=mcviolet2, arc=2pt},
}
\newtcolorbox{lab}[1][Laboratorio]{
  enhanced, breakable, colback=mcgreen2!60, colframe=mcgreen, boxrule=0.8pt,
  left=6pt, right=6pt, top=5pt, bottom=5pt, arc=3pt,
  title={\faFlask\ #1}, fonttitle=\bfseries, coltitle=mcgreen,
  attach boxed title to top left={yshift=-2pt, xshift=6pt},
  boxed title style={boxrule=0pt, colback=mcgreen2, arc=2pt},
}

% entornos matemáticos auxiliares (demostraciones numeradas)
\newcommand{\demostracion}[1]{\par\smallskip\noindent\textbf{\color{mcgreen}Demostración#1.}\ }
\newcommand{\findem}{\hfill$\square$\par\smallskip}

% ======================================================================
%  MACROS MATEMÁTICOS
% ======================================================================
\newcommand{\F}{\mathbb{F}}
\newcommand{\GF}[1]{\mathbb{F}_{#1}}
\newcommand{\wt}{\operatorname{wt}}
\newcommand{\supp}{\operatorname{supp}}
\newcommand{\rank}{\operatorname{rank}}
\newcommand{\Span}{\operatorname{span}}
\newcommand{\im}{\operatorname{im}}
\newcommand{\Kern}{\operatorname{ker}}
\newcommand{\dist}{\operatorname{d}}
\newcommand{\HashH}{\mathsf{H}}         % función hash
\newcommand{\Hmat}{\mathbf{H}}          % matriz de paridad (genérica)
\newcommand{\Gmat}{\mathbf{G}}          % matriz generadora
\newcommand{\Tmat}{\mathbf{T}}          % clave pública de Classic McEliece
\newcommand{\Smat}{\mathbf{S}}
\newcommand{\Pmat}{\mathbf{P}}
\newcommand{\one}{\mathbf{1}}
\newcommand{\ve}{\mathbf{e}}
\newcommand{\vx}{\mathbf{x}}
\newcommand{\vc}{\mathbf{c}}
\newcommand{\vu}{\mathbf{u}}
\newcommand{\vz}{\mathbf{z}}
\newcommand{\vv}{\mathbf{v}}
\newcommand{\vs}{\mathbf{s}}
\newcommand{\vf}{\mathbf{f}}
\newcommand{\va}{\boldsymbol{\alpha}}
\newcommand{\bit}{\{0,1\}}
\newcommand{\Ftwo}{\F_2}
\newcommand{\FF}[1]{\F_{2^{#1}}}
% \FF4 funciona como \FF{4}: en TeX, el nombre de una secuencia de control
% termina en el primer carácter no alfabético, y \FF toma el «4» como argumento.
\newcommand{\Syn}{\operatorname{Synd}}
\newcommand{\OWCPA}{\mathsf{OW\text{-}CPA}}
\newcommand{\INDCCA}{\mathsf{IND\text{-}CCA2}}
\newcommand{\Adv}{\operatorname{Adv}}
\DeclareMathOperator{\lcm}{lcm}
\DeclareMathOperator{\Tr}{Tr}
\DeclareMathOperator{\ord}{ord}
\DeclareMathOperator{\MLD}{MLD}
\DeclareMathOperator{\ISD}{ISD}

% grados y otras marcas tipográficas
\newcommand{\grad}{^{\circ}}

% dígitos hexadecimales y cadenas de bits dentro de fórmulas
\newcommand{\hx}[1]{\text{\texttt{#1}}}

% celdas de matrices de bits (figuras)
\newcommand{\Bcell}[2]{\draw[bit!35,line width=0.15pt] (#1*0.155,-#2*0.155) rectangle ++(0.155,0.155);}
\newcommand{\Bfill}[2]{\fill[mcblue] (#1*0.155,-#2*0.155) rectangle ++(0.155,0.155);}
\newcounter{bitpos}
\makeatletter
\newcommand{\bitsrow}[2]{%
  \setcounter{bitpos}{-1}%
  \@for\bit:=#2\do{%
    \stepcounter{bitpos}%
    \Bcell{\value{bitpos}}{#1}%
    \ifnum\bit=1 \Bfill{\value{bitpos}}{#1}\fi
  }%
}
\makeatother
\definecolor{bit}{HTML}{9FB3C8}

% ======================================================================
%  ALGORITMOS Y CÓDIGO
% ======================================================================
\SetKwInOut{Entrada}{Entrada}
\SetKwInOut{Salida}{Salida}
\SetKw{devolver}{devolver}
\SetKw{fallo}{fallo}
\SetKw{repetir}{repetir}

\lstdefinestyle{pystyle}{
  language=Python,
  basicstyle=\ttfamily\scriptsize,
  keywordstyle=\color{mcblue}\bfseries,
  commentstyle=\color{mcgreen}\itshape,
  stringstyle=\color{mcorange},
  numberstyle=\ttfamily\scriptsize\color{mcgray},
  numbers=left, numbersep=8pt,
  breaklines=true, breakatwhitespace=false,
  showstringspaces=false,
  frame=single, rulecolor=\color{mcgray!50},
  backgroundcolor=\color{mclight},
  columns=fullflexible, keepspaces=true,
  tabsize=4,
}
\lstset{style=pystyle}

% ======================================================================
%  REFERENCIAS Y METADATOS
% ======================================================================
\hypersetup{
  pdftitle={Classic McEliece: apuntes sobre el cifrado de clave pública basado en códigos de Goppa},
  pdfauthor={Apuntes de estudio},
  pdfsubject={Criptografía post-cuántica, códigos de Goppa, Classic McEliece},
  pdfkeywords={McEliece, Goppa, criptografía, post-cuántica, KEM, ISO, NIST},
}
\makeindex[intoc,title=Índice alfabético]

% ======================================================================
%  PORTADA
% ======================================================================
\newcommand{\portada}{%
\begin{titlepage}
\thispagestyle{empty}
\begin{tikzpicture}[remember picture,overlay]
  % banda superior
  \fill[mcblue] (current page.north west) rectangle ([yshift=-5.2cm]current page.north east);
  \fill[mcblue2] ([yshift=-5.2cm]current page.north west) rectangle ([yshift=-5.45cm]current page.north east);
  \node[anchor=north west, text=white, font=\small\bfseries] at ([xshift=2cm,yshift=-1.4cm]current page.north west)
      {\faLock\ \ CLASSIC McELIECE};
  \node[anchor=north east, text=white!85, font=\small] at ([xshift=-2cm,yshift=-1.4cm]current page.north east)
      {Apuntes de estudio --- octubre de 2026};
\end{tikzpicture}

\vspace*{4.4cm}
{\color{mcblue}\Huge\bfseries Classic McEliece\par}
\vspace{0.5em}
{\color{mcgray}\LARGE El cifrado de clave pública basado en\\[0.2em] códigos de Goppa\par}
\vspace{1.4cm}

% esquema decorativo de matriz de bits
\begin{center}
\begin{tikzpicture}[scale=1]
  \foreach \r/\row in {0/{1,0,1,1,0,1,0,0,1,1,0,1,0,1,1,0},
                       1/{0,1,1,0,0,1,1,0,1,0,1,0,0,1,0,1},
                       2/{1,1,0,1,0,0,1,1,1,0,0,1,1,0,0,1},
                       3/{0,1,0,0,1,1,0,1,0,1,1,0,1,1,0,0},
                       4/{1,0,0,1,1,0,1,0,0,1,0,1,1,0,1,1},
                       5/{1,1,1,0,0,0,1,1,1,1,0,0,1,0,1,0},
                       6/{0,1,0,1,0,1,1,0,0,1,1,1,0,0,1,1},
                       7/{1,0,1,1,1,0,0,1,0,0,1,0,1,1,0,0},
                       8/{1,1,0,0,1,1,0,1,1,0,0,1,1,0,0,1},
                       9/{0,0,1,1,0,1,1,0,0,1,0,1,0,0,1,1},
                      10/{1,0,0,1,1,0,0,1,1,1,0,0,1,1,0,0},
                      11/{0,1,1,0,0,1,0,1,0,0,1,1,0,0,1,1},
                      12/{1,1,0,1,0,0,1,0,1,1,0,1,1,0,0,1},
                      13/{0,1,1,0,1,1,0,1,1,0,1,0,0,1,1,0},
                      14/{1,0,0,1,1,0,1,1,0,1,0,1,1,0,0,1},
                      15/{0,0,1,1,0,1,0,1,1,0,1,1,0,1,0,0}}
    {\bitsrow{\r}{\row}}
  \node[right=0.8cm, align=center, text=mcblue, font=\small] at (2.6,-1.24)
       {\faArrowRight[regular]\ \ $C=\Hmat\ve$\\[2pt] cifrado compacto\\ (síndrome)};
\end{tikzpicture}
\end{center}

\vspace{0.8cm}
{\color{mcblue}\large De McEliece (1978) al estándar ISO/IEC 18033-2 (2026):\\
teoría, laboratorio numérico y estado del arte\par}
\vfill
\begin{center}
\small\color{mcgray}
Documento didáctico con ejemplos verificables, figuras y ejercicios.\\[2pt]
\emph{Dedicado a la memoria de Robert J. McEliece (1942--2019).}
\end{center}
\vspace{1cm}
\end{titlepage}
}

\begin{document}

\portada

% ======================================================================
%  PREFACIO
% ======================================================================
\chapter*{Prefacio: cómo usar estos apuntes}
\addcontentsline{toc}{chapter}{Prefacio: cómo usar estos apuntes}
\markboth{PREFACIO}{}

Estas notas explican, con el mayor rigor que hemos podido alcanzar, \textbf{cómo funciona
Classic McEliece}: el cifrado de clave pública basado en \emph{códigos binarios de Goppa},
propuesto por Robert J. McEliece en 1978 y estandarizado por ISO/IEC en junio de 2026
(enmienda a la norma ISO/IEC 18033-2). No es un manual histórico: es un recorrido
completo que empieza en la teoría de códigos, construye el esquema pieza a pieza, lo
ejecuta sobre un ejemplo miniatura \emph{totalmente verificable}, y termina analizando el
estado del arte criptoanalítico \textbf{a fecha de octubre de 2026} --- que incluye una
sacudida reciente que ningún texto de hace un año recoge.

\section*{Para quién es este documento}
\addcontentsline{toc}{section}{Para quién es este documento}
\begin{itemize}[leftmargin=1.4em,itemsep=2pt]
  \item Estudiantes de matemáticas, informática o ingeniería con interés en criptografía.
  \item Profesionales que necesitan \emph{entender} (de verdad) qué se despliega en sus
        sistemas, más allá de la casilla «post-cuántico».
  \item Docentes que busquen material listo para clase: figuras, tablas, ejemplos
        numéricos y problemas con solución.
\end{itemize}
Se asume una base de álgebra lineal (matrices, rango, espacios vectoriales sobre $\F_2$)
y cierta soltura con notación matemática. Todo lo demás --- cuerpos finitos, códigos,
decodificación --- se desarrolla desde cero. No se requieren conocimientos previos de
criptografía.

\section*{Tres recorridos posibles}
\addcontentsline{toc}{section}{Tres recorridos posibles}
\begin{itemize}[leftmargin=1.4em,itemsep=2pt]
  \item \textbf{Exprés (2--3 horas):} capítulos \ref{ch:pqc}, \ref{ch:goppa},
        \ref{ch:1978}, \ref{ch:kem}, \ref{ch:keygen}, \ref{ch:encapdecap} y \ref{ch:2026},
        más el laboratorio del capítulo \ref{ch:labkem}.
  \item \textbf{Estándar (un curso corto):} el documento en orden, saltando los
        apartados marcados como «avanzado» y las letras pequeñas.
  \item \textbf{Profundo (investigación):} todo, incluyendo la parte \ref{part:seguridad},
        el debate metodológico de costes del capítulo \ref{ch:isd} y las referencias.
\end{itemize}

\section*{Convenciones y avisos de rigor}
\addcontentsline{toc}{section}{Convenciones y avisos de rigor}
\begin{itemize}[leftmargin=1.4em,itemsep=2pt]
  \item \textbf{Las fuentes primarias mandan.} Todo dato cuantitativo se contrasta con:
        el artículo original de McEliece (1978)~\cite{mceliece1978}; la especificación
        oficial de Classic McEliece del 23 de octubre de
        2022~\cite{cmspec}; su documento de racional de diseño~\cite{cmrationale} y su
        guía para evaluadores de seguridad~\cite{cmsecurity}; el informe NIST IR~8545
        (marzo de 2025)~\cite{nist8545}; y la literatura de ataques de 2026 citada en el
        capítulo~\ref{ch:2026}.
  \item \textbf{Los números llevan fecha.} En 2026 la criptoanalítica de McEliece se
        movió más rápido que en las dos décadas anteriores. Cada estimación de seguridad
        de este documento indica su fuente, su modelo de coste y su fecha. Donde hay
        debate abierto (y lo hay), se muestra el debate, no una conclusión empaquetada.
  \item \textbf{Distinguimos siempre} entre resultados \emph{demostrados},
        \emph{heurísticos} y \emph{estimaciones concretas} --- y entre ataques
        \emph{ejecutables} y \emph{teóricos}. Los recuadros
        \textcolor{mcred}{\faBalanceScale\ «Rigor»} señalan precisamente esos puntos.
  \item \textbf{Convención de vectores.} En las partes \ref{part:fundamentos} y
        \ref{part:mceliece} los vectores son fila (como en el artículo de 1978 y en casi
        todos los libros de códigos). La especificación de Classic McEliece usa vectores
        \emph{columna}; el capítulo~\ref{ch:kem} explica la traducción y a partir de ahí
        seguimos fielmente la especificación. Este cambio está señalizado y no es un
        descuido.
  \item \textbf{El laboratorio miniatura no es seguro.} Los códigos de los capítulos
        \ref{ch:labgoppa} y \ref{ch:labkem} (longitud 16, ``clave'' de 8 bytes) existen
        para enseñar los mecanismos con números reales; cualquier parecido con la
        seguridad real es casualidad. No los use fuera de estas notas. El código Python
        que reproduce \emph{todas} las tablas y cifras aparece en el
        apéndice~\ref{app:codigo}.
  \item \textbf{Bit vs byte.} Los tamaños de claves y cifrados se dan en bytes
        (redondeando hacia arriba cuando hay bits de relleno); los tamaños
        criptográficos y de seguridad, en bits de operación o bits de seguridad, según se
        indique. Usamos el espacio fino como separador de millares (p.\,ej.\ \num{261120}).
\end{itemize}

\begin{roadmap}[Mapa del documento]
La \textbf{parte~\ref{part:fundamentos}} construye las herramientas (códigos lineales,
cuerpos finitos, códigos de Goppa) y termina con un laboratorio sobre un código de Goppa
real de juguete. La \textbf{parte~\ref{part:mceliece}} presenta el esquema de 1978, su
versión dual de Niederreiter y la transformación a KEM que se estandariza. La
\textbf{parte~\ref{part:classic}} recorre la especificación de Classic McEliece al
detalle (generación de claves, encapsulación, desencapsulación, parámetros y
rendimiento) con un segundo laboratorio \emph{de extremo a extremo}. La
\textbf{parte~\ref{part:seguridad}} analiza la seguridad: reducciones, decodificación por
conjuntos de información, ataques estructurales y el terremoto criptoanalítico
de 2026. La \textbf{parte~\ref{part:cierre}} cierra con preguntas frecuentes, ejercicios,
soluciones, glosario y referencias.
\end{roadmap}

\clearpage
\tableofcontents
\clearpage
\listoffigures
\listoftables
\clearpage

% ======================================================================
%  PARTE I — FUNDAMENTOS
% ======================================================================
\part{Fundamentos: de los bits a los códigos de Goppa}\label{part:fundamentos}

\chapter{La amenaza cuántica y el giro post-cuántico}\label{ch:pqc}

\begin{roadmap}
Este capítulo sitúa a Classic McEliece en su contexto y responde a tres
preguntas: \emph{por qué} hace falta criptografía post-cuántica, \emph{qué}
se está estandarizando y \emph{dónde encaja} un esquema de 1978 en los
sistemas de 2026. Es un capítulo de orientación: la maquinaria matemática
empieza en el capítulo~\ref{ch:lineales}.
\end{roadmap}

Toda la criptografía de clave pública que protege internet hoy ---RSA,
Diffie--Hellman, las curvas elípticas--- descansa sobre dos problemas que
llevamos medio siglo sin resolver con métodos clásicos: factorizar enteros
grandes y calcular logaritmos discretos. En 1994, Peter Shor demostró que un
ordenador cuántico suficientemente grande resolvería \emph{ambos} en tiempo
polinómico~\cite{shor1994}. Dos años después, Lov Grover mostró una
aceleración cuadrática para la búsqueda no estructurada~\cite{grover1996},
que debilita ---que no anula--- la criptografía simétrica. La conclusión
de aquellos dos resultados define el campo: \textbf{necesitamos cifrado de
clave pública que siga siendo seguro en un mundo con ordenadores cuánticos},
y que además funcione hoy en nuestros ordenadores clásicos. Eso es la
criptografía \emph{post-cuántica} (o \emph{crypto resiliente ante lo
cuántico}).

\begin{roadmap}[Dos matices de vocabulario]
«Post-cuántico» no significa «cuántico»: los algoritmos post-cuánticos son
clásicos (corren en máquinas normales); lo que cambia es \emph{contra qué}
adversario prometen seguridad. Y «seguro contra Shor» no es el único
criterio: un esquema debe resistir a toda la criptoanalítica clásica
conocida, que para los esquemas de códigos lleva acumulada casi medio siglo
de trabajo.
\end{roadmap}

\section{La amenaza: Shor, Grover y el adversario que espera}

\subsection*{Dos algoritmos que cambian las reglas}

Shor (1994) resuelve en tiempo polinómico la factorización y el logaritmo
discreto ---dos problemas que sostienen RSA, Diffie--Hellman y ECC---, con
una única herramienta: la transformada cuántica de Fourier. En cambio, no
conoce ningún atajo análogo para los problemas \emph{NP-duros} típicos de la
criptografía basada en códigos o retículos: frente a ellos solo ofrece la
aceleración de Grover, cuadrática y genérica. Para cifrado simétrico eso
implica que AES-128 protege, en el mejor modelo cuántico, como AES de 64
bits; de ahí la recomendación de claves simétricas de 256 bits
(seguridad efectiva $\approx 128$ bits). La tabla~\ref{tab:cuanticos}
resume el efecto.

\begin{table}[htbp]
\centering
\caption{Efecto de los algoritmos cuánticos conocidos sobre las familias
criptográficas clásicas. «Seguridad efectiva» es una estimación de libro
mayor; los detalles de los modelos de coste importan (véase el
capítulo~\ref{ch:isd}).}
\label{tab:cuanticos}
\small
\begin{tabularx}{\textwidth}{@{}l l l X@{}}
\toprule
\textbf{Primitiva} & \textbf{Problema} & \textbf{Algoritmo cuántico} & \textbf{Efecto} \\
\midrule
RSA & factorización & Shor (1994) & \textbf{rota} (polinómico) \\
Diffie--Hellman / ECC & log. discreto & Shor (1994) & \textbf{rota} \\
AES / hashes & búsqueda, colisiones & Grover (1996), BHT & clave de $2\lambda$ bits para $\lambda$ de seguridad \\
Códigos (McEliece) & decodificación por síndrome & Grover / BHT & sin atajo conocido; exponente a lo sumo a la mitad (cap.~\ref{ch:isd}) \\
\bottomrule
\end{tabularx}
\end{table}

\begin{rigor}[¿Cuándo llega el ordenador cuántico?]
Nadie ha construido un ordenador cuántico tolerante a fallos capaz de
ejecutar Shor sobre claves reales, y \emph{no} es este documento el lugar
para zanjar cuándo llegará. Lo que sí está documentado: los recursos
estimados crecen y se han refinado a la baja (un análisis de 2021 para
RSA-2048 habla de $\approx 20$ millones de qubits físicos ruidosos y
$\approx 8$ horas~\cite{gidney2021}); los criptogramas \emph{de hoy} pueden
grabarse y descifrarse \emph{mañana} (\emph{harvest now, decrypt later}),
fenómeno decisivo para datos con larga vida útil; y migrar infraestructura
criptográfica lleva una década o más, como ilustra el propio calendario de
NIST (IR 8547): retirada de los algoritmos vulnerables recomendada para
2030 y prohibición desde 2035~\cite{nist8547}. La prudencia se apoya en
esos tres hechos, no en una profecía.
\end{rigor}

\subsection*{Treinta segundos de historia hasta hoy}

La figura~\ref{fig:timeline} recorre los hitos que desembocan en estas
notas. La cronología no está a escala: los años recientes se han espaciado
para poder rotularlos.

\begin{figure}[htbp]
\centering
\begin{tikzpicture}[
    every node/.style={font=\footnotesize, align=center},
    hit/.style={circle, inner sep=1.7pt},
    lab/.style={text width=2.32cm},
  ]
  \draw[mcgray!55, line width=1.3pt] (-0.7,0) -- (11.9,0);

  % --- 1978 ---
  \node[hit, fill=mcsepia] (m1) at (0,0) {};
  \node[lab, above, text=mcsepia] at (0,0.22) {\textbf{1978}\\ \textcolor{black}{McEliece: cifrado de clave pública con códigos de Goppa}};
  \draw[mcsepia!50] (m1) -- (m1|-0,0);
  % --- 1986 ---
  \node[hit, fill=mcsepia] (m2) at (1.3,0) {};
  \node[lab, below, text=mcsepia] at (1.3,-0.22) {\textbf{1986}\\ \textcolor{black}{Niederreiter: la versión dual}};
  % --- 1994 ---
  \node[hit, fill=mcblue2] (m3) at (2.6,0) {};
  \node[lab, above, text=mcblue2] at (2.6,0.22) {\textbf{1994}\\ \textcolor{black}{Shor: factorización y log. discreto en tiempo polinómico}};
  % --- 1996 ---
  \node[hit, fill=mcblue2] (m4) at (3.9,0) {};
  \node[lab, below, text=mcblue2] at (3.9,-0.22) {\textbf{1996}\\ \textcolor{black}{Grover: aceleración cuadrática de la búsqueda}};
  % --- 2016 ---
  \node[hit, fill=mcgreen] (m5) at (5.2,0) {};
  \node[lab, above, text=mcgreen] at (5.2,0.22) {\textbf{2016--17}\\ \textcolor{black}{NIST convoca la estandarización post-cuántica}};
  % --- 2022 ---
  \node[hit, fill=mcgreen] (m6) at (6.5,0) {};
  \node[lab, below, text=mcgreen] at (6.5,-0.22) {\textbf{2022}\\ \textcolor{black}{Se eligen Kyber, Dilithium, FALCON y SPHINCS+; ese año cae SIKE}};
  % --- 2024 ---
  \node[hit, fill=mcgreen] (m7) at (7.8,0) {};
  \node[lab, above, text=mcgreen] at (7.8,0.22) {\textbf{2024}\\ \textcolor{black}{FIPS 203, 204 y 205 (agosto)}};
  % --- 2025 ---
  \node[hit, fill=mcgreen] (m8) at (9.1,0) {};
  \node[lab, below, text=mcgreen] at (9.1,-0.22) {\textbf{2025}\\ \textcolor{black}{NIST elige HQC (IR 8545, marzo)}};
  % --- 2026a ---
  \node[hit, fill=mcviolet] (m9) at (10.4,0) {};
  \node[lab, above, text=mcviolet] at (10.4,0.22) {\textbf{junio 2026}\\ \textcolor{black}{ISO estandariza Classic McEliece (18033-2 Amd.~2)}};
  % --- 2026b ---
  \node[hit, fill=mcred] (m10) at (11.7,0) {};
  \node[lab, below, text=mcred] at (11.7,-0.22) {\textbf{sep--oct 2026}\\ \textcolor{black}{Sacudida criptoanalítica; el BSI desaconseja despliegues nuevos}};
\end{tikzpicture}
\caption{De 1978 a octubre de 2026. En sepia, la prehistoria del cifrado con
códigos; en azul, los algoritmos cuánticos que motivaron el giro; en verde,
el proceso del NIST; en violeta, la estandarización ISO; en rojo, los
acontecimientos de 2026 que se analizan en el capítulo~\ref{ch:2026}.}
\label{fig:timeline}
\end{figure}

\section{Qué se está estandarizando: el proceso del NIST}

En 2016 el NIST abrió una convocatoria pública para elegir algoritmos
post-cuánticos. Recibió 82 propuestas (69 aceptadas como completas y
correctas), que fueron cribándose en rondas sucesivas: 26 candidatos en la
segunda ronda (2019) y 15 en la tercera (2020: 7 finalistas y 8
alternativos). En julio de 2022 seleccionó cuatro algoritmos y en agosto de
2024 publicó las tres primeras normas federales: FIPS~203 (ML-KEM, ex
CRYSTALS-Kyber), FIPS~204 (ML-DSA, ex Dilithium) y FIPS~205 (SLH-DSA, ex
SPHINCS+)~\cite{fips203,fips204,fips205}. A la ronda final de KEMs llegaron
tres candidatos alternativos: BIKE, HQC y Classic McEliece. En marzo de
2025, tras la cuarta ronda de análisis, el NIST eligió \emph{solo} a
HQC~\cite{nist8545}. Classic McEliece quedó fuera de la selección del
NIST ---no por un fallo criptoanalítico, sino porque sus claves públicas de
kilobytes a megabytes encajan mal en los protocolos de internet--- y tomó
otro camino: ISO/IEC.

\begin{table}[htbp]
\centering
\caption{Estado de la estandarización post-cuántica en octubre de 2026
(KEMs y firmas principales). Fuentes: FIPS~203--205, NIST IR~8545 y la
página ISO del equipo Classic McEliece~\cite{cmiso}.}
\label{tab:estado}
\small
\begin{tabularx}{\textwidth}{@{}l l l X@{}}
\toprule
\textbf{Nombre estándar} & \textbf{Antecesor} & \textbf{Familia} & \textbf{Estado (oct.\ 2026)} \\
\midrule
ML-KEM & CRYSTALS-Kyber & retículos & norma (FIPS 203, ago.\ 2024) \\
ML-DSA & CRYSTALS-Dilithium & retículos & norma (FIPS 204, ago.\ 2024) \\
SLH-DSA & SPHINCS+ & funciones hash & norma (FIPS 205, ago.\ 2024) \\
FN-DSA & FALCON & retículos & borrador en curso \\
HQC & HQC & códigos & elegido (mar.\ 2025); norma prevista para 2027 \\
Classic McEliece & McEliece (1978) & códigos & \textbf{ISO/IEC 18033-2 Amd.~2 (jun.\ 2026)} \\
\bottomrule
\end{tabularx}
\end{table}

\begin{historia}[SIKE: cuando la diversidad de hipótesis importa]
En agosto de 2022, semanas después de que el NIST anunciara su primera
selección, Castryck y Decru rompieron SIKE ---candidato basado en
isogenias--- con un ataque que ejecutaba en minutos sobre una clave que la
comunidad consideraba sana. El episodio explica por qué las instituciones
recomiendan \emph{híbridos} (combinar un KEM post-cuántico con ECDH) y por
qué la diversidad de hipótesis matemáticas ---retículos, códigos,
hashes--- se considera un valor en sí misma.
\end{historia}

\section{Dónde encaja Classic McEliece}

El esquema de estas notas es, cronológicamente, el más veterano de todos los
que aparecen en la tabla~\ref{tab:estado}: McEliece lo publicó en
1978~\cite{mceliece1978} y, con la disciplina de usar siempre código
binario de Goppa y parámetros muy holgados, ha sobrevivido sin rupturas a
\emph{cuarenta y ocho años} de criptoanálisis clásico. Sus señas de
identidad:

\begin{itemize}[leftmargin=1.4em,itemsep=2pt]
  \item \textbf{Suposición muy conservadora.} No depende de estructuras
        nuevas ni de problemáticas recientes: su dureza se apoya en el
        problema de la decodificación por síndrome (capítulo~\ref{ch:lineales}),
        para el cual no se ha encontrado ni un atajo estructural en
        código binario de Goppa (capítulo~\ref{ch:estructura}).
  \item \textbf{Claves públicas enormes, cifrados diminutos.} La clave
        pública es de 255 KiB a 1.3 MiB según el conjunto de parámetros,
        pero un cifrado ocupa apenas 96--208 bytes: el esquema «comprime»
        el mensaje de forma agresiva. La figura~\ref{fig:tamanos} compara
        estos tamaños con los de sus vecinos estandarizados.
  \item \textbf{Rendimiento asimétrico pero rápido.} La generación de
        claves es costosa (décimas de segundo a segundos); la
        encapsulación y desencapsulación tardan microsegundos o decenas de
        microsegundos. En cifrado de sesiones, donde la clave pública se
        reutiliza, el coste amortizado es bueno.
  \item \textbf{Es un KEM, no un cifrado.} Desde 2019 el esquema se
        especifica como mecanismo de encapsulación de claves con rechazo
        implícito (capítulo~\ref{ch:kem}), que es como se estandarizó
        en~\cite{cmspec} y, en junio de 2026, en el apéndice a la norma
        ISO/IEC 18033-2~\cite{iso18033,cmiso}.
  \item \textbf{Qué conjuntos estandariza ISO.} La enmienda ISO recoge los
        conjuntos \texttt{mceliece460896}, \texttt{mceliece6688128},
        \texttt{mceliece6960119} y \texttt{mceliece8192128}, cada uno en
        cuatro variantes (\texttt{f}, \texttt{pc} y \texttt{pcf} además de
        la básica); el conjunto más pequeño (\texttt{mceliece348864})
        quedó fuera: no alcanza los 128 bits de seguridad en el modelo
        cuántico (con aceleración de Grover) que exige
        ISO~\cite{cmiso}. Los detalles, en el capítulo~\ref{ch:parametros}.
\end{itemize}

\begin{figure}[htbp]
\centering
\begin{tikzpicture}
\begin{axis}[
  ybar=2.4pt, bar width=7.5pt,
  width=0.93\textwidth, height=6.4cm,
  ymode=log, ymin=90, ymax=8e6,
  ylabel={bytes (escala logarítmica)},
  symbolic x coords={ML-KEM-768, HQC-128, CM-460896, CM-8192128},
  xtick=data, enlarge x limits=0.16,
  ytick={100,1000,10000,100000,1000000},
  yticklabels={$10^2$,$10^3$,$10^4$,$10^5$,$10^6$},
  legend style={at={(0.5,1.03)}, anchor=south, legend columns=2,
                draw=none, font=\footnotesize},
  grid=major, major grid style={draw=mcgray!25},
  axis line style={draw=mcgray!70},
]
\addplot[fill=mcblue2, draw=mcblue] coordinates {
  (ML-KEM-768,1184) (HQC-128,2249) (CM-460896,524160) (CM-8192128,1357824)};
\addplot[fill=mcorange2, draw=mcorange] coordinates {
  (ML-KEM-768,1088) (HQC-128,4497) (CM-460896,156) (CM-8192128,208)};
\legend{clave pública, cifrado}
\end{axis}
\end{tikzpicture}
\caption{Tamaño de clave pública y de cifrado en tres familias
estandarizadas (octubre de 2026). La escala es logarítmica: cada
división vertical es un factor~10. Classic McEliece intercambia claves de
centenares de KiB por cifrados de un centenar de bytes; los esquemas de
retículos y HQC son más equilibrados. Datos: FIPS~203, especificación de
HQC y especificación de Classic McEliece~\cite{fips203,hqcspec,cmspec}.}
\label{fig:tamanos}
\end{figure}

\begin{practica}[Híbridos: cómo se despliega esto mañana]
Un KEM post-cuántico rara vez sustituye a ECDH en solitario: lo habitual es
el modo híbrido (X25519 + ML-KEM-768, por ejemplo), que sobrevive si
\emph{cualquiera} de los dos resultara roto. Classic McEliece no es el
candidato natural para el apretón de manos de TLS (su clave pública no
cabe en un datagrama), pero sí encaja donde el almacenamiento es barato y
la criptografía es de larga vida: firmware firmado, archivos cifrados,
raíces de confianza. Los tamaños de la figura~\ref{fig:tamanos} se
entienden mejor con esa lógica de despliegue.
\end{practica}

\section{Cómo leer el resto del documento}

La \textbf{parte~\ref{part:fundamentos}} construye las tres herramientas
previas: códigos lineales (cap.~\ref{ch:lineales}), cuerpos finitos
(cap.~\ref{ch:cuerpos}) y códigos de Goppa (cap.~\ref{ch:goppa}), con un
laboratorio numérico totalmente verificable al final
(cap.~\ref{ch:labgoppa}). La \textbf{parte~\ref{part:mceliece}} presenta el
criptosistema original de 1978 (cap.~\ref{ch:1978}), la variante dual de
Niederreiter (cap.~\ref{ch:niederreiter}) y la transformación a KEM
(cap.~\ref{ch:kem}). La \textbf{parte~\ref{part:classic}} recorre la
especificación de Classic McEliece al detalle
(caps.~\ref{ch:keygen}--\ref{ch:parametros}) y ejecuta un mini-KEM de
extremo a extremo (cap.~\ref{ch:labkem}). La
\textbf{parte~\ref{part:seguridad}} analiza qué significa romper McEliece:
reducciones (cap.~\ref{ch:seguridad}), decodificación por conjuntos de
información (cap.~\ref{ch:isd}), ataques estructurales
(cap.~\ref{ch:estructura}) y el terremoto criptoanalítico de 2026
(cap.~\ref{ch:2026}). La \textbf{parte~\ref{part:cierre}} reúne preguntas
frecuentes, ejercicios, soluciones, glosario y referencias.

\begin{idea}[El hilo conductor]
Todo el documento responde a una sola pregunta: \emph{¿por qué publicar una
matriz generadora «barajada» de un código con decodificador eficiente da un
criptosistema?} Veremos la idea, su historia, su estandarización y ---en
2026--- sus grietas y sus límites.
\end{idea}

\section{Ejercicios de calentamiento}

\begin{enumerate}[leftmargin=1.6em,itemsep=2pt]
  \item Sin mirar la tabla~\ref{tab:cuanticos}: ¿por qué AES-256 se
        considera suficiente en el mundo post-cuántico y AES-128 no?
  \item Ordene por tamaño de \emph{cifrado} (de menor a mayor):
        ML-KEM-768, HQC-128, mceliece460896. Después, ordene por tamaño de
        \emph{clave pública}.
  \item Explique en dos frases por qué «harvest now, decrypt later» obliga
        a actuar antes de que exista un ordenador cuántico.
  \item \textnormal{(}Investigación\textnormal{)} Busque el comunicado del
        BSI del 1 de octubre de 2026 citado en el
        capítulo~\ref{ch:2026} y resuma en cinco líneas sus dos argumentos
        principales.
\end{enumerate}

\chapter{Códigos lineales: el lenguaje de la corrección de errores}\label{ch:lineales}

\begin{roadmap}
Este capítulo introduce el vocabulario mínimo imprescindible: códigos
lineales binarios, matrices generadora y de paridad, síndromes y
decodificación. Al final aparecen las dos piezas que conectan la teoría de
códigos con la criptografía: el problema de la decodificación por síndrome
y la idea de esconder un decodificador eficiente bajo una matriz
«barajada». Todo lo demás del documento se apoya en este capítulo.
\end{roadmap}

\section{El espacio de Hamming}

La corrección de errores empieza con una decisión de modelado: los datos
son palabras de $n$ bits y el canal puede \emph{cambiar} algunos bits.
Medimos el daño contando posiciones distintas.

\begin{definicion}{Peso y distancia de Hamming}{hamming}
Para $x,y\in\Ftwo^n$:
\begin{itemize}[leftmargin=1.2em,itemsep=1pt]
  \item El \textbf{peso de Hamming} $\wt(x)$ es el número de posiciones no
        nulas de $x$.
  \item La \textbf{distancia de Hamming} $\dist(x,y)=\wt(x-y)$
        \textnormal{(}en $\Ftwo$, $x-y=x+y$\textnormal{)} es el número de
        posiciones en que difieren.
\end{itemize}
$\dist$ es una distancia en el sentido métrico: simétrica, nula solo
entre iguales y con desigualdad triangular.
\end{definicion}

\begin{ejemplo}{A mano}{hamming}
En $\Ftwo^5$: $x=10110$, $y=00111$; $x+y=10001$, luego
$\dist(x,y)=\wt(10001)=2$.
\end{ejemplo}

\begin{definicion}{Código lineal, parámetros}{codigo}
Un \textbf{código lineal binario} $C$ de longitud $n$ es un subespacio
vectorial de $\Ftwo^n$. Si $\dim C=k$ y la distancia mínima
$d=\min\{\dist(c,c')\colon c\neq c'\in C\}$ vale $d$, se dice que $C$ es un
código $[n,k,d]$. La \textbf{tasa} del código es $k/n$.
\end{definicion}

En un código lineal la distancia mínima se lee sin recorrer todos los
pares:

\begin{proposicion}{Distancia mínima y pesos}{dmin}
Si $C$ es lineal y no nulo, $d=\min_{c\in C,\,c\neq 0}\wt(c)$.
\end{proposicion}
\begin{proof}
$\dist(c,c')=\wt(c-c')$ y $c-c'\in C$; recíprocamente todo $c\in C$ es
$c-0$.
\end{proof}

\begin{definicion}{Matriz generadora}{generadora}
Una \textbf{matriz generadora} de $C$ es una matriz $G\in\Ftwo^{k\times n}$
cuyas filas forman una base de $C$; entonces
$C=\{uG\colon u\in\Ftwo^k\}$. El mapa $u\mapsto uG$ es inyectivo y codifica
$k$ bits en $n$ bits. Cuando $G=(\mathrm{I}_k\mid A)$ se dice que el código
es \textbf{sistemático}: los $k$ bits de mensaje aparecen tal cual en las
primeras $k$ posiciones de la palabra.
\end{definicion}

\begin{definicion}{Matriz de paridad}{paridad}
Una \textbf{matriz de paridad} de $C$ es una matriz
$\Hmat\in\Ftwo^{r\times n}$ cuyas filas generan el ortogonal
$C^{\perp}=\{v\in\Ftwo^n\colon v\cdot c=0\ \forall c\in C\}$; entonces
$C=\{\vx\in\Ftwo^n\colon \Hmat\,\vx^{\top}=0\}$ y $r=n-k$ (si las filas de
$\Hmat$ son independientes).
\end{definicion}

Generar y verificar son las dos caras de la misma moneda:

\begin{proposicion}{$G\Hmat^{\top}=0$}{gh}
Con la convención de vectores \emph{fila}: $G\Hmat^{\top}=0$, y
$C$ es exactamente el núcleo de $\Hmat$ (visto como aplicación lineal
$\Ftwo^n\to\Ftwo^{n-k}$, $\vx\mapsto\vx\Hmat^{\top}$).
\end{proposicion}
\begin{proof}
Cada fila de $G$ es un elemento de $C$, ortogonal a cada fila de $\Hmat$,
luego $G\Hmat^{\top}=0$. El núcleo tiene dimensión $n-\rank(\Hmat)=k$ y
contiene a las filas de $G$ (dimensión $k$), así que coincide con $C$.
\end{proof}

La matriz de paridad también certifica la distancia mínima, con un criterio
que usaremos constantemente:

\begin{teorema}{Columnas de $\Hmat$ y distancia mínima}{columnas}
$d$ es el menor número de columnas de $\Hmat$ linealmente dependientes.
En particular $d\ge d_0$ si y solo si cualesquiera $d_0-1$ columnas de
$\Hmat$ son independientes.
\end{teorema}
\begin{proof}
Si $c\in C$ con $\supp(c)=\{i_1,\dots,i_w\}$, entonces
$0=\Hmat c^{\top}=\sum_{j}\text{col}_{i_j}(\Hmat)$: las columnas del
soporte son dependientes, luego $d\ge{}$(menor tamaño dependiente).
Recíprocamente, una dependencia $\sum_{j\in J}\text{col}_j=0$ con
$|J|=w$ produce el codeword $c=\sum_{j\in J}e_j$ de peso $w$ (o menor si
hay cancelaciones), luego (menor tamaño dependiente)$\ \ge d$ para el
mínimo.
\end{proof}

\begin{ejemplo}{El código de repetición}{repeticion}
$C=\{000,111\}\subset\Ftwo^3$ es un $[3,1,3]$: una matriz generadora es
$G=(1\ 1\ 1)$ y una de paridad
$\Hmat=\begin{pmatrix}1&1&0\\1&0&1\end{pmatrix}$. Toda palabra con a lo
sumo un error se corrige por unanimidad: $y=101$ se «vota» a $111$.
\end{ejemplo}

\begin{ejemplo}{El código de Hamming $[7,4,3]$}{hamming74}
Con
\[
G=\begin{pmatrix}
1&0&0&0&1&1&0\\
0&1&0&0&0&1&1\\
0&0&1&0&1&1&1\\
0&0&0&1&1&0&1
\end{pmatrix},
\qquad
\Hmat=\begin{pmatrix}
1&0&1&1&1&0&0\\
1&1&1&0&0&1&0\\
0&1&1&1&0&0&1
\end{pmatrix},
\]
las columnas de $\Hmat$ son todos los vectores no nulos de $\Ftwo^3$
(¡las representaciones binarias de $1,\dots,7$, en algún orden!): no hay
columnas nulas ni repetidas, luego cualquier suma parcial se anula solo en
el conjunto vacío. Por el teorema~\ref{thm:columnas}, $d\ge 3$; como la
fila 3 de $G$ tiene peso 3, $d=3$.
\end{ejemplo}

\section{Decodificación: síndromes}

Enviar un codeword $c$ por un canal ruidoso produce $y=c+e$, donde $e$ es
el \emph{vector de error}. Decodificar es, dado $y$, recuperar un
codeword válido «lo más cercana posible». El instrumento central es el
síndrome.

\begin{definicion}{Síndrome}{sindrome}
El \textbf{síndrome} de $y\in\Ftwo^n$ (respecto de $\Hmat$) es
$\Syn(y)=y\,\Hmat^{\top}\in\Ftwo^{n-k}$.
\end{definicion}

\begin{proposicion}{Propiedades del síndrome}{sindrome}
\begin{enumerate}[leftmargin=1.6em,itemsep=1pt]
  \item $\Syn$ es lineal: $\Syn(y+y')=\Syn(y)+\Syn(y')$.
  \item $\Syn(y)=\Syn(y')$ si y solo si $y-y'\in C$.
  \item $\Syn(y)=\Syn(e)$: el síndrome solo ve el error.
  \item $y$ admite corrección con error $e$ de peso $w$ si y solo si
        $\Syn(y)=\Syn(e)$ para algún $e$ con $\wt(e)=w$.
\end{enumerate}
\end{proposicion}
\begin{proof}
(1) y (3) salen de la linealidad de $x\mapsto x\Hmat^{\top\!}$. Para (2),
$\Syn$ tiene núcleo $C$ (proposición~\ref{prop:gh}); para (4), si
$\Syn(y)=\Syn(e)$ entonces $y-e\in C$, es decir, $y$ es el codeword $y-e$
más el error $e$.
\end{proof}

El conjunto de vectores con un síndrome dado es una \emph{clase
(lateral)} $y+C$; el representante de peso mínimo se llama
\textbf{líder de clase} y es la mejor corrección posible en el sentido de
máxima verosimilitud (canal simétrico binario). Con $d$ como distancia
mínima, si $\wt(e)\le t$ con
\begin{equation}
t=\Bigl\lfloor\frac{d-1}{2}\Bigr\rfloor,
\label{eq:radio}
\end{equation}
la corrección es además \emph{única}: las bolas de radio $t$ centradas en
codewords no se solapan (figura~\ref{fig:esferas}). A $t$ se le llama
capacidad de corrección del código.

\begin{figure}[htbp]
\centering
\begin{tikzpicture}
  \draw[mcblue2!35, rounded corners=8pt, fill=mcsky!25] (-0.4,-2.6) rectangle (9.6,2.35);
  \node[mcgray, font=\small] at (9.15,2.05) {$\Ftwo^{n}$};
  \coordinate (c1) at (2.6,0);
  \coordinate (c2) at (7.0,0);
  \draw[mcblue, densely dashed] (c1) circle (1.5);
  \draw[mcblue, densely dashed] (c2) circle (1.5);
  \fill[mcblue] (c1) circle (2.4pt);
  \node[mcblue, font=\small, below=1pt] at (c1) {$c_1$};
  \fill[mcblue] (c2) circle (2.4pt);
  \node[mcblue, font=\small, below=1pt] at (c2) {$c_2$};
  \draw[<->, mcteal] (2.6,0) -- (1.1,0)
        node[midway, above=0.5pt, font=\footnotesize, text=mcteal] {$t$};
  \fill[mcorange] (3.5,0.5) circle (2.4pt);
  \node[mcorange, font=\small, right=1pt] at (3.5,0.5) {$y$};
  \node[mcorange, font=\footnotesize, align=center] at (4.85,1.9)
       {$\dist(y,c_1)\le t$\\[1pt] $\Rightarrow$ se decodifica a $c_1$};
  \draw[-{Stealth[length=2mm]}, mcorange] (4.32,1.62) -- (3.64,0.66);
  \draw[<->, mcgray] (2.6,-2.2) -- (7.0,-2.2)
        node[midway, below, font=\footnotesize] {$\dist(c_1,c_2)=d$};
  \draw[mcgray!60, dashed] (2.6,-1.6) -- (2.6,-2.35);
  \draw[mcgray!60, dashed] (7.0,-1.6) -- (7.0,-2.35);
\end{tikzpicture}
\caption{Las bolas de radio $t=\lfloor(d-1)/2\rfloor$ alrededor de cada
codeword no se solapan: cualquier palabra recibida suficientemente cerca de
un codeword se corrige de forma \emph{única}. Fuera de esos radios, la
«mejor» corrección puede ser ambigua.}
\label{fig:esferas}
\end{figure}

\begin{ejemplo}{Decodificación por síndrome en Hamming $[7,4,3]$}{decodificar}
Con las matrices del ejemplo~\ref{ej:hamming74}, $t=1$. Supongamos que se
recibe $y=1110001$. Su síndrome es
\[
y\,\Hmat^{\top}=
(1,1,1,0,0,0,1)
\begin{pmatrix}
1&1&0\\
0&1&1\\
1&1&1\\
1&0&1\\
1&0&0\\
0&1&0\\
0&0&1
\end{pmatrix}
= (0,1,1),
\]
que es exactamente la \emph{segunda columna} de $\Hmat$. El diccionario
«síndrome = columna $j$» identifica el error $\hat e=e_2$ y devuelve
$\hat c=y+\hat e=1010001$. Como el código es sistemático, el mensaje es
$\hat u=1010$. La tabla~\ref{tab:sindrome-hamming} muestra el diccionario
completo.
\end{ejemplo}

\begin{table}[htbp]
\centering
\caption{Diccionario de decodificación del Hamming $[7,4,3]$: cada error de
peso uno tiene por síndrome la columna correspondiente de $\Hmat$. Los
siete síndromes agotan los vectores no nulos de $\Ftwo^3$: es un código
\emph{perfecto}.}
\label{tab:sindrome-hamming}
\small
\begin{tabular}{@{}c c@{\hspace{2.2em}}c c@{\hspace{2.2em}}c c@{\hspace{2.2em}}c c@{}}
\toprule
posición $j$ & síndrome & posición $j$ & síndrome & posición $j$ & síndrome & posición $j$ & síndrome \\
\midrule
1 & $110$ & 2 & $011$ & 3 & $111$ & 4 & $101$ \\
5 & $100$ & 6 & $010$ & 7 & $001$ & --- & --- \\
\bottomrule
\end{tabular}
\end{table}

\section{El problema de la decodificación por síndrome}

Hasta aquí todo es álgebra lineal elemental. La dificultad aparece cuando
\emph{no} se conoce la matriz de paridad secreta que hace fácil corregir,
y solo se dispone de un código «como objeto combinatorio». Formalizamos el
problema central de la criptografía basada en códigos, en sus dos versiones
(seguimos la formulación de la especificación de Classic
McEliece~\cite{cmspec}):

\begin{definicion}{Problema de decodificación por síndrome, SD}{sd}
Sea $\Hmat\in\Ftwo^{r\times n}$, $s\in\Ftwo^{r}$ y $w$ un entero.
\begin{itemize}[leftmargin=1.2em,itemsep=1pt]
  \item \textbf{SD de búsqueda (search SD):} encontrar $e\in\Ftwo^n$ con
        $\wt(e)\le w$ y $e\,\Hmat^{\top}=s$ (si existe).
  \item \textbf{SD de decisión (decision SD):} distinguir con ventaja no
        despreciable los pares $(\Hmat,s)$ en los que tal $e$ existe de
        los pares con $\Hmat$ uniforme y $s$ uniforme.
\end{itemize}
\end{definicion}

\begin{historia}[De Prange a la NP-completitud]
Eugene Prange ya planteó en 1962 la decodificación mediante conjuntos de
información~\cite{prange1962}; en 1978, el mismo año del artículo de
McEliece, Berlekamp, McEliece y van Tilborg demostraron que la versión de
decisión de SD es NP-completa~\cite{bcmv1978}. Es decir: el problema en el
que se apoya McEliece es de los que la teoría de la complejidad considera
duros... en el peor caso. Volvemos sobre este matiz justo debajo.
\end{historia}

\begin{rigor}[«NP-duro» no es una demostración de seguridad]
NP-duro es una garantía de \emph{peor caso}: dice que no habrá un algoritmo
polinómico que resuelva \emph{todas} las instancias (salvo que P$=$NP). Un
criptosistema necesita algo más fuerte: que las instancias \emph{de la
distribución concreta que genera las claves} sean duras \emph{en media} y
\emph{con el tamaño real}. Nadie ha demostrado que SD sea duro en media
para las distribuciones relevantes; la confianza en McEliece es
\emph{empírica e histórica}: casi medio siglo de ataques no ha doblado la
respuesta. Además, los códigos de un criptosistema no son aleatorios (tienen
estructura algebraica, la que permite decodificar con la clave secreta), y
esa estructura es justo lo que atacan los ataques «estructurales» de la
parte~\ref{part:seguridad}. Por último, NP-duro no detiene a Shor: los
algoritmos cuánticos conocidos para SD (Grover, paseos cuánticos) recuperan
a lo sumo una raíz cuadrada del exponente, con sobrecostes grandes; véase
la sección~\ref{sec:isd-cuantico}.
\end{rigor}

\begin{ejemplo}{Resolver SD «a lo bruto»}{sd-bruto}
Para el código de Hamming $[7,4,3]$, dado un síndrome se prueban los
vectores de peso 0 y 1 hasta acertar: a lo sumo 8 candidatos
(tabla~\ref{tab:sindrome-hamming}). Para el conjunto
\texttt{mceliece460896} ($n=4608$, $t=96$), en cambio, probar errores de
peso $\le 96$ a lo bruto es del orden de $\binom{4608}{96}\approx
2^{668.6}$ candidatos. Ese abismo ---8 frente a $2^{668}$--- es,
comprimido, todo el argumento de seguridad de este
documento.
\end{ejemplo}

\section{La idea de McEliece: esconder un decodificador}

\begin{figure}[htbp]
\centering
\begin{tikzpicture}[
    node distance=0.7cm,
    box/.style={draw=mcblue, fill=mcsky!40, rounded corners=2.5pt,
                align=center, font=\footnotesize, inner sep=4pt,
                text width=2.75cm, minimum height=1.55cm},
    boxg/.style={draw=mcgreen, fill=mcgreen2, rounded corners=2.5pt,
                align=center, font=\footnotesize, inner sep=4pt,
                text width=2.75cm, minimum height=1.55cm},
    arr/.style={-{Stealth[length=2.4mm]}, thick, mcgray!80}
  ]
  \node[box] (enc) {\textbf{Codificador}\\ $c=u\,G$\\ {\scriptsize ($k$ bits $\to$ $n$ bits)}};
  \node[box, right=of enc] (ch) {\textbf{Canal}\\ $y=c+e$\\ {\scriptsize $\wt(e)\le t$}};
  \node[box, right=of ch] (syn) {\textbf{Síndrome}\\ $s=y\,\Hmat^{\top}$\\ {\scriptsize ($n$ bits $\to$ $n-k$ bits)}};
  \node[boxg, right=of syn] (dec) {\textbf{Decodificador}\\ $s\longmapsto \hat e$\\ {\scriptsize eficiente solo si el\\ código tiene estructura\\ conocida}};
  \draw[arr] (enc) -- (ch);
  \draw[arr] (ch) -- (syn);
  \draw[arr] (syn) -- (dec);
  \node[font=\scriptsize, mcgray, above=0.06cm of enc] {mensaje $u$};
\end{tikzpicture}
\caption{La cadena completa: codificar, atravesar el canal, calcular el
síndrome y decodificar. El público ve $y$ y, en su caso, la matriz $\Hmat$;
la dificultad no está en el álgebra lineal del síndrome, sino en encontrar
un error de peso pequeño que lo explique.}
\label{fig:pipeline}
\end{figure}

Un código con decodificador eficiente es una \emph{puerta secreta}: quien
conoce el código puede corregir errores de peso hasta $t$ en tiempo
polinómico; quien solo ve su matriz generadora, no. La idea fundacional de
McEliece (1978) es:

\begin{idea}[La puerta secreta de McEliece]
Tómese un código $C$ con un decodificador eficiente ---p.~ej.\ un código
de Goppa, capítulo~\ref{ch:goppa}--- y publíquese una versión
\emph{barajada} de sus matrices:
\[
G' = S\,G\,P,
\]
con $S$ invertible y $P$ una matriz de permutación. La matriz $G'$ genera
un código \emph{equivalente} a $C$ (el mismo código salvo permutar las
coordenadas y cambiar de base), pero sin la estructura que hace eficiente
la decodificación. Cifrar es codificar con $G'$ y sumar un error de peso
$\le t$; descifrar exige deshacer $P$, decodificar con la puerta secreta y
deshacer $S$. Para un adversario sin la clave, no queda más que resolver SD
para un código que parece aleatorio.
\end{idea}

Los tres capítulos siguientes construyen ese código especial (los
\emph{códigos de Goppa}); la parte~\ref{part:mceliece} convierte la idea en
criptosistema y la parte~\ref{part:seguridad} pregunta con dureza qué pasa
cuando el adversario intenta \emph{deshacer el barajado} en lugar de
decodificar.

\begin{practica}[Por qué elegir una permutación «cara»]
Reordenar las coordenadas de un código no cambia su distancia mínima ni sus
propiedades algebraicas, pero destruye la relación directa entre
posiciones y elementos del cuerpo finito que hace eficiente la
decodificación. Esa es toda la «trampa»: un isomorfismo oculto. La
contrapartida es que el atacante puede intentar recuperar el isomorfismo
(support splitting, ataques algebraicos, \ldots{}), y ahí se libra la
batalla real de Classic McEliece (capítulos~\ref{ch:estructura}
y~\ref{ch:2026}).
\end{practica}

\subsection*{Notación a partir de aquí}

La tabla~\ref{tab:notacion-lineales} fija los símbolos que usaremos. En
las partes~\ref{part:fundamentos} y~\ref{part:mceliece} los vectores son
\emph{fila}, como en el artículo original de 1978; en el
capítulo~\ref{ch:kem} la especificación de Classic McEliece impone
vectores \emph{columna} y señalaremos la traducción sin drama.

\begin{table}[htbp]
\centering
\caption{Notación básica de códigos lineales.}
\label{tab:notacion-lineales}
\small
\begin{tabular}{@{}ll@{\hspace{2em}}ll@{}}
\toprule
\textbf{Símbolo} & \textbf{Significado} & \textbf{Símbolo} & \textbf{Significado} \\
\midrule
$n$ & longitud de las palabras & $u$ & mensaje ($k$ bits) \\
$k$ & dimensión del código & $c$ & codeword ($n$ bits) \\
$d$ & distancia mínima & $e$ & vector de error \\
$t$ & $\lfloor(d-1)/2\rfloor$: errores corregibles & $y$ & palabra recibida $c+e$ \\
$G$ & matriz generadora & $s$ & síndrome \\
$\Hmat$ & matriz de paridad & $\Syn$ & aplicación síndrome \\
$C$ & el código & $C^\perp$ & código dual (ortogonal) \\
\bottomrule
\end{tabular}
\end{table}

\section{Ejercicios}

\begin{enumerate}[leftmargin=1.6em,itemsep=2pt]
  \item Demuestre la desigualdad triangular para la distancia de Hamming.
  \item Sea $C=\{0000,1100,0011,1111\}$. Calcule una matriz $G$, una
        matriz $\Hmat$ y la distancia mínima $d$.
  \item Para el Hamming $[7,4,3]$ del ejemplo~\ref{ej:hamming74}, decodifique
        $y=0101010$ y $y'=1100110$ e indique el mensaje de 4 bits en cada
        caso.
  \item ¿Cuántas palabras de peso exactamente $t$ tiene el cubo
        $\Ftwo^n$? Compare $n=7,\,t=1$ con $n=3488,\,t=64$ dando el
        exponente en base 2 del segundo caso.
  \item Explique por qué $\Syn(y)=\Syn(e)$ implica que todo codeword de la
        clase $y+C$ dista de $y$ como mucho $\wt(e'')$, donde $e''$ es el
        mínimo peso de un vector con ese síndrome.
  \item \textnormal{(}Síntesis\textnormal{)} Formule en tres líneas qué
        propiedad debe tener un código para que la construcción «publicar
        $G'=SGP$» sea un cifrado razonable, y qué dos recursos (uno
        computacional, uno algebraico) hacen falta para descifrar.
\end{enumerate}

\clearpage

\chapter{Cuerpos finitos: el terreno donde vive Goppa}\label{ch:cuerpos}

\begin{roadmap}
Los códigos de Goppa se construyen sobre cuerpos finitos: el polinomio $g$
y los puntos de evaluación $\alpha_i$ viven en $\GF{q}$ con $q=2^m$. Este
capítulo construye $\GF{2^m}$ desde cero ---aritmética, potencias, traza,
raíces cuadradas--- y fija el modelo de representación que usa la
especificación de Classic McEliece. El ejemplo de referencia es
$\FF4=\F_2[z]/(z^4+z+1)$, el mismo del laboratorio.
\end{roadmap}

\section{¿Por qué cuerpos finitos?}

En el capítulo anterior, todo ocurría sobre $\F_2$: bits, matrices, síndromes.
El salto de McEliece consiste en construir un código binario \emph{a través}
de un cuerpo mayor. La definición de los códigos de Goppa
(capítulo~\ref{ch:goppa}) necesita dos ingredientes:

\begin{itemize}[leftmargin=1.4em,itemsep=2pt]
  \item un polinomio $g$ de grado $t$ con coeficientes en un cuerpo
        \emph{grande} (no en $\F_2$), y
  \item una lista de $n$ elementos distintos $\alpha_0,\dots,\alpha_{n-1}$
        de ese cuerpo, con $g(\alpha_i)\neq 0$.
\end{itemize}
El código resultante es binario, pero sus ecuaciones de paridad están
escritas sobre $\GF{q}$, $q=2^m$. Por eso necesitamos saber calcular en
$\GF{2^m}$ con soltura: sumar (trivial), multiplicar, invertir, elevar al
cuadrado, sacar raíces cuadradas y evaluar polinomios.

\section{Construcción de $\GF{2^m}$}

\begin{idea}[De la aritmética modular a los cuerpos finitos]
Para quien no haya estudiado teoría de anillos, la construcción de $\GF{2^m}$ es
exactamente análoga a la aritmética modular de enteros ($\mathbb{Z}/p\mathbb{Z}$):
en $\mathbb{Z}/p\mathbb{Z}$ los elementos son los restos posibles
$\{0,1,\dots,p-1\}$ de dividir entre un número primo $p$, y la irreducibilidad de
$p$ asegura que todo elemento no nulo tenga inverso. En $\GF{2^m}$, el papel de
$p$ lo asume un polinomio \emph{irreducible} $f(z)$ de grado $m$: los elementos
del cuerpo son todos los polinomios con coeficientes en $\F_2$ de grado menor que
$m$ (los posibles restos al dividir entre $f(z)$), y «módulo $f(z)$» significa
simplemente quedarse con el resto de la división euclídea de polinomios.
\end{idea}

\begin{definicion}{Cuerpo finito $\GF{2^m}$}{gf2m}
Sea $f(z)$ un polinomio \emph{irreducible} de grado $m$ con coeficientes en
$\F_2$. El anillo cociente
\[
\GF{2^m} \;=\; \F_2[z]/(f(z))
\]
es un cuerpo con $2^m$ elementos. Cada elemento se escribe de forma única
como un polinomio de grado $<m$,
\[
a \;=\; c_0 + c_1 z + \cdots + c_{m-1}z^{m-1},\qquad c_i\in\F_2,
\]
y se representa como el vector de bits $(c_0,c_1,\dots,c_{m-1})$
\emph{(}el bit $i$ es el coeficiente de $z^i$\emph{)}. La suma es el XOR
de bits; el producto es la multiplicación de polinomios seguida de
reducción módulo $f(z)$.
\end{definicion}

\begin{observacion}{Existencia y unicidad}{gf-existe}
Para cada $m\ge 1$ existe un $f$ irreducible de grado $m$ sobre $\F_2$, y
todos los cuerpos de $2^m$ elementos son isomorfos entre sí (se denotan
$\GF{2^m}$ sin ambigüedad). La irreducibilidad es esencial: si $f$ se
factoriza, el anillo tiene divisores de cero y no es un cuerpo. La
demostración estándar (que $\F_2[z]/(f)$ es cuerpo si y solo si $f$ es
irreducible, y que $\GF{2^m}$ existe) puede verse en Lidl y
Niederreiter~\cite{lidl}.
\end{observacion}

\begin{ejemplo}{El cuerpo más pequeño con extensión: $\GF{4}$}{gf4}
Con $f(z)=z^2+z+1$ (el único polinomio irreducible de grado 2 sobre
$\F_2$), los elementos son $0,1,z,1+z$. Multiplicación:
$z\cdot z=z^2\equiv z+1$ (pues $z^2=z+1$ módulo $f$), y por tanto
$(1+z)^2=(1+z)+(z+z^2)=1+z^2=1+(z+1)=z$. Todos los productos no nulos
están en $\{1,z,1+z\}$: no hay divisores de cero.
\end{ejemplo}

\subsection*{El elemento primitivo y los logaritmos}

Un hecho estructural clave: el grupo multiplicativo
$\GF{2^m}^\times=\GF{2^m}\setminus\{0\}$ es \emph{cíclico} de orden
$2^m-1$~\cite{lidl}. Un generador $z$ se llama \emph{elemento primitivo}:
sus potencias $z^0,z^1,\dots,z^{2^m-2}$ recorren todo el cuerpo salvo el
cero. Multiplicar se vuelve sumar exponentes (tablas de logaritmos), que es
como trabajaremos.

\begin{ejemplo}{$\FF4$ de referencia}{gf16}
Tomamos en toda la parte I el cuerpo
\[
\FF4=\F_2[z]/(z^4+z+1),
\]
que es la construcción usada por el laboratorio y por los conjuntos de
parámetros reales con $m=12,13$ (allí, $z^4+z+1$ se sustituye por un
polinomio irreducible de grado $12$ o $13$). Aquí $z$ es primitivo: sus
potencias recorren los $15$ elementos no nulos (tabla~\ref{tab:gf16} y
figura~\ref{fig:ciclo-z}). La suma de bits es XOR; el producto, polinómico
con reducción $z^4=z+1$.
\end{ejemplo}

\begin{table}[htbp]
\centering
\caption{Potencias de $z$ en $\FF4$, con $z^4=z+1$. Cada elemento se
escribe como polinomio, como vector de bits $(b_3 b_2 b_1 b_0)$ (el bit
$i$-ésimo es el coeficiente de $z^i$, como manda la
especificación~\cite{cmspec}) y como dígito hexadecimal.}
\label{tab:gf16}
\small
\begin{minipage}[t]{0.49\textwidth}\centering
\begin{tabular}{@{}cclc@{}}
\toprule
$i$ & $z^i$ & bits $b_3b_2b_1b_0$ & hex \\
\midrule
0 & $1$ & 0001 & 0x1 \\
1 & $z$ & 0010 & 0x2 \\
2 & $z^2$ & 0100 & 0x4 \\
3 & $z^3$ & 1000 & 0x8 \\
4 & $1+z$ & 0011 & 0x3 \\
5 & $z+z^2$ & 0110 & 0x6 \\
6 & $z^2+z^3$ & 1100 & 0xC \\
7 & $1+z+z^3$ & 1011 & 0xB \\
\bottomrule
\end{tabular}
\end{minipage}\hfill
\begin{minipage}[t]{0.49\textwidth}\centering
\begin{tabular}{@{}cclc@{}}
\toprule
$i$ & $z^i$ & bits $b_3b_2b_1b_0$ & hex \\
\midrule
8 & $1+z^2$ & 0101 & 0x5 \\
9 & $z+z^3$ & 1010 & 0xA \\
10 & $1+z+z^2$ & 0111 & 0x7 \\
11 & $z+z^2+z^3$ & 1110 & 0xE \\
12 & $1+z+z^2+z^3$ & 1111 & 0xF \\
13 & $1+z^2+z^3$ & 1101 & 0xD \\
14 & $1+z^3$ & 1001 & 0x9 \\
\phantom{0} & $0$ & 0000 & 0x0 \\
\bottomrule
\end{tabular}
\end{minipage}
\end{table}

\begin{figure}[htbp]
\centering
\begin{tikzpicture}[scale=1]
  \draw[mcblue2!30, dashed] (0,0) circle (2.3);
  \foreach \i in {0,...,14} {
    \fill[mcblue] ({2.3*cos(90-24*\i)},{2.3*sin(90-24*\i)}) circle (1.9pt);
    \node[font=\scriptsize, mcgray] at
      ({2.72*cos(90-24*\i)},{2.72*sin(90-24*\i)}) {$z^{\i}$};
  }
  \foreach \i in {1,2,4,8} {
    \fill[mcorange] ({2.3*cos(90-24*\i)},{2.3*sin(90-24*\i)}) circle (2.7pt);
    \node[font=\scriptsize, mcorange] at
      ({2.72*cos(90-24*\i)},{2.72*sin(90-24*\i)}) {$z^{\i}$};
  }
  \node[align=center, font=\footnotesize, mcgray] at (0,0.25)
       {las $15$ potencias\\ de $z$};
  \node[align=center, font=\footnotesize, mcorange] at (0,-0.45)
       {en naranja:\\ $z \to z^2 \to z^4 \to z^8$};
\end{tikzpicture}
\caption{El grupo multiplicativo de $\FF4$ es cíclico de orden $15$: las
potencias de $z$ recorren todos los elementos no nulos, y $z^{15}=1$. En
naranja, la órbita de $z$ bajo el automorfismo de Frobenius
$x\mapsto x^2$ (véase la sección~\ref{sec:frobenius}).}
\label{fig:ciclo-z}
\end{figure}

\begin{ejemplo}{Aritmética en $\FF4$ con logaritmos}{aritmetica}
\emph{Producto:} $(z^2+z)(1+z^3)$. Con la tabla:
$z^2+z=z^5$ y $1+z^3=z^{14}$; luego el producto es $z^{19}=z^{4}=1+z$.
\emph{Inverso:} $(1+z)^{-1}$: como $1+z=z^4$, el inverso es
$z^{15-4}=z^{11}=z+z^2+z^3$, y en efecto
$(1+z)(z+z^2+z^3)=z^{15}=1$.
\end{ejemplo}

\begin{trampa}[Sobre la representación y el orden de bits]
Conviene distinguir con cuidado dos formas de escribir un elemento
$a = c_0 + c_1 z + \cdots + c_{m-1}z^{m-1} \in \GF{2^m}$:
\begin{itemize}[leftmargin=1.4em,itemsep=2pt]
  \item \textbf{Como vector de coordenadas:} $(c_0, c_1, \dots, c_{m-1})$,
        con el término independiente $c_0$ en la primera posición (índice $0$),
        que es como se disponen las columnas y filas de las matrices en los
        capítulos~\ref{ch:labgoppa} y~\ref{ch:labkem}.
  \item \textbf{Como entero o hexadecimal:} $b_{m-1}\dots b_1 b_0$ con el bit más
        significativo a la izquierda ($b_i = c_i$). Por ejemplo,
        $z = 0\cdot z^3 + 0\cdot z^2 + 1\cdot z + 0\cdot 1$ se escribe en binario
        como \texttt{0010} ($\hx 2$), y su cuadrado $z^2$ como \texttt{0100}
        ($\hx 4$). Al elevar un monomio al cuadrado, su bit se desplaza hacia la
        izquierda en la notación numérica, pero hacia la derecha en el vector de
        coordenadas.
\end{itemize}
Toda la especificación de Classic McEliece~\cite{cmspec} y los códigos de este
libro respetan estrictamente esta correspondencia.
\end{trampa}

\section{Frobenius, traza y conjugados}\label{sec:frobenius}

\begin{proposicion}{El automorfismo de Frobenius}{frobenius}
La aplicación $\varphi\colon \GF{2^m}\to\GF{2^m}$, $\varphi(x)=x^2$, es un
automorfismo de cuerpos que fija exactamente a $\F_2$. Los
\emph{conjugados} de un elemento $\beta$ son
$\beta,\beta^2,\beta^4,\dots,\beta^{2^{m-1}}$ (con repeticiones si la
órbita es más corta), y $\varphi^m(\beta)=\beta^{2^m}=\beta$: tras $m$
pasos el ciclo se cierra. El producto de los $m$ conjugados (con
multiplicidades) es la norma $\beta^{1+2+\cdots+2^{m-1}}=\beta^{2^m-1}$,
que vale $1$ para $\beta\neq0$.
\end{proposicion}
\begin{proof}
$(x+y)^2=x^2+y^2$ y $(xy)^2=x^2y^2$ en característica~2, luego $\varphi$ es
morfismo inyectivo y, por ser finito, biyectivo. Los elementos fijos de
$\varphi$ son raíces de $X^2-X$, es decir, $0$ y $1$. La órbita de
$\beta\neq0$ bajo $\varphi$ es un ciclo porque $\beta^{2^m}=\beta$
(grupo multiplicativo de orden $2^m-1$).
\end{proof}

La traza linealiza el cuerpo hasta $\F_2$ y será la herramienta para
decidir irreducibilidad de cuadráticas:

\begin{definicion}{Traza}{traza}
$\Tr(\beta)=\beta+\beta^2+\beta^4+\cdots+\beta^{2^{m-1}}\in\F_2$.
\end{definicion}

\begin{proposicion}{Propiedades de la traza}{traza}
$\Tr$ es $\F_2$-lineal; $\Tr(\beta^2)=\Tr(\beta)$; y para $a\in\F_2$,
$\Tr(a)=ma \bmod 2$.
\end{proposicion}
\begin{proof}
Ejercicio para el lector; la linealidad sale de que cada término
$\beta\mapsto\beta^{2^i}$ es lineal en $\F_2$.
\end{proof}

\begin{teorema}{Criterio de irreducibilidad para cuadráticas}{cuadratica}
Sea $\beta\in\GF{2^m}$. El polinomio $x^2+x+\beta$ es irreducible sobre
$\GF{2^m}$ si y solo si $\Tr(\beta)=1$ \textnormal{(}si $\beta=0$ es
reducible: $x(x+1)$\textnormal{)}.
\end{teorema}
\begin{proof}[Bosquejo]
Si existe raíz $\gamma$ (es decir, $\gamma^2+\gamma=\beta$), entonces
$\Tr(\beta)=\Tr(\gamma^2)+\Tr(\gamma)=0$ por~\ref{prop:traza}. Recíprocamente,
si $\Tr(\beta)=0$ y $\beta\neq0$, la ecuación $x^2+x=\beta$ tiene solución
en $\GF{2^m}$ (argumento clásico de conteo: la traza es la suma de los
conjugados y el mapa $x\mapsto x^2+x$ tiene núcleo $\F_2$; véase
\cite{lidl}).
\end{proof}

\begin{ejemplo}{Traza en $\FF4$}{traza}
$\Tr(z^3)=z^3+z^6+z^{12}+z^{9}$; con la tabla,
$z^3{=}0\text{x}8$, $z^6{=}0\text{x}C$, $z^{12}{=}0\text{x}F$,
$z^9{=}0\text{x}A$, y $0\text{x}8\oplus 0\text{x}C\oplus 0\text{x}F\oplus
0\text{x}A=0\text{x}1=1$. Por el teorema~\ref{thm:cuadratica}, $x^2+x+z^3$
es irreducible: es el $g$ del juguete.
\end{ejemplo}

\section{Raíces cuadradas: cuadrar es fácil, y también su inversa}

\begin{proposicion}{Raíz cuadrada en característica 2}{sqrt}
El mapa $x\mapsto x^2$ es $\F_2$-lineal y biyectivo. Su inversa es la raíz
cuadrada, que se calcula como
\[
\sqrt{x}\;=\;x^{\,2^{m-1}} .
\]
\end{proposicion}
\begin{proof}
$(x^2)^{2^{m-1}}=x^{2^m}=x$ para todo $x$ (el grupo multiplicativo tiene
orden $2^m-1$, y el cero se comprueba aparte).
\end{proof}

\begin{ejemplo}{Raíz cuadrada en $\FF4$}{sqrt}
$\sqrt{1+z}$: como $1+z=z^4$, la raíz es $z^2$, y
$(z^2)^2=z^4=1+z$. Con la fórmula: $(z^4)^{8}=z^{32}=z^2$ (los exponentes
se reducen módulo $15$).
\end{ejemplo}

Este detalle, que parece anecdótico, es el corazón del \emph{decodificador
de Patterson} (capítulo~\ref{ch:labgoppa}): en característica 2 sacar
raíces cuadradas de polinomios es tan barato como calcular cuadrados, y eso
hace que la decodificación de Goppa sea eficiente.

\section{Cómo se muestrean cuerpos finitos en la práctica}

La semántica abstracta («elige un polinomio irreducible al azar, un orden
al azar de los elementos») se implementa en Classic McEliece con reglas
precisas que conviene conocer, porque reaparecen en la generación de claves
(capítulo~\ref{ch:keygen}):

\begin{itemize}[leftmargin=1.4em,itemsep=3pt]
  \item \textbf{Elementos a partir de bits.} Los $m$ bits de la convención
        anterior se extraen directamente del flujo pseudoaleatorio (PRG
        basado en SHAKE256 y la semilla $\delta$).
  \item \textbf{Polinomio $g$ (algoritmo \textsf{Irreducible}).} Se genera un
        elemento $\beta$ de una extensión de grado $t$ y se calcula su
        polinomio mínimo sobre $\GF{q}$; si el grado es $t$, es irreducible
        y se devuelve $g$. Si no, se reintenta. La probabilidad de éxito en
        cada intento es mayor que $0{,}99$ para $t\ge16$
        (proposición~\ref{prop:irreducible-uniforme}), así que en la
        práctica se acierta al primer intento.
  \item \textbf{Orden de la lista $L$ (algoritmo \textsf{FieldOrdering}).}
        Se sortean $q$ enteros de $\sigma_2=32$ bits (uno por elemento),
        se ordenan los pares $(a_i,i)$ y resulta la permutación $\pi$ con
        $a_{\pi(0)}<\cdots<a_{\pi(q-1)}$; se define
        $\alpha_i=\sum_j \pi(i)_j\,z^{m-1-j}$, donde $\pi(i)_j$ es el bit
        $j$-ésimo del índice $\pi(i)$. La permutación se guarda comprimida
        como los bits de control de una red de Benes
        (sección~\ref{sec:benes}, en el capítulo~\ref{ch:keygen}).
\end{itemize}

\begin{observacion}{El sentido del «exponente invertido»}{orden}
El uso de $z^{m-1-j}$ (bit $j$ del entero $\mapsto$ exponente $m-1-j$) hace
que $\alpha_i$ \emph{no} sea la representación numérica del índice
$\pi(i)$: los bits se invierten (si $\pi$ fuera la identidad, $L$
empezaría en el cero y seguiría con el elemento de representación
$2^{m-1}$). Este detalle está en la
especificación~\cite{cmspec} y afecta a la reproducibilidad exacta de las
claves; en nuestros ejemplos de juguete usamos el orden natural, lo que
señalaremos donde importe.
\end{observacion}

\begin{practica}[Por qué importa el muestreo]
Dos implementaciones que «hagan lo mismo» pero muestreen el orden de $L$ de
forma distinta producen claves \emph{incompatibles}. La especificación
dedica secciones enteras (representación de polinomios, orden de cuerpos,
redes de Benes) a eliminar cualquier ambigüedad: la reproducibilidad byte a
byte de las claves es un requisito de interoperabilidad, no una manía.
\end{practica}

\section{Ejercicios}

\begin{enumerate}[leftmargin=1.6em,itemsep=2pt]
  \item Demuestre que en $\FF4$, $z^5=z+z^2$ y calcule $z^{-5}$.
  \item Verifique con la tabla que $\Tr(1)=0$ y $\Tr(z)=z+z^2+z^4+z^8$.
        ¿Cuántos elementos de $\FF4$ tienen traza 1?
  \item Calcule $\sqrt{z^3+z^2}$ y $\sqrt{z^2+z+1}$ en $\FF4$.
  \item \label{ej:traza-ej} Demuestre que $\Tr(\beta^2)=\Tr(\beta)$ y que
        $\Tr$ es $\F_2$-lineal, usando la definición.
  \item Explique por qué $x^2+x+1$ es irreducible sobre $\F_2$ pero $x^2+x$
        no. ¿Qué dice el criterio de la traza en cada caso?
  \item \textnormal{(}Síntesis\textnormal{)} Enumere las cuatro operaciones
        que Patterson necesitará sobre $\FF{m}$: suma, producto,
        inverso y raíz cuadrada. ¿Cuál es más cara de implementar en
        hardware y por qué?
\end{enumerate}

\chapter{Códigos de Goppa: el corazón del esquema}\label{ch:goppa}

\begin{roadmap}
Este capítulo define los códigos de Goppa binarios, demuestra que la
matriz de evaluación $t\times n$ sobre $\GF{2^m}$ es una matriz de
paridad, fija los parámetros ($k\ge n-mt$, $d\ge 2t+1$ para $g$ libre de
cuadrados) y explica por qué esta familia concreta ---y no otras--- ha
resistido desde 1970. El capítulo cierra con el puente hacia el
criptosistema: la matriz sistemática $T$, que será la clave pública.
\end{roadmap}

\section{La definición de Goppa}

La idea de V.\,D. Goppa (1970) fue construir un código binario mediante
condiciones de divisibilidad sobre un cuerpo mayor. Fijemos:

\begin{definicion}{Código de Goppa}{goppa}
Sea $q=2^m$, sea $g\in\GF{q}[x]$ mónico, irreducible de grado $t$, y sea
$L=(\alpha_0,\dots,\alpha_{n-1})$ una lista de $n\le q$ elementos
\emph{distintos} de $\GF{q}$ con $g(\alpha_i)\neq0$ para todo $i$. Sea
\[
A(x)=\prod_{j=0}^{n-1}(x-\alpha_j).
\]
El \textbf{código de Goppa $q$-ario} es
\[
\Gamma_q(L,g)=\Bigl\{c\in\GF{q}^{\,n}\;:\;
\sum_{i=0}^{n-1} c_i\,\frac{A(x)}{x-\alpha_i}\;\equiv\;0 \pmod{g(x)}
\Bigr\},
\]
y el \textbf{código de Goppa binario} es la parte binaria,
$\Gamma_2(L,g)=\Gamma_q(L,g)\cap\Ftwo^n$.
\end{definicion}

\begin{idea}[Cómo leer la definición]
El polinomio $A(x)/(x-\alpha_i)$ es el «selector» de la posición $i$:
vale $A'(\alpha_i)\neq0$ en $\alpha_i$ y $0$ en todos los demás $\alpha_j$.
La suma $\sum_i c_i A/(x-\alpha_i)$ es, por tanto, un polinomio de grado
$\le n-1$ que codifica $c$; la condición dice que $g$ lo divide. Dividir
por $x-\alpha_i$ es legítimo porque $g(\alpha_i)\neq0$: el denominador
nunca se anula \emph{módulo} $g$.
\end{idea}

\section{La matriz de paridad: evaluar en los $\alpha_i$}

La definición anterior es elegante pero incómoda para calcular. El primer
resultado sustituye la condición polinómica por $t$ ecuaciones lineales
concretas sobre $\GF{q}$.

\begin{proposicion}{La matriz de evaluación es de paridad}{paridad-goppa}
Con las hipótesis de la definición~\ref{def:goppa}, el código $\Gamma_q$
es el núcleo de la matriz $t\times n$
\[
\tilde{\Hmat}=\Bigl(\,\tilde h_{j,i}\,\Bigr)_{0\le j<t,\;0\le i<n},
\qquad
\tilde h_{j,i}=\frac{\alpha_i^{\,j}}{g(\alpha_i)} .
\]
En particular, expandiendo cada entrada $\tilde h_{j,i}\in\GF{q}$ en sus
$m$ bits se obtiene una matriz binaria $\hat{\Hmat}$ ($mt\times n$) de
paridad para $\Gamma_2$.
\end{proposicion}
\begin{proof}[Demostración]
Como $\gcd(A,g)=1$ (los $\alpha_i$ no son raíces de $g$), multiplicar por
$A$ es invertible módulo $g$, así que la condición de la
definición~\ref{def:goppa} equivale a
\begin{equation}
R_c(x):=\sum_{i=0}^{n-1}\frac{c_i}{x-\alpha_i}\;\equiv\;0\pmod{g(x)}.
\label{eq:goppa-racional}
\end{equation}
Para cada $i$, define $Q_i(x)=\dfrac{g(x)-g(\alpha_i)}{x-\alpha_i}
=\sum_{k=1}^{t} g_k\bigl(x^{k-1}+x^{k-2}\alpha_i+\cdots+\alpha_i^{k-1}\bigr)$,
polinomio de grado $\le t-1$. Entonces
$(x-\alpha_i)Q_i(x)\equiv-g(\alpha_i)\pmod{g(x)}$ y, por tanto,
\[
\frac{1}{x-\alpha_i}\;\equiv\;-\frac{Q_i(x)}{g(\alpha_i)}
\pmod{g(x)} .
\]
Sustituyendo en~\eqref{eq:goppa-racional}: $R_c\equiv0$ si y solo si el
polinomio $\sum_i c_i Q_i(x)/g(\alpha_i)$ de grado $\le t-1$ es
\emph{idénticamente} cero. Mirando el coeficiente de $x^{t-1-j}$ se obtiene
la forma lineal $\sum_i c_i\,v_j(\alpha_i)/g(\alpha_i)$, donde
$v_j(y)=\sum_{l=0}^{j}g_{t-j+l}\,y^{l}$ es un polinomio mónico de grado
exactamente $j$ (en particular $v_0=g_t=1$ y
$v_{t-1}=\bigl(g(y)-g(0)\bigr)/y$). El sistema en las $t$ formas $v_j$ es triangular respecto de las
potencias $\alpha_i^{0},\dots,\alpha_i^{t-1}$ (porque $g_t=1$), luego ambas
familias de ecuaciones tienen el mismo núcleo. La última afirmación es
inmediata: cada ecuación sobre $\GF{q}$ son $m$ ecuaciones binarias una vez
elegida la base (la representación en bits de la definición~\ref{def:gf2m}).
\end{proof}

\noindent La figura~\ref{fig:goppa-check} resume el mecanismo, y la
matriz del juguete (capítulo~\ref{ch:labgoppa}) muestra los números.

\begin{figure}[htbp]
\centering
\begin{tikzpicture}[
    node distance=0.6cm,
    box/.style={draw=mcblue, fill=mcsky!40, rounded corners=2.5pt,
                align=center, font=\footnotesize, inner sep=4pt,
                text width=3.0cm, minimum height=1.5cm},
    boxg/.style={draw=mcgreen, fill=mcgreen2, rounded corners=2.5pt,
                align=center, font=\footnotesize, inner sep=4pt,
                text width=3.0cm, minimum height=1.5cm},
    arr/.style={-{Stealth[length=2.4mm]}, thick, mcgray!80}
  ]
  \node[box] (c) {codeword candidato\\ $c=(c_0,\dots,c_{n-1})\in\Ftwo^n$};
  \node[box, right=of c] (r) {función racional\\ $R_c(x)=\displaystyle\sum_i\frac{c_i}{x-\alpha_i}$};
  \node[box, right=of r] (m) {¿$R_c\equiv 0 \pmod{g(x)}$?\\ {\scriptsize(equivalente: $t$ ecuaciones lineales sobre $\GF{q}$)}};
  \node[boxg, right=of m] (y) {sí: $c\in\Gamma_2$\\ no: fuera del código\\ {\scriptsize(error corregible si $\wt(c-y)\le t$)}};
  \draw[arr] (c) -- (r);
  \draw[arr] (r) -- (m);
  \draw[arr] (m) -- (y);
\end{tikzpicture}
\caption{El código de Goppa como «máquina de verificación»: $c$ pertenece
al código si su función racional asociada es divisible por $g$. La
proposición~\ref{prop:paridad-goppa} convierte esta condición en $t$
ecuaciones lineales sobre $\GF{q}$ ($mt$ binarias).}
\label{fig:goppa-check}
\end{figure}

\section{Parámetros: dimensión y distancia}

\begin{teorema}{Parámetros de un código de Goppa binario}{parametros-goppa}
Sea $\Gamma_2(L,g)=[n,k,d]$ como en la definición~\ref{def:goppa}, con
$g$ libre de cuadrados \textnormal{(}en particular, $g$ irreducible como en
Classic McEliece\textnormal{)}. Entonces
\[
k\;\ge\;n-mt,
\qquad
d\;\ge\;2t+1 .
\]
\end{teorema}

Presentamos la cota de dimensión y la mitad «fácil» de la cota de
distancia, y explicamos con honestidad qué es clásico y qué se demuestra.

\demostracion{(dimensión)}
$\hat{\Hmat}$ tiene $mt$ filas, luego su núcleo tiene dimensión al menos
$n-mt$; ese núcleo \emph{es} $\Gamma_2$ por la
proposición~\ref{prop:paridad-goppa}. (Para $n=q$ y $g$ irreducible, se
tiene de hecho $k=n-mt$; en el juguete, $8=16-2\cdot4$.)\findem

\demostracion{(distancia, paso $t+1$)}
Sea $c\in\Gamma_2$ no nula, $I=\supp(c)$, $w=|I|$, y
$V_I(x)=\prod_{i\in I}(x-\alpha_i)$. Multiplicando~\eqref{eq:goppa-racional}
por $V_I$,
\[
\sum_{i\in I}\frac{V_I(x)}{x-\alpha_i}=V_I'(x)\;\equiv\;0 \pmod{g(x)},
\]
donde $V_I'$ es la derivada formal (definida algebraicamente por
$(x^k)'=kx^{k-1}$, cumpliendo la regla del producto), es decir, $g\mid V_I'$.
Pero $V_I$ es un producto de factores simples, así que $V_I'\neq0$ y
$\deg V_I'\le w-1$ (en característica~2, la derivada de un producto simple
baja un grado más). Si $w\le t$ tendríamos
$0\neq\deg V_I'<t=\deg g$, imposible que $g$ lo divida. Luego $w\ge t+1$
para todo codeword no nulo: $d\ge t+1$. \findem

\begin{teorema}{La versión clásica completa}{goppa-2t}
Con las hipótesis del teorema~\ref{thm:parametros-goppa}, de hecho $d\ge
2t+1$.
\end{teorema}

La mejora de $t+1$ a $2t+1$ es el hecho profundo de la teoría de códigos
de Goppa binarios. La razón es una identidad clásica de característica~2
\emph{(identidad de van Lint)}: si $g$ es libre de cuadrados,
\begin{equation}
\Gamma_2(L,g)=\Gamma_2(L,g^2).
\label{eq:vanlint}
\end{equation}
El código de la derecha tiene módulo de grado $2t$, y el «paso fácil»
aplicado a $g^2$ ---una vez verificado que $g^2$ factoriza en factores
simples \emph{para este propósito}--- da $d\ge 2t+1$. La identidad no es
trivial: expresa que en característica 2 el cuadrado del síndrome racional
de un codeword binario vuelve a ser un síndrome del mismo código con
$g$ reemplazado por $g^2$. En el laboratorio la comprobaremos
exhaustivamente (es exacta para el juguete: $256$ palabras por cada lado);
la demostración estándar puede verse en MacWilliams y
Sloane~\cite{msbook} o van Lint~\cite{vlint}, y la usamos tal cual.

\begin{corolario}{Se corrige cualquier error de peso $\le t$}{corrige-t}
Todo vector de $\Ftwo^n$ está a distancia $\le t$ de, a lo sumo, un
codeword de $\Gamma_2(L,g)$.
\end{corolario}
\begin{proof}
Si $\dist(y,c_1),\dist(y,c_2)\le t$ con $c_1\neq c_2\in\Gamma_2$, entonces
$\wt(c_1-c_2)\le 2t<d$, contradicción con $d\ge2t+1$.
\end{proof}

\begin{rigor}[Qué es teorema y qué es costumbre]
La cota $d\ge 2t+1$ es un teorema con 50 años de antigüedad (Goppa, van
Lint; véanse los libros citados), no una conjetura. La afirmación
«$k=n-mt$» en cambio es \emph{típica}, no universal: depende de que
$\hat{\Hmat}$ tenga rango máximo, cosa que la especificación maneja
explícitamente en \textsf{MatGen} (y que en la práctica siempre ocurre con
los parámetros elegidos). Distinguimos ambas cosas porque en la
parte~\ref{part:seguridad} algunos ataques explotan precisamente
irregularidades estructurales de los códigos, no violaciones de estas
cotas.
\end{rigor}

\begin{ejemplo}{El juguete, anunciado}{juguete-anuncio}
Con $m=4$ ($q=16$), $t=2$, $g=x^2+x+z^3$ y $L=(0,1,z,1+z,\dots)$:
\[
n=16,\qquad k=16-2\cdot4=8,\qquad d\ge 5,
\]
y una comprobación exhaustiva da $d=5$ exactamente, $|\Gamma_2|=256$ y
unicidad de la decodificación para todos los errores de peso $\le2$. Es el
código que estudiaremos pieza a pieza en el capítulo~\ref{ch:labgoppa}.
\end{ejemplo}

\section{¿Por qué Goppa binario y no otra cosa?}

Conviene decirlo pronto, porque es la pregunta que todo el mundo hace:
existen muchas familias de códigos con decodificador eficiente, y casi
todas las que se probaron en esquemas tipo McEliece acabaron
rompiéndose o perdiendo sus ventajas (Reed--Solomon, códigos AG,
quasi-cíclicos, rango, \ldots). Los códigos de Goppa \emph{binarios}
tienen tres propiedades que explican su resistencia:

\begin{enumerate}[leftmargin=1.6em,itemsep=3pt]
  \item \textbf{Densidad de códigos buenos.} Para $n=q$ y casi cualquier
        $g$ irreducible, los parámetros alcanzan la cota de
        Gilbert--Varshamov: hay muchos códigos con distancia relativa
        excelente. El atacante no puede «adivinar» una estructura débil:
        no la hay.
  \item \textbf{La identidad~\eqref{eq:vanlint} y la rigidez en
        característica 2.} Los ataques algebraicos que derriban variantes
        sobre $\GF{q}$ con $q>2$ (capítulo~\ref{ch:estructura}) se estrellan
        en el caso binario: el «cuadrado de Schur» del código no escapa de
        la familia de Goppa, así que no delata el polinomio secreto.
  \item \textbf{Casi 50 años de escrutinio.} Desde 1978 se han publicado
        cientos de artículos intentando romper McEliece con Goppa binario;
        ninguno ha degradado la seguridad más allá de factores polinómicos
        o mejoras de decodificación (capítulos~\ref{ch:isd}
        y~\ref{ch:2026}).
\end{enumerate}

\begin{historia}[Goppa, antes de McEliece]
Los códigos de Goppa (1970) son un poco mayores que el criptosistema: V.\,D.
Goppa los introdujo como códigos correctores «algebraicos» con una cota de
distancia inusualmente buena para su tasa~\cite{goppa1970}. McEliece tuvo
la ocurrencia, ocho años después, de que un código con decodificador
eficiente \emph{y} matriz generadora indistinguible de aleatoria valía
como clave pública. Los códigos de Goppa cumplían ambas cosas entonces y
las siguen cumpliendo hoy; el resto de la historia es de la
parte~\ref{part:mceliece}.
\end{historia}

\section{Ejercicios}

\begin{enumerate}[leftmargin=1.6em,itemsep=2pt]
  \item ¿Por qué la definición exige $g(\alpha_i)\neq 0$? Muestre que si
        algún $g(\alpha_i)=0$, la definición de $\Gamma_q$ se
        \emph{degenera} (todas las condiciones se anulan en la posición
        $i$).
  \item Compruebe que $A(x)/(x-\alpha_i)$ evaluado en $\alpha_i$ vale
        $A'(\alpha_i)=\prod_{j\neq i}(\alpha_i-\alpha_j)$.
  \item Para el juguete ($t=2$), escriba la matriz $\tilde{\Hmat}$ de
        $2\times16$ sobre $\FF4$ como en la especificación: filas
        $\alpha_i^0/g(\alpha_i)$ y $\alpha_i^1/g(\alpha_i)$. (La solución
        está en la salida del laboratorio, pero intente calcular la
        columna $i=0$ y la columna $i=1$ a mano.)
  \item Deduzca del teorema~\ref{thm:parametros-goppa} que \emph{ningún}
        codeword no nulo tiene peso $\le 2t$. Explique por qué esto no
        contradice que $d$ pueda ser mayor que $2t+1$.
  \item \textnormal{(}Conceptual\textnormal{)} Un estudiante propone
        cambiar $g$ irreducible por $g$ con raíces repetidas
        ($g=h^2$). ¿Qué le respondió el teorema~\ref{thm:goppa-2t} y por
        qué McEliece nunca lo intentó?
  \item \textnormal{(}Síntesis\textnormal{)} Escriba en cinco líneas la
        ruta completa: bits $\to$ $g,L$ $\to$ $\Gamma_2$ $\to$ $\hat{\Hmat}$
        $\to$ forma sistemática $(\mathrm{I}\,|\,T)$. ¿Qué pieza de todas
        ellas se publica y cuál se mantiene secreta?
\end{enumerate}

\chapter{Laboratorio 1: el juguete completo}\label{ch:labgoppa}

\begin{roadmap}
En este capítulo ejecutamos, con todos sus números, el código de Goppa
miniatura anunciado en los capítulos~\ref{ch:cuerpos} y~\ref{ch:goppa}:
$m=4$, $n=16$, $t=2$ y $g(x)=x^2+x+z^3$. Construimos las matrices
$\tilde{\Hmat}$, $\hat{\Hmat}$ y la forma sistemática
$(\mathrm{I}_8\,|\,\Tmat)$; calculamos el síndrome de un error de dos
posiciones; lo decodificamos paso a paso con Patterson; verificamos
exhaustivamente la corrección ($120/120$ errores de peso $2$) y la
identidad de van Lint sobre las \num{65536} palabras de $\Ftwo^{16}$; y
terminamos cifrando y descifrando un mensaje con el McEliece de 1978 en
miniatura. Cada cifra de este capítulo es reproducible con el script del
apéndice~\ref{app:codigo}.
\end{roadmap}

\begin{lab}[Lectura y reproducibilidad]
El script \texttt{toy.py} (apéndice~\ref{app:codigo}) imprime todos los
números de este capítulo, en el mismo orden en que aparecen. Tres
convenciones de lectura, avisadas una sola vez:
\begin{itemize}[leftmargin=1.3em,itemsep=2pt]
  \item Los \textbf{elementos de $\GF{16}$} se escriben con un
        \emph{dígito} hexadecimal, como en la tabla de potencias del
        capítulo~\ref{ch:cuerpos}: $9=1+z^3$, $a=z+z^3$,
        $d=1+z^2+z^3$, $f=1+z+z^2+z^3$.
  \item Los \textbf{vectores de bits} se escriben con la \emph{posición
        $0$ a la izquierda}, en el mismo orden que las columnas $j$ de
        las matrices y que la lista $L$: el error con unos en las
        posiciones $10$ y $13$ es $\hx{0000000000100100}$.
  \item Los \textbf{polinomios} sobre $\GF{16}$ se escriben con sus
        coeficientes en ese mismo dígito hexadecimal:
        $S(x)=\hx{4}+\hx{9}x$ significa $S(x)=z^2+(1+z^3)x$.
\end{itemize}
\end{lab}

\section{El montaje: un código de Goppa de longitud 16}

El escenario es el del capítulo~\ref{ch:cuerpos}:
\[
\GF{16}=\F_2[z]/(z^4+z+1),
\]
con $z$ primitivo: sus potencias recorren el grupo multiplicativo, de
orden $15$. El polinomio de Goppa del juguete es
\[
g(x)=x^2+x+z^3,
\]
irreducible porque $\Tr(z^3)=z^3+z^6+z^{12}+z^9=1$
(ejemplo~\ref{ej:traza} y criterio~\ref{thm:cuadratica}). El script
comprueba además que $g(\alpha)\neq0$ para todo $\alpha\in\GF{16}$, que
es la condición de la definición~\ref{def:goppa}. Como puntos de
evaluación tomamos la lista completa del cuerpo en orden natural,
\[
L=(0,\;1,\;z,\;1+z,\;z^2,\;1+z^2,\;z+z^2,\;1+z+z^2,\;z^3,\;\ldots,\;1+z+z^2+z^3),
\]
y los parámetros son
\[
n=q=16,\qquad t=2,\qquad mt=8,\qquad k=n-mt=8,\qquad d\ge 2t+1=5 .
\]
La tabla~\ref{tab:lab-columnas} contiene, columna a columna, todo lo
necesario para reconstruir las matrices: cada fila es una columna de
$\tilde{\Hmat}$, es decir, el par $\bigl(1/g(\alpha_j),\ \alpha_j/g(\alpha_j)\bigr)$.

\begin{table}[htbp]
\centering
\footnotesize
\setlength{\tabcolsep}{4.5pt}
\begin{minipage}[t]{0.49\textwidth}\centering
\begin{tabular}{rrrr}
\toprule
$j$ & $\alpha_j$ & $g(\alpha_j)$ & $(1/g,\ \alpha_j/g)$\\
\midrule
0 & $0$ & $z^3$ & $(\hx f,\hx 0)$\\
1 & $1$ & $z^3$ & $(\hx f,\hx f)$\\
2 & $z$ & $z+z^2+z^3$ & $(\hx 3,\hx 6)$\\
3 & $1+z$ & $z+z^2+z^3$ & $(\hx 3,\hx 5)$\\
4 & $z^2$ & $1+z+z^2+z^3$ & $(\hx 8,\hx 6)$\\
5 & $1+z^2$ & $1+z+z^2+z^3$ & $(\hx 8,\hx e)$\\
6 & $z+z^2$ & $1+z^3$ & $(\hx 2,\hx c)$\\
7 & $1+z+z^2$ & $1+z^3$ & $(\hx 2,\hx e)$\\
\bottomrule
\end{tabular}
\end{minipage}\hfill
\begin{minipage}[t]{0.49\textwidth}\centering
\begin{tabular}{rrrr}
\toprule
$j$ & $\alpha_j$ & $g(\alpha_j)$ & $(1/g,\ \alpha_j/g)$\\
\midrule
8 & $z^3$ & $z^2+z^3$ & $(\hx a,\hx f)$\\
9 & $1+z^3$ & $z^2+z^3$ & $(\hx a,\hx 5)$\\
10 & $z+z^3$ & $z+z^3$ & $(\hx c,\hx 1)$\\
11 & $1+z+z^3$ & $z+z^3$ & $(\hx c,\hx d)$\\
12 & $z^2+z^3$ & $1+z+z^3$ & $(\hx 5,\hx 9)$\\
13 & $1+z^2+z^3$ & $1+z+z^3$ & $(\hx 5,\hx c)$\\
14 & $z+z^2+z^3$ & $1+z^2+z^3$ & $(\hx 4,\hx d)$\\
15 & $1+z+z^2+z^3$ & $1+z^2+z^3$ & $(\hx 4,\hx 9)$\\
\bottomrule
\end{tabular}
\end{minipage}
\caption{Las $16$ columnas del código de Goppa del juguete: valores
$g(\alpha_j)$ y los dos síndromes sobre $\GF{16}$,
$1/g(\alpha_j)$ y $\alpha_j/g(\alpha_j)$. Todas las cifras provienen de
\texttt{toy.py}.}
\label{tab:lab-columnas}
\end{table}

\begin{ejemplo}{Dos columnas a mano}{lab-columnas-mano}
Comprobemos la columna $j=0$ ($\alpha_0=0$) y la $j=1$ ($\alpha_1=1$),
como pedía el ejercicio del capítulo~\ref{ch:goppa}. Para $j=0$:
\[
g(0)=z^3,\qquad
\frac{1}{g(0)}=z^{-3}=z^{12}=1+z+z^2+z^3=\hx{f},\qquad
\frac{\alpha_0}{g(\alpha_0)}=0 .
\]
Para $j=1$, en característica~$2$, $g(1)=1+1+z^3=z^3$, así que
$\frac{1}{g(1)}=\hx f$ y $\frac{\alpha_1}{g(\alpha_1)}=1\cdot\hx f=\hx f$.
La tabla~\ref{tab:lab-columnas} muestra, en efecto, primero la columna
$(\hx f,\hx 0)$ y después $(\hx f,\hx f)$.
\end{ejemplo}

\begin{observacion}{Una simetría de la tabla}{lab-g-simetria}
Como $g(x+1)=(x+1)^2+(x+1)+z^3=x^2+x+z^3=g(x)$, los pares
$\alpha_{2k},\alpha_{2k+1}$ (que difieren en $1$) comparten $g$ y sus
columnas difieren solo en la segunda fila:
$\frac{\alpha+1}{g(\alpha)}=\frac{\alpha}{g(\alpha)}+\frac1{g(\alpha)}$.
Compruébelo en la tabla, por ejemplo en $j=2,3$:
$\hx 6\oplus\hx 3=\hx 5$.
\end{observacion}

\section{Las matrices: $\tilde{\Hmat}$, $\hat{\Hmat}$ y la forma sistemática}

Con la convención de la especificación, la fila $i$ de $\tilde{\Hmat}$
($i=0,1$) tiene entradas $\alpha_j^i/g(\alpha_j)$: exactamente lo que
lista la tabla~\ref{tab:lab-columnas}. En hexadecimal (columnas
$j=0,\dots,15$):
\[
{\small
\tilde{\Hmat}=
\begin{array}{c|cccccccccccccccc}
 & 0&1&2&3&4&5&6&7&8&9&10&11&12&13&14&15\\\hline
i=0 & \hx{f}&\hx{f}&\hx{3}&\hx{3}&\hx{8}&\hx{8}&\hx{2}&\hx{2}&\hx{a}&\hx{a}&\hx{c}&\hx{c}&\hx{5}&\hx{5}&\hx{4}&\hx{4}\\
i=1 & \hx{0}&\hx{f}&\hx{6}&\hx{5}&\hx{6}&\hx{e}&\hx{c}&\hx{e}&\hx{f}&\hx{5}&\hx{1}&\hx{d}&\hx{9}&\hx{c}&\hx{d}&\hx{9}
\end{array}}
\]
Cada entrada es un elemento de $\GF{16}$ y por tanto cuatro bits, con el
bit $\ell$ = coeficiente de $z^\ell$ (la convención del
capítulo~\ref{ch:cuerpos}). Expandiendo cada columna en sus bits y
apilando las filas de $i=0$ encima de las de $i=1$ se obtiene
$\hat{\Hmat}$, la matriz binaria de $8\times16$ que muestra la
figura~\ref{fig:lab-hat}. El rango de $\hat{\Hmat}$ es $8$: el código
$\Gamma_2=\ker\hat{\Hmat}$ tiene dimensión $k=8$, como manda el
teorema~\ref{thm:parametros-goppa}.

\begin{figure}[htbp]
\centering
\begin{tikzpicture}
  % ---------- panel A: columnas 0-7 ----------
  \begin{scope}
    \begin{scope}[scale=1.55]
      \bitsrow{0}{1,1,1,1,0,0,0,0}
      \bitsrow{1}{1,1,1,1,0,0,1,1}
      \bitsrow{2}{1,1,0,0,0,0,0,0}
      \bitsrow{3}{1,1,0,0,1,1,0,0}
      \bitsrow{4}{0,1,0,1,0,0,0,0}
      \bitsrow{5}{0,1,1,0,1,1,0,1}
      \bitsrow{6}{0,1,1,1,1,1,1,1}
      \bitsrow{7}{0,1,0,0,0,1,1,1}
    \end{scope}
    \foreach \r/\lab in {0/{(0,0)},1/{(0,1)},2/{(0,2)},3/{(0,3)},4/{(1,0)},5/{(1,1)},6/{(1,2)},7/{(1,3)}}{
      \node[anchor=east,font=\tiny,inner sep=1pt] at (-0.06,{-\r*0.24025-0.12013}) {\lab};
    }
    \foreach \c in {0,...,7}{
      \node[font=\tiny,inner sep=1pt] at ({\c*0.24025+0.12013},-2.0) {\c};
    }
    \node[font=\scriptsize,color=mcblue] at (0.96,0.33) {$\hat{\Hmat}$ (cols.\ $0$--$7$)};
  \end{scope}
  % ---------- panel B: columnas 8-15 ----------
  \begin{scope}[shift={(2.75,0)}]
    \begin{scope}[scale=1.55]
      \bitsrow{0}{0,0,0,0,1,1,0,0}
      \bitsrow{1}{1,1,0,0,0,0,0,0}
      \bitsrow{2}{0,0,1,1,1,1,1,1}
      \bitsrow{3}{1,1,1,1,0,0,0,0}
      \bitsrow{4}{1,1,1,1,1,0,1,1}
      \bitsrow{5}{1,0,0,0,0,0,0,0}
      \bitsrow{6}{1,1,0,1,0,1,1,0}
      \bitsrow{7}{1,0,0,1,1,1,1,1}
    \end{scope}
    \foreach \r/\lab in {0/{(0,0)},1/{(0,1)},2/{(0,2)},3/{(0,3)},4/{(1,0)},5/{(1,1)},6/{(1,2)},7/{(1,3)}}{
      \node[anchor=east,font=\tiny,inner sep=1pt] at (-0.06,{-\r*0.24025-0.12013}) {\lab};
    }
    \foreach \c in {0,...,7}{
      \node[font=\tiny,inner sep=1pt] at ({\c*0.24025+0.12013},-2.0) {\pgfmathparse{int(\c+8)}\pgfmathresult};
    }
    \node[font=\scriptsize,color=mcblue] at (0.96,0.33) {$\hat{\Hmat}$ (cols.\ $8$--$15$)};
  \end{scope}
  % ---------- flecha ----------
  \draw[-{Stealth[length=2.2mm]},thick,mcgray] (5.05,-0.96) -- (6.45,-0.96)
    node[midway,above,font=\tiny]{reducción};
  % ---------- panel C: (I8 | T) ----------
  \begin{scope}[shift={(6.9,0)}]
    \fill[mcsky!55] (0,0) rectangle ({8*0.24025},-1.922);
    \begin{scope}[scale=1.55]
      \bitsrow{0}{1,0,0,0,0,0,0,0,1,1,0,1,1,0,0,1}
      \bitsrow{1}{0,1,0,0,0,0,0,0,1,1,1,0,0,1,1,0}
      \bitsrow{2}{0,0,1,0,0,0,0,0,0,0,1,0,1,1,1,0}
      \bitsrow{3}{0,0,0,1,0,0,0,0,0,0,0,1,1,1,0,1}
      \bitsrow{4}{0,0,0,0,1,0,0,0,0,1,1,1,1,0,1,0}
      \bitsrow{5}{0,0,0,0,0,1,0,0,1,0,1,1,0,1,0,1}
      \bitsrow{6}{0,0,0,0,0,0,1,0,0,1,0,0,1,0,1,1}
      \bitsrow{7}{0,0,0,0,0,0,0,1,1,0,0,0,0,1,1,1}
    \end{scope}
    \draw[mcblue,thick] ({8*0.24025},0.05) rectangle ({16*0.24025},-1.972);
    \foreach \r in {0,...,7}{
      \node[anchor=east,font=\tiny,inner sep=1pt] at (-0.06,{-\r*0.24025-0.12013}) {\r};
    }
    \foreach \c in {0,4,8,12,15}{
      \node[font=\tiny,inner sep=1pt] at ({\c*0.24025+0.12013},-2.0) {\c};
    }
    \node[font=\scriptsize,color=mcblue] at (0.96,0.33) {$\mathrm{I}_8$};
    \node[font=\scriptsize,color=mcblue] at ({12*0.24025},0.33) {$\Tmat$};
  \end{scope}
\end{tikzpicture}
\caption{La matriz de paridad binaria $\hat{\Hmat}$ del juguete, en dos
mitades (la fila $(i,\ell)$ contiene el bit $\ell$ —el coeficiente de
$z^\ell$— de la entrada $\alpha_j^i/g(\alpha_j)$), y su forma
sistemática $(\mathrm{I}_8\,|\,\Tmat)$ obtenida por operaciones de fila,
que no cambian el núcleo. La mitad derecha $\Tmat$ ($8\times8$) es la
misma que aparecerá como clave pública en el laboratorio~2.}
\label{fig:lab-hat}
\end{figure}

De aquí en adelante identificamos el código con su versión sistemática:
un vector $\vc=(w\,|\,v)$ con $w,v\in\Ftwo^8$ pertenece a $\Gamma_2$ si y
solo si $w=\Tmat v$ (vectores columna), o, en convención de fila, si
$\vc=\vu\,\Gmat$ con $\Gmat=(\Tmat^{\top}\,|\,\mathrm{I}_8)$
(ecuación~\eqref{eq:lab-G} más abajo).

Para terminar de conocer el juguete, el script verifica por fuerza bruta
las \num{65536} palabras de $\Ftwo^{16}$:
\begin{itemize}[leftmargin=1.4em,itemsep=2pt]
  \item $|\Gamma_2|=\num{256}=2^8$: la dimensión efectiva es $k=8$.
  \item $d_{\min}=5$ exactamente, igual que la cota $2t+1$; la
        figura~\ref{fig:lab-pesos} muestra el histograma completo de
        pesos. No hay palabras de peso $1,\dots,4$, y el peso mínimo se
        alcanza $24$ veces.
  \item Los síndromes de los $1+16+120=137$ errores de peso $\le2$ son
        todos distintos entre sí: en este código, cada error de peso
        $\le t$ es la \emph{única} solución de su síndrome. Solo $137$
        de los $256$ síndromes posibles vienen de un error de peso
        $\le2$ (esta asimetría reaparecerá en el laboratorio~2).
\end{itemize}

\begin{figure}[htbp]
\centering
\begin{tikzpicture}
\begin{axis}[
  width=12.6cm, height=5.6cm,
  ybar, bar width=9pt,
  xlabel={peso $w$ de la palabra},
  ylabel={\# palabras},
  xmin=-0.7, xmax=12.7, ymin=0, ymax=52,
  xtick={0,1,2,3,4,5,6,7,8,9,10,11,12},
  ytick={0,10,20,30,40,50},
  ymajorgrids, grid style={mcgray!25},
  axis lines=left, tick align=outside,
  ylabel style={at={(axis description cs:-0.09,0.5)},rotate=90,anchor=south},
]
\addplot[fill=mcblue!35, draw=mcblue] coordinates {
(0,1) (1,0) (2,0) (3,0) (4,0) (5,24) (6,44) (7,40)
(8,45) (9,40) (10,28) (11,24) (12,10)};
\addplot[draw=none, fill=none, nodes near coords, nodes near coords style={font=\tiny,color=mcblue}]
  coordinates {(5,24) (6,44) (7,40) (8,45) (9,40) (10,28) (11,24) (12,10)};
\draw[dashed,mcred] (axis cs:4.5,0) -- (axis cs:4.5,50);
\node[anchor=west,font=\scriptsize,color=mcred] at (axis cs:4.8,47) {$d_{\min}=5$};
\end{axis}
\end{tikzpicture}
\caption{Histograma de pesos de los $256$ codewords del juguete
(\num{65536} palabras examinadas). El hueco de pesos $1$--$4$ refleja
$d_{\min}=5=2t+1$: el código corrige $t=2$ errores.}
\label{fig:lab-pesos}
\end{figure}

\section{El síndrome de un error de dos posiciones}

Elegimos el error que nos acompañará durante todo el capítulo:
\[
e=e_{10}+e_{13}=\hx{0000000000100100}
\qquad(\text{unos en las posiciones }10\text{ y }13).
\]
Sus dos columnas en $\tilde{\Hmat}$ son, según la
tabla~\ref{tab:lab-columnas},
\[
\text{col}_{10}=\binom{\hx c}{\hx 1},\qquad
\text{col}_{13}=\binom{\hx 5}{\hx c},
\]
y sumando coordenada a coordenada sobre $\GF{16}$:
\[
\tilde{\Hmat} e=\binom{\hx c\oplus\hx 5}{\hx 1\oplus\hx c}
=\binom{\hx 9}{\hx d}
=\bigl(1+z^3,\;1+z^2+z^3\bigr).
\]
En bits, descomponiendo según la base $\{1,z,z^2,z^3\}$ de menor a mayor
grado ($1+z^3 \mapsto (1,0,0,1)$ y $1+z^2+z^3 \mapsto (1,0,1,1)$), el
síndrome binario es el vector $\hx{10011011}$, en el orden de filas
$(0,0),(0,1),(0,2),(0,3),(1,0),\dots$ que fija la
figura~\ref{fig:lab-hat} (nótese que al escribir el entero $\hx d$ en binario
habitual con el bit más significativo a la izquierda se lee $1101_2$, mientras
que como vector de coordenadas el coeficiente de $z^0$ ocupa la posición inicial
$(1,0,1,1)$).

Para decodificar, Patterson no usa $\tilde{\Hmat} e$ directamente sino el
\emph{síndrome racional}
\[
S(x)=\sum_{j\in\{10,13\}}\frac{1}{x+\alpha_j}\pmod{g(x)},
\]
la misma cantidad de la ecuación~\eqref{eq:goppa-racional} (en
característica $2$, $x-\alpha=x+\alpha$). Para calcularla, obsérvese que
\[
g(x)+g(\alpha)=(x^2+x+z^3)+(\alpha^2+\alpha+z^3)
=(x+\alpha)^2+(x+\alpha)=(x+\alpha)(x+\alpha+1),
\]
y como $g(x)\equiv0\pmod g$, resulta
$g(\alpha)\equiv(x+\alpha)(x+\alpha+1)\pmod g$; dividiendo por
$(x+\alpha)\,g(\alpha)$:
\begin{equation}
\frac{1}{x+\alpha}\equiv\frac{x+\alpha+1}{g(\alpha)}\pmod{g(x)}.
\label{eq:lab-inversa}
\end{equation}
Sumando sobre las dos posiciones del error:
\[
S(x)=x\underbrace{\sum_{j\in E}\frac1{g(\alpha_j)}}_{=:\,s_0}
+\underbrace{\sum_{j\in E}\frac{\alpha_j+1}{g(\alpha_j)}}_{=:\,s_0+s_1},
\]
es decir, con $s_0=\hx9$ y $s_1=\hx d$ (las componentes de
$\tilde{\Hmat} e$):
\begin{equation}
S(x)=s_0x+(s_0+s_1)=\hx9x+(\hx9\oplus\hx d)
=\hx9x+\hx4=z^2+(1+z^3)x.
\label{eq:lab-S}
\end{equation}

\begin{trampa}[El síndrome no es $S(x)$]
$\tilde{\Hmat} e=(s_0,s_1)=(\hx9,\hx d)$ y $S(x)=\hx4+\hx9x$ son dos
representaciones del mismo dato. La primera es el vector de síndrome
($8$ bits; en el laboratorio~2 será el criptograma); la segunda es el
polinomio módulo $g$ que necesita Patterson. Se relacionan por una
transformación lineal fija —aquí, $S(x)=s_0x+(s_0+s_1)$—, de modo que
los coeficientes de $S$ \emph{no} son $(s_0,s_1)$: confundirlos daría una
decodificación sin sentido.
\end{trampa}

\section{Decodificación de Patterson, paso a paso}

\begin{idea}[La derivada formal en característica $2$ y la descomposición $A^2+xB^2$]
Sobre cualquier cuerpo finito, la derivada de un polinomio se define formalmente
mediante la regla algebraica $(x^k)'=kx^{k-1}$, extendida por linealidad.
En característica $2$ ($2=0$), las potencias pares se anulan al derivar:
$(x^{2k})'=2kx^{2k-1}=0$, mientras que las impares dan $(x^{2k+1})'=(2k+1)x^{2k}=x^{2k}=(x^k)^2$.

Además, por el «sueño del novato» en característica~$2$, elevar al cuadrado conmuta
con la suma: $(u+v)^2=u^2+v^2$. Agrupando cualquier polinomio $\sigma(x)$ en sus términos
de grado par y de grado impar, y extrayendo la raíz cuadrada de sus coeficientes
(que en $\GF{2^m}$ siempre existe y es única por la proposición~\ref{prop:sqrt}),
se obtienen dos polinomios únicos $A(x)$ y $B(x)$ tales que $\sigma(x) = A(x)^2 + x B(x)^2$.
Al derivar:
\[
\sigma'(x) \;=\; (A(x)^2)' + (x B(x)^2)' \;=\; 0 + 1\cdot B(x)^2 + x\cdot (B(x)^2)' \;=\; B(x)^2.
\]
Esta notable simplificación es la base del algoritmo de Patterson.
\end{idea}

Sea $E=\{10,13\}$ y sea $\sigma(x)=\prod_{j\in E}(x+\alpha_j)$ el
\emph{localizador del error}. En característica $2$, derivar un producto
elimina cada factor y deja el producto de los demás:
\[
\sigma'(x)=\sum_{j\in E}\ \prod_{k\neq j}(x+\alpha_k)
=\sum_{j\in E}\frac{\sigma(x)}{x+\alpha_j}.
\]
Multiplicando la definición de $S$ por $\sigma$ y reduciendo módulo $g$
aparece la \emph{ecuación clave}:
\begin{equation}
\sigma(x)\,S(x)\equiv\sigma'(x)\pmod{g(x)}.
\label{eq:lab-keyeq}
\end{equation}
Patterson la resuelve así: en característica $2$ todo polinomio se
escribe como $\sigma=A^2+xB^2$ (separando coeficientes de grado par e
impar con la proposición~\ref{prop:sqrt}), y entonces $\sigma'=B^2$; con
$R:=S^{-1}$, la ecuación~\eqref{eq:lab-keyeq} equivale a
\[
A^2+xB^2\equiv R\,B^2
\;\Longrightarrow\;
A^2\equiv B^2\,(x+R)
\;\Longrightarrow\;
A\equiv B\,W,
\]
donde $W^2\equiv x+R\pmod g$. Con $\tau:=W^{-1}$, queda por encontrar el
par $(A,B)$ tal que $\deg A\le t/2$, $\deg B\le(t-1)/2$ y
$B\equiv A\,\tau\pmod g$; el algoritmo de Euclides parcial lo produce
(véase~\cite{patterson1975}). En el juguete todo se reduce a siete
operaciones, resumidas en la figura~\ref{fig:lab-patterson} y detalladas
aquí con \emph{certificados exactos}: los múltiplos de $g$ que el script
calcula sin dejar resto.

\begin{observacion}{Patterson general: raíces cuadradas y Euclides con parada temprana}{patterson-general}
En el juguete ($t=2$) el proceso termina en una sola división, pero en Classic McEliece
($t=64, 96, 128$) las etapas clave operan así:
\begin{itemize}[leftmargin=1.4em,itemsep=2pt]
  \item \textbf{Raíz cuadrada $W=\sqrt{x+R}\bmod g(x)$:} Como $g$ es irreducible de grado $t$
        sobre $\GF{q}$, el anillo $\GF{q}[x]/(g(x))$ es el cuerpo $\GF{q^t}$. En característica $2$,
        la aplicación $P\mapsto P^2\bmod g$ es un automorfismo lineal invertible; por tanto,
        la raíz cuadrada módulo $g$ siempre existe y es única. En las implementaciones se
        calcula mediante una multiplicación matriz--vector precalculada o por exponenciación.
  \item \textbf{Euclides extendido con parada temprana:} Para hallar $(A,B)$ con
        $\deg A \le t/2$ y $\deg B \le (t-1)/2$ tales que $A \equiv B W \pmod g$, se aplica el
        algoritmo de Euclides extendido a $g(x)$ y $W(x)$. Se detiene en el primer resto
        $r_k$ cuyo grado cumple $\deg r_k \le t/2$. El resto $r_k$ y su correspondiente
        coeficiente de Bézout proporcionan directamente los polinomios $A$ y $B$.
\end{itemize}
\end{observacion}

\begin{enumerate}[leftmargin=1.8em,itemsep=3pt]
  \item $S(x)=\hx4+\hx9x$ (ecuación~\eqref{eq:lab-S}).
  \item $R=S^{-1}=\hx a+\hx7x$;\quad certificado $SR=1+(z+z^3)\,g$.
  \item $x+R=\hx a+\hx6x$.
  \item $W=\sqrt{x+R}=\hx c+\hx7x$;\quad certificado
        $W^2=(x+R)+(z+z^2)\,g$.
  \item $\tau=W^{-1}=\hx3+\hx b x$;\quad certificado
        $W\tau=1+z^2\,g$.
  \item Dividiendo $g$ entre $\tau$: cociente $A=\hx3+\hx5x$ y resto
        $B=\hx d$, de modo que $g=\tau A+B$ \emph{exactamente}; y
        $\deg A=1\le t/2$, $\deg B=0\le(t-1)/2$.
  \item $\sigma=A^2+xB^2=(\hx3+\hx5x)^2+x\,(\hx d)^2
        =(\hx5+\hx2x^2)+\hx e x=\hx5+\hx e x+\hx2x^2$, usando
        $\hx3^2=\hx5$, $\hx5^2=\hx2$ y $\hx d^2=\hx e$: cuadrar en
        característica $2$ duplica los exponentes y eleva al cuadrado
        los coeficientes.
\end{enumerate}

\begin{figure}[htbp]
\centering
\begin{tikzpicture}[
  pbox/.style={draw=mcgreen, fill=mcgreen2!60, rounded corners=2.5pt,
               align=center, font=\footnotesize, inner sep=4pt,
               text width=2.85cm, minimum height=1.2cm},
  arr/.style={-{Stealth[length=2.2mm]}, thick, mcgray}
]
\node[pbox] (b1) {síndrome\\ $S=\hx4+\hx9x$};
\node[pbox, right=4mm of b1] (b2) {$R=S^{-1}=\hx a+\hx7x$\\ {\tiny $SR\equiv1$}};
\node[pbox, right=4mm of b2] (b3) {$x+R=\hx a+\hx6x$};
\node[pbox, right=4mm of b3] (b4) {$W=\sqrt{x+R}=\hx c+\hx7x$\\ {\tiny $W^2\equiv x+R$}};
\node[pbox, below=9mm of b1] (b5) {$\tau=W^{-1}=\hx3+\hx b x$\\ {\tiny $W\tau\equiv1$}};
\node[pbox, right=4mm of b5] (b6) {$g=\tau A+B$\\ $A=\hx3+\hx5x$, $B=\hx d$};
\node[pbox, right=4mm of b6] (b7) {$\sigma=A^2+xB^2$\\ $=\hx5+\hx e x+\hx2x^2$};
\node[pbox, right=4mm of b7] (b8) {raíces $\alpha_{10},\alpha_{13}$\\ $\Rightarrow e'=e$};
\draw[arr] (b1)--(b2); \draw[arr] (b2)--(b3); \draw[arr] (b3)--(b4);
\draw[arr] (b4.south) -- ++(0,-0.30) -| (b5.north);
\draw[arr] (b5)--(b6); \draw[arr] (b6)--(b7); \draw[arr] (b7)--(b8);
\end{tikzpicture}
\caption{La decodificación de Patterson del error $e_{10}+e_{13}$, paso
a paso, con los números del juguete. Cada caja lleva su certificado
exacto (múltiplo de $g$) en letra pequeña.}
\label{fig:lab-patterson}
\end{figure}

La comprobación final es la ecuación~\eqref{eq:lab-keyeq}, que se cumple
con holgura:
\[
\sigma S=\sigma'+(x+\hx e)\,g,\qquad \sigma'=\hx e .
\]
Evaluando $\sigma$ en la lista $L$ (tabla~\ref{tab:lab-sigma}), los
únicos ceros son $\sigma(\alpha_{10})=\sigma(\alpha_{13})=0$; el error
reconstruido es
\[
e'=e_{10}+e_{13}=\hx{0000000000100100}=e .
\]

\begin{table}[htbp]
\centering
\footnotesize
\begin{minipage}[t]{0.48\textwidth}\centering
\begin{tabular}{rrr}
\toprule
$j$ & $\alpha_j$ & $\sigma(\alpha_j)$\\
\midrule
0 & $0$ & $1+z^2$\\
1 & $1$ & $1+z^3$\\
2 & $z$ & $z$\\
3 & $1+z$ & $z+z^2+z^3$\\
4 & $z^2$ & $z+z^2+z^3$\\
5 & $1+z^2$ & $z$\\
6 & $z+z^2$ & $1+z^3$\\
7 & $1+z+z^2$ & $1+z^2$\\
\bottomrule
\end{tabular}
\end{minipage}\hfill
\begin{minipage}[t]{0.48\textwidth}\centering
\begin{tabular}{rrr}
\toprule
$j$ & $\alpha_j$ & $\sigma(\alpha_j)$\\
\midrule
8 & $z^3$ & $1+z+z^2$\\
9 & $1+z^3$ & $1+z+z^3$\\
10 & $z+z^3$ & \textcolor{mcred}{\textbf{$0$}}\\
11 & $1+z+z^3$ & $z^2+z^3$\\
12 & $z^2+z^3$ & $z^2+z^3$\\
13 & $1+z^2+z^3$ & \textcolor{mcred}{\textbf{$0$}}\\
14 & $z+z^2+z^3$ & $1+z+z^3$\\
15 & $1+z+z^2+z^3$ & $1+z+z^2$\\
\bottomrule
\end{tabular}
\end{minipage}
\caption{Evaluación del localizador $\sigma(x)=\hx5+\hx e x+\hx2x^2$ en
los $16$ puntos de $L$. Sus raíces (en rojo) señalan las posiciones del
error: $10$ y $13$.}
\label{tab:lab-sigma}
\end{table}

\begin{observacion}{La simetría de la tabla de $\sigma$}{lab-palindrome}
Los valores $\sigma(\alpha_j)$ se repiten en espejo dentro de cada mitad
de la tabla~\ref{tab:lab-sigma}. No es casualidad: la suma de las raíces
es $\alpha_{10}+\alpha_{13}=1+z+z^2=\alpha_7$, así que el apoyo del
error, $\{10,13\}=\{10,\,10\oplus7\}$, es invariante bajo
$j\mapsto j\oplus7$; y como en esta lista
$\alpha_{j\oplus7}=\alpha_j+\alpha_7$, resulta
$\sigma(x+\alpha_7)=\sigma(x)$ (compruébelo con los coeficientes:
$2c^2+\hx e\,c=c+c=0$ para $c=\alpha_7$). Por eso
$\sigma(\alpha_j)=\sigma(\alpha_{j\oplus7})$ y la tabla es palindrómica.
\end{observacion}

Con esto queda verificado un error. El script repite el procedimiento
sobre los $\binom{16}{2}=120$ errores de peso $2$ posibles y verifica el
resultado: $\num{120}/\num{120}$ correctos. Los casos límite también
están claros: si $S=0$ no hay error ($e'=0$), y si el vector recibido no
está a distancia $\le t$ de ningún codeword, las raíces no reconstruyen
un error consistente; el decodificador del esquema real lo detecta
verificando el resultado (capítulo~\ref{ch:encapdecap}).

\begin{lab}[La promesa del capítulo~\ref{ch:goppa}: van Lint]
Para comprobar la identidad $\Gamma_2(L,g)=\Gamma_2(L,g^2)$ —la
ecuación~\eqref{eq:vanlint}— el script construye la matriz sobre
$\GF{16}$ de $4$ filas con entradas $\alpha_j^i/g(\alpha_j)^2$
($i=0,\dots,3$; $16$ filas al expandir en bits) y contrasta las dos
pertenencias sobre las \num{65536} palabras de $\Ftwo^{16}$: ambas dan
exactamente $\num{256}$ palabras y coinciden palabra a palabra. Es una
corroboración finita de un teorema clásico (véanse
\cite{msbook,vlint}), no una demostración alternativa; lo que sí
muestra es que la construcción del juguete es coherente con la teoría
del capítulo~\ref{ch:goppa}.
\end{lab}

\section{Vista previa: el juguete cifra (McEliece 1978 en miniatura)}

Para cerrar la parte~\ref{part:fundamentos}, montamos con estas piezas el
criptosistema de 1978 en miniatura (el esquema general se estudia en el
capítulo~\ref{ch:1978}). El código es el núcleo de
$(\mathrm{I}_8\,|\,\Tmat)$: un vector $\vc=(w\,|\,v)$ pertenece a
$\Gamma_2$ si y solo si $w=\Tmat v$ (vectores columna). En convención de
fila, esto es $\vc=\vu\,\Gmat$ con
\begin{equation}
\Gmat=(\,\Tmat^{\top}\,|\,\mathrm{I}_8\,),
\label{eq:lab-G}
\end{equation}
cuyas filas son $\Gmat_i=(\text{columna }i\text{ de }\Tmat,\ \ve_i)$:
\[
\Gmat=
\begin{pmatrix}
\hx{1100010110000000}\\
\hx{1100101001000000}\\
\hx{0110110000100000}\\
\hx{1001110000010000}\\
\hx{1011101000001000}\\
\hx{0111010100000100}\\
\hx{0110101100000010}\\
\hx{1001011100000001}
\end{pmatrix}.
\]
A continuación se eligen una matriz $S$ invertible aleatoria y una
permutación $\sigma$ de las posiciones:
\[
\begin{array}{r|cccccccccccccccc}
j&0&1&2&3&4&5&6&7&8&9&10&11&12&13&14&15\\\hline
\sigma(j)&14&2&7&10&4&8&12&6&9&15&11&5&0&3&13&1
\end{array}
\]
con la convención $(\vv P)_j=\vv_{\sigma(j)}$: la coordenada antigua
$\sigma(j)$ pasa a la posición $j$. La clave pública del juguete es
$\Gmat'=S\,\Gmat\,P$, en la figura~\ref{fig:lab-SPG}.

\begin{figure}[htbp]
\centering
\begin{tikzpicture}
  % ---------- S ----------
  \begin{scope}
    \begin{scope}[scale=1.55]
      \bitsrow{0}{0,1,1,0,1,0,1,0}
      \bitsrow{1}{1,0,0,1,0,0,0,1}
      \bitsrow{2}{0,1,0,1,0,0,0,1}
      \bitsrow{3}{0,1,0,1,0,1,1,0}
      \bitsrow{4}{0,1,1,1,0,1,1,0}
      \bitsrow{5}{0,0,1,1,1,1,1,0}
      \bitsrow{6}{1,0,0,1,1,1,0,0}
      \bitsrow{7}{0,0,0,1,0,1,0,1}
    \end{scope}
    \foreach \r in {0,...,7}{
      \node[anchor=east,font=\tiny,inner sep=1pt] at (-0.06,{-\r*0.24025-0.12013}) {\r};
    }
    \foreach \c in {0,...,7}{
      \node[font=\tiny,inner sep=1pt] at ({\c*0.24025+0.12013},-2.0) {\c};
    }
    \node[font=\scriptsize,color=mcblue] at (0.96,0.33) {$S$};
  \end{scope}
  % ---------- G' ----------
  \begin{scope}[shift={(2.7,0)}]
    \begin{scope}[scale=1.55]
      \bitsrow{0}{1,1,1,1,0,0,1,1,1,0,0,1,0,1,0,1}
      \bitsrow{1}{0,0,0,0,1,1,0,1,0,1,1,1,1,0,0,1}
      \bitsrow{2}{0,0,1,0,0,0,0,0,1,1,1,0,1,0,0,1}
      \bitsrow{3}{1,0,0,0,1,0,0,0,1,0,1,0,0,0,1,1}
      \bitsrow{4}{1,1,0,1,0,0,0,0,1,0,1,1,0,0,1,0}
      \bitsrow{5}{1,0,0,1,0,0,1,0,0,0,1,1,0,1,1,1}
      \bitsrow{6}{0,0,0,0,0,1,1,1,0,0,1,1,1,1,1,0}
      \bitsrow{7}{0,1,0,0,1,0,0,1,0,1,1,1,0,1,1,1}
    \end{scope}
    \draw[mcblue,thick] ({8*0.24025},0.06) -- ({8*0.24025},-1.98);
    \foreach \r in {0,...,7}{
      \node[anchor=east,font=\tiny,inner sep=1pt] at (-0.06,{-\r*0.24025-0.12013}) {\r};
    }
    \foreach \c in {0,2,4,6,8,10,12,14}{
      \node[font=\tiny,inner sep=1pt] at ({\c*0.24025+0.12013},-2.0) {\c};
    }
    \node[font=\scriptsize,color=mcblue] at (0.96,0.33) {$\Gmat'$ (cols.\ $0$--$7$)};
    \node[font=\scriptsize,color=mcblue] at ({12*0.24025},0.33) {(cols.\ $8$--$15$)};
  \end{scope}
\end{tikzpicture}
\caption{La matriz aleatoria invertible $S$ ($8\times8$) y la clave
pública del juguete $\Gmat'=S\,\Gmat\,P$ ($8\times16$). La línea
gruesa separa las $8$ primeras coordenadas de las $8$ últimas.}
\label{fig:lab-SPG}
\end{figure}

\textbf{Cifrado.} Tomamos el mensaje $\vu=\hx{01001101}$ (unos en
$i=1,4,5,7$) y el error $\vz$ de peso $2$ con apoyo $\{3,12\}$:
\[
\vz=\hx{0001000000001000}.
\]
La primera mitad es una suma de filas de $\Gmat'$,
$\vu\Gmat'=\bigoplus_{i\in\{1,4,5,7\}}\Gmat'_i$:
\[
\Gmat'_1=\hx{0000110101111001},\qquad
\Gmat'_4=\hx{1101000010110010},
\]
\[
\Gmat'_5=\hx{1001001000110111},\qquad
\Gmat'_7=\hx{0100100101110111},
\]
y sumando las cuatro (equivalentemente, sumando por columnas) más el
error:
\[
\vu\Gmat'=\hx{0000011010001011},
\qquad
\vx=\vu\Gmat'+\vz=\hx{0001011010000011}.
\]
La figura~\ref{fig:lab-cifrado} muestra las tiras de bits de cada paso.

\begin{figure}[htbp]
\centering
\begin{tikzpicture}
  \foreach \r in {0,1,2,3}{
    \foreach \j in {0,...,15}{
      \draw[bit!70,line width=0.2pt] ({\j*0.3},{-\r*0.62-0.3}) rectangle ++(0.3,0.3);
    }
  }
  \foreach \j in {5,6,8,12,14,15}{\fill[mcblue] ({\j*0.3},-0.3) rectangle ++(0.3,0.3);}
  \foreach \j in {3,12}{\fill[mcblue] ({\j*0.3},-0.92) rectangle ++(0.3,0.3);}
  \foreach \j in {3,5,6,8,14,15}{\fill[mcblue] ({\j*0.3},-1.54) rectangle ++(0.3,0.3);}
  \foreach \j in {1,8,9,10,12,13}{\fill[mcblue] ({\j*0.3},-2.16) rectangle ++(0.3,0.3);}
  \foreach \j in {3,12}{\draw[mcred,thick] ({\j*0.3},-1.54) rectangle ++(0.3,0.3);}
  \foreach \j in {0,10}{\draw[mcred,thick] ({\j*0.3},-2.16) rectangle ++(0.3,0.3);}
  \node[anchor=east,font=\scriptsize] at (-0.1,-0.15) {$\vu\Gmat'$};
  \node[anchor=east,font=\scriptsize] at (-0.1,-0.77) {$\vz$};
  \node[anchor=east,font=\scriptsize] at (-0.1,-1.39) {$\vx$};
  \node[anchor=east,font=\scriptsize] at (-0.1,-2.01) {$\vx P^{-1}$};
  \foreach \j in {0,...,15}{
    \node[font=\tiny,inner sep=0pt] at ({\j*0.3+0.15},-2.62) {\j};
  }
\end{tikzpicture}
\caption{El cifrado en tiras: en azul, las posiciones con bit $1$; en
rojo, las posiciones alteradas por el error. De arriba abajo:
$\vu\Gmat'$ (codeword sin error), el error $\vz$, el criptograma
$\vx=\vu\Gmat'+\vz$ y, tras deshacer $P$, la palabra recibida
$\vx P^{-1}$, con el error ya en las posiciones $0$ y $10$.}
\label{fig:lab-cifrado}
\end{figure}

\textbf{Descifrado.} Se deshace la permutación,
$(\vv P^{-1})_j=\vv_{\sigma^{-1}(j)}$:
\[
\vx P^{-1}=\hx{0100000011101100},
\qquad
\supp\bigl(\vz P^{-1}\bigr)=\{\sigma(3),\sigma(12)\}=\{10,0\}.
\]
Ahora el problema es el de decodificar una palabra a distancia $2$ de un
codeword. Patterson (la misma rutina de la
sección anterior; aquí $S(x)=\hx2+\hx3x$ y $\sigma(x)=\hx c\,x+\hx f
x^2$) devuelve
\[
\vz'=\hx{1000000000100000}\qquad(\text{apoyo }\{0,10\}),
\qquad
\vc=\vx P^{-1}+\vz'=\hx{1100000011001100}.
\]
Los últimos $8$ bits de $\vc$ son el mensaje «escalado»
$\vu S=\hx{11001100}$; deshaciendo $S$:
\[
\vu=(\vu S)\,S^{-1}=\hx{01001101},
\]
el mensaje original. Fin del juguete: en el
laboratorio~2 (capítulo~\ref{ch:labkem}) el mismo montaje se convierte
en un KEM completo con SHAKE256, y en el capítulo~\ref{ch:1978} el
esquema de 1978 se presenta en toda su generalidad. Advertencia: que
este juguete cifre \emph{no} significa que proteja algo; cuantificar por
qué es el ejercicio~6 y, en serio, la
parte~\ref{part:seguridad}.

\begin{practica}[Del juguete al esquema real]
La tabla mental del capítulo:
$\GF{16}\leftrightarrow\GF{2^m}$ con $m=12$ o $13$; $\deg g=2
\leftrightarrow t=64,\dots,128$; los $16$ puntos de $L$
$\leftrightarrow$ las $n\approx\num{8000}$ posiciones del orden de
cuerpo (\emph{FieldOrdering}); $\tilde{\Hmat},\hat{\Hmat}$ de $2\times16$ y
$8\times16$ $\leftrightarrow$ $t\times n$ y $mt\times n$; la forma
sistemática $(\mathrm{I}_8\,|\,\Tmat)$ $\leftrightarrow$
$(\mathrm{I}_{mt}\,|\,\Tmat)$ con $\Tmat$ de $mt\times k$; Patterson,
idéntico; el error de peso $2$ de esta sección $\leftrightarrow$ un
error de peso $t$ entre $\binom{n}{t}\approx2^{456}$ posibilidades para
\texttt{mceliece348864}; y el mini-esquema $\leftrightarrow$ el KEM
completo de los capítulos~\ref{ch:kem}--\ref{ch:encapdecap}, con
rechazo implícito y confirmación de texto claro.
\end{practica}

\section{Ejercicios}

\begin{enumerate}[leftmargin=1.6em,itemsep=4pt]
  \item \textbf{La columna $j=2$ a mano.} Con $z^4=z+1$ y la tabla de
        potencias del capítulo~\ref{ch:cuerpos}, verifique que
        $g(z)=z+z^2+z^3$, que $1/g(z)=\hx3$ y que $z/g(z)=\hx6$;
        compare con la tabla~\ref{tab:lab-columnas}.
  \item \textbf{Linealidad del síndrome.} Calcule $\tilde{\Hmat} e_5$,
        $\tilde{\Hmat} e_9$ y $\tilde{\Hmat}(e_5+e_9)$: debe obtener
        $(\hx2,\hx b)$. ¿Por qué $e_5+e_9$ también se decodifica de
        forma única?
  \item \textbf{Un certificado a mano.} Compruebe, usando una y otra vez
        $z^4=z+1$, que $(\hx4+\hx9x)(\hx a+\hx7x)=1+(z+z^3)\,g(x)$; es
        el certificado del paso~2 de Patterson.
  \item \textbf{El caso de un solo error.} Si $E=\{j\}$, demuestre que
        $\sigma(x)=x+\alpha_j$ y compruebe la
        ecuación~\eqref{eq:lab-keyeq} (aquí $\sigma'=1$).
        ¿Qué devuelve Patterson?
  \item \textbf{La permutación, paso a paso.} Con
        $(\vv P)_j=\vv_{\sigma(j)}$, compruebe que
        $(\vx P^{-1})_1=\vx_{\sigma^{-1}(1)}=\vx_{15}=1$ y que
        $(\vx P^{-1})_3=\vx_{\sigma^{-1}(3)}=\vx_{13}=0$. ¿En qué
        posiciones quedan los dos errores de $\vz$?
  \item \textbf{¿Por qué el juguete no es seguro?} Dé dos ataques
        concretos: uno enumerando palabras de $\Ftwo^{16}$ y otro
        usando los $256$ síndromes. Explique en una línea qué cambia al
        pasar a $n\approx\num{8000}$ (parte~\ref{part:seguridad}).
\end{enumerate}

\begin{center}\rule{0.55\textwidth}{0.4pt}\end{center}
\medskip
\noindent\textit{Cierre de la parte~\ref{part:fundamentos}.} Con esto
termina la parte de fundamentos: los códigos lineales y sus parámetros
(capítulos~\ref{ch:pqc} y~\ref{ch:lineales}), la aritmética de cuerpos
finitos (capítulo~\ref{ch:cuerpos}), los códigos de Goppa
(capítulo~\ref{ch:goppa}) y un juguete completo, verificado número a
número, que ya construye matrices, decodifica con Patterson y cifra. La
parte~\ref{part:mceliece} convierte todo esto en el criptosistema de
clave pública: el McEliece de 1978, su dual de Niederreiter y la
transformación a KEM que lleva a Classic McEliece.

\clearpage

% ======================================================================
%  PARTE II — EL CRIPTOSISTEMA DE MCELIECE
% ======================================================================
\part{El criptosistema de McEliece}\label{part:mceliece}

\chapter{El McEliece de 1978}\label{ch:1978}

\begin{roadmap}
Este capítulo presenta el criptosistema original de Robert McEliece
(1978), con el que empezó todo: la construcción, la demostración de que
el descifrado funciona, el argumento de seguridad que lo acompaña y sus
tres debilidades prácticas (tamaño de clave, cifrado sin garantías
modernas y ausencia de firmas). El lector reconocerá de inmediato el
juguete del capítulo~\ref{ch:labgoppa}: este capítulo es aquel
laboratorio con parámetros grandes, y el resto de la
parte~\ref{part:mceliece} convierte el esquema en un KEM.
\end{roadmap}

\begin{historia}[Dos páginas y media que esperaron cuarenta años]
El esquema apareció en el \emph{DSN Progress Report 42-44} del Jet
Propulsion Laboratory (enero--febrero de 1978, pp.~114--116), casi al
mismo tiempo que RSA, y McEliece lo cierra con dos observaciones que la
historia confirmó palabra por palabra: permite «tasas de comunicación
cercanas a $10^6$ bits/s o más», y su algoritmo de descifrado \emph{no}
sirve para firmar. Durante décadas fue una curiosidad teórica, sin
estándares ni implementaciones masivas; la amenaza cuántica
(capítulo~\ref{ch:pqc}) lo devolvió al primer plano, y su descendiente
directo, Classic McEliece, es el esquema que ISO/IEC estandarizó en
junio de 2026 como enmienda a la norma 18033-2~\cite{iso18033,cmiso}.
\end{historia}

\section{La idea: una trampa amable}

En 1976, Diffie y Hellman habían enseñado a construir canales seguros
sin clave compartida, y en 1978 existían ya dos «trampas» concretas
—problemas fáciles con información secreta y difíciles sin ella—: la
factorización de enteros (RSA) y el problema del mochilero
(Merkle--Hellman). McEliece propuso una tercera, que no viene de la
teoría de números sino de los \emph{códigos correctores de errores}. La
asimetría que explota está demostrada en las partes anteriores:

\begin{itemize}[leftmargin=1.4em,itemsep=3pt]
  \item \textbf{Con} la descripción secreta —un polinomio de Goppa $g$
        de grado $t$ y su soporte $L$— existe un decodificador
        eficiente: el algoritmo de Patterson corrige hasta $t$ errores
        en tiempo $O(nt)$~\cite{patterson1975}.
  \item \textbf{Sin} ella, el problema de decodificar un código lineal
        arbitrario a partir de una matriz generadora es NP-completo
        (Berlekamp, McEliece y van Tilborg, 1978)~\cite{bcmv1978}.
\end{itemize}

El truco consiste en esconder el código de Goppa a la vista de todos: se
multiplica su matriz generadora por una matriz invertible densa $S$ y se
permutan las coordenadas con $P$. Ninguna de las dos operaciones cambia
la dimensión, la distancia ni la decodificabilidad del código, pero
$G'=SGP$ pierde toda la apariencia de «matriz de Goppa»: el adversario
solo ve lo que parece una matriz aleatoria $k\times n$.

\begin{idea}[La asimetría en una línea]
Con la clave secreta, descifrar es decodificar un código de Goppa
(fácil). Sin ella, es decodificar un código lineal cualquiera
(conjeturadamente difícil).
\end{idea}

\section{El criptosistema, paso a paso}

\begin{definicion}{El criptosistema de McEliece (1978)}{1978-esquema}
Parámetros: enteros $m$ y $t$; longitud $n=2^m$; dimensión $k=n-mt$.
Tres algoritmos:
\begin{enumerate}[label=\textsc{\Alph*.},leftmargin=1.9em,itemsep=3pt]
  \item \textsf{Generación de claves.} Elíjase $g\in\FF{m}[x]$
        irreducible de grado $t$ y tómese como soporte la lista $L$ de
        todos los elementos de $\FF{m}$. Sea $\Gamma(L,g)$ el código de
        Goppa correspondiente y sea $G$ una matriz generadora
        \emph{canónica} ($k\times n$; por ejemplo, en forma
        sistemática). Sortéense una matriz invertible $S\in
        \mathrm{GL}(k,\Ftwo)$ y una permutación $P$ de las $n$
        coordenadas. La \textbf{clave pública} es
        \begin{equation}
          G' = S\,G\,P \qquad\text{junto con el número }t ;
          \label{eq:1978-gp}
        \end{equation}
        la \textbf{clave secreta} es $(g,L,S,P)$.
  \item \textsf{Cifrado.} Para cifrar $\vu\in\Ftwo^k$ se sortea un
        error $\vz\in\Ftwo^n$ de peso \emph{exactamente} $t$ y se envía
        \begin{equation}
          \vx = \vu\,G' + \vz .
          \label{eq:1978-enc}
        \end{equation}
  \item \textsf{Descifrado.} Dado $\vx$: (1)~deshágase la permutación,
        $\vx P^{-1}$; (2)~decodifíquese con Patterson, lo que produce
        el codeword $(\vu S)\,G$; (3)~léase $\vu S$ en las coordenadas
        sistemáticas de $G$ y devuélvase $\vu=(\vu S)\,S^{-1}$.
\end{enumerate}
\end{definicion}

\begin{proposicion}{El descifrado recupera el mensaje}{1978-correcto}
Con la notación de la definición~\ref{def:1978-esquema}, si
$\vx=\vu G'+\vz$ y $\wt(\vz)\le t$, el descifrado devuelve $\vu$.
\end{proposicion}

\begin{proof}
Como $G'=SGP$ y $PP^{-1}=\mathrm{I}_n$,
\begin{equation}
  \vx P^{-1}
  = \vu S\,G\,(\underbrace{P\,P^{-1}}_{\mathrm{I}_n}) + \vz P^{-1}
  = \underbrace{(\vu S)\,G}_{\text{codeword de }\Gamma_2} + \vz P^{-1},
  \label{eq:1978-xpinv}
\end{equation}
y $\wt(\vz P^{-1})=\wt(\vz)\le t$, porque permutar coordenadas no
cambia el peso. El corolario~\ref{cor:corrige-t} garantiza que ese
codeword es el único a distancia $\le t$, y Patterson lo encuentra. La
matriz $G$ es canónica, así que sus coordenadas sistemáticas son
$\vu S$; multiplicando por $S^{-1}$ se recupera $\vu$.
\end{proof}

\noindent Quien cifra solo necesita $(G',t)$; quien descifra necesita la
clave secreta. La figura~\ref{fig:1978-esquema} resume el circuito
completo y la tabla~\ref{tab:1978-claves} separa lo público de lo
secreto, con sus tamaños para los parámetros originales.

\begin{figure}[htbp]
\centering
\begin{tikzpicture}[
    pbox/.style={draw=mcblue,fill=mcsky!45,rounded corners=2.5pt,
      align=center,font=\footnotesize,inner sep=3.5pt,
      text width=2.7cm,minimum height=1.35cm},
    sbox/.style={draw=mcorange,fill=mcorange2!60,rounded corners=2.5pt,
      align=center,font=\footnotesize,inner sep=3.5pt,
      text width=2.7cm,minimum height=1.35cm},
    arr/.style={-{Stealth[length=2.2mm]},thick,mcgray}
  ]
  % ---- fila 1: generación de claves ----
  \node[sbox] (gL) {$g(x)$ y $L$ (secreto)\\ $\deg g=t$};
  \node[sbox,right=5mm of gL] (Gm) {$G$ canónica\\ $k\times n$ (secreto)};
  \node[sbox,right=5mm of Gm] (SP) {$S$, $P$ (secreto)\\ $S$ invertible};
  \node[pbox,right=5mm of SP] (Gp) {$G'=SGP$ (público)\\ con el parámetro $t$};
  \draw[arr] (gL)--(Gm); \draw[arr] (Gm)--(SP); \draw[arr] (SP)--(Gp);
  % ---- fila 2: cifrado ----
  \node[pbox,below=10mm of gL] (u) {mensaje $\vu$\\ $k$ bits};
  \node[pbox,right=5mm of u] (x) {criptograma\\ $\vx=\vu G'+\vz$, $\wt(\vz)=t$};
  \draw[arr] (u)--(x);
  \draw[arr,dashed,mcblue] (Gp.south) -- node[font=\tiny,color=mcblue,pos=0.55,above,sloped]{se usa} (x.north);
  % ---- fila 3: descifrado ----
  \node[sbox,below=10mm of x] (xp) {$\vx P^{-1}$, y Patterson\\ (secreto)};
  \node[sbox,right=5mm of xp] (us) {codeword $(\vu S)G$;\\ se lee $\vu S$};
  \node[pbox,right=5mm of us,text width=1.7cm] (u2) {$\vu$};
  \draw[arr] (x)--(xp); \draw[arr] (xp)--(us); \draw[arr] (us)--(u2);
\end{tikzpicture}
\caption{El circuito del esquema de 1978. Arriba, la generación de
claves: la clave pública es la matriz de Goppa disfrazada $G'=SGP$.
Abajo, el uso: el emisor suma un error de peso $t$ a un codeword y el
receptor deshace el disfraz y decodifica. Los recuadros naranjas
contienen material secreto.}
\label{fig:1978-esquema}
\end{figure}

\begin{table}[htbp]
\centering
\footnotesize
\setlength{\tabcolsep}{6pt}
\begin{tabular}{lll}
\toprule
 & \textbf{Contenido} & \textbf{Tamaño (parámetros originales: $n=1024$, $k=524$, $t=50$)}\\
\midrule
Público & $G'=SGP$, de $k\times n$ & $524\cdot1024=\num{536576}$ bits $=64$~KiB\\
        & $t$ & despreciable\\
Secreto & $g$ (grado $t$) & $mt=\num{500}$ bits\\
        & $L$ & implícita ($L=\FF{10}$ completo)\\
        & $S$ & $k^2=\num{274576}$ bits ($\approx 34$~KiB)\\
        & $P$ & $n\log_2 n=\num{10240}$ bits ($\approx 1{,}3$~KiB)\\
\bottomrule
\end{tabular}
\caption{Reparto de material y tamaños en el McEliece de 1978. La
clave pública domina; reducirla será el primer problema a resolver
(capítulos~\ref{ch:niederreiter} y~\ref{ch:keygen}).}
\label{tab:1978-claves}
\end{table}

\section{El juguete del capítulo~\ref{ch:labgoppa} es este esquema}

El capítulo~\ref{ch:labgoppa} fue, literalmente, una ejecución de la
definición~\ref{def:1978-esquema} con $m=4$, $t=2$: allí $n=16$, $k=8$
y $g(x)=x^2+x+z^3$. La matriz $S$ de $8\times8$ y la permutación están
en la figura~\ref{fig:lab-SPG}; la clave pública $G'=SGP$ también; el
criptograma $\vx=\hx{0001011010000011}$ en la
figura~\ref{fig:lab-cifrado}; y el descifrado con Patterson, la lectura
de $\vu S$ y el producto por $S^{-1}$, con nombres y números, en todo
el capítulo. La única diferencia es la escala —y en criptografía la
escala lo es todo: el capítulo~\ref{ch:isd} mostrará por qué $n=16$ se
rompe en microsegundos y $n=3488$ resiste a todos los algoritmos
conocidos.

\section{El argumento de seguridad de 1978}

Un adversario que conoce $G'$ (y $t$) y ve $\vx$ tiene, como el propio
artículo señala, dos caminos (figura~\ref{fig:1978-ataques}):

\begin{enumerate}[leftmargin=1.6em,itemsep=3pt]
  \item \textbf{Ataques estructurales:} recuperar $(g,L,S,P)$, o al
        menos \emph{una} descripción equivalente del código, para
        decodificar a gusto. El artículo lo considera «desesperado»
        por el número astronómico de descomposiciones
        $S^{-1}G'P^{-1}$ posibles; la parte~\ref{part:seguridad}
        mostrará que la intuición era correcta en el recuento pero
        incompleta: la especialidad de Goppa deja huellas.
  \item \textbf{Ataques decodificadores:} resolver directamente
        \begin{equation}
          \text{dados } G' \text{ y } \vx,\quad
          \text{hallar } \vu \text{ con }
          \wt(\vx-\vu G')\le t,
          \label{eq:1978-mld}
        \end{equation}
        sin aprender nada de la estructura. Es el problema de la
        decodificación de máxima verosimilitud (MLD) restringido a la
        clase de códigos que $G'$ genera.
\end{enumerate}

El segundo camino se apoya en un teorema: el MLD para códigos lineales
binarios arbitrarios es NP-completo~\cite{bcmv1978}. Pero un teorema
así no cierra la cuestión.

\begin{rigor}[Qué prueba y qué no la NP-dureza]
La NP-dureza es un enunciado del \emph{peor caso}: «existen instancias
difíciles». La criptografía necesita más: que las instancias que el
esquema \emph{genera} sean difíciles. Y la instancia de McEliece no es
un código lineal arbitrario, sino un código de Goppa disfrazado, una
subclase muy especial sobre la que el teorema no dice nada. La
seguridad de McEliece es, desde 1978, una \textbf{conjetura}
cuantificada una y otra vez por la comunidad: el coste del mejor
algoritmo conocido (capítulos~\ref{ch:isd} y~\ref{ch:2026}) acota la
confianza, pero no la demuestra.
\end{rigor}

\begin{figure}[htbp]
\centering
\begin{tikzpicture}[
    cbox/.style={draw=mcgray,fill=mclight,rounded corners=2.5pt,align=center,
      font=\footnotesize,inner sep=3.5pt,text width=3.6cm,minimum height=1.0cm},
    abox/.style={draw=mcred,fill=mcred2!55,rounded corners=2.5pt,align=center,
      font=\footnotesize,inner sep=3.5pt,text width=4.4cm,minimum height=1.7cm},
    arr/.style={-{Stealth[length=2.4mm]},thick,mcgray}
  ]
  \node[cbox] (known) at (0,0) {el adversario conoce $G'$ y $t$, y ve $\vx$};
  \node[abox] (A) at (-3.4,-2.35)
    {\textbf{Ataque estructural}\\ recuperar $(g,L,S,P)$\\[1pt]
     {\scriptsize\itshape «¿cuántas descomposiciones $S^{-1}G'P^{-1}$ hay?»}\\[1pt]
     {\scriptsize $\to$ capítulo~\ref{ch:estructura}}};
  \node[abox] (B) at (3.4,-2.35)
    {\textbf{Ataque decodificador}\\ resolver~\eqref{eq:1978-mld}\\[1pt]
     {\scriptsize\itshape «MLD es NP-completo»}\\[1pt]
     {\scriptsize $\to$ capítulo~\ref{ch:isd}}};
  \draw[arr] (known.south) -- ++(0,-0.55) -| (A.north);
  \draw[arr] (known.south) -- ++(0,-0.55) -| (B.north);
\end{tikzpicture}
\caption{Las dos familias de ataques que ya identifica el artículo de
1978. La seguridad práctica moderna se discute toda dentro de este
dibujo, con dos precisiones que veremos en la
parte~\ref{part:seguridad}: los ataques estructurales también
\emph{distinguen} códigos de Goppa escondidos de códigos aleatorios, y
los decodificadores también pueden llegar a recuperar la clave.}
\label{fig:1978-ataques}
\end{figure}

\section{Los números del artículo, recalculados}

Con $(m,t)=(10,50)$ se tiene $n=1024$, $k=524$ y
$d\ge2t+1=101$. McEliece sopesó los ataques evidentes; la
tabla~\ref{tab:1978-numeros} compara sus cifras con el recálculo de hoy.

\begin{table}[htbp]
\centering
\small
\begin{tabular}{lcc}
\toprule
Magnitud & Artículo (1978) & Recálculo exacto\\
\midrule
Codewords posibles ($2^k$) & $2^{524}\approx10^{158}$ & igual\\
Clases laterales ($2^{n-k}$) & --- & $2^{500}\approx10^{150.5}$\\
Ataque de las $k$ coordenadas & $10^{19}\ (\approx2^{63})$ & $2^{80.7}\approx2\cdot10^{24}$\\
$\Pr[\text{decodificar un vector al azar}]$ & $2^{-215.59}$ & $2^{-215.89}$\\
\bottomrule
\end{tabular}
\caption{Estimaciones del artículo de 1978 para $n=1024$, $t=50$,
$k=524$, y recálculo con aritmética exacta (caja siguiente). Coinciden
en lo esencial y las tres favorecen al defensor.}
\label{tab:1978-numeros}
\end{table}

\begin{rigor}[Recalculando un artículo de 1978]
Dos de las cifras del artículo se pueden recalcular hoy sin margen de
interpretación.
\begin{enumerate}[label=(\arabic*),leftmargin=1.8em,itemsep=3pt]
  \item \textbf{Ataque de las $k$ coordenadas.} La probabilidad de
        que $k$ coordenadas elegidas al azar eviten las $t$ posiciones
        del error (sin reemplazo) es
        \[
          \frac{\binom{n-t}{k}}{\binom{n}{k}}
          \;=\;\prod_{i=0}^{k-1}\frac{974-i}{1024-i}
          \;=\;2^{-53.61},
        \]
        de modo que el trabajo esperado es
        $2^{53.61}\cdot k^{3}\approx 2^{80.7}\approx 2\cdot10^{24}$.
        El artículo usó la aproximación \emph{con} reemplazo
        $(1-t/n)^k=2^{-37.8}$, que da $1{,}9\cdot10^{19}$. La
        corrección \emph{aumenta} el coste del ataque: la estimación
        de 1978 era conservadora a favor del atacante.
  \item \textbf{La (no) firma.} El número de vectores de longitud $n$
        que están a distancia $\le t$ de algún codeword es
        $2^{k}\cdot\sum_{i=0}^{t}\binom{n}{i}$, y
        \[
          \sum_{i=0}^{50}\binom{1024}{i}=2^{284.11},
          \qquad
          \Pr=\frac{2^{524}\cdot2^{284.11}}{2^{1024}}=2^{-215.89}.
        \]
        El artículo escribe $2^{284.41}$ y $2^{-215.59}$: una
        diferencia de $0{,}3$ bits sin consecuencias (firmar exigiría
        acertar un suceso de probabilidad $2^{-215}$).
\end{enumerate}
En ambos casos la discrepancia es de aritmética de la época, no de
concepto, y las dos conclusiones —atacar es prohibitivamente caro,
«firmar» es imposible— quedan intactas.
\end{rigor}

La velocidad seguía siendo el argumento de venta: «cercanas a $10^6$
bits/s o más». Sigue siéndolo: cifrar es una multiplicación de bits
$k\times n$ más un xor, y descifrar añade Patterson. Cifrar es
prácticamente gratis; descifrar cuesta del orden de un milisegundo en
software actual, con aritmética de cuerpos finitos en vez de aritmética
modular sobre enteros enormes.

\section{Dónde falla (y dónde no)}

Tres grietas prácticas —todas arreglables o evitables— y una conjetura
que, por ahora, aguanta.

\smallskip
\noindent\textbf{1. No hay firmas.} El propio autor lo advierte: el
descifrado solo está definido para vectores a distancia $\le t$ de un
codeword, y esa fracción es $2^{-215.89}$ (caja anterior). Un esquema
de firma necesita que «descifrar» esté definido en todo el espacio;
«firmar» cifrando con la clave secreta produce basura casi siempre.
Existen esquemas de firma basados en códigos desde 2001, pero son otra
familia: Classic McEliece, como veremos, es \emph{solo} un KEM.

\smallskip
\noindent\textbf{2. Cifrado sin garantías modernas.} El cifrado de
1978 no aspira a IND-CPA ni a IND-CCA: es la capa que reduce la
recuperación del mensaje a un problema de decodificación. Dos
observaciones elementales lo muestran:
\begin{itemize}[leftmargin=1.4em,itemsep=2pt]
  \item Dado $\vx=\vu G'+\vz$, el error $\vz$ está
        \emph{unívocamente} determinado (ejercicio~3): el cifrado no
        esconde «ruido», esconde un problema MLD.
  \item El cifrado es \emph{maleable}: si se modifica $\vx$ en
        posiciones del error, el descifrado sigue devolviendo el
        mismo $\vu$: $\vx\oplus\ve'$ con
        $\supp(\ve')\subseteq\supp(\vz)$ decodifica a $\vu$.
\end{itemize}
Un adversario con acceso al oráculo de descifrado explota lo segundo
como \emph{ataque de reacción}: pregunta por $\vx\oplus\ve_j$ para
cada $j$ y aprende, por las reacciones, dónde está el error
(ejercicio~6). El remedio moderno —rechazo implícito— llega con
Classic McEliece (capítulos~\ref{ch:kem} y~\ref{ch:encapdecap}).

\begin{trampa}[«El descifrado siempre devuelve el mensaje»]
Falso: Patterson solo garantiza la decodificación si $\wt(\vz)\le t$.
Si el canal introduce tres errores en el juguete ($t=2$), o si un
adversario retoca un bit de más, el descifrador produce basura o falla;
y en el esquema de 1978 ese \emph{fallo es observable}. En un KEM
moderno el fallo se oculta (rechazo implícito,
capítulo~\ref{ch:encapdecap}).
\end{trampa}

\smallskip
\noindent\textbf{3. El tamaño de la clave pública.} $G'$ tiene
$k\times n$ bits: en 1978, $524\cdot1024$ bits $=64$~KiB; con
parámetros actuales, $2720\cdot3488$ bits $\approx1{,}13$~MB. Dos
remedios históricos: publicar la matriz de paridad del código dual
(Niederreiter, capítulo~\ref{ch:niederreiter}) —más pequeña cuando
$k>n/2$— o publicar solo la parte sistemática
$(\mathrm{I}_{mt}\,|\,\Tmat)$ y regenerar $\Tmat$ de una semilla corta,
que es lo que hace Classic McEliece (capítulos~\ref{ch:keygen}
y~\ref{ch:parametros}).

\smallskip
\noindent\textbf{4. Lo que no ha fallado.} En casi medio siglo no se ha
roto la conjetura central: no se conoce ningún ataque que recupere el
mensaje o la clave de los parámetros recomendados en tiempo factible.
Pero todo lo demás ha cambiado: los parámetros de 1978 están hoy por
debajo de cualquier margen razonable; la seguridad se mide con
algoritmos ISD cada vez mejores (capítulo~\ref{ch:isd}); hay ataques
estructurales que explotan la especialidad de las variantes mal
elegidas (capítulo~\ref{ch:estructura}); y 2025--2026 trajo resultados
nuevos que obligan a revisar supuestos (capítulo~\ref{ch:2026}). La
próxima parada es la reformulación dual de Niederreiter,
capítulo~\ref{ch:niederreiter}.

\section{Ejercicios}

\begin{enumerate}[leftmargin=1.6em,itemsep=4pt]
  \item \textbf{Corrección.} Repase la
        proposición~\ref{prop:1978-correcto} comprobando cada igualdad.
        ¿Dónde se usa $\wt(\vz)\le t$ (y no la igualdad)? ¿Qué fallaría
        si el emisor eligiera errores de peso $t+1$ «a veces»?
  \item \textbf{Tamaños.} Calcule los bits de $G'$ para
        $(n,k)=(1024,524)$ y para \texttt{mceliece348864}
        ($k=2720$, $n=3488$, clave pública de \num{261120} bytes).
        ¿Cuántos bytes ocupa la parte sistemática
        $(\mathrm{I}_{2720}\,|\,\Tmat)$? ¿Qué factor de ahorro supone
        frente a $G'$ completa?
  \item \textbf{El error es único.} Pruebe que, con
        $\vx=\vu G'+\vz$ y $\wt(\vz)\le t$, no existe otro par
        $(\vu',\vz')$ con $\vu'\ne\vu$, $\wt(\vz')\le t$ y
        $\vx=\vu'G'+\vz'$. (Use $d\ge2t+1$.)
  \item \textbf{Los números del artículo.} Con una calculadora o
        Python, reproduzca los tres valores no triviales de la
        tabla~\ref{tab:1978-numeros}: $2^{500}$, el cociente
        $\binom{974}{524}/\binom{1024}{524}$ y
        $\sum_{i\le50}\binom{1024}{i}$. Pista: trabaje con logaritmos
        en base $2$.
  \item \textbf{Por qué no basta la NP-dureza.} Explique en cinco
        líneas por qué «MLD es NP-completo» no demuestra que McEliece
        sea seguro.
  \item \textbf{Ataque de reacción.} Suponga que el descifrado «avisa»
        cuando su entrada no es decodificable. Demuestre que enviando
        $\vx\oplus\ve_j$ para cada posición $j$ se aprenden todas las
        posiciones del error y, con ellas, $\vu$: ¿cuántas consultas
        hacen falta? ¿Cómo se comporta el KEM real ante esa consulta
        (capítulo~\ref{ch:encapdecap})?
\end{enumerate}

\chapter{La reformulación de Niederreiter}\label{ch:niederreiter}

\begin{roadmap}
Este capítulo da la vuelta al esquema: en lugar de publicar una matriz
generadora y enviar un codeword con ruido, publicaremos una
\emph{matriz de paridad} y enviaremos el \emph{síndrome} de un vector de
peso fijo. El resultado (Niederreiter, 1986) tiene criptogramas más
cortos y un espacio de mensajes distinto, con la misma dificultad de
fondo. No es un capítulo de arqueología: el KEM real,
Classic McEliece, es niederreiteriano de principio a fin.
\end{roadmap}

\begin{historia}[1986: publicar la otra matriz]
Harald Niederreiter propuso en 1986~\cite{niederreiter1986} cambiar
\emph{qué} se publica y \emph{qué} se transmite: en vez de la matriz
generadora de 1978, la matriz de paridad; y en vez de un codeword con
error, el síndrome del error. Casi cuarenta años después, la
especificación de Classic McEliece hace exactamente esto: su clave
pública $T$ define la matriz de paridad $H=(\mathrm{I}_{mt}\,|\,T)$ y
sus criptogramas son síndromes de $mt$ bits~\cite{cmspec}.
\end{historia}

\section{La idea dual: mensajes de peso fijo}

Recordemos el cuadro de los capítulos~\ref{ch:lineales}
y~\ref{ch:goppa}: un código lineal $\Gamma$ de parámetros $[n,k,d]$ se
describe indiferentemente por una matriz generadora $G$ ($k\times n$)
o por una matriz de paridad $H_0$ ($(n-k)\times n$), con
\[
\vc\in\Gamma \;\Longleftrightarrow\;
\vc\,H_0^{\top}=\mathbf{0} .
\]
En el McEliece de 1978 se publica $G'=SGP$ disfrazada y se envía
$\vx=\vu G'+\vz$: $n$ bits, de los que $k$ son mensaje. Niederreiter
observó que se puede cifrar con el código «al revés»:

\begin{itemize}[leftmargin=1.4em,itemsep=3pt]
  \item El \textbf{mensaje} es un vector $\ve\in\Ftwo^n$ de peso
        \emph{exactamente} $t$ (no una palabra arbitraria). Hay
        $\binom{n}{t}$ mensajes posibles: para $(n,t)=(1024,50)$,
        $\binom{1024}{50}\approx2^{284}$, es decir $\approx284$ bits
        útiles por bloque. Generar un mensaje es «codificar en peso
        fijo» (\emph{FixedWeight} en la especificación; algoritmos en
        el capítulo~\ref{ch:keygen}).
  \item El \textbf{criptograma} es su síndrome:
        $\vc=\ve\,(H')^{\top}$, un vector de $n-k=mt$ bits,
        donde $H'$ es la matriz de paridad pública.
\end{itemize}

Con los parámetros de 1978, $n=1024$ y $t=50$: el criptograma baja de
$1024$ a $n-k=500$ bits, mientras que el mensaje baja de $k=524$ bits
libres a $\log_2\binom{1024}{50}\approx284$ bits. La ganancia es real
pero modesta; con los parámetros de hoy, en cambio, la balanza se
inclina: $k\approx0{,}78\,n$ y el criptograma cae por debajo de la
cuarta parte de $n$ (tabla~\ref{tab:nied-tamanos}). Esa es una de las
razones por las que la forma dual esperó a los KEM: su ventaja se
materializa cuando $n-k\ll n$.

\begin{figure}[htbp]
\centering
\begin{tikzpicture}[
    pbox/.style={draw=mcblue,fill=mcsky!45,rounded corners=2.5pt,align=center,
      font=\footnotesize,inner sep=3.5pt,text width=3.3cm,minimum height=1.05cm},
    mbox/.style={draw=mcorange,fill=mcorange2!55,rounded corners=2.5pt,align=center,
      font=\footnotesize,inner sep=3.5pt,text width=3.3cm,minimum height=1.05cm},
    arr/.style={-{Stealth[length=2.2mm]},thick,mcgray}
  ]
  % --- panel izquierdo: McEliece ---
  \node[font=\small\bfseries,color=mcblue] at (1.65,0.75) {McEliece (1978)};
  \node[mbox] (u) at (1.65,0) {mensaje $\vu$\\ $k$ bits libres};
  \node[pbox] (xc) at (1.65,-1.75) {criptograma\\ $\vx=\vu G'+\vz$, $\wt(\vz)=t$};
  \node[align=center,font=\scriptsize,color=mcgray] at (1.65,-3.0)
    {longitud $n$ (por ejemplo, $1024$ bits)};
  \draw[arr] (u)--(xc) node[midway,right,font=\tiny]{codifica y añade ruido};
  % --- panel derecho: Niederreiter ---
  \begin{scope}[shift={(7.0,0)}]
  \node[font=\small\bfseries,color=mcblue] at (1.65,0.75) {Niederreiter (1986)};
  \node[mbox] (e) at (1.65,0) {mensaje $\ve$\\ peso exacto $t$};
  \node[pbox] (c) at (1.65,-1.75) {criptograma\\ $\vc=\ve\,(H')^{\top}$, $n-k$ bits};
  \node[align=center,font=\scriptsize,color=mcgray] at (1.65,-3.0)
    {longitud $n-k=mt$ (por ejemplo, $768$ bits)};
  \draw[arr] (e)--(c) node[midway,right,font=\tiny]{toma el síndrome};
  \end{scope}
  % flecha dual
  \draw[arr,dashed,mcred] (4.35,-0.9) -- (6.1,-0.9)
    node[midway,above,font=\tiny,color=mcred]{visión dual};
\end{tikzpicture}
\caption{Las dos miradas al mismo código. En McEliece se publica
$G'$ y se envían $n$ bits (codeword con ruido); en Niederreiter se
publica $H'$ y se envían $n-k$ bits (síndrome del mensaje-error).
Classic McEliece usará la segunda columna.}
\label{fig:nied-dual}
\end{figure}

\section{El criptosistema (1986), en nuestras convenciones}

\begin{definicion}{El criptosistema de Niederreiter (1986)}{nied-esquema}
Parámetros $m$, $t$; $n=2^m$; $k=n-mt$. Sea $\Gamma(L,g)$ el código de
Goppa secreto y sea $H_0$ una matriz de paridad suya de rango completo.
Las operaciones de fila sobre $H_0$ no cambian el código que comprueba,
y una permutación de columnas se absorbe en $P$; así pues, sin pérdida
de generalidad suponemos $H_0$ reducida:
$H_0=(\mathrm{I}_{n-k}\,|\,\Tmat_0)$. (Es la misma forma
$(\mathrm{I}_8\,|\,\Tmat)$ que el juguete destiló en la
figura~\ref{fig:lab-hat}.) Tres algoritmos:
\begin{enumerate}[label=\textsc{\Alph*.},leftmargin=1.9em,itemsep=3pt]
  \item \textsf{Generación de claves.} Se sortean $S\in
        \mathrm{GL}(n-k,\Ftwo)$ y una permutación $P$ de las $n$
        coordenadas. La \textbf{clave pública} es
        \begin{equation}
          H' = S\,H_0\,P \qquad\text{(y el número } t),
          \label{eq:nied-pk}
        \end{equation}
        y la \textbf{clave secreta} es $(g,L,S,P)$.
  \item \textsf{Cifrado.} Para cifrar $\ve\in\Ftwo^n$ con
        $\wt(\ve)=t$, envíese
        \begin{equation}
          \vc = \ve\,(H')^{\top}\in\Ftwo^{\,n-k}
          \qquad
          \bigl(\text{en notación de columnas: }\vc^{\top}=H'\ve^{\top}\bigr).
          \label{eq:nied-enc}
        \end{equation}
  \item \textsf{Descifrado.} Dado $\vc$: (1)~calcúlese
        $\gamma=\vc\,(S^{-1})^{\top}$; (2)~encuéntrese el único
        $\vf\in\Ftwo^n$ con $\vf H_0^{\top}=\gamma$ y $\wt(\vf)\le t$
        (con Patterson: basta decodificar la palabra
        $(\gamma,\mathbf{0})$, cuyo síndrome es $\gamma$);
        (3)~devuélvase $\ve=\vf P$.
\end{enumerate}
\end{definicion}

La convención de permutación es la del capítulo~\ref{ch:labgoppa}:
$(\vv P)_j=\vv_{\sigma(j)}$, la coordenada antigua $\sigma(j)$ pasa a
la posición $j$; recordemos también que $P^{\top}=P^{-1}$.

\begin{proposicion}{El descifrado recupera el mensaje-error}{nied-correcto}
Con la notación de la definición~\ref{def:nied-esquema}, si
$\vc=\ve(H')^{\top}$ y $\wt(\ve)=t$, el descifrado devuelve $\ve$.
\end{proposicion}

\begin{proof}
Como $H'=SH_0P$,
\[
  \vc\,(S^{-1})^{\top}
  = \ve\,(SH_0P)^{\top}(S^{-1})^{\top}
  = \ve\,P^{\top}H_0^{\top}S^{\top}(S^{-1})^{\top}
  = \underbrace{(\ve P^{-1})}_{\text{peso }t}H_0^{\top}
  =: \gamma .
\]
Por tanto $\vf=\ve P^{-1}$ es una solución de peso $t$ del paso (2)
(de hecho, $(\ve P^{-1})H_0^{\top}=\gamma$ es la identidad anterior), y
es la única: si $\vf'$ fuera otra con $\wt(\vf')\le t$, entonces
$\vf+\vf'\in\ker H_0^{\top}=\Gamma$, con
$\wt(\vf+\vf')\le2t<d$, imposible. Patterson encuentra ese único $\vf$
(el capítulo~\ref{ch:labgoppa} mostró el mecanismo con todas sus
cuentas). Y $\vf P=\ve P^{-1}P=\ve$.
\end{proof}

\begin{observacion}{El mensaje es el error}{nied-mensaje-error}
En McEliece, el mensaje $\vu$ viaja «escalado» dentro de $\vu S$ y hay
que deshacer $S$; el ruido $\vz$ es solo eso, ruido. En Niederreiter no
hay mensaje separado del error: el vector de peso $t$ \emph{es} el
mensaje, y el descifrado de un síndrome es recuperarlo. La frase «el
descifrado de Niederreiter es un problema de decodificación por
síndromes» es literal.
\end{observacion}

\section{Equivalencia con el esquema de 1978}

Los dos esquemas resuelven el mismo problema con distinto embalaje.
Formalmente:

\begin{proposicion}{Un ataque niederreiteriano rompe McEliece}{nied-reduccion}
Sean $G'=S_1GP$ y $H'=S_2H_0P$ dos disfraces del mismo código secreto,
con la \emph{misma} permutación $P$ en ambos ($S_1$ y $S_2$ pueden ser
distintas). Un algoritmo que, dados $H'$ y $\vc=\ve(H')^{\top}$ con
$\wt(\ve)=t$, recupera $\ve$, recupera también el mensaje de todo
criptograma de McEliece $\vx=\vu G'+\vz$.
\end{proposicion}

\begin{proof}
Sea $\vx=\vu G'+\vz$ un criptograma de McEliece. Lo primero es que
$G'$ y $H'$ son ortogonales aunque los disfraces $S_1$, $S_2$ difieran:
\[
  G'\,(H')^{\top}
  = S_1\,G\,\underbrace{P\,P^{\top}}_{\mathrm{I}_n}\,H_0^{\top}\,S_2^{\top}
  = S_1\,\underbrace{(G\,H_0^{\top})}_{=\,0}\,S_2^{\top}
  = 0 ,
\]
porque $G$ genera el mismo código que $H_0$ comprueba. Por tanto
\[
  \vx\,(H')^{\top}
  = \vu\,G'(H')^{\top}+\vz\,(H')^{\top}
  = \vz\,(H')^{\top},
\]
es decir: el síndrome público de $\vx$ es el síndrome de $\vz$, y se
calcula sin ninguna clave. Se lo damos al recuperador niederreiteriano:
devuelve el único $\ve$ de peso $t$ con
$\ve(H')^{\top}=\vz(H')^{\top}$. Pero entonces
$(\vz+\ve)(H')^{\top}=0$, así que $\vz+\ve$ es un codeword de un código
equivalente a $\Gamma$ (sus coordenadas solo están permutadas, y eso no
cambia el peso) y tiene peso $\le2t<d$: la única posibilidad es
$\ve=\vz$, por el mismo argumento de unicidad de la
proposición~\ref{prop:nied-correcto}. Recuperado el error,
$\vx+\vz=\vu G'$, y $\vu$ se despeja por álgebra lineal (o se lee en
las coordenadas sistemáticas de $G'$).
\end{proof}

La dirección demostrada es la que se usa: los ataques modernos contra
Classic McEliece se describen casi siempre como decodificación por
síndromes contra $H'$ (parte~\ref{part:seguridad}). El recíproco —que un
atacante capaz de descifrar McEliece rompa también Niederreiter— fue
establecido por Li, Deng y Wang en 1994~\cite{lidw1994}, que probaron
la equivalencia polinómica de ambos esquemas con parámetros
correspondientes. La guía de los diseñadores lo resume así: la variante
de Niederreiter «cambia la seguridad de forma
insignificante»~\cite{cmrationale}. Equivalencia, con matices: los
\emph{problemas de inversión} son esencialmente el mismo; los esquemas
completos (tamaños de clave, longitud de criptograma y las
transformaciones a KEM del capítulo~\ref{ch:kem}) no son idénticos, y
ahí está toda la diferencia práctica.

\section{Los números: cuándo gana la forma dual}

En 1978 se publicaba una matriz $k\times n$; en la forma dual se publica
una matriz $(n-k)\times n$. El factor de ahorro es $k/(n-k)$: con los
parámetros originales, $k\approx n/2$ y el ahorro resultaba mísero
($5\%$); con los parámetros reales, $k\approx0{,}78\,n$ y el ahorro es
de $3{,}5$ veces. La tabla~\ref{tab:nied-tamanos} muestra los números
del estándar, que ya usan la compresión sistemática que veremos en el
capítulo~\ref{ch:keygen}.

\begin{table}[htbp]
\centering
\small
\setlength{\tabcolsep}{5pt}
\begin{tabular}{lccccccc}
\toprule
Conjunto & $m$ & $n$ & $t$ & $k$ & $mt$ & clave pública & criptograma\\
\midrule
juguete (cap.~\ref{ch:labgoppa}) & 4 & 16 & 2 & 8 & 8 & 8~B & 1~B\\
\texttt{mceliece348864} & 12 & 3488 & 64 & 2720 & 768 & \num{261120}~B & 96~B\\
\texttt{mceliece460896} & 13 & 4608 & 96 & 3360 & 1248 & \num{524160}~B & 156~B\\
\texttt{mceliece6688128} & 13 & 6688 & 128 & 5024 & 1664 & \num{1044992}~B & 208~B\\
\texttt{mceliece6960119} & 13 & 6960 & 119 & 5413 & 1547 & \num{1047319}~B & 194~B\\
\texttt{mceliece8192128} & 13 & 8192 & 128 & 6528 & 1664 & \num{1357824}~B & 208~B\\
\bottomrule
\end{tabular}
\caption{Los cinco tamaños de Classic McEliece (sin las variantes
\texttt{f} y \texttt{pc}; detalles en el capítulo~\ref{ch:parametros})
frente al juguete. La clave pública guarda la matriz $T$, de
$mt\times k$ bits, en $mt\lceil k/8\rceil$ bytes; el criptograma es un
síndrome de $\lceil mt/8\rceil$ bytes (en \texttt{mceliece6960119},
$mt=1547$ no es múltiplo de $8$ y ambas cosas se rellenan). A modo de
comparación: una matriz generadora completa de \texttt{mceliece348864}
ocuparía $k\,n/8=\num{1185920}$ bytes, más de cuatro veces lo que ocupa
$T$.}
\label{tab:nied-tamanos}
\end{table}

\section{Por qué Classic McEliece es niederreiteriano}

En la especificación, la función \textsf{Encode} recibe el vector $\ve$
y la clave pública $T$ —una matriz de $mt\times k$ bits— y devuelve
\[
  C=(\mathrm{I}_{mt}\,|\,T)\,\ve^{\top},
\]
un síndrome de $n-k=mt$ bits; la función \textsf{Decode} reconstruye la
palabra $v=(C,\mathbf{0})$ (los $mt$ bits del síndrome y $k$ ceros) y
busca el codeword único a distancia $\le t$, del que recupera
$\ve=v+c$~\cite{cmspec}. Son los pasos (2) y (3) de la
definición~\ref{def:nied-esquema}, con dos simplificaciones de
ingeniería que estudiaremos en los capítulos~\ref{ch:keygen}
y~\ref{ch:encapdecap}: (i)~la clave pública no lleva una $S$ explícita,
porque $T$ ya está en forma canónica y las operaciones de fila quedan
absorbidas en ella; (ii)~la permutación $P$ no se sortea como matriz
aparte: el soporte de Goppa se elige en un orden aleatorio (la
\emph{FieldOrdering} de la especificación) y ese orden vive en la clave
secreta. La matemática es la del capítulo; el formato, el del KEM:
clave pública más de cuatro veces más corta y criptogramas de $mt$ bits.

\noindent El capítulo~\ref{ch:kem} convierte esta forma dual en un KEM;
los capítulos~\ref{ch:keygen} y~\ref{ch:encapdecap} abren la
especificación línea a línea.

\section{Ejercicios}

\begin{enumerate}[leftmargin=1.6em,itemsep=4pt]
  \item \textbf{El truco del relleno de ceros.} Verifique que, con
        $H_0=(\mathrm{I}_{n-k}\,|\,\Tmat_0)$, la palabra
        $v=(\gamma,\mathbf{0})$ tiene síndrome
        $vH_0^{\top}=\gamma$; explique por qué eso convierte el paso
        (2) del descifrado en exactamente el problema de decodificación
        del capítulo~\ref{ch:goppa}, con Patterson aplicable tal cual.
  \item \textbf{Unicidad del síndrome.} Pruebe que si
        $\vf H_0^{\top}=\vf' H_0^{\top}$ con
        $\wt(\vf),\wt(\vf')\le t$, entonces $\vf=\vf'$. (Pista: la
        diferencia está en $\Gamma$ y su peso es $\le2t<d$.)
  \item \textbf{La reducción, con lápiz.} Reproduzca la demostración de
        la proposición~\ref{prop:nied-reduccion} e indique en qué paso
        se usa $PP^{\top}=\mathrm{I}_n$ y en cuál la ortogonalidad
        $GH_0^{\top}=0$. ¿Por qué el argumento falla si $G'$ y $H'$
        usaran permutaciones distintas?
  \item \textbf{Factores de ahorro.} Calcule el factor $k/(n-k)$ de los
        cinco conjuntos de la tabla~\ref{tab:nied-tamanos} y compárelo
        con el del juguete (donde vale $1$). ¿Por qué el esquema
        original de $(1024,524)$ apenas ganaba?
  \item \textbf{Peso exacto.} El descifrado también funcionaría con
        $\wt(\ve)<t$; dé dos razones por las que la especificación fija
        el peso en $t$ exacto. (Pista: tamaño del espacio de mensajes y
        uniformidad: las transformaciones a KEM del
        capítulo~\ref{ch:kem} exigen que $\ve$ sea uniforme en el
        conjunto de vectores de peso $t$.)
  \item \textbf{Mensaje y error.} Escriba los dos problemas de inversión
        uno al lado del otro,
        \[
          \text{(McEliece)}\quad \vx=\vu G'+\vz\mapsto \vu ;
          \qquad
          \text{(Niederreiter)}\quad \vc=\ve(H')^{\top}\mapsto \ve ,
        \]
        y explique en qué sentido son «el mismo problema» y en qué se
        parecen a la decodificación por síndromes del
        capítulo~\ref{ch:goppa}.
\end{enumerate}

\chapter{Del cifrado a un KEM}\label{ch:kem}

\begin{roadmap}
Los esquemas de los capítulos~\ref{ch:1978} y~\ref{ch:niederreiter}
cierran la parte~\ref{part:mceliece}: sabemos cifrar, y sabemos que
romper el cifrado es decodificar. Pero un protocolo real no necesita
«cifrar mensajes»: necesita \emph{acordar claves de sesión} que un
adversario activo no pueda distinguir. Este capítulo presenta ese
objetivo (un KEM IND-CCA2), muestra por qué el esquema desnudo no lo
alcanza y describe la transformación ---tipo Fujisaki--Okamoto, con
\emph{rechazo implícito}--- que adoptó Classic McEliece. Es el puente
conceptual hacia la especificación de la parte~\ref{part:classic}.
\end{roadmap}

\begin{historia}[De «descifrar» a «no distinguir»]
Los objetivos de seguridad evolucionaron junto con los esquemas:
recuperar el mensaje era el listón de 1978; a mediados de los noventa
llegaron la indistinguibilidad y, con ella, el cifrado elegido
(IND-CCA2). En 1999 Fujisaki y Okamoto publicaron su transformación
clásica~\cite{foj1999}; en 2001 Shoup propuso empaquetar las claves de
sesión en KEMs~\cite{shoup2001}; el rechazo implícito apareció en 2012,
en la tesis de Persichetti, para esquemas basados en
códigos~\cite{persichetti2012}, y Hofheinz, Hövelmanns y Kiltz le
dieron una prueba modular en 2017~\cite{hhk2017}. Classic McEliece usa
la versión de 2018 de Bernstein y Persichetti~\cite{bp2018}, la adapta
y añade mejoras propias.
\end{historia}

\section{De filas a columnas (y por qué)}

Hasta aquí escribimos los vectores como \emph{filas}; la especificación
de Classic McEliece los escribe como \emph{columnas}. Es un cambio
puramente tipográfico ---transponer---, pero conviene fijarlo antes de
seguir, porque las fórmulas de los capítulos~\ref{ch:keygen}
a~\ref{ch:labkem} estarán en la segunda convención. El diccionario
completo es la tabla~\ref{tab:kem-columnas}; la matemática no cambia
en absoluto.

\begin{table}[htbp]
\centering
\small
\begin{tabular}{lcc}
\toprule
Objeto & en filas (caps.~1--7) & en columnas (desde aquí)\\
\midrule
vector & $v$ de tamaño $1\times n$ & $v$ de tamaño $n\times1$\\
código & $c\,H_0^{\top}=0$ & $H_0\,c=0$\\
cifrado de McEliece & $x=uG'+z$ & $x=(G')^{\top}u+z$\\
cifrado de Niederreiter & $C=e\,(H')^{\top}$ & $C=H'\,e$\\
decodificación & Patterson, igual & Patterson, igual\\
\bottomrule
\end{tabular}
\caption{Diccionario entre las dos convenciones. Transponer no cambia
ningún resultado: es la misma álgebra con las transposiciones
explícitas. La matriz de un código de Goppa (las $t$ filas de
$\alpha_i^{j}/g(\alpha_i)$, $j=0,\dots,t-1$) funciona igual en ambos
mundos.}
\label{tab:kem-columnas}
\end{table}

\section{Qué es un KEM y qué le pedimos}

Un \emph{mecanismo de encapsulación de claves} (KEM) no cifra mensajes:
produce un criptograma y, a la vez, una clave de sesión. Formalmente:

\begin{definicion}{KEM: sintaxis y corrección}{kem-def}
Un KEM es una terna de algoritmos:
\begin{itemize}[leftmargin=1.5em,itemsep=2pt]
  \item $\textsf{Gen}$: sin entrada más que los parámetros, produce una
        pareja $(\mathrm{pk},\mathrm{sk})$;
  \item $\textsf{Encap}(\mathrm{pk})$: con la clave pública, produce un
        criptograma $C$ y una clave de sesión $K$ (en Classic McEliece,
        $32$ bytes);
  \item $\textsf{Decap}(\mathrm{sk},C)$: con la clave secreta y el
        criptograma, reproduce la clave $K$ (o, en diseños con
        \emph{rechazo explícito}, el símbolo $\bot$).
\end{itemize}
El KEM es \emph{correcto} si para toda pareja $(\mathrm{pk},\mathrm{sk})$
y todo $(C,K)\leftarrow\textsf{Encap}(\mathrm{pk})$ se tiene
$\textsf{Decap}(\mathrm{sk},C)=K$.
\end{definicion}

\begin{definicion}{Seguridad IND-CCA2 de un KEM}{kem-indcca}
El juego entre un adversario $\mathcal{A}$ y un retador:
\begin{enumerate}[leftmargin=1.8em,itemsep=2pt]
  \item el retador calcula
        $(\mathrm{pk},\mathrm{sk})\leftarrow\textsf{Gen}$ y
        $(C^*,K_0^*)\leftarrow\textsf{Encap}(\mathrm{pk})$; sortea
        $b\leftarrow\{0,1\}$ y una clave uniforme $K_1^*$;
  \item $\mathcal{A}$ recibe $(\mathrm{pk},C^*,K_b^*)$ y puede
        consultar un oráculo de descifrado que, para cualquier
        $C\ne C^*$, responde $\textsf{Decap}(\mathrm{sk},C)$;
  \item $\mathcal{A}$ termina con un bit $b'$ y gana si $b'=b$.
\end{enumerate}
La \emph{ventaja} de $\mathcal{A}$ es
$\bigl|\Pr[b'=b]-\tfrac12\bigr|$. El KEM es IND-CCA2 si todo
adversario eficiente tiene ventaja despreciable.
\end{definicion}

Los objetivos de seguridad se ordenan: OW-CPA (recuperar el mensaje),
IND-CPA (distinguirlo sin oráculo) e IND-CCA2 (distinguirlo con oráculo).
En este libro el trabajo duro se concentra en OW-CPA: la
parte~\ref{part:seguridad} mide, muy en serio, la dificultad de
recuperar el mensaje; todo lo demás lo aporta la transformación de este
capítulo. ¿Y por qué KEMs y no cifrados de clave pública? Porque los
protocolos de transporte (TLS, SSH) necesitan exactamente esto: una
clave simétrica mutuamente acordada. El patrón estándar es
KEM--DEM~\cite{shoup2001}: el KEM encapsula la clave $K$ y un cifrado
simétrico autenticado (AES-GCM, por ejemplo) cifra los datos con ella.

\begin{trampa}[«Un KEM no cifra»]
La confusión inversa también es un error: \textsf{Encap} \emph{no}
recibe un mensaje, y \textsf{Decap} \emph{no} devuelve un mensaje. Si su
protocolo necesita mandar datos, no intente «pasar el mensaje por el
KEM»: encapsule una clave y cifre los datos con un DEM. Todos los KEMs
modernos (ML-KEM, HQC, Classic McEliece) funcionan así.
\end{trampa}

\section{Por qué el esquema desnudo no basta}

Lo que el esquema de Niederreiter \emph{sí} ofrece es OW-CPA: recuperar
$e$ de $C=H'e$ (o $u$ de $x=uG'+z$) es exactamente el problema de
decodificación que estudia la parte~\ref{part:seguridad}. Pero como
interfaz criptográfica falla ante un adversario activo:

\begin{itemize}[leftmargin=1.5em,itemsep=3pt]
  \item \textbf{Relaciones entre criptogramas.} En McEliece, dos
        cifrados del mismo mensaje difieren en $z_1+z_2$, un vector de
        peso $\le2t$: un atacante que recibe dos reenvíos ya no ve dos
        problemas independientes, sino uno con un vector de peso
        pequeño conocido de antemano (ataque de reenvío de Berson,
        1997~\cite{berson1997}). En Niederreiter, repetir la
        aleatoriedad $e$ repite el criptograma exacto.
  \item \textbf{Oráculos de descifrado.} En el capítulo~\ref{ch:1978}
        (ejercicio~6) vimos que, si un descifrador responde de forma
        distinguible al éxito y al fracaso, cada consulta sobre
        $C+\delta$ revela información sobre el error y el atacante
        acaba reconstruyendo $z$ entero, consulta a consulta.
  \item \textbf{Ciclos del generador aleatorio.} Si la aleatoriedad se
        repite entre usuarios o entre sesiones (ciclos de RNG, claves
        mal generadas), los criptogramas se relacionan y el esquema se
        rompe de forma trivial~\cite{cmsecurity}.
\end{itemize}

La conclusión es importante y sutil: \emph{el problema matemático no es
el culpable}; lo que hay que blindar es la interfaz. Eso es,
exactamente, lo que hace la transformación que viene a continuación.

\section{La transformación tipo FO con rechazo implícito}

La idea, depurada por veinte años de literatura, cabe en tres frases:
(1)~que la clave de sesión dependa de la semilla aleatoria $e$
\emph{y} del criptograma $C$; (2)~que el descifrado verifique, sin
ambigüedad, que el criptograma es legítimo; (3)~que el fallo no sea
distinguible: en vez de $\bot$, se devuelve una clave pseudoaleatoria
derivada de una clave secreta $s$.

\begin{definicion}{Transformación tipo FO con rechazo implícito}{kem-transform}
Sea $(\textsf{Gen}',\textsf{Enc},\textsf{Dec})$ una PKE correcta y
\emph{determinista} con espacio de mensajes $P$ ---aquí, los vectores
de peso $t$--- y sea $H$ una función de hash con separación de dominio.
El KEM se construye así:
\begin{enumerate}[leftmargin=1.8em,itemsep=2pt]
  \item $\textsf{Gen}$: $(\mathrm{pk}',\mathrm{sk}')\leftarrow
        \textsf{Gen}'$; se sortea la clave de rechazo $s$; la clave
        pública es $\mathrm{pk}=\mathrm{pk}'$ y la clave secreta es
        $(\mathrm{sk}',\mathrm{pk}',s)$. (La versión genérica guarda
        $\mathrm{pk}'$ para reencriptar; veremos que Classic McEliece
        puede prescindir de ella.)
  \item $\textsf{Encap}(\mathrm{pk})$: se sortea $e\leftarrow P$
        uniforme y se calcula
        \[
          C=\textsf{Enc}(\mathrm{pk},e),
          \qquad
          K=H(1,e,C).
          \label{eq:kem-encap}
        \]
  \item $\textsf{Decap}(\mathrm{sk},C)$: se descifra
        $e'\leftarrow\textsf{Dec}(\mathrm{sk}',C)$ y la clave es
        \[
          K=
          \begin{cases}
            H(1,e',C), & \text{si $\textsf{Enc}(\mathrm{pk},e')=C$},\\[2pt]
            H(0,s,C),  & \text{en caso contrario.}
          \end{cases}
          \label{eq:kem-decap}
        \]
\end{enumerate}
\end{definicion}

\begin{figure}[htbp]
\centering
\begin{tikzpicture}[
    pbox/.style={draw=mcblue,fill=mcsky!45,rounded corners=2.5pt,align=center,
      font=\footnotesize,inner sep=3.5pt,text width=3.6cm,minimum height=0.95cm},
    gbox/.style={draw=mcgreen,fill=mcgreen2!60,rounded corners=2.5pt,align=center,
      font=\footnotesize,inner sep=3.5pt,text width=3.6cm,minimum height=0.95cm},
    rbox/.style={draw=mcred,fill=mcred2!60,rounded corners=2.5pt,align=center,
      font=\footnotesize,inner sep=3.5pt,text width=3.5cm,minimum height=0.95cm},
    arr/.style={-{Stealth[length=2.2mm]},thick,mcgray}
  ]
  % --- panel izquierdo: Encap ---
  \node[font=\small\bfseries,color=mcblue] at (1.95,0.95) {\textsf{Encap}$(pk)$};
  \node[pbox] (e) at (1.95,0) {$e\leftarrow P$ uniforme\\ (peso $t$: la «semilla»)};
  \node[pbox] (c) at (1.95,-1.6) {$C=\textsf{Enc}(pk,e)$};
  \node[gbox] (k) at (1.95,-3.2) {$K=H(1,e,C)$};
  \draw[arr] (e)--(c);
  \draw[arr] (c)--(k);
  % --- panel derecho: Decap ---
  \begin{scope}[shift={(6.2,0)}]
  \node[font=\small\bfseries,color=mcblue] at (1.95,0.95) {\textsf{Decap}$(sk,C)$};
  \node[pbox] (c2) at (1.95,0) {$C$ (llega del canal)};
  \node[pbox] (d) at (1.95,-1.6) {$e'\leftarrow\textsf{Dec}(sk',C)$};
  \node[gbox] (k2) at (0.0,-3.2) {si es válido:\\ $K=H(1,e',C)$};
  \node[rbox] (kb) at (3.9,-3.2) {si no:\\ $K=H(0,s,C)$};
  \draw[arr] (c2)--(d);
  \draw[arr] (d.south) -- (k2.north);
  \draw[arr] (d.south) -- (kb.north);
  \end{scope}
  % --- el criptograma cruza el canal ---
  \draw[arr,dashed,mcred] (3.95,-1.6) -- (6.15,-1.6)
    node[midway,above,font=\tiny,color=mcred]{$C$};
\end{tikzpicture}
\caption{La transformación, en una imagen. \textsf{Encap} sortea la
semilla $e$, la cifra y deriva la clave de sesión del par $(e,C)$.
\textsf{Decap} descifra; si el criptograma es válido reconstruye la
misma clave (verde); si no, devuelve una clave pseudoaleatoria derivada
de $s$ (rojo), indistinguible de una clave legítima: el
\emph{rechazo implícito}. El adversario ve $(\mathrm{pk},C,K)$ y puede
consultar \textsf{Decap} con cualquier $C'\ne C$.}
\label{fig:kem-transform}
\end{figure}

La seguridad de esta construcción está demostrada en dos modelos:

\begin{teorema}{IND-CCA2 en el modelo del oráculo aleatorio (ROM)}{kem-rom}
Sea $\mathcal{A}$ un adversario IND-CCA2 contra el KEM de la
definición~\ref{def:kem-transform} que hace $q_D$ consultas al oráculo
de descifrado y $q_H$ a $H$. Existe un adversario $B$ contra la
propiedad OW-CPA de la PKE, con tiempo y recursos comparables a los de
$\mathcal{A}$ (más el costo de reencriptar cada consulta a $H$), tal que
\[
  \bigl|\Pr[b'=b]-\tfrac12\bigr|
  \;\le\;
  \Pr[\text{$B$ gana OW-CPA}]
  \;+\;\frac{q_D}{|P|}
  \;+\;\frac{2\,q_H}{|S|},
\]
donde $|P|$ es el tamaño del espacio de mensajes de la PKE y $|S|$ el
de la clave de rechazo. \emph{(Bernstein--Persichetti, 2018~\cite{bp2018},
teorema~14.3.)}
\end{teorema}

Los tres términos admiten lectura directa: el primero es la conjetura
dura (la que mide la parte~\ref{part:seguridad}); el segundo es la
ventaja de adivinar un criptograma válido y consultarlo al oráculo; el
tercero, la probabilidad de adivinar la clave de rechazo. Con los
números reales del capítulo~\ref{ch:parametros} son minúsculos: para
\texttt{mceliece348864}, $|P|>2^{456}$, mientras que la clave de rechazo
$s$ mide $n=3488$ bits ($|S|=2^{3488}$, con entropía efectiva acotada por
la semilla $\delta$ de 256 bits, $|S|_{\text{ef}}\ge 2^{256}$). Incluso con
la cota conservadora $|S|=2^{256}$, con $q_D=q_H=2^{64}$ consultas los dos
últimos términos no pasan de $2^{-392}$ y $2^{-191}$.

\begin{teorema}{IND-CCA2 en el modelo cuántico (QROM)}{kem-qrom}
En el modelo del oráculo aleatorio cuántico, con las mismas
notaciones,
\[
  \bigl|\Pr[b'=b]-\tfrac12\bigr|
  \;\le\;
  2\sqrt{\Pr[\text{$B$ gana OW-CPA}]}
  \;+\;2\,\varepsilon_{\mathrm{PRF}}
  \;+\;\varepsilon_{\mathrm{iny}},
\]
donde $\varepsilon_{\mathrm{iny}}=0$ para una PKE correcta
(«inyectividad») y $\varepsilon_{\mathrm{PRF}}$ es la ventaja contra la
pseudoaleatoriedad de $H(0,s,\cdot)$.
\emph{(Bindel--Hamburg--Hövelmanns--Hülsing--Persichetti,
2019~\cite{bindel2019}, teorema~2.)}
\end{teorema}

El término $\varepsilon_{\mathrm{PRF}}$ corresponde a una búsqueda
tipo Grover sobre $s$: con $|S|=2^{256}$ es inalcanzable. La novedad
del QROM es la \emph{raíz cuadrada} del primer término: si se quiere
mantener la ventaja IND-CCA2 por debajo de $1/1000$, el teorema exige
que ningún ataque OW-CPA con recursos comparables tenga éxito por
encima de $1/4\,000\,000$. Dicho de otro modo: \emph{en el QROM, la
seguridad probada de un KEM tipo FO es, en bits, la mitad de la
seguridad OW-CPA subyacente}. Es una pérdida de ajuste real (la
``tightness''), bien conocida, y una de las razones por las que los
parámetros de Classic McEliece se eligen con márgenes
holgados~\cite{cmsecurity}.

\begin{rigor}[La seguridad demostrable tiene grietas y errores]
La propia guía de seguridad de Classic McEliece lo advierte con
honestidad ejemplar: «la seguridad demostrable, como el
criptoanálisis, puede tener lagunas y errores». Bernstein y Persichetti
documentan contraejemplos a teoremas de seguridad IND-CCA2 publicados
antes de 2018~\cite{bp2018}: conviene tratar estos teoremas como
herramientas a revisar, no como cajas negras. Además, una prueba en el
QROM no es una prueba en el mundo real, y \emph{ninguna} prueba
descarta ataques mucho más rápidos contra OW-CPA: la confianza última
en Classic McEliece no viene de estos teoremas, sino del criptoanálisis
de la parte~\ref{part:seguridad}.
\end{rigor}

\section{El KEM de Classic McEliece: tres simplificaciones}

Classic McEliece no usa la transformación genérica tal cual: usa una
versión específica ``para mayor eficiencia'', en palabras de sus
diseñadores~\cite{cmrationale}, con tres simplificaciones. En la
práctica, su \textsf{Encap} y su \textsf{Decap} se leen así
(pk$=T$, la clave pública de $mt\times k$ bits;
sk$=(\delta,c,g,\alpha,s)$, los datos del código y la clave de rechazo;
capítulos~\ref{ch:keygen}--\ref{ch:encapdecap}):

\begin{practica}[El KEM de Classic McEliece, en dos columnas]
\begin{minipage}[t]{0.47\textwidth}\raggedright
\textbf{\textsf{Encap}$(\mathrm{pk})$}
\begin{enumerate}[leftmargin=1.4em,itemsep=2pt,topsep=2pt]
  \item $e\leftarrow$ vector uniforme de peso $t$
        ($\textsf{FixedWeight}$).
  \item $C\leftarrow(\mathrm{I}_{mt}\,|\,T)\,e$.
  \item $K\leftarrow H(1,e,C)$.
  \item salida: $(C,K)$.
\end{enumerate}
\end{minipage}\hfill
\begin{minipage}[t]{0.47\textwidth}\raggedright
\textbf{\textsf{Decap}$(\mathrm{sk},C)$}
\begin{enumerate}[leftmargin=1.4em,itemsep=2pt,topsep=2pt]
  \item $e'\leftarrow\textsf{Decode}(C)$
        ($\textsf{Patterson}$; puede fallar).
  \item si $e'$ es válido, $\wt(e')=t$ y
        $(\mathrm{I}_{mt}\,|\,T)\,e'=C$:
        $K\leftarrow H(1,e',C)$.
  \item si no: $K\leftarrow H(0,s,C)$ (rechazo implícito).
  \item salida: $K$.
\end{enumerate}
\end{minipage}
\end{practica}

Las tres diferencias con la definición~\ref{def:kem-transform} son:

\begin{enumerate}[leftmargin=1.8em,itemsep=3pt]
  \item \textbf{Sin reencriptar.} La comprobación
        $\textsf{Enc}(\mathrm{pk},e')=C$ no hace falta: el propio
        $\textsf{Decode}$ ya exige $\wt(e')=t$ y
        $(\mathrm{I}_{mt}|T)e'=C$. La guía de seguridad llama a esto la
        \emph{rigidez} de la PKE: si un criptograma se descifra, se
        descifra a un único $e'$ verificado~\cite{cmsecurity}.
  \item \textbf{Sin clave pública dentro de la clave secreta.} Como el
        descifrado no reencripta, no necesita $\mathrm{pk}'$; se
        elimina~\cite{cmsecurity}.
  \item \textbf{La clave de rechazo $s$ es una cadena uniforme de $n$
        bits} (no un vector de peso $t$): es más simple de generar, y
        los teoremas posteriores de 2017--2019 permiten claves de
        rechazo más grandes que el espacio de
        mensajes~\cite{hhk2017,bp2018,bindel2019}.
\end{enumerate}

\begin{table}[htbp]
\centering
\small
\begin{tabular}{p{3.4cm} p{5.1cm} p{5.3cm}}
\toprule
 & transformación genérica & Classic McEliece \\
 & (definición~\ref{def:kem-transform}) & (especificación) \\
\midrule
clave secreta del KEM & $(\mathrm{sk}',\mathrm{pk}',s)$
  & $(\mathrm{sk}',s)$, sin $\mathrm{pk}'$ \\
verificación al descifrar & reencriptar:
  $\textsf{Enc}(\mathrm{pk},e')=C$
  & ya la hace $\textsf{Decode}$: $\wt(e')=t$ y $He'=C$ \\
clave de rechazo $s$ & $256$ bits uniformes
  & $n$ bits uniformes ($n\ge3488$) \\
semilla $e$ & uniforme en $P$
  & $\textsf{FixedWeight}$, uniforme en peso $t$ \\
\bottomrule
\end{tabular}
\caption{La transformación genérica y el KEM real, lado a lado. La
distribución de $e$ y la corrección del esquema son idénticas; lo que
cambia es cuánto material hay que guardar y recalcular.}
\label{tab:kem-cm}
\end{table}

Queda un detalle de la especificación que usaremos en el
capítulo~\ref{ch:encapdecap}: el hash $H$ tiene \emph{separación de
dominio} por etiquetas. La etiqueta $1$ deriva claves legítimas, la
$0$ claves de rechazo y la $2$ se reserva para la confirmación de texto
claro de las variantes \texttt{pc} (capítulo~\ref{ch:parametros}). Es
la versión concreta de la idea del capítulo: distintas entradas,
distintos espacios de salida, sin ambigüedades.

\section{Ejercicios}

\begin{enumerate}[leftmargin=1.6em,itemsep=4pt]
  \item \textbf{Corrección del KEM.} Recorra
        $\textsf{Decap}(\mathrm{sk},\textsf{Encap}(\mathrm{pk}))$ en la
        definición~\ref{def:kem-transform} y compruebe que las dos
        partes calculan la misma clave. Señale el único punto donde se
        usa que la PKE es \emph{correcta}.
  \item \textbf{¿Por qué $H(1,e,C)$ y no $H(e)$?} En una
        implementación real, dos cadenas de bytes distintas pueden
        decodificar al mismo criptograma $C$ (bits de relleno; véase el
        capítulo~\ref{ch:parametros}): el oráculo de descifrado acepta
        ambas. Demuestre que un esquema con $K=H(e)$ sería vulnerable
        en ese escenario: el adversario consulta el oráculo con la
        segunda codificación de $C^*$ y obtiene $K^*$. ¿Por qué la
        misma consulta es inútil con $K=H(1,e,C)$?
  \item \textbf{Rechazo implícito.} En muchos protocolos, un
        $\textsf{Decap}$ que devuelve $\bot$ obliga a manejar el error
        (alertas TLS, reintentos\ldots), y esa reacción es observable.
        Explique por qué devolver siempre una clave ---pseudoaleatoria
        si hace falta--- es una ventaja de seguridad, y en qué parte
        de los teoremas~\ref{thm:kem-rom} y~\ref{thm:kem-qrom} aparece
        la clave $s$.
  \item \textbf{Los números de la cota ROM.} Para
        \texttt{mceliece348864} se tiene $|P|>2^{456}$ y la entropía efectiva
        de la clave de rechazo $s\in\Ftwo^n$ está fijada por la semilla
        $\delta$ ($|S|\ge 2^{256}$). Con $q_D=q_H=2^{64}$, calcule los valores
        de los términos segundo y tercero del teorema~\ref{thm:kem-rom}. ¿Qué haría
        falta para que el término $\frac{2q_H}{|S|}$ dejara de ser
        despreciable?
  \item \textbf{La raíz cuadrada del QROM.} Usando el
        teorema~\ref{thm:kem-qrom} (ignore los términos de $s$):
        (a)~¿qué éxito OW-CPA máximo tolera el teorema si se quiere
        ventaja IND-CCA2 menor que $1/1000$?; (b)~¿y si se quieren
        $128$ bits de ventaja, es decir, $2^{-128}$? Comente por qué
        esto obliga a elegir parámetros con margen amplio.
  \item \textbf{¿Por qué no Fujisaki--Okamoto a secas?} El rationale
        dice que la conversión de Classic McEliece «no es FO
        genérico»~\cite{cmrationale}. Explique, con la
        tabla~\ref{tab:kem-cm} delante, qué ahorra cada una de las
        simplificaciones (material de clave, reencriptaciones) y por
        qué la verificación de $\textsf{Decode}$ bastaría para
        justificarlas.
\end{enumerate}


\begin{center}\rule{0.55\textwidth}{0.4pt}\end{center}
\medskip
\noindent\textit{Cierre de la parte~\ref{part:mceliece}.} La idea de
McEliece ---esconder un código con decodificador eficiente tras un
isomorfismo--- quedó convertida en criptosistema
(capítulo~\ref{ch:1978}), reformulada en su forma dual con mensajes
de peso fijo (capítulo~\ref{ch:niederreiter}) y transformada en un
KEM con rechazo implícito (capítulo~\ref{ch:kem}), con sus
reducciones a la dureza de la decodificación y sus dos modelos, ROM y
QROM. La parte~\ref{part:classic} abre ahora la especificación de
Classic McEliece palabra por palabra: semillas, matrices,
laboratorios y parámetros.


% ======================================================================
%  PARTE III — CLASSIC MCELIECE
% ======================================================================
\part{Classic McEliece}\label{part:classic}

\chapter{La generación de claves}\label{ch:keygen}

\begin{roadmap}
Aquí empieza el Classic McEliece real. Este capítulo abre la
especificación por su primera mitad ---la generación de claves--- y
muestra algo notable: todo el proceso es determinista a partir de una
semilla de $256$ bits. Veremos cómo se sortea el orden del soporte
($\textsf{FieldOrdering}$), cómo se sortea el polinomio de Goppa
($\textsf{Irreducible}$) y cómo se reduce la matriz de paridad a forma
sistemática ($\textsf{MatGen}$), con sus probabilidades de éxito y sus
tamaños exactos. Es el capítulo que convierte la teoría de la parte~I
en ingeniería auditable.
\end{roadmap}

\begin{historia}[De una propuesta a un estándar]
Classic McEliece se presentó al concurso del NIST en 2017 y llegó a la
ronda final de KEMs; en 2022 el NIST eligió Kyber (ML-KEM) y en marzo
de 2025 cerró su proceso con HQC: Classic McEliece quedó fuera de la
selección (capítulo~\ref{ch:parametros}). El veredicto no fue un juicio
sobre su seguridad ---el criptoanálisis no encontró grietas en sus
parámetros--- sino sobre sus claves de cientos de kilobytes, incómodas
para los usos que allí se buscaban. En junio de 2026, ISO/IEC publicó
su estandarización en la enmienda~2 de la norma de cifrado asimétrico
(ISO/IEC 18033-2:2006/Amd.~2:2026), junto a ML-KEM y FrodoKEM, con los
$16$ conjuntos de parámetros que veremos en el
capítulo~\ref{ch:parametros}.
\end{historia}

\section{La idea: una semilla manda}

La especificación define la generación de claves en dos niveles. El
algoritmo $\textsf{KeyGen}$ no recibe más que los parámetros: sortea una
cadena uniforme $\delta$ de $\ell=256$ bits ---la \emph{semilla}--- y
llama a $\textsf{SeededKeyGen}(\delta)$. Y $\textsf{SeededKeyGen}$ es
completamente \emph{determinista}: la misma $\delta$ produce siempre la
misma clave. Sus pasos son los siguientes:

\begin{enumerate}[leftmargin=1.8em,itemsep=3pt]
  \item Sea $E=G(\delta)$ la cadena de
        $n+\sigma_2q+\sigma_1t+\ell$ bits que resulta de aplicar el
        generador
        $G=\textsf{SHAKE256}(64,\delta)$ (primero la etiqueta $64$, que
        lo separa del resto de usos del hash).
  \item Los primeros $n$ bits de $E$ son $s$, la clave de rechazo
        implícito del capítulo~\ref{ch:kem}.
  \item Con los siguientes $\sigma_2q=32q$ bits, $\textsf{FieldOrdering}$
        produce $q$ elementos distintos
        $\alpha_0,\dots,\alpha_{q-1}$ de $\GF q$. Si falla, se
        reinicia con $\delta\leftarrow\delta'$.
  \item Con los siguientes $\sigma_1t=16t$ bits, $\textsf{Irreducible}$
        produce un polinomio mónico irreducible $g$ de grado $t$. Si
        falla, se reinicia con $\delta\leftarrow\delta'$.
  \item Sea $\Gamma=(g,\alpha_0,\alpha_1,\dots,\alpha_{n-1})$: el
        soporte del código son los primeros $n$ elementos del
        ordenamiento ($\alpha_n,\dots,\alpha_{q-1}$ no se usan).
  \item Con $\textsf{MatGen}(\Gamma)$ se calculan la clave pública $T$
        y los datos $\Gamma_0$ y $c$ (la lista de columnas permutadas,
        para las variantes \texttt{f}). Si falla, se reinicia con
        $\delta\leftarrow\delta'$.
  \item La clave pública es $T$; la clave secreta es
        $(\delta,c,g,\alpha,s)$, donde $\alpha$ es el ordenamiento
        $\alpha_0,\dots,\alpha_{n-1}$ tras los intercambios de
        $\textsf{MatGen}$, seguido de $\alpha_n,\dots,\alpha_{q-1}$.
\end{enumerate}

El truco del reinicio merece un comentario: los $256$ bits finales de
$E$ son $\delta'$, la semilla del intento siguiente. \emph{El azar no
se pide de nuevo al generador externo}: todo el proceso, incluidos los
reintentos, es una función determinista de la semilla original. Por eso
una clave privada de Classic McEliece podría comprimirse a sus $32$
bytes de semilla y regenerarse por completo; la especificación guarda
la versión expandida solo por eficiencia.

\begin{figure}[htbp]
\centering
\begin{tikzpicture}[
    sbox/.style={draw=mcblue,fill=mcsky!45,rounded corners=2.5pt,align=center,
      font=\footnotesize,inner sep=3.5pt,text width=2.55cm,minimum height=0.9cm},
    fbox/.style={draw=mcorange,fill=mcorange2!55,rounded corners=2.5pt,align=center,
      font=\footnotesize,inner sep=3.5pt,text width=2.8cm,minimum height=0.9cm},
    arr/.style={-{Stealth[length=2.2mm]},thick,mcgray}
  ]
  % semilla y PRG
  \node[sbox] (delta) at (0,0) {$\delta$: $256$ bits\\(semilla)};
  \draw[arr] (delta.south) -- (0,-2.05);
  \node[align=center,font=\scriptsize,color=mcgray] at (3.1,-1.2)
    {$E=G(\delta)=\textsf{SHAKE256}(64,\delta)$\\ (primeros $n+32q+16t+256$ bits)};
  % particion de E
  \node[sbox,text width=1.9cm] (sb) at (-4.3,-2.6) {$s$\\ $n$ bits};
  \node[sbox,text width=2.0cm] (ab) at (-1.5,-2.6) {$\alpha^\circ$\\ $32q$ bits};
  \node[sbox,text width=1.9cm] (gb) at (1.5,-2.6) {$g^\circ$\\ $16t$ bits};
  \node[sbox,text width=1.9cm] (db) at (4.3,-2.6) {$\delta'$\\ $256$ bits};
  % procesos
  \node[fbox] (fo) at (-1.5,-4.2) {$\textsf{FieldOrdering}$};
  \node[fbox] (ir) at (1.5,-4.2) {$\textsf{Irreducible}$};
  \draw[arr] (ab)--(fo);
  \draw[arr] (gb)--(ir);
  % salidas
  \node[align=center,font=\footnotesize] (alpha) at (-1.5,-5.6)
    {$\alpha_0,\dots,\alpha_{q-1}$\\ orden uniforme\\(y secreto)};
  \node[align=center,font=\footnotesize] (gg) at (1.5,-5.6)
    {$g$ mónico irreducible\\de grado $t$};
  \draw[arr] (fo)--(alpha);
  \draw[arr] (ir)--(gg);
  % s y reinicio
  \draw[arr] (sb)--(-4.3,-3.3);
  \node[align=center,font=\scriptsize,color=mcgray] at (-4.3,-3.75)
    {clave de rechazo\\(capítulo~\ref{ch:kem})};
  \draw[arr,dashed,mcred] (db.east) -- (5.9,-2.6) -- (5.9,0) -- (1.42,0);
  \node[align=center,font=\tiny,color=mcred] at (5.35,-1.3)
    {si algo falla:\\ $\delta\leftarrow\delta'$};
  % Gamma
  \node[fbox,text width=6.4cm] (gamma) at (0,-7.2)
    {$\Gamma=(g,\alpha_0,\alpha_1,\dots,\alpha_{n-1})$};
  \draw[arr] (alpha)--(gamma);
  \draw[arr] (gg)--(gamma);
\end{tikzpicture}
\caption{La expansión de la semilla. De $\delta$ sale todo: $s$, la
materia prima de $\alpha$ y $g$, y la semilla $\delta'$ de un eventual
reintento. Los cuerpos son $\GF q=\GF{2^m}$ con $q=2^m$; el campo de
Goppa es $\GF q[y]/F(y)$ con $F$ fijado por los parámetros
(tabla~\ref{tab:keygen-pol}).}
\label{fig:keygen-semilla}
\end{figure}

\section{\textsf{FieldOrdering}: los $\alpha_i$}

El soporte de un código de Goppa es una lista de $q$ elementos
distintos de $\GF q$: ordenados de una manera y no de otra, definen
---como vimos en el capítulo~\ref{ch:niederreiter}--- la permutación
secreta que esconde la estructura del código. La especificación lo
sortea así:

\begin{enumerate}[leftmargin=1.8em,itemsep=2pt]
  \item Los $\sigma_2q=32q$ bits de entrada se leen como $q$ enteros de
        $32$ bits $a_0,\dots,a_{q-1}$.
  \item Si los $a_i$ no son todos distintos, se falla ($\bot$).
  \item Se ordenan los pares $(a_i,i)$ lexicográficamente; resulta una
        permutación $\pi$ con $a_{\pi(0)}<a_{\pi(1)}<\cdots<
        a_{\pi(q-1)}$.
  \item Se define $\alpha_i=\sum_{j=0}^{m-1}\pi(i)_j\,z^{m-1-j}$, donde
        $\pi(i)_j$ es el $j$-ésimo bit menos significativo de
        $\pi(i)$.
\end{enumerate}

\begin{proposicion}{$\textsf{FieldOrdering}$ produce un orden uniforme}{field-ordering}
Si los $32q$ bits de entrada son uniformes e independientes y el
algoritmo no falla, $(\alpha_0,\dots,\alpha_{q-1})$ es un ordenamiento
uniforme de $\GF q$. La probabilidad de no fallar es
\[
  \prod_{i=1}^{q-1}\Bigl(1-\frac{i}{2^{32}}\Bigr),
\]
que supera $0{,}99$ para $q\le2^{13}$ y $0{,}6$ para $q\le2^{16}$.
\end{proposicion}

\begin{proof}
Condicionados a ser distintos, los $a_i$ son un vector uniforme de
enteros distintos (simetría de las $q!$ posibilidades). Ordenarlos da
una permutación uniforme $\pi$, y el mapa $\pi(i)\mapsto\alpha_i$ es una
biyección fija de $\{0,\dots,q-1\}$ en $\GF q$; componer una biyección
con una permutación uniforme da un ordenamiento uniforme. La
probabilidad de que $q$ enteros uniformes sean distintos es el producto
de las probabilidades condicionales $\bigl(1-\tfrac{i}{2^{32}}\bigr)$;
las cotas numéricas están en la razón de
diseño~\cite{cmrationale}.
\end{proof}

Un ejemplo minúsculo: con $m=2$ ($q=4$) y enteros $(7,1,9,3)$, el orden
de $\pi$ es $1<3<7<9$, luego $\pi=(1,3,0,2)$ y, leyendo los bits en el
orden de la definición, $(\alpha_0,\alpha_1,\alpha_2,\alpha_3)
=(z,\,z+1,\,0,\,1)$. Con todos los $a_i$ iguales, el algoritmo fallaría:
la distinción estricta es lo que garantiza que $\pi$ sea una permutación
de verdad.

La clave privada no guarda la lista como $q$ enteros de $m$ bits, sino
la permutación $\pi$ comprimida como los bits de control de una
\emph{red de Beneš}: una red de interconexión con $2m-1$ etapas de $2^{m-1}$
conmutadores elementales de $2\times2$. Cada conmutador toma dos entradas y,
según un bit de control $b\in\{0,1\}$, las deja pasar directas ($b=0$) o las
intercambia ($b=1$). En total se requieren $(2m-1)2^{m-1}$ bits de control,
empaquetados en $\lceil(2m-1)2^{m-4}\rceil$ bytes. Para $m=12$, $\alpha$ ocupa
$5888$ bytes (la columna $\alpha$ de la tabla~\ref{tab:keygen-tamanos});
decapsular recorre la red en sentido inverso y reconstruye el orden
exacto. El formato es canónico ---los mismos bytes en toda
implementación--- porque la reproducibilidad de la clave es un
requisito de interoperabilidad del estándar, no una manía de sus
autores.
\label{sec:benes}

\begin{idea}[¿Por qué una red de Beneš en vez de una tabla?]
Si la desencapsulación reconstruyera el orden secreto consultando una tabla
$\pi[i]$ en memoria, los accesos con índices secretos provocarían variaciones
temporales medibles a través de la memoria caché (ataques de canal lateral).
En una red de Beneš, la permutación se aplica mediante una secuencia fija de
operaciones lógicas sobre los cables guiadas por los bits de control: las
direcciones de memoria y el tiempo de ejecución son \emph{estrictamente independientes}
del secreto, garantizando tiempo constante por diseño.
\end{idea}

\section{\textsf{Irreducible}: el polinomio $g$}

El polinomio de Goppa debe ser mónico, irreducible y de grado $t$; y,
para que el sistema sea seguro, debe ser \emph{impredecible}. La
especificación lo genera de una forma elegante:

\begin{enumerate}[leftmargin=1.8em,itemsep=2pt]
  \item Los $\sigma_1t=16t$ bits de entrada se parten en $t$ bloques de
        $16$ bits; de cada bloque se usan los primeros $m$ bits para
        formar $\beta_j=\sum_{i=0}^{m-1}d_{\sigma_1 j+i}z^i$, y se
        ignora el resto.
  \item Se construye $\beta=\beta_0+\beta_1y+\cdots+\beta_{t-1}y^{t-1}$
        en el cuerpo $\GF q[y]/F(y)$ (de tamaño $q^t$).
  \item Se calcula el polinomio minimal $g$ de $\beta$ sobre $\GF q$.
        Si su grado es $t$, se devuelve; si no, se falla ($\bot$).
\end{enumerate}

\begin{idea}[El polinomio minimal vía álgebra lineal]
El \emph{polinomio minimal} de un elemento $\beta\in\GF{q^t}$ sobre $\GF{q}$ es el
polinomio mónico de menor grado no nulo con coeficientes en $\GF{q}$ que satisface
$g(\beta)=0$. Dado que $\GF{q^t}$ tiene dimensión $t$ como espacio vectorial sobre
$\GF{q}$, los $t+1$ vectores $\{1,\beta,\beta^2,\dots,\beta^t\}$ son necesariamente
linealmente dependientes. Encontrar la combinación lineal mónica
$c_0+c_1\beta+\cdots+\beta^t=0$ es simplemente resolver un sistema lineal homogéneo
sobre $\GF{q}$ mediante eliminación gaussiana. Si no existe dependencia lineal
de grado menor que $t$, el polinomio resultante $g$ tiene grado $t$ y es
automáticamente irreducible sobre $\GF{q}$.
\end{idea}

\begin{proposicion}{$\textsf{Irreducible}$ produce un polinomio uniforme}{irreducible-uniforme}
Si los $16t$ bits de entrada son uniformes, la salida $g$ (cuando no se
falla) es un polinomio mónico irreducible de grado $t$ uniforme, y la
probabilidad de éxito es mayor que $0{,}99$ para $t\ge16$.
\end{proposicion}

\begin{proof}
$\beta$ es un elemento uniforme del cuerpo $\GF{q^t}$. Cada polinomio
mónico irreducible $g$ de grado $t$ tiene exactamente $t$ raíces en
$\GF{q^t}$ (divide a $y^{q^t}-y$, que se escinde con raíces simples), y
son exactamente los $\beta$ cuyo minimal es $g$: cada $g$ aparece por
tanto $t$ de las $q^t$ elecciones de $\beta$, es decir, con la misma
probabilidad. El algoritmo falla únicamente cuando $\beta$ cae en un
subcuerpo propio de $\GF{q^t}$; el número de tales elementos está
dominado por $q^{t/2+\varepsilon}$ (suma sobre divisores propios de
$t$), de modo que, para $t\ge16$, el fallo es inferior al $1\%$ (en
nuestros parámetros, muchísimo menos)~\cite{cmrationale}.
\end{proof}

\section{\textsf{MatGen}: de $\Gamma$ a la clave pública}

Con $g$ y el ordenamiento $\alpha$ a mano, $\textsf{MatGen}$ fabrica la
clave pública. En el caso sistemático ($(\mu,\nu)=(0,0)$, los conjuntos
sin sufijo \texttt{f}):

\begin{enumerate}[leftmargin=1.8em,itemsep=2pt]
  \item Se construye la matriz $\tilde H$, de $t\times n$ sobre
        $\GF q$, con entradas $h_{i,j}=\alpha_j^i/g(\alpha_j)$ para
        $i=0,\dots,t-1$ y $j=0,\dots,n-1$. Es la matriz de paridad
        clásica del código de Goppa de la definición del
        capítulo~\ref{ch:goppa}.
  \item Se expande cada entrada (un elemento de $\GF q$, es decir, $m$
        bits $u_0,\dots,u_{m-1}$) a una columna de $m$ bits: resulta
        $\hat H$, de $mt\times n$ sobre $\Ftwo$.
  \item Se reduce $\hat H$ a forma sistemática $(\mathrm{I}_{mt}\,|\,T)$
        mediante eliminación gaussiana. Si no es posible, se falla
        ($\bot$).
  \item Se devuelve $T$: la clave pública, de $mt\times k$ bits.
\end{enumerate}

\begin{figure}[htbp]
\centering
\begin{tikzpicture}[
    sbox/.style={draw=mcblue,fill=mcsky!45,rounded corners=2.5pt,align=center,
      font=\footnotesize,inner sep=3.5pt,minimum height=0.95cm},
    fbox/.style={draw=mcorange,fill=mcorange2!55,rounded corners=2.5pt,align=center,
      font=\footnotesize,inner sep=3.5pt,minimum height=0.95cm},
    arr/.style={-{Stealth[length=2.2mm]},thick,mcgray}
  ]
  \node[sbox,text width=2.7cm] (G) at (0,0) {$\Gamma=(g,\alpha_0,\dots,\alpha_{n-1})$};
  \node[fbox,text width=3.3cm] (Ht) at (4.7,0)
    {$\tilde H$: $t\times n$ sobre $\GF q$\\ $h_{i,j}=\alpha_j^i/g(\alpha_j)$};
  \node[fbox,text width=3.5cm] (Hh) at (9.6,0)
    {$\hat H$: $mt\times n$ sobre $\Ftwo$\\ (cada entrada $\to$ $m$ bits)};
  \draw[arr] (G)--(Ht);
  \draw[arr] (Ht)--(Hh);
  % matriz sistematica
  \begin{scope}[shift={(9.6,-2.9)}]
    \draw[fill=mcblue!75,draw=mcblue] (-2.35,-0.65) rectangle (-1.05,0.65);
    \draw[fill=mcsky!45,draw=mcblue] (-1.05,-0.65) rectangle (1.45,0.65);
    \node[font=\scriptsize,color=white] at (-1.7,0) {$\mathrm{I}_{mt}$};
    \node[font=\scriptsize] at (0.2,0) {$T$};
    \draw[decorate,decoration={brace,amplitude=4pt},mcgray] (-2.35,-0.85)--(-1.05,-0.85)
      node[midway,below=2pt,font=\tiny,color=mcgray]{$mt$};
    \draw[decorate,decoration={brace,amplitude=4pt},mcgray] (-1.05,-0.85)--(1.45,-0.85)
      node[midway,below=2pt,font=\tiny,color=mcgray]{$k$};
  \end{scope}
  \draw[arr] (Hh)--(9.6,-2.1);
  \node[align=center,font=\scriptsize,color=mcgray,text width=4.6cm] at (4.6,-2.9)
    {eliminación gaussiana; si no se puede, $\bot$: nuevo intento
     ($\approx29\%$ de éxito; $\sim3{,}4$ intentos de media)};
  \node[fbox,text width=3.2cm] (T) at (9.6,-5.2)
    {clave pública: $T$ ($mt\times k$ bits)};
  \draw[arr] (9.6,-3.55)--(T);
\end{tikzpicture}
\caption{$\textsf{MatGen}$ en tres actos. La matriz de Goppa $t\times
n$ sobre $\GF q$ se expande bit a bit a una matriz binaria $mt\times
n$; la eliminación gaussiana la lleva a la forma sistemática
$(\mathrm{I}_{mt}\,|\,T)$ y publica el bloque $T$. La permutación
secreta del capítulo~\ref{ch:niederreiter} es aquí el ordenamiento
$\alpha$: las columnas de $\tilde H$ están etiquetadas por elementos
del cuerpo, en un orden que solo conoce la clave secreta.}
\label{fig:keygen-matgen}
\end{figure}

¿Por qué puede fallar el paso 3? La eliminación por filas no puede
mover columnas: existe la forma $(\mathrm{I}_{mt}\,|\,T)$ exactamente
cuando las primeras $mt$ columnas de $\hat H$ son linealmente
independientes. Para una matriz binaria «genérica» eso ocurre con
probabilidad
\[
  \prod_{i=0}^{mt-1}\Bigl(1-2^{\,i-mt}\Bigr)\approx0{,}289,
\]
es decir, falla en torno al $71\%$ de los intentos: se necesitan
$\sim3{,}4$ de media, cada uno con su semilla nueva
$\delta\leftarrow\delta'$~\cite{cmrationale}. Las variantes \texttt{f}
admiten la forma \emph{semi-sistemática} $(\mu,\nu)=(32,64)$: se
permiten hasta $\mu=32$ intercambios de columna dentro de una ventana
de $\nu=64$, lo que reduce el fallo por intento a menos de $2^{-30}$
(prácticamente, un solo intento); los intercambios quedan registrados
en la lista $c$ de la clave secreta~\cite{cmrationale}. La conversión
cambia la seguridad en, a lo sumo, $2$ bits, y la estudiaremos en el
capítulo~\ref{ch:parametros}.

\begin{observacion}{El código es público; el secreto es la descripción}{code-public}
Tras $\textsf{MatGen}$, cualquiera puede calcular el código $\ker H$
con $H=(\mathrm{I}_{mt}|T)$: el subespacio no es secreto. Lo que guarda
la clave privada es la \emph{descripción eficiente} del código ---el
polinomio $g$ y el orden de las $\alpha_i$---, la única que permite
decodificar con Patterson. Esta distinción se confunde a menudo y es
esencial: un atacante que descubriera $g$ y el ordenamiento tendría un
decodificador; un atacante con solo el código en la mano tiene el
problema duro de la parte~\ref{part:seguridad}.
\end{observacion}

\section{Los tamaños reales}

Los cinco conjuntos principales fijan el cuerpo y los polinomios de la
tabla~\ref{tab:keygen-pol}; la tabla~\ref{tab:keygen-tamanos} da los
tamaños de las claves.

\begin{table}[htbp]
\centering
\small
\begin{tabular}{lccccc}
\toprule
Conjunto & $m$ & $n$ & $t$ & $k=n-mt$ & polinomios $f(z)$ (cuerpo) y $F(y)$ (Goppa)\\
\midrule
\texttt{mceliece348864} & 12 & 3488 & 64 & 2720
  & $z^{12}\!+\!z^3\!+\!1$; \quad $y^{64}\!+\!y^3\!+\!y\!+\!z$\\
\texttt{mceliece460896} & 13 & 4608 & 96 & 3360
  & $z^{13}\!+\!z^4\!+\!z^3\!+\!z\!+\!1$; \quad $y^{96}\!+\!y^{10}\!+\!y^9\!+\!y^6\!+\!1$\\
\texttt{mceliece6688128} & 13 & 6688 & 128 & 5024
  & $z^{13}\!+\!z^4\!+\!z^3\!+\!z\!+\!1$; \quad $y^{128}\!+\!y^7\!+\!y^2\!+\!y\!+\!1$\\
\texttt{mceliece6960119} & 13 & 6960 & 119 & 5413
  & $z^{13}\!+\!z^4\!+\!z^3\!+\!z\!+\!1$; \quad $y^{119}\!+\!y^8\!+\!1$\\
\texttt{mceliece8192128} & 13 & 8192 & 128 & 6528
  & $z^{13}\!+\!z^4\!+\!z^3\!+\!z\!+\!1$; \quad $y^{128}\!+\!y^7\!+\!y^2\!+\!y\!+\!1$\\
\bottomrule
\end{tabular}
\caption{Los cinco conjuntos sin sufijo (véanse las variantes
\texttt{f}, \texttt{pc} y \texttt{pcf} en el
capítulo~\ref{ch:parametros}). $f$ define el cuerpo
$\GF q=\Ftwo[z]/f(z)$; $F$ define $\GF q[y]/F(y)$, el cuerpo donde
vive $\beta$ para $\textsf{Irreducible}$.}
\label{tab:keygen-pol}
\end{table}

\begin{table}[htbp]
\centering
\small
\setlength{\tabcolsep}{6pt}
\begin{tabular}{lrrrrrrr}
\toprule
Conjunto & $pk$ & $\delta$ & $c$ & $g$ & $\alpha$ & $s$ & $sk$\\
 & (bytes) & (bytes) & (bytes) & (bytes) & (bytes) & (bytes) & (bytes)\\
\midrule
\texttt{mceliece348864} & \num{261120} & 32 & 8 & 128 & 5888 & 436 & \num{6492}\\
\texttt{mceliece460896} & \num{524160} & 32 & 8 & 192 & 12800 & 576 & \num{13608}\\
\texttt{mceliece6688128} & \num{1044992} & 32 & 8 & 256 & 12800 & 836 & \num{13932}\\
\texttt{mceliece6960119} & \num{1047319} & 32 & 8 & 238 & 12800 & 870 & \num{13948}\\
\texttt{mceliece8192128} & \num{1357824} & 32 & 8 & 256 & 12800 & 1024 & \num{14120}\\
\bottomrule
\end{tabular}
\caption{Tamaños de las claves. La pública ocupa $mt\lceil k/8\rceil$
bytes: fila a fila, la matriz $T$ con $\lceil k/8\rceil$ bytes por fila.
La secreta suma la semilla $\delta$ ($32$), la lista $c$ ($8$ en ambos
tipos de conjunto: $\lceil\nu/8\rceil$ o la codificación fija de
$2^{32}-1$ para $(\mu,\nu)=(0,0)$), el polinomio $g$
($t\lceil m/8\rceil$), el ordenamiento $\alpha$ (la red de Benes,
$\lceil(2m-1)2^{m-4}\rceil$ bytes: $2m-1$ etapas de $2^{m-1}$
conmutadores de un bit) y la clave de rechazo $s$ ($\lceil n/8\rceil$).}
\label{tab:keygen-tamanos}
\end{table}

Dos observaciones sobre la tabla. La primera: la clave pública tiene el
tamaño $k(n-k)$ que prometía el capítulo~\ref{ch:niederreiter}
($mt\cdot k$ bits); en los conjuntos \texttt{mceliece6*} ronda el
megabyte, y el equipo de Classic McEliece recomienda no bajar de
$n=6688$ para seguridad a largo plazo. La segunda: en la clave secreta,
el ordenamiento $\alpha$ pesa
más que todo lo demás junto ---es una permutación de $q$ elementos---,
y aun así podemos regenerarlo entero desde los $32$ bytes de $\delta$.

\section{Ejercicios}

\begin{enumerate}[leftmargin=1.6em,itemsep=4pt]
  \item \textbf{Reproducibilidad.} Suponga que, en el paso 3,
        $\textsf{FieldOrdering}$ falla por un choque entre enteros.
        Explique por qué el proceso completo sigue siendo determinista
        y cómo las pruebas de vectores de la especificación pueden
        ejecutar (y comparar) los reintentos sin generar azar nuevo.
  \item \textbf{$\textsf{FieldOrdering}$ con el juguete.} Con $m=2$
        ($q=4$) y el vector de enteros de 32 bits $(7,1,9,3)$,
        reproduzca $\pi$ y deduzca $(\alpha_0,\dots,\alpha_3)$. Repita
        con $(0,1,2,3)$ y con $(0,0,0,0)$; explique qué ocurre en el
        último caso.
  \item \textbf{Uniformidad de $\textsf{Irreducible}$.} Demuestre con
        detalle que cada polinomio mónico irreducible de grado $t$
        aparece con exactamente $t$ de las $q^t$ elecciones de
        $\beta$; concluya la uniformidad de $g$.
  \item \textbf{La probabilidad del $29\%$.} Deduzca
        $\prod_{i=0}^{mt-1}(1-2^{i-mt})\approx0{,}289$ como
        probabilidad de que una matriz binaria $mt\times n$
        «genérica» tenga sus primeras $mt$ columnas independientes.
        ¿Cuántos intentos de media se necesitan? ¿Qué compra la forma
        semi-sistemática $(\mu,\nu)=(32,64)$?
  \item \textbf{Tamaños.} Verifique con las fórmulas del pie de la
        tabla~\ref{tab:keygen-tamanos} los números de
        \texttt{mceliece8192128}: $pk=1357824$ bytes y $sk=14120$
        bytes. ¿Qué componente domina la clave secreta, y cuál la
        pública?
  \item \textbf{El secreto es la descripción.} Argumente que el código
        $\ker(\mathrm{I}_{mt}|T)$ es público y calcule, para
        \texttt{mceliece348864}, su dimensión y su número de
        codewords. ¿Qué gana un atacante que descubra $g$ y el
        ordenamiento $\alpha$? Relacione la respuesta con la
        observación~\ref{obs:code-public} y con la
        proposición~\ref{prop:nied-reduccion} del
        capítulo~\ref{ch:niederreiter}.
\end{enumerate}


\chapter{Cifrado y descifrado}\label{ch:encapdecap}

\begin{roadmap}
Este capítulo recorre las rutinas que ejecuta el KEM en cada
operación ---$\textsf{FixedWeight}$, $\textsf{Encode}$, $\textsf{Decode}$
y el hash final--- en el orden exacto de la especificación. Veremos por
qué el mensaje de un KEM niederreiteriano es un vector de peso exacto,
cómo se descifra con Patterson sobre el soporte secreto y por qué un
último chequeo (peso y síndrome) es imprescindible para que la
seguridad IND-CCA2 sea la del capítulo~\ref{ch:kem}. Al terminar, el
lector podrá leer la especificación línea a línea sin sorpresas.
\end{roadmap}

\begin{historia}[Por qué «Classic»]
El nombre no es nostalgia. A partir de 2017, el NIST recibió docenas
de propuestas basadas en códigos ---cuasicíclicas, de paridad moderada,
convolucionales\ldots--- que ahorraban bytes añadiendo estructura. El
equipo de McEliece presentó lo contrario: el McEliece clásico, con
códigos de Goppa binarios, de 1978, sin estructura adicional. Frente a
los trucos que ahorran espacio a costa de nueva superficie de ataque,
la apuesta fue pagar claves grandes y conservar casi medio siglo de
criptoanálisis limpio. La apuesta le valió un estándar ISO en 2026 y
sigue sin roturas algorítmicas conocidas.
\end{historia}

\section{\textsf{FixedWeight}: la semilla del cifrador}

En un cifrado de ElGamal se sortea un exponente; en un KEM
niederreiteriano, el azar que hay que sortear es un \emph{vector de
peso exacto}. La especificación genera $e$ con $\textsf{FixedWeight}$,
que no recibe nada y devuelve $e\in\Ftwo^{n}$ con $\wt(e)=t$ uniforme. El
algoritmo usa lotes de $\tau$ enteros, con $\tau=t$ si $n=q$ y
$\tau=2t$ si $q/2\le n<q$ (todos nuestros conjuntos cumplen una de las
dos cosas):

\begin{enumerate}[leftmargin=1.8em,itemsep=2pt]
  \item Se generan $\sigma_1\tau=16\tau$ bits uniformes
        $b_0,\dots,b_{16\tau-1}$.
  \item Cada bloque de $16$ bits produce un entero
        $d_j=\sum_{i=0}^{m-1}b_{16j+i}2^i$ para $j=0,\dots,\tau-1$
        (solo se usan los primeros $m$ bits de cada bloque).
  \item De la lista $d_0,\dots,d_{\tau-1}$ se toman las primeras $t$
        entradas \emph{en el rango} $\{0,\dots,n-1\}$. Si hay menos de
        $t$, se reinicia.
  \item Si esas $t$ entradas no son todas distintas, se reinicia.
  \item Se define $e$ como el vector de peso $t$ cuyo soporte son esas
        $t$ posiciones.
\end{enumerate}

\begin{proposicion}{$\textsf{FixedWeight}$ es exactamente uniforme}{fixedweight-uniform}
La salida $e$ de $\textsf{FixedWeight}$ (cuando no se reinicia) está
uniformemente distribuida sobre los $\binom{n}{t}$ vectores de peso
$t$. La probabilidad de que cada intento supere la prueba de
distinción es $\prod_{i=1}^{t-1}\bigl(1-\tfrac{i}{n}\bigr)$, mayor que
$1/4$ en todos los conjuntos; y la de que el lote aporte $t$ entradas
en rango supera $0{,}96$ en el peor caso,
\texttt{mceliece460896}~\cite{cmrationale}.
\end{proposicion}

\begin{proof}
Los $d_j$ son independientes y uniformes en $\{0,\dots,q-1\}$; la
pertenencia al rango $\{0,\dots,n-1\}$ es así independiente para cada
$j$, y condicionada a qué posiciones caen en rango, sus valores son
uniformes en el rango (el submuestreo de una sucesión uniforme por un
criterio independiente de los valores lo es). La primera vez que
aparecen $t$ entradas en rango no es distinto: se toma la subsucesión
de entradas en rango, que hereda la uniformidad. Finalmente, exigir
que las $t$ entradas sean distintas, por rechazo, deja una sucesión
ordenada uniforme de $t$ valores distintos, es decir, un soporte
uniforme entre los $\binom{n}{t}$ posibles. Las probabilidades son las
del argumento de cumpleaños, con las cotas de la razón de
diseño~\cite{cmrationale}.
\end{proof}

\begin{figure}[htbp]
\centering
\begin{tikzpicture}[
    sbox/.style={draw=mcblue,fill=mcsky!45,rounded corners=2.5pt,align=center,
      font=\footnotesize,inner sep=3.5pt,text width=2.9cm,minimum height=0.95cm},
    fbox/.style={draw=mcorange,fill=mcorange2!55,rounded corners=2.5pt,align=center,
      font=\footnotesize,inner sep=3.5pt,text width=3.3cm,minimum height=0.95cm},
    arr/.style={-{Stealth[length=2.2mm]},thick,mcgray}
  ]
  \node[sbox] (bits) at (1.4,0) {$16\tau$ bits\\ al azar};
  \node[sbox] (ints) at (5.3,0) {$\tau$ enteros $d_j$\\ de $m$ bits};
  \node[fbox] (filt) at (9.4,0) {descarta $d_j\ge n$;\\ toma los $t$ primeros distintos};
  \node[sbox] (e) at (13.3,0) {$e$: unos en esos\\ $t$ índices (peso $t$)};
  \draw[arr] (bits)--(ints);
  \draw[arr] (ints)--(filt);
  \draw[arr] (filt)--(e);
  \node[align=center,font=\tiny,color=mcgray] at (5.3,-1.15)
    {$\tau=t$ si $n=q$;\\ $\tau=2t$ si $q/2\le n<q$};
  \node[align=center,font=\tiny,color=mcred] at (9.4,-1.15)
    {si hay menos de $t$\\ o repetidos: reiniciar};
  \node[align=center,font=\tiny,color=mcgray] at (1.4,-1.15)
    {RNG (o semilla)};
\end{tikzpicture}
\caption{$\textsf{FixedWeight}$: de bits al azar a un soporte de
tamaño $t$ exacto. El lote de $\tau$ enteros abarata el muestreo y
facilita la vectorización; el rechazo garantiza uniformidad exacta
(proposición~\ref{prop:fixedweight-uniform}). Los valores $\tau$ son
$128,192,256,238,128$ para los cinco conjuntos de la
tabla~\ref{tab:keygen-pol}.}
\label{fig:fixedweight}
\end{figure}

\section{\textsf{Encode}: el criptograma es un síndrome}

Con $e$ en la mano, encapsular es casi gratis:
\[
  C=(\mathrm{I}_{mt}\,|\,T)\,e \in \Ftwo^{mt}.
  \label{eq:encap-c}
\]
Recordemos el diccionario del capítulo~\ref{ch:kem}: todo en columnas;
$T$ es la clave pública de $mt\times k$; escribir $e=(e_{\mathrm{alto}},
e_{\mathrm{bajo}})$ con $e_{\mathrm{alto}}\in\Ftwo^{mt}$ y
$e_{\mathrm{bajo}}\in\Ftwo^{k}$ permite ver que la multiplicación real
usa solo $T$:
\[
  C=e_{\mathrm{alto}}+T\,e_{\mathrm{bajo}}.
\]
Ese es todo el cifrado: $mt\cdot k$ bits de multiplicación
matriz--vector, unos $0{,}25$~MB de operaciones de XOR para el conjunto
más pequeño. La ecuación~\eqref{eq:encap-c} es también la
razón por la que $C$ no se puede «descifrar» sin la clave secreta:
encontrar $e$ de peso $t$ con $He=C$ es \emph{exactamente} el problema
duro de la parte~\ref{part:seguridad} (proposición~\ref{prop:nied-reduccion}).

\section{\textsf{Decode}: Patterson con el soporte secreto}

Descifrar es lo caro, y es donde vive la trampa de Goppa. La entrada es
$C$ y la clave secreta; $\textsf{Decode}$ procede así:

\begin{enumerate}[leftmargin=1.8em,itemsep=2pt]
  \item Se extiende $C$ con $k$ ceros: $v=(C,\mathbf{0},\dots,\mathbf{0})
        \in\Ftwo^{n}$.
  \item Se busca el único $c\in\Ftwo^{n}$ con (1)~$Hc=0$ y (2)~la distancia
        de Hamming entre $c$ y $v$ es $\le t$. Si no existe, se
        devuelve $\bot$.
  \item Se define $e=v+c$.
  \item Si $\wt(e)=t$ y $C=He$, se devuelve $e$; si no, $\bot$.
\end{enumerate}

El paso 2 es el descifrado de Patterson del capítulo~\ref{ch:goppa},
aplicado al código $\Gamma_0$ que definen $g$ y el ordenamiento
$\alpha'_0,\dots,\alpha'_{n-1}$ (la versión «ya permutada» que guarda
la clave secreta). El truco del relleno de ceros que ya usamos en el
capítulo~\ref{ch:niederreiter} es justo el que permite aplicarlo tal
cual: como $v+e$ es un codeword ($H(v+e)=C+C=0$), el síndrome de
Goppa de $v$ coincide con el de $e$, y Patterson trabaja directamente
sobre la información recibida.
Para un criptograma legítimo, $c=v+e$ es un codeword a distancia $t$
de $v$, y la unicidad del paso 2 se debe a la distancia mínima:
si $c_1\ne c_2$ estuvieran ambas a distancia $\le t$ de $v$, entonces
$\wt(c_1+c_2)\le2t<d$, imposible en un código de distancia $d$.

El paso 4 no es decorativo, y merece un párrafo. Si $C$ es un
criptograma legítimo, el chequeo pasa siempre (lo comprobamos en los
ejercicios). Pero $C$ también admite valores «tramposos»: por ejemplo,
$C=He'$ con $\wt(e')=t-1$. Entonces $v+e'$ \emph{sí} es un codeword,
Patterson lo encuentra, y el paso 2 devuelve una distancia menor que
$t$: exactamente lo que el chequeo rechaza. Sin él, el descifrador
aceptaría mensajes de peso distinto de $t$, la imagen de
$\textsf{Encode}$ dejaría de ser el conjunto de los vectores de peso
exacto ---la uniformidad que exige el capítulo~\ref{ch:kem}--- y la
«rigidez» que sustituye a la reencriptación dejaría de valer. Este
chequeo es, en miniatura, toda la transformación FO: verificación
exacta de que el criptograma es legítimo, y ningún fallo distinguible
(el rechazo implícito del paso final).

\section{Una vuelta completa}

Juntando todo, la especificación ejecuta esto:

\begin{center}
\small
\begin{tabular}{p{6.6cm}p{6.6cm}}
\toprule
$\textsf{Encap}(\mathrm{pk}=T)$ & $\textsf{Decap}(\mathrm{sk},C)$\\
\midrule
1. $e\leftarrow\textsf{FixedWeight()}$ &
  1. $b\leftarrow1$\\
2. $C\leftarrow(\mathrm{I}_{mt}|T)e$ &
  2. extraer $s$ y $\Gamma_0=(g,\alpha'_0,\dots,\alpha'_{n-1})$
    de la clave secreta\\
3. $K\leftarrow H(1,e,C)$ &
  3. $e'\leftarrow\textsf{Decode}(C,\Gamma_0)$;\newline
    si $e'=\bot$: $e'\leftarrow s$, $b\leftarrow0$\\
4. salida $(C,K)$ &
  4. $K\leftarrow H(b,e',C)$\\
 & 5. salida $K$\\
\bottomrule
\end{tabular}
\end{center}

\begin{figure}[htbp]
\centering
\begin{tikzpicture}[
    pbox/.style={draw=mcblue,fill=mcsky!45,rounded corners=2.5pt,align=center,
      font=\footnotesize,inner sep=3.5pt,text width=2.7cm,minimum height=0.95cm},
    fbox/.style={draw=mcorange,fill=mcorange2!55,rounded corners=2.5pt,align=center,
      font=\footnotesize,inner sep=3.5pt,text width=3.1cm,minimum height=0.95cm},
    gbox/.style={draw=mcgreen,fill=mcgreen2!60,rounded corners=2.5pt,align=center,
      font=\footnotesize,inner sep=3.5pt,text width=4.2cm,minimum height=0.9cm},
    rbox/.style={draw=mcred,fill=mcred2!60,rounded corners=2.5pt,align=center,
      font=\footnotesize,inner sep=3.5pt,text width=4.2cm,minimum height=0.9cm},
    arr/.style={-{Stealth[length=2.2mm]},thick,mcgray}
  ]
  \node[pbox] (fw) at (1.5,0) {$\textsf{FixedWeight}$\\ $e$ de peso $t$};
  \node[pbox] (en) at (5.5,0) {$\textsf{Encode}$\\ $C=(\mathrm{I}_{mt}|T)e$};
  \node[fbox] (de) at (9.7,0) {$\textsf{Decode}$ (Patterson)\\ con $g,\alpha'$: $e'$ o $\bot$};
  \node[fbox] (ch) at (13.4,0) {chequeo\\ $\wt(e')=t$, $He'=C$};
  \draw[arr] (fw)--(en) node[midway,above,font=\tiny]{$e$};
  \draw[arr] (en)--(de) node[midway,above,font=\tiny]{$C$ (canal)};
  \draw[arr] (de)--(ch);
  \node[align=center,font=\tiny,color=mcgray] at (1.5,-0.95) {azar (RNG)};
  \node[align=center,font=\tiny,color=mcgray] at (5.5,-0.95) {clave pública $T$};
  \node[align=center,font=\tiny,color=mcgray] at (9.7,-0.95) {clave secreta: $g,\alpha'$};
  \node[gbox] (kg) at (11.4,-2.1) {si es válido: $K=H(1,e',C)$};
  \node[rbox] (kr) at (11.4,-3.7) {si $\bot$ o no valida: $K=H(0,s,C)$};
  \draw[arr] (ch.south) -- (kg.north);
  \draw[arr,mcred] (ch.east) -- (15.2,0) -- (15.2,-3.7) -- (kr.east);
\end{tikzpicture}
\caption{El recorrido completo. En azul, lo que hace
$\textsf{Encap}$ con la clave pública; en naranja, lo que hace
$\textsf{Decap}$ con la secreta; en verde y rojo, las dos formas de
terminar. La especificación unifica el hash final como
$K=H(b,v,C)$: con $b=1$ y $v=e'$ si el descifrado valida, y con $b=0$
y $v=s$ (los $n$ bits secretos) si no.}
\label{fig:encapdecap}
\end{figure}

Sobre el hash, la especificación es precisa:
\[
  H(b,v,C)=\text{primeros }256\text{ bits de}\;\,
  \textsf{SHAKE256}\bigl(\mathrm{byte}(b)\,\|\,\mathrm{bytes}(v)\,\|\,\mathrm{bytes}(C)\bigr),
\]
con $v$ codificado en $\lceil n/8\rceil$ bytes y $C$
en $\lceil mt/8\rceil$ bytes. El byte inicial $0$ o $1$ es la
separación de dominio del capítulo~\ref{ch:kem}: las claves legítimas y
las de rechazo viven en «mundos» distintos del hash, y ningún
criptograma puede hacer colisionar ambos. Para
\texttt{mceliece348864}, la entrada del hash ocupa
$1+\lceil3488/8\rceil+\lceil768/8\rceil=533$ bytes; su salida, $32$.

Dos detalles finos que el lector encontrará en los archivos de prueba
de la especificación: (1)~$b$ es un bit, pero se escribe como byte
entero; (2)~en los conjuntos donde $n$ o $mt$ no son múltiplos de $8$
(\texttt{mceliece6960119}: $mt=1547$), los bits sobrantes del último
byte son \emph{relleno}. La especificación define dos modos de lectura:
\emph{Simply Decoded} (ignora el relleno) y \emph{Narrowly Decoded}
(rechaza criptogramas con relleno distinto de cero). Volveremos a esta
distinción ---que es una forma benigna de maleabilidad--- en el
capítulo~\ref{ch:parametros}.

\section{¿Cuánto cuesta?}

Los costos, en orden de magnitud y por operación, son los siguientes.

\begin{itemize}[leftmargin=1.5em,itemsep=3pt]
  \item $\textsf{Encap}$: un muestreo de peso $t$ y una multiplicación
        binaria $mt\times k$: para \texttt{mceliece348864}, unos
        $261120$ bytes de lectura y XOR con un acumulador de $96$
        bytes; para \texttt{mceliece8192128}, unos $1{,}3$~MB. Es
        comparable a una exponenciación modular pequeña, pero con
        operaciones de bits: decenas de microsegundos en hardware
        moderno.
  \item $\textsf{Decap}$: el síndrome en $\GF q$ cuesta del orden de
        $tn$ operaciones en el cuerpo (unos $223\,000$ productos para
        \texttt{mceliece348864}); Patterson añade $O(t^2)$ y la
        búsqueda de raíces un coste similar. Sigue siendo, con
        diferencia, la operación más lenta ---milisegundos o menos en
        implementaciones optimizadas---, pero muy lejos de la
        decodificación genérica de la parte~\ref{part:seguridad}: la
        trampa hace su trabajo~\cite{cmimpl}.
  \item Memoria: $\textsf{Decap}$ toca la clave secreta y el
        criptograma; nada más. La semilla $\delta$ no se usa en línea
        (solo para regenerar la clave), y $T$ no se usa para descifrar: la
        eliminación de la clave pública que vimos en el
        capítulo~\ref{ch:kem}.
\end{itemize}

\begin{observacion}{Patterson no es opcional}{patterson-no-opcional}
El descifrado de un código de Goppa \emph{no} es «buscar el codeword
más cercano» de forma exhaustiva: es un algoritmo algebraico
(Patterson, 1975) que usa el polinomio $g$ y el soporte para resolver
el problema con álgebra de polinomios. Ese algoritmo es la trampa, y
por eso su descripción (la clave secreta) es lo único que no puede
salir de la caja. La distancia mínima garantiza además que, cuando el
descifrado existe, es único: la trampa no es solo eficiente, es
inequívoca.
\end{observacion}

\section{Ejercicios}

\begin{enumerate}[leftmargin=1.6em,itemsep=4pt]
  \item \textbf{La ruta de $v$ (laboratorio).} En el juguete del
        capítulo~\ref{ch:labgoppa} ($n=16$, $k=8$, $t=2$, soporte
        $\{10,13\}$), escriba $v=(C,\mathbf{0})$ y compruebe con las
        matrices del laboratorio que $c=v+e$ satisface $Hc=0$ y que su
        distancia a $v$ es $t$. Repita el cálculo para el síndrome
        tramposo $C'=He'$ con $\wt(e')=1$ y describa qué devolvería
        cada paso de $\textsf{Decode}$.
  \item \textbf{Uniformidad de $\textsf{FixedWeight}$.} Repita la
        demostración de la proposición~\ref{prop:fixedweight-uniform}
        señalando (a)~dónde se usa que $16\ge m$;
        (b)~dónde se usa la independencia; (c)~por qué el rechazo por
        repeticiones no sesga el resultado. Calcule
        $\prod_{i=1}^{t-1}(1-i/n)$ para el juguete ($n=16$, $t=2$) y
        para \texttt{mceliece348864} (use $e^{-t(t-1)/(2n)}$ como
        aproximación).
  \item \textbf{Unicidad del candidato.} Demuestre con detalle que en
        el paso 2 no puede haber dos codewords distintos $c_1,c_2$
        con distancia $\le t$ a $v$. (Pista:
        $\wt(c_1+c_2)\le2t$ y la distancia mínima del código de Goppa
        es $\ge2t+1$.)
  \item \textbf{El chequeo imprescindible.} Con el juguete, fabrique
        $C'=He'$ con $\wt(e')=1$ y explique qué devolverían el paso 2
        y el paso 4. ¿Qué pasaría con la uniformidad de los mensajes
        aceptados si suprimiéramos el chequeo? ¿Cómo se relaciona esto
        con la rigidez del capítulo~\ref{ch:kem}?
  \item \textbf{Cuentas de costos.} Para \texttt{mceliece8192128}
        ($m=13$, $n=8192$, $t=128$, $k=6528$): calcule los bits que
        toca $\textsf{Encode}$ ($mt\cdot k$) y las operaciones de
        cuerpo de un síndrome de $\textsf{Decode}$ ($tn$). ¿Cuál de
        las dos operaciones domina, y por qué eso es exactamente lo
        que busca un KEM (cifrar barato, descifrar con clave)?
  \item \textbf{La entrada del hash.} Para \texttt{mceliece6960119}:
        calcule los bytes de la entrada $H(b,v,C)$ y explique para qué
        sirve cada pieza (etiqueta, vector, criptograma). ¿Por qué el
        hash no recibe la lista de posiciones del soporte de $e$ en
        vez del vector entero?
\end{enumerate}


\chapter{Laboratorio 2: el KEM completo}\label{ch:labkem}

\begin{roadmap}
El juguete del capítulo~\ref{ch:labgoppa} decodificaba; los
capítulos~\ref{ch:kem} a~\ref{ch:encapdecap} definieron el KEM. Este
laboratorio enchufa las dos cosas: genera una clave del juguete desde
una semilla (como el capítulo~\ref{ch:keygen}), encapsula, desencapsula
y comprueba la igualdad exacta de las claves de $32$ bytes. Después
rompe el descifrado a propósito, cambiando un solo bit del criptograma,
y comprueba que el rechazo implícito no deja rastro. Todos los números
de este capítulo se reproducen ejecutando el programa del
apéndice~\ref{app:codigo}.
\end{roadmap}

\begin{practica}[Los vectores de prueba del NIST son esto]
La especificación de Classic McEliece distribuye archivos
\emph{known-answer test} (KAT): para cada conjunto y cada semilla
$\delta$, dan la clave pública, la privada, el criptograma y la clave
de sesión esperados. La razón de ser es exactamente la de este
laboratorio: que cualquier implementación independiente reproduzca
\emph{bit a bit} los mismos resultados. Nuestra tabla del juguete es un
KAT en miniatura, escrito a mano.
\end{practica}

\section{La clave del juguete (KeyGen en miniatura)}

Fijamos, como el generador de claves, una semilla
$\delta=\hx{5a5a...5a}$ (el byte $\hx{5a}$ repetido $32$ veces) y
aplicamos la expansión del capítulo~\ref{ch:keygen}:

\begin{itemize}[leftmargin=1.5em,itemsep=2pt]
  \item $G=\textsf{SHAKE256}(64\|\delta)$. Sus primeros $n=16$ bits son
        la clave de rechazo:
        \[
          s=\hx{0111 0101 1101 0011}\quad(\text{bit }0\text{ a la izquierda}),
        \]
        esto es, $s=\hx{CBAE}$ como entero, con bytes
        little-endian $\hx{AE\,CB}$.
  \item El «\textsf{FieldOrdering}» del juguete es el orden natural:
        $\alpha_j$ es el elemento de $\FF4$ cuyo entero vale $j$, para
        $j=0,\dots,15$ (nótese que $s$, en cambio, es una cadena de
        bits \emph{uniforme}, no un vector de peso $t$: exactamente lo
        que ordena la especificación).
  \item El polinomio de Goppa es $g(x)=x^2+x+z^3$, el irreducible del
        capítulo~\ref{ch:labgoppa} (comprobado: $g(\alpha)\ne0$ para
        todo $\alpha\in\FF4$).
  \item \textsf{MatGen} produce, como ya vimos,
        $(\mathrm{I}_8\,|\,T)$ con la $T$ de $8\times8$ bits del
        capítulo~\ref{ch:labgoppa}.
\end{itemize}

Así, la clave pública del juguete es $T$ (¡$8$ bytes!) y la secreta es
$(\delta,g,\alpha,s)$. Todo lo demás del KEM son cuentas.

\section{Encap: de la semilla al criptograma}

\textsf{FixedWeight} sortea el error
$e=\hx{0000000000100100}$ (posiciones $10$ y $13$, como en todas las
cuentas del juguete). Para cifrar, separamos $e$ en dos mitades de $8$
bits:

\[
  e_{\text{alto}}=\hx{00000000},\qquad
  e_{\text{bajo}}=\hx{00100100},
\]

y calculamos $C=e_{\text{alto}}+T\,e_{\text{bajo}}$. La mitad alta es
cero, así que basta combinar las columnas $2$ y $5$ de $T$, que son las
que selecciona $e_{\text{bajo}}$:

\[
  C=\mathrm{col}_2(T)\oplus\mathrm{col}_5(T)
  =\hx{36}\oplus\hx{AE}=\hx{98}
  \quad\Longrightarrow\quad
  C=\hx{00011001}.
\]

El hash final recibe la etiqueta y las codificaciones en bytes
little-endian de $e$ y $C$: $e=9216=\hx{2400}$ aporta los bytes
$\hx{00\,24}$, y $C$ aporta $\hx{98}$:

\[
  \begin{gathered}
    K=H(1,e,C)=\textsf{SHAKE256}(\hx{01}\,\|\,\hx{00\,24}\,\|\,\hx{98}),\\
    K=\hx{597f46fbba325c68f0b0b45c6ea984961461eb220d1f24024ed888912b336278}.
  \end{gathered}
\]

El criptograma es, entonces, $C=\hx{98}$ (un solo byte) y la clave de
sesión, $32$ bytes.

\section{Decap: la vuelta, paso a paso}

Del otro lado, con la clave secreta, el descifrado repite el ejercicio
del capítulo~\ref{ch:labgoppa} pero sobre
$v=(C,\mathbf{0})=\hx{0001100100000000}$ (el relleno de ceros del que
hablaba el capítulo~\ref{ch:encapdecap}). Patterson calcula el
localizador:

\[
  S(x)=\hx4+\hx9x,\qquad \sigma(x)=\hx5+\hx e x+\hx2x^2,
  \qquad \text{raíces: }\{10,13\},
\]

las mismas cuentas que la traza del capítulo~\ref{ch:labgoppa}
(§~«Patterson, paso a paso»), como debe ser: el síndrome de $v=e+c$,
con $c\in\Gamma$, es el de $e$, porque para todo codeword
$\sum_j 1/(x+\alpha_j)\equiv0\pmod g$. Encontrado
$e'=\hx{0000000000100100}=e$, el chequeo final pasa
($\wt(e')=2$ y $(\mathrm{I}_8|T)e'=C$) y la clave es la misma:

\[
  K'=H(1,e',C)=\textsf{SHAKE256}(\hx{01}\|\hx{00\,24}\|\hx{98})=K.
\]

La figura~\ref{fig:labkem-pipeline} resume la vuelta completa con todos
los valores a la vista.

\begin{figure}[htbp]
\centering
\begin{tikzpicture}[
    pbox/.style={draw=mcblue,fill=mcsky!45,rounded corners=2.5pt,align=center,
      font=\footnotesize,inner sep=3.5pt,text width=5.0cm,minimum height=0.9cm},
    fbox/.style={draw=mcorange,fill=mcorange2!55,rounded corners=2.5pt,align=center,
      font=\footnotesize,inner sep=3.5pt,text width=5.6cm,minimum height=0.9cm},
    gbox/.style={draw=mcgreen,fill=mcgreen2!60,rounded corners=2.5pt,align=center,
      font=\footnotesize,inner sep=3.5pt,text width=5.6cm,minimum height=0.9cm},
    rbox/.style={draw=mcred,fill=mcred2!60,rounded corners=2.5pt,align=center,
      font=\footnotesize,inner sep=3.5pt,text width=11.4cm,minimum height=0.9cm},
    arr/.style={-{Stealth[length=2.2mm]},thick,mcgray}
  ]
  \node[font=\small\bfseries,color=mcblue] at (2.0,1.0) {Encap$(pk=T)$};
  \node[font=\small\bfseries,color=mcorange] at (9.2,1.0) {Decap$(sk,C)$};
  % columna izquierda
  \node[pbox] (fw) at (2.0,0) {$\textsf{FixedWeight}$: $e\mapsto\{10,13\}$.};
  \node[pbox] (en) at (2.0,-1.7)
    {$\textsf{Encode}$: $C=\mathrm{col}_2(T)\oplus\mathrm{col}_5(T)
     =\hx{36}\oplus\hx{AE}=\hx{98}$;\\ $C=\hx{00011001}$.};
  \node[gbox] (k) at (2.0,-3.4)
    {hash: $H(1,e,C)$ con bytes $\hx{01}\,\hx{00\,24}\,\hx{98}$;\\
     $K=\hx{597f46\ldots6278}$.};
  \draw[arr] (fw)--(en);
  \draw[arr] (en)--(k);
  % columna derecha
  \node[fbox] (v) at (9.2,-1.7)
    {llega $C$; se rellena: $v=(C,\mathbf{0})=\hx{0001100100000000}$.};
  \node[fbox] (pa) at (9.2,-3.4)
    {Patterson: $S(x)=\hx4+\hx9x$; $\sigma(x)=\hx5+\hx e x+\hx2x^2$;\\
     raíces $\{10,13\}$.};
  \node[fbox] (ch) at (9.2,-5.1)
    {chequeo: $\wt(e')=2$ y $(\mathrm{I}_8|T)e'=C$: pasa.};
  \node[gbox] (kp) at (9.2,-6.8)
    {$K'=H(1,e',C)=\hx{597f46\ldots6278}$.};
  \draw[arr] (v)--(pa);
  \draw[arr] (pa)--(ch);
  \draw[arr] (ch)--(kp);
  % canal C
  \draw[arr] (en.east) -- (v.west) node[midway,above,font=\tiny]{$C$ ($1$ byte)};
  % igualdad de claves
  \draw[mcgreen,dashed,thick,{Stealth[length=2.2mm]}-{Stealth[length=2.2mm]}]
    (k.east) -- (kp.west) node[midway,sloped,above,font=\tiny,color=mcgreen]{$=$};
  % rechazo implicito
  \node[rbox] (r) at (5.6,-8.6)
    {con $\tilde{C}=\hx{10011001}$ (un bit cambiado): $\tilde{\sigma}$ no tiene raíces
     $\Rightarrow\bot$; \textsf{Decap} devuelve\\ $K''=H(0,s,\tilde{C})$ con bytes
     $\hx{00}\,\hx{AE\,CB}\,\hx{99}$, es decir,
     $K''=\hx{ba1a72\ldots1494}$ (rechazo implícito).};
\end{tikzpicture}
\caption{El KEM del juguete, con todos los números. A la izquierda,
\textsf{Encap}: del error al criptograma y a la clave. A la derecha,
\textsf{Decap}: del criptograma al error y a la misma clave (el signo
$=$ verde). Abajo, el caso manipulado: un bit distinto y una clave
completamente distinta, indistinguible de una legítima.}
\label{fig:labkem-pipeline}
\end{figure}

\section{Un criptograma manipulado}

Cambiemos un solo bit del criptograma: $\tilde{C}=\hx{10011001}$, el
resultado de $C\oplus\hx{01}$. La especificación no responde con un
error: intenta descifrar y, al fallar, recurre a la clave de rechazo.
Los números son los siguientes:

\begin{itemize}[leftmargin=1.5em,itemsep=2pt]
  \item El síndrome cambia: $\tilde{S}(x)=\hx b+\hx6x$, y Patterson llega a
        $\tilde{\sigma}(x)=\hx5+\hx8x+\hx a x^2$, que \emph{no tiene raíces} en
        $\FF4$: no hay error de peso $\le2$ que explique $\tilde{C}$.
  \item Con el fallo, \textsf{Decap} pone $b\leftarrow0$ y
        $e\leftarrow s$, y devuelve
        \[
          \begin{gathered}
            K''=H(0,s,\tilde{C})=\textsf{SHAKE256}(\hx{00}\,\|\,\hx{AE\,CB}\,\|\,\hx{99}),\\
            K''=\hx{ba1a72ac26b0fcf20c55c108bb681e27aa9c270c60802a821601130b7f1e1494}.
          \end{gathered}
        \]
  \item Nadie distingue $K''$ de una clave legítima: es la salida de un
        hash sobre entradas distintas, de $32$ bytes como todas. Ese es
        el rechazo implícito del capítulo~\ref{ch:kem} en acción.
\end{itemize}

Hay una estadística del juguete que conviene mirar de cerca: de los
$2^8=256$ síndromes posibles, solo $137$ son alcanzables por errores de
peso $\le2$:
\[
  1+16+120=137
  \qquad(\text{peso }0,1,2).
\]
Patterson resuelve esos $137$, pero \textsf{Decap} solo acepta los
$120$ de peso exactamente $2$: los $17$ de peso $0$ o $1$ se rechazan
por el chequeo de peso, igual que los $119$ síndromes inalcanzables.
Es decir, un criptograma manipulado al azar se rechaza con
probabilidad $136/256\approx53\%$: más de la mitad de las veces. Sin
rechazo implícito, esa mayoría de fallos sería una señal detectable;
con él, \textsf{Decap} es una función total que siempre devuelve $32$
bytes.

\section{Del juguete al estándar (y qué no prueba el juguete)}

La tabla~\ref{tab:labkem-toy-real} pone los dos mundos al lado.

\begin{table}[htbp]
\centering
\small
\begin{tabular}{lcc}
\toprule
 & juguete & \texttt{mceliece348864}\\
\midrule
$m$ (bits por elemento) & 4 & 12\\
$n$ (longitud) & 16 & 3488\\
$t$ (peso del error) & 2 & 64\\
$k$ (dimensión) & 8 & 2720\\
clave pública $T$ & $8$ bytes & $261\,120$ bytes ($255$~KiB)\\
criptograma $C$ & $1$ byte & $96$ bytes\\
semilla $\delta$ & $32$ bytes & $32$ bytes\\
clave de sesión $K$ & $32$ bytes & $32$ bytes\\
seguridad de la decodificación & nula ($2^{16}$ palabras) & $\approx141$ bits
  (capítulo~\ref{ch:seguridad})\\
\bottomrule
\end{tabular}
\caption{El juguete y el primer conjunto real. Los elementos
criptográficos «gruesos» ---semilla y clave de sesión--- ya tienen el
tamaño definitivo; lo que crece es la matemática: código, matriz
pública y criptograma.}\label{tab:labkem-toy-real}
\end{table}

El juguete demuestra que la \emph{tubería} funciona: las mismas
etiquetas, el mismo relleno, el mismo rechazo implícito, las mismas
cuentas de bytes. Lo que no demuestra ---ni puede--- es seguridad: con
$n=16$, cualquiera decodifica por fuerza bruta, y el problema «duro» no
es duro. Esa parte vive en la
parte~\ref{part:seguridad}, y es la única que no se puede simular en
pequeño.

\section{Ejercicios}

\begin{enumerate}[leftmargin=1.6em,itemsep=4pt]
  \item \textbf{Reproducir la clave.} Ejecute el programa del
        apéndice~\ref{app:codigo} (o reimplemente sus tres pasos) y
        compruebe que $\textsf{SHAKE256}(\hx{01}\|\hx{00\,24}\|\hx{98})$
        produce exactamente la $K$ de la
        figura~\ref{fig:labkem-pipeline}. Compruebe también
        $s=\hx{CBAE}$ y sus bytes.
  \item \textbf{Bits y bytes.} Explique por qué el entero
        $e=9216=\hx{2400}$ se codifica como $\hx{00\,24}$ (little-endian)
        y cómo se concilia con la convención «posición $0$ a la
        izquierda». Repita el razonamiento para $s$.
  \item \textbf{El síndrome que no cambia.} Demuestre que, si
        $v=(C,\mathbf{0})$ y $C=He$, entonces el síndrome polinomial de
        $v$ coincide con el de $e$. (Pista: recuerde que
        $\sigma_j(x)=1/(x+\alpha_j)$ cumple
        $\sum_jc_j\sigma_j\equiv0\pmod g$ para todo codeword $c$.)
        ¿Por qué este hecho \emph{no} se cumple para $\tilde{C}$?
  \item \textbf{¿Por qué no devolver $\bot$?} Suponga que
        \textsf{Decap} respondiera «error» cuando el descifrado falla.
        Describa cómo un atacante usaría el bit manipulado $\tilde{C}$ para
        distinguir respuestas, y qué papel juegan en ello los $119$
        síndromes inalcanzables de la sección anterior.
  \item \textbf{Contar síndromes.} Verifique $1+16+120=137$ y
        $1+16+120+\binom{16}{3}+\cdots$ (¿por qué los errores de peso
        $3$ ya \emph{no} están garantizados como únicos?). Relacione
        $137/256\approx53{,}5\%$ con la fracción de bytes que Patterson
        resuelve y calcule la probabilidad de que \textsf{Decap}
        rechace un byte manipulado al azar (pista: solo los errores de
        peso exactamente $2$ pasan el chequeo).
  \item \textbf{Escalar.} Con la tabla~\ref{tab:labkem-toy-real}:
        calcule cuántas veces más grande es la clave pública real que la
        del juguete, y el mismo cociente para el criptograma. Enumere
        tres mecanismos del esquema real que el juguete omite (pistas:
        $\sigma_1$, $\sigma_2$, forma semi-sistemática, red de Benes)
        e indique en qué capítulo se estudia cada uno.
\end{enumerate}

\chapter{Los conjuntos y las variantes}\label{ch:parametros}

\begin{roadmap}
El capítulo~\ref{ch:keygen} mostró cómo nace una clave; los
capítulos~\ref{ch:encapdecap} y~\ref{ch:labkem}, cómo se encapsula y
desencapsula. Falta el catálogo: qué conjuntos concretos define Classic
McEliece, qué significan los sufijos \texttt{f}, \texttt{pc} y
\texttt{pcf}, cómo se guardan claves y criptogramas byte a byte, y qué
seguridad se les atribuye. Al terminar, el lector debería poder elegir
un conjunto con criterio y saber cuántos bytes maneja en cada paso.
\end{roadmap}

\section{Los cinco tamaños}\label{sec:param-cinco}

Las tablas~\ref{tab:keygen-pol} y~\ref{tab:keygen-tamanos} fijaron los
cinco pares $(m,n,t)$ y el desglose interno de las claves. La
tabla~\ref{tab:param-base} añade las dos cifras que faltaban para
decidir: el criptograma y la categoría de seguridad.

\begin{table}[htbp]
\centering
\small
\setlength{\tabcolsep}{7pt}
\begin{tabular}{lcrrrc}
\toprule
Conjunto & categoría & $ek$ (bytes) & $ct$ (bytes) & $sk$ (bytes) & ¿ISO?\\
\midrule
\texttt{mceliece348864} & 1 & \num{261120}  & 96  & \num{6492}  & no\\
\texttt{mceliece460896} & 3 & \num{524160}  & 156 & \num{13608} & sí\\
\texttt{mceliece6688128} & 5 & \num{1044992} & 208 & \num{13932} & sí\\
\texttt{mceliece6960119} & 5 & \num{1047319} & 194 & \num{13948} & sí\\
\texttt{mceliece8192128} & 5 & \num{1357824} & 208 & \num{14120} & sí\\
\bottomrule
\end{tabular}
\caption{Los cinco conjuntos base, con
$ek=mt\lceil k/8\rceil$, $ct=\lceil mt/8\rceil$ y $sk$ como en la
tabla~\ref{tab:keygen-tamanos}. La columna «¿ISO?» anticipa la
sección siguiente: el estándar ISO incluye los cuatro tamaños
mayores, cada uno con cuatro variantes; \texttt{mceliece348864} queda
fuera.}\label{tab:param-base}
\end{table}

¿De dónde salen estos números, a primera vista caprichosos? Los
conjuntos se eligieron siempre \emph{maximizando la seguridad sujeta a
un límite de tamaño} \cite{cmrationale}:

\begin{itemize}[leftmargin=1.5em,itemsep=3pt]
  \item \texttt{mceliece6960119}, con $(n,t)=(6960,119)$, maximiza la
        seguridad para claves públicas que quepan en $2^{20}$ bytes.
        La elección de $n=6960$ viene de la optimización clásica de
        2008: para límites de $2^{16}$ a $2^{20}$ bytes, los códigos
        Goppa de longitudes $1744$, $2480$, $3408$, $4624$ y $6960$,
        con $t$ proporcionalmente creciente, maximizan la dificultad
        del ataque; la tasa óptima ronda $0{,}75$. La única diferencia
        con aquella elección es que Classic McEliece añade $119$
        errores en lugar de $121$: sin decodificación por listas no
        puede decodificar de forma única más allá de la capacidad del
        código, y añadir dos errores menos reduce marginalmente la
        seguridad.
  \item \texttt{mceliece6688128} también maximiza dentro de $2^{20}$
        bytes, pero exigiendo que $n$ y $t$ sean múltiplos de $32$;
        \texttt{mceliece460896} y \texttt{mceliece348864} hacen lo
        mismo dentro de $2^{19}$ y $2^{18}$ bytes.
  \item \texttt{mceliece8192128} es el conjunto «clásico»: $n$ y $t$
        potencias de dos, como en el McEliece original de 1978, a
        costa de pasarse un poco de $2^{20}$ bytes
        ($\num{1357824}$ bytes $=1{,}29$~MiB).
\end{itemize}

Maximizar bajo un límite es más robusto que apuntar a un nivel exacto:
las métricas de seguridad varían entre modelos, y un conjunto clavado
a un objetivo puede quedar por debajo en otro modelo distinto. Al
maximizar por tamaño, la incertidumbre cambia poco el orden.

\begin{observacion}{Múltiplos de $32$: por qué}{param-mult32}
Cuando $n$ y $t$ son múltiplos de $32$ ---y por tanto $mt$ y $k$
también lo son, porque $k=n-mt$ y $q=2^{13}$--- no hay \emph{bits de
relleno} en las codificaciones de claves y criptogramas, y la
distinción entre las dos formas de leer un criptograma
(sección~\ref{sec:param-narrowly}) desaparece. Es una pequeña ventaja
de interoperabilidad a cambio de un pequeño empeoramiento en la
relación tamaño-seguridad. Los conjuntos $3488/4608/6688$ y el
$8192$ la tienen; el $6960$, no.
\end{observacion}

\section{Los dieciséis conjuntos de ISO}\label{sec:param-iso}

ISO estandarizó Classic McEliece en junio de 2026 (enmienda 2 de
ISO/IEC 18033-2:2006 \cite{iso18033}), y el estándar incluye
\textbf{dieciséis conjuntos}: los cuatro tamaños de la
tabla~\ref{tab:param-base} (sin \texttt{mceliece348864}) cruzados con
cuatro variantes:

\begin{itemize}[leftmargin=1.5em,itemsep=2pt]
  \item el conjunto \textbf{base}, tal cual lo define el equipo;
  \item el sufijo \textbf{\texttt{f}}: forma semi-sistemática
        $(\mu,\nu)=(32{,}64)$ (sección~\ref{sec:param-f});
  \item el sufijo \textbf{\texttt{pc}}: confirmación de texto claro
        (sección~\ref{sec:param-pc});
  \item el sufijo \textbf{\texttt{pcf}}: ambas cosas a la vez.
\end{itemize}

La tabla~\ref{tab:param-iso} y la figura~\ref{fig:param-iso} muestran
el catálogo completo. Obsérvese lo que cada sufijo hace a los tamaños:
\texttt{f} no cambia nada; \texttt{pc} suma exactamente
$\lceil\ell/8\rceil=32$ bytes al criptograma.

\begin{table}[htbp]
\centering
\small
\setlength{\tabcolsep}{6pt}
\begin{tabular}{lcccccc}
\toprule
tamaño & $(m,n,t,k)$ & $ek$ & $sk$ & $ct$ \texttt{base},\texttt{f} & $ct$ \texttt{pc},\texttt{pcf} & cat.\\
\midrule
\texttt{460896}  & $(13,4608,96,3360)$  & \num{524160}  & \num{13608} & 156 & 188 & 3\\
\texttt{6688128} & $(13,6688,128,5024)$ & \num{1044992} & \num{13932} & 208 & 240 & 5\\
\texttt{6960119} & $(13,6960,119,5413)$ & \num{1047319} & \num{13948} & 194 & 226 & 5\\
\texttt{8192128} & $(13,8192,128,6528)$ & \num{1357824} & \num{14120} & 208 & 240 & 5\\
\bottomrule
\end{tabular}
\caption[Los dieciséis conjuntos de ISO, cada fila con sus cuatro variantes.]{Los dieciséis conjuntos de ISO: cada fila, con los sufijos
---, \texttt{f}, \texttt{pc}, \texttt{pcf}. El nombre completo es
\texttt{mceliece} seguido del número y el sufijo, por ejemplo
\texttt{mceliece6960119pcf}. Los tamaños de cada fila valen para las
cuatro variantes salvo el criptograma de \texttt{pc}/\texttt{pcf}.}
\label{tab:param-iso}
\end{table}

\begin{figure}[htbp]
\centering
\begin{tikzpicture}[
    chip/.style={draw=mcblue,fill=mcsky!45,rounded corners=3pt,align=center,
      font=\footnotesize,inner sep=4pt,minimum width=3.5cm,minimum height=1.05cm},
    pchip/.style={draw=mcorange,fill=mcorange2!55,rounded corners=3pt,align=center,
      font=\footnotesize,inner sep=4pt,minimum width=3.5cm,minimum height=1.05cm},
    head/.style={font=\small\bfseries}
  ]
  \node[head,color=mcblue] at (2.0,0.85) {base};
  \node[head,color=mcblue] at (6.0,0.85) {\texttt{f}};
  \node[head,color=mcorange] at (10.0,0.85) {\texttt{pc}};
  \node[head,color=mcorange] at (14.0,0.85) {\texttt{pcf}};
  \foreach \y/\nm/\ca/\cb in {0/mceliece460896/156/188,-1.4/mceliece6688128/208/240,-2.8/mceliece6960119/194/226,-4.2/mceliece8192128/208/240} {
    \node[chip] at (2.0,\y) {\texttt{\nm}\\[1pt] $ct$: \ca\ bytes};
    \node[chip] at (6.0,\y) {\texttt{\nm}\texttt{f}\\[1pt] $ct$: \ca\ bytes};
    \node[pchip] at (10.0,\y) {\texttt{\nm}\texttt{pc}\\[1pt] $ct$: \cb\ bytes};
    \node[pchip] at (14.0,\y) {\texttt{\nm}\texttt{pcf}\\[1pt] $ct$: \cb\ bytes};
  }
  \draw[decorate,decoration={brace,amplitude=5pt},thick,mcorange]
    (8.2,-4.95) -- (15.8,-4.95);
  \node[font=\scriptsize,color=mcorange,align=center] at (12.0,-5.55)
    {$+\lceil\ell/8\rceil=32$ bytes de confirmación};
  \draw[decorate,decoration={brace,mirror,amplitude=5pt},thick,mcblue]
    (4.2,-4.95) -- (7.8,-4.95);
  \node[font=\scriptsize,color=mcblue,align=center] at (6.0,-5.55)
    {mismos tamaños:\\ otra generación de clave};
\end{tikzpicture}
\caption{Los dieciséis conjuntos de ISO, como familia. Los cuatro
tamaños (filas) y las cuatro variantes (columnas). La confirmación de
texto claro (\texttt{pc}, \texttt{pcf}) añade $32$ bytes al
criptograma; la variante \texttt{f} no cambia ningún tamaño, solo la
manera de generar la clave (forma semi-sistemática).}
\label{fig:param-iso}
\end{figure}

¿Por qué falta \texttt{mceliece348864}? Por el modelo cuántico de ISO.
ISO pidió a los candidatos \textbf{128 bits de seguridad en su modelo
cuántico}, que asume la aceleración de raíz cuadrada de la búsqueda de
Grover; pedir $128$ bits en ese modelo es mucho más exigente que pedir
$128$ bits clásicos \cite{cmiso}. Con la raíz cuadrada aplicada a los
mejores ataques clásicos sin trucos de memoria:

\begin{center}
\small
\begin{tabular}{lcccc}
\toprule
conjunto & ataque clásico & mitad (modelo cuántico) & ¿$\ge2^{128}$?\\
\midrule
\texttt{mceliece348864} & $2^{151{,}4}$ & $2^{75{,}7}$ & no\\
\texttt{mceliece460896} & $2^{193{,}3}$ & $2^{96{,}7}$ & no\\
\texttt{mceliece6688128} & $2^{268{,}0}$ & $2^{134{,}0}$ & sí\\
\texttt{mceliece6960119} & $2^{268{,}4}$ & $2^{134{,}2}$ & sí\\
\texttt{mceliece8192128} & $2^{303{,}2}$ & $2^{151{,}6}$ & sí\\
\bottomrule
\end{tabular}
\end{center}

Los valores clásicos son los de Stern bajo memoria gratis
(tabla~\ref{tab:param-seguridad}), cuyos estimadores cuentan
operaciones vectoriales como una sola operación, de modo que son cotas
prudentes. Con esa vara, el borrador ISO que el equipo publicó en
2023 incluyó solo los tamaños $6688$, $6960$ y $8192$ ---todos
$\ge2^{134}$ cuánticos---. Pero la aceleración de Grover está, en la
práctica, limitada por la latencia del ataque, así que ISO decidió
estandarizar también \texttt{mceliece4*} (el conjunto $460896$) a
sabiendas de que no alcanza el listón cuántico
\cite{cmiso}. El equipo recomienda los tamaños
\texttt{mceliece6*} para seguridad a largo plazo.

\begin{historia}[Cómo llegó Classic McEliece a ISO]
En octubre de 2022, unas diapositivas de Microsoft revelaron que el
subcomité ISO/IEC JTC1/SC27/WG2 había solicitado comentarios sobre la
inclusión de FrodoKEM, Classic McEliece y CRYSTALS-Kyber en sus
estándares. El equipo de Classic McEliece preparó como borrador el
documento \texttt{iso-mceliece-20230419} y lo publicó para recibir
comentarios. Siguiendo la regla ISO de expresar los requisitos «en
términos de rendimiento, no de diseño o características descriptivas»,
el borrador especificaba las \emph{funciones matemáticas} a computar,
sin constreñir cómo se computan. ISO estandarizó Classic McEliece en
junio de 2026, junto a ML-KEM y FrodoKEM, como enmienda a
ISO/IEC 18033-2 \cite{cmiso}. Mientras tanto, el NIST había
\emph{aplazado} la estandarización de McEliece en 2025, alegando que
«la estandarización concurrente por parte del NIST y de ISO crea el
riesgo de estándares incompatibles» y que «después de que concluya el
proceso de ISO, el NIST podría considerar desarrollar un estándar
basado en el de ISO» (traducción propia, \cite{cmnist}). El NIST, por
tanto, no ha estandarizado McEliece; su única selección de cifrado
post-cuántico adicional ha sido HQC \cite{nist8545}.
\end{historia}

\section{La variante \texttt{f}: forma semi-sistemática}\label{sec:param-f}

Recordemos cómo genera el capítulo~\ref{ch:keygen} la clave pública
(sección~«MatGen»): se reduce la matriz binaria $\hat H$ a forma
sistemática $(\mathrm{I}_{mt}\,|\,T)$ y, si la reducción falla, se
empieza de nuevo con otra clave. Con $(\mu,\nu)=(0,0)$ el intento
tiene éxito con probabilidad $\approx29\%$, así que en promedio hacen
falta unos $3{,}4$ intentos \cite{cmrationale}.

La forma \emph{$(\mu,\nu)$-semi-sistemática} relaja ese requisito. Una
matriz $R$ de rango $r$ en forma escalonada reducida, con pivotes en
las columnas $c_0<\cdots<c_{r-1}$, está en forma
$(\mu,\nu)$-semi-sistemática si no tiene filas nulas, $c_i=i$ para
$0\le i<r-\mu$, y cada pivote posterior admite un margen:
\[
  c_i\;\le\;i-\mu+\nu \qquad (0\le i<r).
\]
Para $(\mu,\nu)=(0,0)$ esto es, exactamente, la forma sistemática. Si
$\nu>\mu$, hay muchas más matrices admisibles, y con
$(\mu,\nu)=(32,64)$ basta con que $R$ sea escalonada salvo por los
últimos $\mu=32$ pivotes, que pueden desplazarse dentro de una ventana
de $\nu=64$ columnas. MatGen (caso general) reduce $\hat H$ hasta esa
forma, y luego intercambia las últimas columnas $i\leftrightarrow c_i$
---arrastrando los $\alpha_j$--- hasta dejar $(\mathrm{I}_{mt}\,|\,T)$.
El «libro de permutaciones» $c=(c_{mt-\mu},\dots,c_{mt-1})$ se guarda
en la clave privada (8 bytes; es la lista $c$ de la
tabla~\ref{tab:keygen-tamanos}, presente también en los conjuntos
$(\mu,\nu)=(0,0)$, donde codifica $2^{32}-1$).
La figura~\ref{fig:param-semisys} resume el mecanismo.

\begin{figure}[htbp]
\centering
\begin{tikzpicture}[
    cell/.style={draw=mcgray!50,minimum size=0.30cm,inner sep=0pt},
    one/.style={cell,fill=mcblue,draw=mcblue},
    piv/.style={cell,fill=mcorange,draw=mcorange},
    ilabel/.style={font=\scriptsize,align=center},
    arr/.style={-{Stealth[length=2.2mm]},thick,mcgray}
  ]
  % ventana sombreada
  \fill[mcorange2!30,rounded corners=2pt] (0.45,0.15) rectangle (2.25,-1.65);
  % rejilla izquierda 6x20
  \foreach \i in {0,...,5} \foreach \j in {0,...,19} {
    \node[cell] at (\j*0.3,-\i*0.3) {};
  }
  \foreach \i/\j in {0/0,1/1} \node[one] at (\j*0.3,-\i*0.3) {};
  \foreach \i/\j in {2/4,3/5,4/6,5/7} \node[piv] at (\j*0.3,-\i*0.3) {};
  \node[ilabel] at (1.35,0.55) {escalonada:};
  \node[ilabel,color=mcorange,align=center] at (4.6,-2.35)
    {ventana de $\nu$ columnas:\\ últimas $\mu$ filas};
  \draw[arr,mcorange] (3.3,-1.9) -- (2.05,-1.0);
  % flecha
  \draw[arr] (5.95,-0.75) -- (7.25,-0.75);
  \node[ilabel] at (6.6,-0.28) {$\le\mu$ intercambios};
  % rejilla derecha 6x12 (I|T)
  \fill[mcsky!45,rounded corners=2pt] (9.35,0.15) rectangle (11.05,-1.65);
  \foreach \i in {0,...,5} \foreach \j in {0,...,11} {
    \node[cell] at (7.6+\j*0.3,-\i*0.3) {};
  }
  \foreach \i in {0,...,5} \node[one] at (7.6+\i*0.3,-\i*0.3) {};
  \node[ilabel] at (9.3,0.55) {sistemática};
  \node[ilabel] at (9.3,-1.95) {$(\mathrm{I}_{mt}\,|\,T)$: clave pública};
  \node[ilabel,color=mcsky] at (11.5,-0.75) {$T$};
\end{tikzpicture}
\caption{La forma $(\mu,\nu)$-semi-sistemática, en esquema reducido
($\mu=4$, $\nu=6$; en el estándar, $(32,64)$). Izquierda: la matriz
escalonada; los primeros $mt-\mu$ pivotes ya están en la diagonal y
los últimos $\mu$ viven dentro de una ventana de $\nu$ columnas.
Basta con $\le\mu$ intercambios de columnas para alcanzar la forma
sistemática de la derecha.}
\label{fig:param-semisys}
\end{figure}

¿Qué se gana? Probabilidad de éxito. Con $(\mu,\nu)=(32,64)$ la
probabilidad de fallo de cada intento cae por debajo de $2^{-30}$: casi
siempre basta \emph{un} intento, frente a los $\approx3{,}4$ de
$(\mu,\nu)=(0,0)$. El precio: la computación en forma escalonada debe
hacerse sin ramificaciones dependientes de secretos (en tiempo
constante), y eso solo es razonable si queda confinada a una ventana
pequeña ---de ahí que $\nu$ se mantenga modesto---. La ganancia de
holgura en la demostración de seguridad también se cuantifica: aceptar
más matrices pierde como mucho $2$ bits de tensión en la reducción
\cite{cmrationale}.

¿Por qué existen entonces conjuntos sin \texttt{f}? Porque hay tres
argumentos para mantener $(\mu,\nu)=(0,0)$ \cite{cmrationale}: (1)~el
software optimizado de Classic McEliece ya acelera los reintentos
saltándose casi toda la eliminación gaussiana en los casos de fallo;
(2)~muchas aplicaciones no tienen la generación de claves como cuello
de botella; (3)~la simplicidad tiene valor. ISO estandarizó las
cuatro variantes para cada tamaño: quien quiera velocidad de
generación elige \texttt{f}; quien prefiera el camino más simple y
auditado, el conjunto base.

\section{La variante \texttt{pc}: confirmación de texto claro}\label{sec:param-pc}

En la ronda 3, la transformación CCA de Classic McEliece usaba
\emph{dos} defensas a la vez: rechazo implícito \emph{y} confirmación
de texto claro. La confirmación añade un segundo componente al
criptograma, $C_1=H(2,e)$, y el descifrado lo comprueba
\cite{cmpc}:

\begin{itemize}[leftmargin=1.5em,itemsep=3pt]
  \item $\textsf{Encap}$: $C_0=\textsf{Encode}(e,T)$,
        $C_1=H(2,e)$; el criptograma es $C=(C_0,C_1)$ y la clave,
        $K=H(1,e,C)$.
  \item $\textsf{Decap}$: se parte $C=(C_0,C_1)$; se descifra
        $C_0$ como siempre; después se recalcula $C_1'=H(2,e)$ y, si
        no coincide con $C_1$, también se toma la rama de rechazo
        ($e\leftarrow s$, $b\leftarrow0$). La clave sale de
        $K=H(b,e,C)$.
\end{itemize}

Nótese el detalle: la confirmación \emph{no} introduce errores
explícitos. Si la comprobación falla, \textsf{Decap} sigue devolviendo
$32$ bytes por la rama de rechazo implícito; el KEM continúa siendo
una función total. Lo que añade es una \emph{demostración}: $C_1$ solo
puede calcularlo quien conoce $e$, así que el criptograma certifica
que su emisor conocía el error. Quien posee la clave privada puede
verificar esa demostración explícitamente (descifra $C_0$, recalcula
$H(2,e')$ y compara con $C_1$), si su aplicación lo necesita.

El costo y la historia: $C_1$ ocupa $\lceil\ell/8\rceil=32$ bytes (los
mismos que una clave de sesión), y el hash pasa a tener tres tipos de
entrada, distinguidas por el primer byte: $(2,v)$, $(1,v,C)$ y
$(0,v,C)$. En la ronda 4, el equipo \emph{quitó} la confirmación del
conjunto por defecto y la dejó como variante opcional
\cite{cmmods3}, por tres razones: (1)~separarse limpiamente de la
patente estadounidense US~9.912.479 ---que no se aplica literalmente
a Classic McEliece, cuyo paquete de cifrado tiene un solo objeto, el
error, y no «el vector de error y el valor del texto claro», pero cuya
sombra legal se prefirió despejar, volviendo a la «forma tradicional
de Niederreiter»---; (2)~ahorrar los $32$ bytes, que en un esquema de
criptogramas diminutos se notan; y (3)~simplificar especificación y
software, ya que para las demostraciones ajustadas de IND-CCA2 en ROM
y cuasi-ajustadas en QROM (teoremas~\ref{thm:kem-rom}
y~\ref{thm:kem-qrom}) el rechazo implícito por sí solo basta.
En el debe: quitar la confirmación abre un ataque eficiente en un
modelo \emph{más allá} de IND-CCA2 en el que un solo bit de la clave
privada almacenada se invierte (por ejemplo, un fallo de DRAM); el
ataque análogo contra NTRU-HRSS \cite{ntrufault} motivó precisamente
este cuidado. La mitigación recomendada no toca el KEM: codificar la
clave privada almacenada con corrección de errores (los sistemas
comerciales ya hacen SECDED sobre la DRAM). ISO, en fin, estandariza
las dos familias de variantes: quien quiera las defensas duales de la
ronda 3 usa \texttt{pc} o \texttt{pcf}.

\section{Simply Decoded y Narrowly Decoded: el relleno}
\label{sec:param-narrowly}

Los conjuntos de parámetros no eligen un múltiplo de $8$ por
casualidad: el criptograma ocupa $\lceil mt/8\rceil$ bytes y la clave
pública $mt\lceil k/8\rceil$ bytes (una fila de $T$ por línea, rellena
cada una a byte completo). Cuando $mt$ o $k$ \emph{no} son múltiplos
de $8$, la última posición contiene bits de relleno que la
especificación deja a cero, y ahí nacen dos formas de leer una cadena
de bytes \cite{cmspec}:

\begin{itemize}[leftmargin=1.5em,itemsep=3pt]
  \item \textbf{Simply Decoded}: los bits de relleno se ignoran al
        interpretar la entrada;
  \item \textbf{Narrowly Decoded}: se rechazan las entradas
        (criptogramas y claves públicas) con relleno distinto de
        cero, devolviendo $\bot$.
\end{itemize}

De los dieciséis conjuntos, \textbf{solo \texttt{mceliece6960119}
tiene relleno}: allí $mt=1547$ (el criptograma tiene $194$ bytes, con
$5$ bits de relleno en el último) y $k=5413$ (cada fila de la clave
pública tiene $3$ bits de relleno). En los demás, $mt$ y $k$ son
múltiplos de $8$ y las dos lecturas coinciden. La
figura~\ref{fig:param-narrowly} ilustra el caso.

\begin{figure}[htbp]
\centering
\begin{tikzpicture}[
    bcell/.style={draw=mcgray!60,minimum size=0.5cm,inner sep=0pt,font=\scriptsize},
    rcell/.style={bcell,fill=mcblue!25},
    pcell/.style={bcell,fill=mcorange2!55},
    arr/.style={-{Stealth[length=2.2mm]},thick,mcgray},
    lbl/.style={font=\scriptsize,align=center}
  ]
  % bytes completos
  \node[draw=mcgray!60,rounded corners=2pt,minimum width=3.4cm,minimum height=0.5cm,
        font=\scriptsize,fill=mcgray!10] at (1.7,0) {bytes $0$--$192$};
  % ultimo byte
  \foreach \b in {0,...,7} {
    \pgfmathsetmacro{\cx}{4.15+\b*0.5}
    \ifnum\b<3 \node[rcell] at (\cx,0) {\b};
    \else \node[pcell] at (\cx,0) {\b};\fi
  }
  \node[lbl] at (5.9,0.55) {byte $193$ (el último): bits $1544$--$1551$};
  \node[rcell,minimum size=0.22cm] at (4.35,-0.75) {};
  \node[lbl,anchor=west] at (4.5,-0.75) {bits reales ($1547$ en total)};
  \node[pcell,minimum size=0.22cm] at (7.3,-0.75) {};
  \node[lbl,anchor=west] at (7.45,-0.75) {relleno ($5$ bits)};
  % interpretaciones
  \node[draw=mcgreen,fill=mcgreen2!45,rounded corners=2pt,font=\scriptsize,
        align=center,inner sep=4pt,text width=4.6cm] (sim) at (2.2,-2.5)
        {\textbf{Simply}: se ignoran los bits de relleno: el criptograma
         es el mismo elemento de $\Ftwo^{1547}$ y la clave no cambia.};
  \node[draw=mcred,fill=mcred2!45,rounded corners=2pt,font=\scriptsize,
        align=center,inner sep=4pt,text width=4.6cm] (nar) at (8.2,-2.5)
        {\textbf{Narrowly}: si algún bit de relleno es $1$, la entrada
         se rechaza ($\bot$).};
  \draw[arr] (5.4,-1.15) -- (3.6,-1.85);
  \draw[arr] (6.4,-1.15) -- (8.2,-1.85);
\end{tikzpicture}
\caption{El relleno de \texttt{mceliece6960119}: el criptograma son
$1547$ bits dentro de $194$ bytes, y a los $5$ bits finales la
especificación les da dos lecturas. Una implementación u otra debe
elegirse de forma consistente entre las dos partes del protocolo;
mezclarlas convierte en divergencia lo que es una maleabilidad
benigna.}
\label{fig:param-narrowly}
\end{figure}

¿Por qué se llama maleabilidad benigna? Porque dos cadenas de bytes
que difieran solo en el relleno representan el \emph{mismo}
criptograma (el mismo elemento de $\Ftwo^{mt}$), y por tanto el mismo
$K$ en modo Simply; en modo Narrowly una se acepta y la otra no. La
distinción no rompe la seguridad ---el teorema de IND-CCA2 cubre las
cadenas codificadas---, pero es un contrato de interoperabilidad que
las partes deben acordar.

\begin{trampa}[Tres bits por fila, $580$ bytes]
Un implementador que empaquete la clave pública «apretada», como una
tira continua de $kmt$ bits, obtendrá para \texttt{mceliece6960119}
una clave de $\lceil kmt/8\rceil=1\,046\,739$ bytes. La cadena
estándar tiene $1\,047\,319$: la diferencia, $580$ bytes, son los
$3$ bits de relleno de cada una de las $1547$ filas
($1547\times3/8=580{,}125$). Ambos empaquetados describen la misma
matriz $T$, pero solo el segundo es válido frente a otra
implementación ---y frente a los vectores de prueba del
estándar---. Con $T$ en la mano el error se detecta enseguida, pero es
un clásico de los primeros días de una implementación.
\end{trampa}

\section{¿Cuánta seguridad?}\label{sec:param-seguridad}

La tabla~\ref{tab:param-seguridad} resume las estimaciones de la guía
de seguridad \cite{cmsecurity}: el mejor ataque de decodificación por
conjuntos de información (capítulo~\ref{ch:isd}) según el estimador de
Esser--Bellini, en dos de sus modelos de memoria, junto a la categoría
que cada conjunto reclama.

\begin{table}[htbp]
\centering
\footnotesize
\setlength{\tabcolsep}{5pt}
\begin{tabular}{lccccc}
\toprule
conjunto & Prange & MO (mem.\ libre) & MO (mem.\ con costo) & categoría & categoría\\
 & & & & (mem.\ realista) & (mem.\ gratis)\\
\midrule
\texttt{mceliece348864} & $173{,}4$ & $140{,}8$ & $158{,}6$ & 1 & 1\\
\texttt{mceliece460896} & $216{,}8$ & $179{,}8$ & $201{,}0$ & 3 & 2\\
\texttt{mceliece6688128} & $295{,}7$ & $246{,}2$ & $278{,}2$ & 5 & 4\\
\texttt{mceliece6960119} & $296{,}8$ & $245{,}7$ & $279{,}2$ & 5 & 4\\
\texttt{mceliece8192128} & $334{,}1$ & $275{,}6$ & $315{,}6$ & 5 & 5\\
\bottomrule
\end{tabular}
\caption{Costos en bits ($\log_2$) de los mejores ataques conocidos
según el estimador de Esser--Bellini: contra Prange (1962) y contra
May--Ozerov (MO), con memoria de acceso gratis y con costo de acceso
en raíz cuadrada. Las dos últimas columnas: la categoría que el
equipo asigna contando costos realistas de memoria, y la que
correspondería si la memoria se contara como gratis.}
\label{tab:param-seguridad}
\end{table}

Tres advertencias sobre esa tabla, todas de la propia guía:

\begin{enumerate}[leftmargin=1.6em,itemsep=3pt]
  \item \textbf{Los números subestiman.} El estimador cuenta cada
        operación vectorial como \emph{una} operación, sea cual sea su
        longitud. El $140{,}8$ de \texttt{mceliece348864}, por ejemplo,
        cuenta más de $2^{143}$ operaciones de bits, y el ataque
        requiere memoria $2^{86{,}6}$ \cite{cmsecurity}.
  \item \textbf{La memoria no es gratis.} La columna «mem.\ libre»
        supone acceso libre a memoria arbitrariamente grande,
        físicamente irreal; ahí los algoritmos con
        «representaciones» (BJMM, MO) brillan. Cuando se cobra la
        memoria a su costo plausible, la ventaja se estrecha.
  \item \textbf{Las categorías dependen del modelo.} Contando
        memoria realista, el equipo asigna $1,3,5,5,5$; contando
        memoria gratis y operaciones de bits, $1,2,4,4,5$. Ambas
        tablas son sostenibles sobre los mismos datos; el número clave
        es que \texttt{mceliece6960119} es categoría 5 si se cuentan
        los accesos a memoria ($2^{279{,}2}>2^{272}$, la referencia de
        AES-256) y categoría 4 si no ($2^{245{,}7}<2^{256}$).
\end{enumerate}

Las referencias del NIST para las categorías 1, 3 y 5 son búsquedas
exhaustivas de AES-128, AES-192 y AES-256, estimadas en $2^{143}$,
$2^{207}$ y $2^{272}$ operaciones de bits. De ahí las tensiones
numéricas reseñadas: $140{,}8$, el peor valor, roza por abajo la
referencia de la categoría 1; $201{,}0$ cae bajo $2^{207}$; y
$245{,}7$, bajo $2^{256}$. El equipo responde, con razón documentada,
que estas comparaciones mezclan métricas: no es lo mismo «operaciones
de bits contadas por un estimador que trata un XOR de mil bytes como
una operación» que operaciones reales. La discusión completa, y por
qué el récord de estabilidad del esquema es el argumento más fuerte,
llega en los capítulos~\ref{ch:seguridad} y~\ref{ch:isd}.

\section{Elegir un conjunto}\label{sec:param-elegir}

Con todo lo anterior, la guía práctica es corta:

\begin{itemize}[leftmargin=1.5em,itemsep=3pt]
  \item \textbf{Tamaño.} Para seguridad a largo plazo, el equipo
        recomienda los tamaños \texttt{mceliece6*}
        ($6688128$, $6960119$) \cite{cmiso}; añaden un margen
        generoso sobre la categoría~5 clásica. Si el entorno aprieta
        (clave pública de $512$~KiB), el candidato es
        \texttt{mceliece460896}; \texttt{mceliece348864} existe en la
        especificación, pero queda fuera de ISO y no es la elección
        para décadas.
  \item \textbf{Variante \texttt{f}.} Úsese cuando la generación de
        claves sea cuello de botella o deba ser rápida en tiempo
        constante; añade la forma semi-sistemática sin cambiar
        ningún tamaño.
  \item \textbf{Variante \texttt{pc}/\texttt{pcf}.} Úsese cuando el
        protocolo quiera la confirmación de texto claro de la ronda 3
        (por ejemplo, para verificar el conocimiento de $e$ o para
        defensas en modelos con fallos de clave privada); cuesta
        $32$ bytes por criptograma.
\end{itemize}

Para calibrar los compromisos, la tabla~\ref{tab:param-contexto}
compara los tamaños con los otros KEM que un despliegue real barajará.
La silueta de Classic McEliece es inconfundible: frente a ML-KEM-768,
una clave pública entre cuatrocientas y mil doscientas veces mayor
(según el conjunto), un criptograma entre cinco y ocho veces menor, y
una clave privada moderada. La clave de HQC, en
cambio, es de $40$ bytes: es una semilla, regenerada al vuelo, la
misma filosofía que la $\delta$ de Classic McEliece pero llevada al
secreto entero.

\begin{table}[htbp]
\centering
\small
\setlength{\tabcolsep}{7pt}
\begin{tabular}{lrrr}
\toprule
esquema & $ek$ (bytes) & $ct$ (bytes) & $sk$ (bytes)\\
\midrule
ML-KEM-768 & 1184 & 1088 & 2400\\
ML-KEM-1024 & 1568 & 1568 & 3168\\
HQC-128 & 2249 & 4497 & 40\\
HQC-256 & 7245 & 14485 & 40\\
\texttt{mceliece460896} & \num{524160} & 156 & \num{13608}\\
\texttt{mceliece6960119} & \num{1047319} & 194 & \num{13948}\\
\bottomrule
\end{tabular}
\caption{Tamaños de los KEM post-cuánticos relevantes (ML-KEM según
FIPS~203 \cite{fips203}; HQC según el informe del NIST \cite{nist8545};
Classic McEliece según la tabla~\ref{tab:param-base}). El criptograma
de McEliece es el más pequeño con diferencia; la clave pública, la más
grande.}
\label{tab:param-contexto}
\end{table}

\begin{practica}[¿Y en un protocolo real?]
La clave pública de Classic McEliece no viaja en cada conexión: viaja
una vez y se cachea (es estable durante años). El criptograma, en
cambio, viaja en cada intercambio: ahí el tamaño importa, y es donde
los $96$--$240$ bytes juegan a favor. En TLS~1.3, por ejemplo, ML-KEM
y McEliece pueden convivir como grupos híbridos; la página de
implementación del equipo documenta los detalles
\cite{cmimpl}.
\end{practica}

\section{Ejercicios}\label{sec:param-ejercicios}

\begin{enumerate}[leftmargin=1.6em,itemsep=4pt]
  \item \textbf{Aritmética de bytes.} Con $ct=\lceil mt/8\rceil$ y
        $ek=mt\lceil k/8\rceil$, verifique las columnas de la
        tabla~\ref{tab:param-base}. ¿Por qué para cuatro conjuntos
        $ek=\lceil kmt/8\rceil$ pero para \texttt{mceliece6960119} no?
        Calcule la diferencia exacta.
  \item \textbf{La variante \texttt{pc}.} Para
        \texttt{mceliece460896pc}, complete la
        tabla~\ref{tab:param-iso} para los dieciséis conjuntos y
        calcule después cuántos bytes recibe cada hash: $H(2,e)$ por
        un lado, y $H(1,e,C)$ y $H(0,s,C)$ por otro.
  \item \textbf{El listón cuántico de ISO.} Usando el modelo de raíz
        cuadrada, calcule la mitad del logaritmo del ataque de Stern
        (tabla~\ref{tab:param-seguridad}) para los cinco conjuntos y
        compruebe qué conjuntos superan $2^{128}$. ¿Coincide con los
        que eligió el borrador ISO? ¿Y con lo que ISO decidió
        estandarizar al final? Explique la diferencia.
  \item \textbf{Intentos de generación.} Con $(\mu,\nu)=(0,0)$ el
        éxito de MatGen ronda el $29\%$: calcule el número esperado de
        intentos y la probabilidad de necesitar más de dos. Con
        $(\mu,\nu)=(32,64)$ el fallo por intento es $<2^{-30}$:
        calcule el número esperado de intentos. ¿Cuánto trabajo de
        generación se ahorra la variante \texttt{f}?
  \item \textbf{Relleno benigno.} Para \texttt{mceliece6960119},
        construya dos cadenas de $194$ bytes que difieran solo en los
        $5$ bits de relleno. ¿Cuántos $K$ distintos producen en modo
        Simply? ¿Y en modo Narrowly? ¿Por qué esta diferencia debe
        resolverse por acuerdo de protocolo y no dejarse al azar de
        la implementación?
  \item \textbf{Elegir con criterio.} Un servicio de mensajería
        cachea claves públicas de servidor y envía millones de
        criptogramas pequeños; una etiqueta IoT de bajo consumo
        genera su propia clave una vez y no la cachea. ¿Qué conjunto
        y qué variante propondría para cada caso? Justifique con las
        tablas~\ref{tab:param-base} y~\ref{tab:param-contexto}.
\end{enumerate}

\begin{center}\rule{0.55\textwidth}{0.4pt}\end{center}
\medskip
\noindent\textit{Cierre de la parte~\ref{part:classic}.} Los cuatro
capítulos de esta parte recorrieron la especificación completa: la
generación determinista de claves (capítulo~\ref{ch:keygen}), la
encapsulación y desencapsulación con Patterson y rechazo implícito
(capítulos~\ref{ch:encapdecap} y~\ref{ch:labkem}) y el catálogo de
los dieciséis conjuntos ISO con sus variantes
(capítulo~\ref{ch:parametros}). El mecanismo está entero sobre la
mesa; queda la pregunta que decide si merece la pena: su seguridad.
La parte~\ref{part:seguridad} la analiza sin adornos.


\clearpage

% ======================================================================
%  PARTE IV — SEGURIDAD
% ======================================================================
\part{Seguridad}\label{part:seguridad}

\chapter{¿Qué significa romper McEliece?}\label{ch:seguridad}

\begin{roadmap}
La especificación ya está contada: sabemos generar claves, encapsular
y desencapsular. Falta la pregunta que lo decide todo: ¿cuánto cuesta
de verdad romper un conjunto de parámetros? Esta parte la responde en
cuatro capítulos. Este fija el vocabulario: qué objetivos puede
perseguir un atacante (el mensaje, la clave, distinguir), qué está
demostrado y qué se conjetura. Los capítulos~\ref{ch:isd}
y~\ref{ch:estructura} recorren las dos grandes familias de
criptoanálisis ---decodificación y ataques estructurales--- y
el~\ref{ch:2026} fecha el estado del arte: los resultados de 2025 y
2026 que obligaron a revisar supuestos.
\end{roadmap}

\section{Tres objetivos (y no todos rompen)}

Un atacante contra Classic McEliece ve una clave pública
$\mathrm{pk}$ y, en el modelo activo
(definición~\ref{def:kem-indcca}), puede además preguntar al oráculo de
desencapsulación. A grandes rasgos puede perseguir tres objetivos, de
dificultad y significado muy distintos:

\begin{enumerate}[leftmargin=1.7em,itemsep=3pt]
  \item \textbf{Romper el mensaje.} Dados $\mathrm{pk}=T$ y un
        criptograma $C$, encontrar la semilla $e$ con la que se
        encapsuló: la única de peso $t$ con $\Hmat e=C$, donde
        $\Hmat=(\mathrm{I}_{mt}\,|\,\Tmat)$.
  \item \textbf{Romper la clave.} Encontrar \emph{alguna} descripción
        eficiente del código público: el par $(g,\alpha)$ secreto, o
        más en general cualquier algoritmo que decodifique $t$
        errores para el código de $\mathrm{pk}$ (una \emph{clave
        equivalente}, definición~\ref{def:seg-equivalente}). Con ella
        se descifra todo, hoy y en el futuro de esa clave.
  \item \textbf{Distinguir.} Decidir si $\mathrm{pk}$ es una clave
        legítima o una matriz aleatoria del mismo tamaño. Es
        información \emph{estructural}: revela que el código lleva
        huellas de Goppa, aunque no descifra nada por sí sola.
\end{enumerate}

\begin{definicion}{Clave equivalente}{seg-equivalente}
Sea $\mathrm{pk}=T$ con $\Hmat=(\mathrm{I}_{mt}\,|\,\Tmat)$. Una
\emph{clave equivalente} es cualquier algoritmo $\mathcal{D}$ que,
dado $C\in\Ftwo^{mt}$, devuelve un $e\in\Ftwo^{n}$ con $\wt(e)\le t$
y $\Hmat e=C$ siempre que tal $e$ exista, y el símbolo de fallo en
caso contrario.
\end{definicion}

El nombre es exacto: quien tiene una clave equivalente descifra
directamente el criptograma $C$ ---basta pedirle $\mathcal{D}(C)$:
eso es exactamente lo que hace $\textsf{Decode}$--- sin saber nada
de $g$ ni del orden de las $\alpha_i$. Los ataques estructurales del
capítulo~\ref{ch:estructura} buscan exactamente esto, o más: a veces
se conforman con distinguir.

La figura~\ref{fig:seg-mapa} ordena los tres objetivos y sus
relaciones. Recuperar la clave implica romper el mensaje: una clave
equivalente decodifica cualquier criptograma. El recíproco no se
conoce: un algoritmo que resuelva \emph{el juego} del objetivo~1 solo
tiene que ganar con los criptogramas que el juego le entrega, y nadie
ha reducido a partir de ahí un decodificador general. Y distinguir
está, hoy por hoy, al margen: un buen distinguidor de claves públicas
no se convierte ---con lo que se sabe--- en un ataque contra el
mensaje. Cada flecha de esa figura tiene, por tanto, una dirección
demostrada o un interrogante, y la parte~\ref{part:seguridad} entera
se organiza alrededor de ellas.

\begin{figure}[htbp]
\centering
\begin{tikzpicture}[
    obj/.style={draw=mcblue,fill=mcsky!40,rounded corners=3pt,align=center,
      font=\footnotesize,inner sep=5pt,text width=4.0cm,minimum height=1.9cm},
    objg/.style={obj,draw=mcgreen,fill=mcgreen2!45},
    objv/.style={obj,draw=mcviolet,fill=mcviolet!15},
    arr/.style={-{Stealth[length=2.4mm]},thick},
    darr/.style={dashed,thick,mcgray}
  ]
  \node[obj]  (m1) at (2.2,0)
    {\textbf{1.~Romper el mensaje}\\[2pt]
     {\scriptsize hallar $e$ con $\Hmat e=C$; es el problema SD sobre
      la clave pública (capítulo~\ref{ch:isd})}};
  \node[objg] (m2) at (8.0,0)
    {\textbf{2.~Romper la clave}\\[2pt]
     {\scriptsize hallar una clave equivalente; no hace falta el $g$
      de Goppa (capítulo~\ref{ch:estructura})}};
  \node[objv] (m3) at (13.8,0)
    {\textbf{3.~Distinguir}\\[2pt]
     {\scriptsize decidir si $\mathrm{pk}$ parece aleatoria
      (capítulo~\ref{ch:estructura})}};
  \draw[arr,mcgreen] ([yshift=1.5mm]m2.west) -- ([yshift=1.5mm]m1.east);
  \draw[darr] ([yshift=-1.5mm]m1.east) -- ([yshift=-1.5mm]m2.west);
  \draw[darr,{Stealth}-{Stealth}] ([yshift=1.5mm]m2.east)
    -- node[above=3pt,font=\tiny,color=mcgray]{sin reducción conocida}
    ([yshift=1.5mm]m3.west);
  % aspas de "no se conoce" sobre la flecha discontinua 1<-2
  \draw[mcred,thick] (4.78,-0.05) -- (4.98,-0.25);
  \draw[mcred,thick] (4.78,-0.25) -- (4.98,-0.05);
  % aspas sobre la doble flecha 2<->3
  \draw[mcred,thick] (10.80,0.05) -- (11.00,0.25);
  \draw[mcred,thick] (10.80,0.25) -- (11.00,0.05);
  \node[font=\scriptsize,align=center] at (8.0,-1.55)
    {\textcolor{mcgreen}{$\longrightarrow$}~quien tiene una clave
     equivalente decodifica el mensaje \quad
     \textcolor{mcred}{$\times$}~sin implicación demostrada};
\end{tikzpicture}
\caption{Los tres objetivos de un atacante, con la dirección
demostrada (clave $\Rightarrow$ mensaje) y los interrogantes
abiertos. La seguridad del KEM descansa, en último término, en el
objetivo~1: los teoremas del capítulo~\ref{ch:kem} reducen su
IND-CCA2 a la OW-CPA del esquema desnudo, que es exactamente este
problema.}
\label{fig:seg-mapa}
\end{figure}

\section{El juego OW-CPA y el problema SD}

Para medir el primer objetivo se usa la noción clásica de las PKE
deterministas.

\begin{definicion}{One-wayness bajo ataque de texto claro elegido (OW-CPA)}{seg-ow}
Sea $(\textsf{Gen}',\textsf{Enc},\textsf{Dec})$ una PKE determinista
y correcta, con espacio de mensajes $P$. El juego:
\begin{enumerate}[leftmargin=1.8em,itemsep=2pt]
  \item el retador genera
        $(\mathrm{pk}',\mathrm{sk}')\leftarrow\textsf{Gen}'$ y sortea
        $e\leftarrow P$ uniforme;
  \item calcula $C=\textsf{Enc}(\mathrm{pk}',e)$ y entrega
        $(\mathrm{pk}',C)$ al adversario~$\mathcal{A}$;
  \item $\mathcal{A}$ termina con $e'$ y gana si $e'=e$.
\end{enumerate}
La PKE es OW-CPA si todo adversario eficiente gana con probabilidad
despreciable.
\end{definicion}

En Classic McEliece, $\textsf{Enc}$ es $\textsf{Encode}(\cdot,T)$ y
$P$ son los vectores de peso exactamente $t$ ---que
$\textsf{FixedWeight}$ recorre de forma uniforme
(proposición~\ref{prop:fixedweight-uniform})---. El lema siguiente
traduce el juego al lenguaje de la teoría de códigos.

\begin{lema}{OW-CPA de Classic McEliece es SD sobre la clave pública}{seg-ow-sd}
Fijemos un conjunto de parámetros y sea
$(\textsf{KeyGen},\textsf{Encode},\textsf{Decode})$ el esquema
subyacente. Ganar el juego OW-CPA de la definición~\ref{def:seg-ow}
es, exactamente, resolver SD (definición~\ref{def:sd}, versión de
búsqueda) para el código público: dados $\Tmat$ y
$C=\textsf{Encode}(e,T)$, hallar la palabra de error $e$ de peso $t$.
En particular, un atacante OW-CPA es un algoritmo de decodificación
por síndromes para la distribución de matrices que genera
$\textsf{KeyGen}$, y todo algoritmo así gana el juego.
\end{lema}

\begin{proof}
Por linealidad de $\textsf{Encode}$: $\textsf{Encode}(e',T)=\Hmat e'$
con $\Hmat=(\mathrm{I}_{mt}\,|\,\Tmat)$. Si $\mathcal{A}$ gana con
$e'$, entonces $\Hmat(e+e')=0$: el vector $e+e'$ pertenece al código
público, que tiene distancia mínima $\ge 2t+1$
(teorema~\ref{thm:goppa-2t}), y su peso es a lo sumo $2t$. Por tanto
$e+e'=0$, es decir, $e'=e$: no hay soluciones «falsas» de peso
$\le t$, y el juego coincide tarea a tarea con SD. En sentido
contrario, quien resuelve SD para $C=\Hmat e$ devuelve $e$ y gana el
juego.
\end{proof}

Obsérvese lo que el lema \emph{no} dice: nada sobre la dificultad de
resolver SD en la distribución de matrices que produce
$\textsf{KeyGen}$. Traducir un problema a otro no lo vuelve difícil;
solo identifica el enemigo con nombre y apellidos. Eso permite medir
la seguridad con las herramientas del capítulo~\ref{ch:isd}, pero
también recuerda que la confianza última es una conjetura sobre esa
distribución concreta, no sobre «SD en general».

\section{La cadena de reducciones y la conjetura central}

La figura~\ref{fig:seg-cadena} ordena lo que sabemos. La seguridad
IND-CCA2 del KEM está demostrada (capítulo~\ref{ch:kem}) a partir de
dos insumos: la OW-CPA del esquema desnudo y suposiciones simétricas
(el hash y la pseudoaleatoriedad de la clave de rechazo). La OW-CPA
es, por el lema~\ref{lem:seg-ow-sd}, la dificultad de SD sobre la
clave pública. Y ahí se detiene la cadena demostrativa: la dificultad
de SD \emph{en esa distribución} es una conjetura.

\begin{figure}[htbp]
\centering
\begin{tikzpicture}[
    conjbox/.style={draw=mcorange,fill=mcorange2!45,rounded corners=3pt,
      align=center,font=\footnotesize,inner sep=4pt,text width=4.6cm,
      minimum height=1.0cm},
    statebox/.style={draw=mcblue,fill=mcsky!40,rounded corners=3pt,
      align=center,font=\footnotesize,inner sep=4pt,text width=4.6cm,
      minimum height=0.9cm},
    thmbox/.style={draw=mcgreen,fill=mcgreen2!45,rounded corners=3pt,
      align=center,font=\footnotesize,inner sep=4pt,text width=4.6cm,
      minimum height=0.9cm},
    sidebox/.style={draw=mcviolet,fill=mcviolet!12,rounded corners=3pt,
      align=center,font=\scriptsize,inner sep=3.5pt,text width=3.3cm},
    implbox/.style={draw=mcgray,dashed,fill=mcgray!8,rounded corners=3pt,
      align=center,font=\scriptsize,inner sep=3.5pt,text width=3.3cm},
    arr/.style={-{Stealth[length=2.4mm]},thick,mcgray}
  ]
  \node[conjbox] (sd) at (0,0)
    {\textbf{SD sobre la distribución de claves}\\
     {\scriptsize (matrices que genera $\textsf{KeyGen}$) ---
      \textbf{conjetura}}};
  \node[statebox] (ow) at (0,-2.9)
    {\textbf{OW-CPA} del esquema desnudo\\
     {\scriptsize (hallar el error $e$ del criptograma)}};
  \node[thmbox] (ind) at (0,-5.8)
    {\textbf{IND-CCA2} del KEM\\
     {\scriptsize (clave de sesión indistinguible)}};
  \node[sidebox] (hash) at (6.4,-5.0)
    {suposiciones simétricas: el hash $H$ y la PRF de $\delta$};
  \node[implbox] (impl) at (6.4,-1.55)
    {implementación real: canales laterales, fallos de hardware,
     aleatoriedad --- fuera del modelo};
  \draw[arr] (sd.south) -- node[left=3pt,font=\scriptsize,align=right]
    {lema~\ref{lem:seg-ow-sd}: son la misma tarea\\ (implicación trivial)}
    (ow.north);
  \draw[arr] (ow.south) -- node[left=3pt,font=\scriptsize,align=right]
    {teoremas~\ref{thm:kem-rom} y~\ref{thm:kem-qrom}\\ (ROM; en QROM,
     con pérdida)} (ind.north);
  \draw[arr] (hash.west) -- (ind.east);
  \draw[dashed,mcgray] (impl.west) -- (0.4,-1.55)
    node[midway,above,font=\tiny,color=mcgray]{no cubierto};
\end{tikzpicture}
\caption{La cadena de reducciones de Classic McEliece. Arriba, la
conjetura dura; en medio, la propiedad que se mide con el
criptoanálisis (capítulos~\ref{ch:isd} y~\ref{ch:estructura}); abajo,
la garantía del KEM, que además necesita el hash y la PRF. La
implementación real queda fuera del modelo, y se vigila aparte.}
\label{fig:seg-cadena}
\end{figure}

\begin{rigor}[La conjetura central, sin adornos]
Para cada conjunto de parámetros, la conjetura es esta: no existirá
un algoritmo esencialmente mejor que los mejores conocidos hoy para
encontrar $e$ a partir de $(\Tmat,\textsf{Encode}(e,T))$ con $e$
uniforme de peso $t$. La historia admite lectura doble. Es la
conjetura que ha sostenido casi medio siglo de escrutinio, desde
1978, y ninguna de las mejoras publicadas ha pasado de recortes del
exponente (capítulo~\ref{ch:isd}). Pero también es una apuesta
inductiva: nada demuestra que la curva de ataques se vaya a quedar
quieta, y los parámetros se eligen con margen ---y se revisan cuando
el margen se estrecha--- precisamente por eso. Lo que sí está
demostrado: la corrección del esquema
(proposiciones~\ref{prop:1978-correcto} y~\ref{prop:nied-correcto}),
la equivalencia del lema~\ref{lem:seg-ow-sd}, la NP-dureza de SD en
el peor caso~\cite{bcmv1978} y las reducciones del
capítulo~\ref{ch:kem}. Lo conjetural es solo ---y nada menos que---
la dureza en la distribución real.
\end{rigor}

\section{Las familias de ataque de esta parte}

El criptoanálisis de McEliece se organiza en dos grandes familias,
más su versión cuántica; la tabla~\ref{tab:seg-familias} las sitúa y
las reparte entre los capítulos siguientes.

\begin{table}[htbp]
\centering
\small
\setlength{\tabcolsep}{6pt}
\begin{tabular}{@{}p{3.2cm}p{3.4cm}p{7.0cm}l@{}}
\toprule
familia & objetivo & herramienta principal & dónde\\
\midrule
decodificación clásica & el mensaje (y, con trabajo extra, la
clave) & conjuntos de información: de Prange (1962) a May--Ozerov,
con estimadores modernos & cap.~\ref{ch:isd}\\
decodificación cuántica & el mensaje & Grover y paseos cuánticos: a
lo sumo, la raíz del exponente clásico & cap.~\ref{ch:isd}\\
estructural & la clave o distinguir & particiones de soporte,
álgebra sobre $\GF{2^m}$, distinguidores & cap.~\ref{ch:estructura}\\
\bottomrule
\end{tabular}
\caption{Las familias de ataque contra Classic McEliece. El
capítulo~\ref{ch:2026} añade la cuarta dimensión: el tiempo, con los
resultados de 2025 y 2026 y las respuestas del equipo.}
\label{tab:seg-familias}
\end{table}

Las fronteras entre familias son porosas, en las dos direcciones. Un
decodificador \emph{es} una clave equivalente ---por definición--- y
un ataque estructural que recupere una descripción decodificante
rompe el mensaje. Y en sentido inverso: un algoritmo que resolviera
SD para síndromes \emph{elegidos} por el atacante ---no solo los que
le da el juego OW-CPA--- podría reconstruir una matriz de paridad del
código público; informalmente, pidiendo la solución del síndrome
$\Hmat_j+\Hmat_k$, que es la palabra de peso $2\le t$ única
$e_j+e_k$, y despejando columnas. Por eso los ataques de
decodificación también producen, con trabajo adicional, recuperación
de clave, y las dos familias se vigilan juntas.

\begin{trampa}[«Distinguen McEliece»: titular y matiz]
Cada pocos años aparece un titular que distingue las claves públicas
de McEliece de matrices aleatorias. Es un resultado sobre la
\emph{estructura} del código, y en sí no descifra nada: no se conoce
ninguna reducción de un distinguidor a un ataque contra el mensaje.
Lo que un buen distinguidor sí hace es erosionar el argumento
---informal--- de que «la clave pública no revela nada», y obliga a
revisar qué garantías dependían de ese parecido. El episodio de 2025
y 2026 se cuenta con números en los
capítulos~\ref{ch:estructura} y~\ref{ch:2026}.
\end{trampa}

\section{Cómo se leen las cifras de seguridad}

Todas las estimaciones de esta parte son \emph{cotas inferiores del
esfuerzo del mejor ataque conocido}: nada más, y nada menos. Un
«$2^b$» significa siempre $2^b$ unidades de un modelo de coste
declarado, en una fecha, contra un objetivo declarado. Los tres
lenguajes de coste más frecuentes se resumen en la
tabla~\ref{tab:seg-costes}.

\begin{table}[htbp]
\centering
\footnotesize
\setlength{\tabcolsep}{6pt}
\begin{tabular}{@{}p{3.2cm}p{4.4cm}p{6.8cm}@{}}
\toprule
lenguaje & una unidad es\ldots{} & se usa para\ldots{} y su
peligro al comparar\\
\midrule
operaciones de bits & un XOR, un AND\ldots{} sobre un bit &
comparar con las referencias tipo NIST (AES-128 $\approx2^{143}$);
engaña si el ataque trabaja con vectores: un XOR de vectores son
miles de operaciones de bits\\
operaciones vectoriales & un XOR entre vectores completos & seguir
los estimadores de decodificación (las unidades de la
tabla~\ref{tab:param-seguridad}); engaña a la baja: es «vender
barato» el trabajo real\\
tiempo en hardware & un segundo medido & validar implementaciones y
ataques reales; no extrapola a décadas ni a máquinas futuras\\
\bottomrule
\end{tabular}
\caption{Tres lenguajes de coste. Toda comparación honesta de cifras
exige fijar las unidades y el modelo de memoria; la
sección~\ref{sec:param-seguridad} ya dio las advertencias para la
tabla~\ref{tab:param-seguridad}.}
\label{tab:seg-costes}
\end{table}

Un ejemplo de lectura, con la tabla~\ref{tab:param-seguridad} delante.
Para \texttt{mceliece348864}, el mejor ataque clásico contra el
mensaje se cotiza en $2^{140{,}8}$ operaciones \emph{vectoriales} con
memoria $2^{86{,}6}$ (modelo «memoria gratis»), o $2^{158{,}6}$ si la
memoria se paga; el más simple, Prange, cuesta $2^{173{,}4}$. La
referencia NIST de categoría~1 es AES-128, estimado en $2^{143}$
operaciones de bits: comparar $140{,}8$ con $143$ sin más sería un
error de unidades (la propia guía señala que ese ataque cuenta más de
$2^{143}$ operaciones de bits, precisamente por eso). Del lado
cuántico, el mensaje es corto: los algoritmos conocidos ofrecen, en
el mejor caso, la raíz del exponente clásico ---con sobrecostes
enormes y supuestos de hardware lejanos~\cite{gidney2021}---, y los
parámetros se eligen con ese margen en cuenta
(capítulo~\ref{ch:parametros}); la
sección~\ref{sec:isd-cuantico} lo cuantifica.

Y cada cifra lleva fecha. Las de la
tabla~\ref{tab:param-seguridad} proceden de la guía de seguridad del
equipo (2022)~\cite{cmsecurity}; el capítulo~\ref{ch:2026} actualiza
el estado del arte a octubre de 2026.

\section{Fuera del modelo: la implementación}

\begin{practica}[El adversario también mira el código]
Los teoremas y las cifras de esta parte hablan del adversario
matemático. El adversario real, además: mide tiempos de respuesta
(las implementaciones de referencia de Classic McEliece están
escritas para no filtrar secretos por el tiempo~\cite{cmimpl});
provoca fallos de hardware (un bit invertido en la clave privada
almacenada puede ser fatal en modelos más allá de IND-CCA2 ---el
capítulo~\ref{ch:parametros} contó el episodio, las variantes
\texttt{pc} y las mitigaciones con corrección de errores---); y
ataca la aleatoriedad (una semilla $\delta$ o un $e$ sesgados rompen
cualquier KEM por muy dura que sea su matemática). Ninguna de las
tres cosas la cubre una reducción, y ninguna se arregla con
parámetros más grandes.
\end{practica}

\section{Ejercicios}

\begin{enumerate}[leftmargin=1.6em,itemsep=4pt]
  \item \textbf{Direcciones de implicación.} Con la
        figura~\ref{fig:seg-mapa} delante: (a)~explique por qué una
        clave equivalente rompe el mensaje; (b)~explique por qué un
        rompemensajes «de juego» no se convierte en decodificador
        general con lo que se sabe; (c)~¿qué implicaciones tiene un
        distinguidor de claves públicas? ¿Por qué su papel es
        erosionar supuestos y no descifrar?
  \item \textbf{SD en el juguete, a mano.} En el juguete del
        capítulo~\ref{ch:labgoppa}, $n=16$ y $t=2$. Enumere las
        $1+16+120=137$ palabras de peso $\le2$ y localice la única
        cuyo síndrome ---en el sentido del
        capítulo~\ref{ch:labgoppa}--- es $(\hx2,\hx b)$. ¿Por qué
        la unicidad estaba garantizada de antemano por la distancia
        mínima? Después compare escalas: enumerar palabras de peso
        $64$ para \texttt{mceliece348864} ronda $2^{456}$ candidatos,
        mientras el ataque de Prange se cotiza en $2^{173{,}4}$.
        ¿De qué vive esa diferencia?
  \item \textbf{Los cuantificadores de la conjetura.} Escriba la
        conjetura del recuadro «Rigor» con cuantificadores
        explícitos: qué algoritmos, sobre qué distribución, con qué
        éxito, con qué recursos. Explique después por qué estos tres
        enunciados \emph{no} son equivalentes a ella:
        (a)~«SD es NP-duro»; (b)~«no existe un algoritmo
        polinómico»; (c)~«el mejor ataque conocido cuesta
        $2^{140{,}8}$». ¿Cuál de los tres es, en la práctica, el que
        sostiene las decisiones?
  \item \textbf{Leer las columnas.} Explique, sin inventar cifras
        nuevas, por qué $2^{140{,}8}$ (tabla~\ref{tab:param-seguridad})
        no contradice la categoría~1 pese a que AES-128 se estima en
        $2^{143}$: unidades, papel de la memoria $2^{86{,}6}$ y
        significado de «operaciones vectoriales». Calcule además el
        factor entre $2^{140{,}8}$ y $2^{158{,}6}$.
  \item \textbf{La mitad cuántica.} Para \texttt{mceliece348864}, el
        mejor ataque clásico sin trucos de memoria se cotiza en
        $2^{151{,}4}$. Calcule su raíz \emph{en el exponente} (la
        «mitad»), compárela con el listón de ISO del
        capítulo~\ref{ch:parametros} ($2^{128}$ en su modelo
        cuántico) y explique, con la historia de ese capítulo, por
        qué ISO estandarizó igualmente el conjunto \texttt{460896}.
  \item \textbf{Taxonomía de titulares.} Clasifique cada titular en
        su familia y capítulo, y diga qué rompe y qué no:
        (a)~«Un algoritmo distingue las claves públicas de McEliece
        de matrices aleatorias»; (b)~«Mejora de un factor $2^{10}$
        en el coste de los ataques de decodificación»; (c)~«Se
        recupera una clave de un despliegue real midiendo tiempos de
        respuesta».
\end{enumerate}


\chapter{Decodificación por conjuntos de información}\label{ch:isd}

\begin{roadmap}
El lema~\ref{lem:seg-ow-sd} del capítulo anterior redujo el objetivo
principal de un atacante ---recuperar la semilla $e$--- al problema
SD sobre la clave pública. Este capítulo mide ese problema. La
familia dominante de ataques son los \emph{decodificadores por
conjuntos de información}: de Prange (1962) a May--Ozerov (2015),
sesenta años de mejoras combinatorias que han bajado el exponente de
$2^{173{,}4}$ a $2^{140{,}8}$ para \texttt{mceliece348864}. Veremos la
idea de cada algoritmo, la tabla de costes de la guía de seguridad
(2022) y, con lupa, el debate que la rodea: qué cuenta un $2^b$, qué
papel juega la memoria y qué significa exactamente ``la mitad'' en el
modelo cuántico. Es el material con el que se asignan las categorías
de seguridad de la tabla~\ref{tab:param-seguridad}.
\end{roadmap}

\section{La idea: un conjunto de información}

Recordemos el montaje exacto. La clave pública de Classic McEliece es
$\mathrm{pk}=\Tmat$ con $\Hmat=(\mathrm{I}_{mt}\,|\,\Tmat)$, el
criptograma es $C=\Hmat e$ y el objetivo es hallar $e$ de peso $t$
(lema~\ref{lem:seg-ow-sd}). Todos los ataques de este capítulo tratan
$\Tmat$ ---o su matriz generadora equivalente--- como una matriz
\emph{sin estructura}, como si fuese aleatoria. Esa apuesta da a los
algoritmos su enorme generalidad: si mañana alguien encuentra una
grieta en las matrices de Goppa concretas (capítulo~\ref{ch:estructura}),
estos ataques no se dan por enterados; su coste depende solo de
$n$, $k$ y $t$.

El algoritmo más viejo sigue siendo el más fácil de contar, y data de
1962~\cite{prange1962}. En la forma del McEliece de 1978 ---clave
pública $G$ de $k\times n$, texto cifrado $y=Ga+e$--- funciona así:

\begin{enumerate}[leftmargin=1.8em,itemsep=2pt]
  \item se eligen $k$ columnas de $G$ al azar ---el
        \emph{conjunto de información}--- y se comprueba que formen
        una matriz invertible $G_I$;
  \item se calcula $a=G_I^{-1}y_I$. Si los $t$ errores de $e$
        cayeran \emph{todos} fuera de $I$, esta $a$ es la buena;
  \item se verifica: si $y-Ga$ tiene peso $t$, el ataque termina.
\end{enumerate}

La figura~\ref{fig:isd-prange} cuenta la aritmética. En cada intento,
la probabilidad de que las $t$ posiciones del error caigan fuera del
conjunto de información es
\[
  \prod_{i=0}^{t-1}\Bigl(1-\frac{k}{n-i}\Bigr)
  \;\approx\;\Bigl(1-\frac{k}{n}\Bigr)^{t}
  \;=\;\Bigl(\frac{n-k}{n}\Bigr)^{t},
\]
y cada intento cuesta la eliminación gaussiana de una matriz de
$k\times k$ ---un orden de $k^{3}$ operaciones--- más la
verificación. Para \texttt{mceliece348864}, con $k=2720$ y $t=64$:
\[
  \Bigl(\frac{3488-2720}{3488}\Bigr)^{64}
  = (0{,}2202\ldots)^{64}
  \approx 2^{-139{,}7},
  \qquad k^{3}\approx 2^{34{,}2},
\]
de modo que el coste total ronda
$2^{139{,}7}\cdot 2^{34}\approx 2^{173}$: ahí está el $173{,}4$ de la
tabla~\ref{tab:param-seguridad}. Prange, con dos multiplicaciones y
una eliminación gaussiana, fija la vara de medir de todo lo que sigue.

\begin{figure}[htbp]
\centering
\begin{tikzpicture}[
    lbl/.style={font=\scriptsize,align=center},
    lblb/.style={font=\footnotesize,align=center}
  ]
  % franja de columnas
  \fill[mcsky!45,rounded corners=2pt] (0,-0.32) rectangle (11.7,0.32);
  \fill[mcgreen2!60,rounded corners=2pt] (11.7,-0.32) rectangle (15.0,0.32);
  \draw[mcblue,thick,rounded corners=2pt] (0,-0.32) rectangle (11.7,0.32);
  \draw[mcgreen,thick,rounded corners=2pt] (11.7,-0.32) rectangle (15.0,0.32);
  % errores (aspas) fuera del conjunto de informacion
  \foreach \x in {11.95,12.23,12.51,12.79,13.07,13.35,13.63,13.91,14.19,14.47}
    \draw[mcred,thick] (\x-0.1,-0.1) -- (\x+0.1,0.1)
      (\x-0.1,0.1) -- (\x+0.1,-0.1);
  % etiquetas de la franja
  \node[lbl] at (5.85,0.78)
    {$k=2720$ columnas: el conjunto de información (se resuelve $a$)};
  \node[lbl] at (13.35,0.78)
    {$n-k=mt=768$};
  \node[lbl] at (4.4,-0.82)
    {si alguna $\times$ cayera aquí,\\ el intento falla};
  \node[lbl,mcred,text width=6.4cm,align=flush center] at (11.2,-0.98)
    {las $t=64$ posiciones del error, todas fuera del conjunto:
     éxito con probabilidad $\approx 0{,}22^{64}\approx 2^{-139{,}7}$};
  % formulas abajo
  \node[lblb] (f1) at (3.4,-2.15)
    {un intento: eliminación $k\times k$\\ $\approx k^{3}\approx 2^{34}$};
  \node[lblb] (f2) at (11.4,-2.15)
    {coste total del ataque:\\ $2^{139{,}7}\cdot2^{34}\approx 2^{173{,}4}$};
  \draw[-{Stealth[length=2.2mm]},thick,mcgray] (f1.east) -- (f2.west);
  \draw[-{Stealth[length=2.2mm]},thick,mcgray] (7.6,-1.2) -- (7.6,-1.75);
\end{tikzpicture}
\caption{Prange (1962): elegir $k$ columnas, resolver el sistema y
repetir hasta que el azar deje el error fuera del conjunto de
información. Los números son los de \texttt{mceliece348864}
($n=3488$, $k=2720$, $t=64$) y reproducen el $2^{173{,}4}$ de la
guía de seguridad~\cite{cmsecurity}.}
\label{fig:isd-prange}
\end{figure}

\begin{historia}[1962 y 1978: el ataque ya tenía dieciséis años]
Prange publicó su decodificación por conjuntos de información en
1962, dieciséis años antes de que McEliece propusiera usar SD como
candado~\cite{prange1962,mceliece1978}. La consecuencia es una virtud
didáctica poco común: cuando McEliece presentó su esquema, el mejor
ataque genérico ya existía y estaba estudiado, y todo lo que la
comunidad ha hecho desde entonces ha sido \emph{afilar} ese ataque.
Cada mejora ---Stern 1988, Dumer 1991, ball-collision 2011,
BJMM 2012, May--Ozerov 2015--- es una página de esa historia, y las
páginas que siguen las recorren en orden.
\end{historia}

\section{Sesenta años de mejoras}

Las mejoras a Prange se entienden mejor como ataques a la condición
«hay que acertar que \emph{ningún} error cae en el conjunto de
información», que es carísima. Todas aflojan esa condición de una
manera u otra:

\begin{itemize}[leftmargin=1.5em,itemsep=3pt]
  \item \textbf{Lee--Brickell (1988) y Leon (1988).}
        En vez de exigir cero errores en las $k$ columnas elegidas,
        se toleran unos pocos: se resuelve el sistema con un pequeño
        conjunto de posiciones presuntamente erróneas y se busca que
        el vector residual tenga el peso correcto~\cite{lee1988,leon1988}.
  \item \textbf{Stern (1988/89): cruzar listas.} El ataque se
        convierte en una búsqueda de tipo \emph{meet-in-the-middle}:
        se parte el conjunto de información en dos mitades, se
        enumeran errores parciales en cada una y se cruzan las listas
        por su síndrome en una ventana fija
        (figura~\ref{fig:isd-stern})~\cite{stern1989}. El cruce de
        dos listas de tamaño $L$ cuesta del orden de $2L$
        operaciones \emph{por iteración}, y a cambio la probabilidad de
        éxito por iteración sube muchísimo.
  \item \textbf{Dumer (1991).} Reutiliza las listas de Stern entre
        permutaciones en lugar de recomputarlas: mismo coste, menos
        trabajo por intento~\cite{dumer1991}.
  \item \textbf{Ball-collision (2011).} Bernstein, Lange y Peters
        convierten el cruce de listas en un problema de colisiones
        de bolas con radios optimizables~\cite{blp2011}.
  \item \textbf{Representaciones: MMT (2011) y BJMM (2012).} La
        idea ``$1+1=0$'': hay muchas maneras de descomponer un error
        en suma de vectores parciales, y se puede pedir que las
        listas contengan solo \emph{representantes}, reduciendo su
        tamaño sin perder soluciones~\cite{mmt2011,bjmm2012}.
  \item \textbf{May--Ozerov (2015).} Sustituye el cruce exacto de
        listas por la búsqueda eficiente de \emph{vecinos cercanos}
        en métrica de Hamming, con las estructuras de datos de
        Indyk--Motwani; es la mejora que fija los mejores números de
        la tabla~\ref{tab:isd-guia}~\cite{mo2015}.
  \item \textbf{Both--May (2018).} Afina el régimen de tasas de
        error altas y las estructuras de vecinos cercanos; su
        variante con peso parcialmente disjunto es una de las
        columnas de la tabla~\ref{tab:isd-guia}~\cite{bothmay2018}.
\end{itemize}

\begin{figure}[htbp]
\centering
\begin{tikzpicture}[
    st/.style={draw=mcblue,fill=mcsky!40,rounded corners=3pt,align=center,
      font=\scriptsize,inner sep=4pt,text width=2.6cm,minimum height=1.05cm},
    st2/.style={st,draw=mcgreen,fill=mcgreen2!45},
    st3/.style={st,draw=mcorange,fill=mcorange2!45},
    arr/.style={-{Stealth[length=2.2mm]},thick,mcgray}
  ]
  \node[st] (p) at (0,0) {se permutan las columnas; el conjunto de
    información se parte en mitades $k_1+k_2$};
  \node[st2] (l) at (3.9,0) {listas $L_1$, $L_2$: errores parciales de
    cada mitad con el mismo síndrome parcial};
  \node[st3] (v) at (7.8,0) {cruce en una ventana: coincidencia de
    $\ell$ bits fijos};
  \node[st] (c) at (11.6,0) {candidato $e$ de peso $t$; si no, nueva
    permutación};
  \draw[arr] (p) -- (l);
  \draw[arr] (l) -- (v);
  \draw[arr] (v) -- (c);
  \draw[arr,mcred] (c.south) -- ++(0,-0.7) -| node[below,pos=0.25,
    font=\tiny,color=mcred]{siguiente iteración} (p.south);
\end{tikzpicture}
\caption{La idea de Stern y de toda la familia moderna: en vez de
exigir un conjunto de información perfecto, enumerar errores
parciales en dos listas y cruzarlas. Las mejoras de 1991--2018
refinan cada caja: cómo se construyen las listas, cómo se cruzan y
qué se reutiliza entre permutaciones.}
\label{fig:isd-stern}
\end{figure}

\begin{practica}[Reproducir la tabla en una tarde]
El estimador que produce todos los números de este capítulo es
público: el \emph{syndrome decoding estimator} de Esser y Bellini,
publicado en PKC 2022~\cite{esserbellini2022} y disponible (con un
sucesor más amplio) en el repositorio
\texttt{Crypto-TII/syndrome\_decoding\_estimator}. Con él y los
parámetros $(n,k,w)=(3488,2720,64)$ ---y análogos para los otros
cuatro conjuntos--- se reproducen la tabla~\ref{tab:isd-guia} y sus
variantes de memoria. Es la mejor forma de comprobar que estos
números son \emph{cotas del mejor algoritmo conocido}, no teoremas:
viven en una herramienta, con fecha y versión.
\end{practica}

\section{Los números: memoria libre y memoria con coste}

La tabla~\ref{tab:isd-guia} reproduce la tabla principal de la guía de
seguridad del equipo~\cite{cmsecurity}: la salida del estimador de
Esser--Bellini para los cinco conjuntos, con el modelo de
\emph{memoria libre} (acceso a memoria arbitrariamente grande sin
coste). La segunda tabla añade los dos modelos con coste de acceso,
que discutiremos en la sección siguiente.

\begin{table}[htbp]
\centering
\footnotesize
\setlength{\tabcolsep}{4pt}
\begin{tabular}{lcccccccc}
\toprule
conjunto & Prange & Stern & Dumer & BC & BJMM & pdw & MO & BM\\
\midrule
\texttt{mceliece348864} & 173{,}4 & 151{,}4 & 151{,}4 & 151{,}5 & 141{,}9 & 143{,}4 & \textbf{140{,}8} & 142{,}7\\
\texttt{mceliece460896} & 216{,}8 & 193{,}3 & 193{,}2 & 193{,}3 & 180{,}8 & 182{,}6 & \textbf{179{,}8} & 182{,}1\\
\texttt{mceliece6688128} & 295{,}7 & 268{,}0 & 267{,}9 & 268{,}0 & 247{,}3 & 248{,}8 & \textbf{246{,}2} & 249{,}6\\
\texttt{mceliece6960119} & 296{,}8 & 268{,}4 & 268{,}3 & 268{,}4 & 246{,}6 & 248{,}4 & \textbf{245{,}7} & 249{,}2\\
\texttt{mceliece8192128} & 334{,}1 & 303{,}2 & 307{,}4 & 303{,}4 & 277{,}7 & 279{,}1 & \textbf{275{,}6} & 281{,}1\\
\bottomrule
\end{tabular}
\caption{El mejor ataque conocido contra cada conjunto (en bits, modelo
de memoria libre), según el estimador de Esser--Bellini, en la
tabla~1 de la guía de seguridad~\cite{cmsecurity}. BC: ball-collision
(Bernstein--Lange--Peters, 2011); pdw: variante de BJMM con peso
parcialmente disjunto; MO: May--Ozerov, 2015; BM: Both--May, 2018. El
mínimo de cada fila está en negrita: May--Ozerov gana en los cinco
conjuntos.}
\label{tab:isd-guia}
\end{table}

\begin{table}[htbp]
\centering
\footnotesize
\setlength{\tabcolsep}{5pt}
\begin{tabular}{llcccc}
\toprule
conjunto & mem & Prange & Stern & BJMM & MO\\
\midrule
\multirow{2}{*}{\texttt{mceliece348864}} & $1/3$ & 176{,}7 & 159{,}3 & 159{,}1 & 156{,}4\\
                                        & $1/2$ & 178{,}3 & 162{,}8 & 162{,}2 & \textbf{158{,}6}\\
\multirow{2}{*}{\texttt{mceliece460896}} & $1/3$ & 220{,}4 & 201{,}7 & 201{,}6 & 198{,}6\\
                                        & $1/2$ & 222{,}1 & 205{,}3 & 204{,}8 & \textbf{201{,}0}\\
\multirow{2}{*}{\texttt{mceliece6688128}} & $1/3$ & 299{,}4 & 279{,}2 & 278{,}9 & 274{,}9\\
                                         & $1/2$ & 301{,}2 & 283{,}0 & 282{,}3 & \textbf{278{,}2}\\
\multirow{2}{*}{\texttt{mceliece6960119}} & $1/3$ & 300{,}4 & 280{,}1 & 279{,}4 & 275{,}8\\
                                         & $1/2$ & 302{,}2 & 283{,}9 & 283{,}3 & \textbf{279{,}2}\\
\multirow{2}{*}{\texttt{mceliece8192128}} & $1/3$ & 337{,}7 & 316{,}8 & 316{,}3 & 311{,}9\\
                                         & $1/2$ & 339{,}5 & 320{,}7 & 320{,}0 & \textbf{315{,}6}\\
\bottomrule
\end{tabular}
\caption{Los mismos ataques con \emph{costes de acceso a memoria}
(modelo $1/3$: coste de raíz cúbica; modelo $1/2$: coste de raíz
cuadrada), tabla~1 de~\cite{cmsecurity}. Compárese una fila con la
otra: las mejoras combinatorias valen mucho menos cuando la memoria se
paga (para \texttt{mceliece348864}, de $178{,}3$ a $158{,}6$ frente a
de $173{,}4$ a $140{,}8$).}
\label{tab:isd-mem}
\end{table}

La figura~\ref{fig:isd-escalera} cuenta la historia de esos números
para \texttt{mceliece348864} como una escalera. Prange cotiza en
$2^{173{,}4}$; el primer salto ---Stern y su cruce de listas--- baja
hasta $2^{151{,}4}$, y treinta años de refinamientos posteriores
rascan otros $10{,}6$ bits hasta el $2^{140{,}8}$ de May--Ozerov. Con
la memoria pagada, la escalera se aplana: de $178{,}3$ a $158{,}6$,
donde las mejoras modernas aportan solo $4$ bits sobre Stern. Leer
las dos escaleras juntas es leer toda la sección siguiente.

\begin{figure}[htbp]
\centering
\begin{tikzpicture}[
    x=0.23cm,
    dot/.style={circle,inner sep=1.4pt},
    lab/.style={font=\scriptsize},
    labt/.style={font=\tiny}
  ]
  % eje
  \draw[mcgray,thick,-{Stealth[length=2.2mm]}] (-1,0) -- (52,0);
  \foreach \x/\v in {0/130,10/140,20/150,30/160,40/170,50/180}
    \draw[mcgray] (\x,-0.15) -- (\x,0.15)
      node[below,font=\tiny,color=mcgray] {$2^{\v}$};
  % carril memoria libre (y = 1.3)
  \node[lab,anchor=east,align=right] at (-1.5,1.3)
    {memoria libre};
  \draw[mcsky!80,dashed] (43.4,1.3) -- (10.8,1.3);
  \fill[mcblue] (43.4,1.3) circle (2pt) node[above,labt]{Prange $173{,}4$};
  \fill[mcteal] (21.4,1.3) circle (2pt) node[above,labt,yshift=6pt]{Stern $151{,}4$};
  \fill[mcorange] (11.9,1.3) circle (2pt) node[above right,labt,xshift=3pt]{BJMM $141{,}9$};
  \fill[mcgreen] (10.8,1.3) circle (2.6pt) node[above left,labt,xshift=-2pt]{MO $140{,}8$};
  % carril memoria 1/2 (y = -1.3)
  \node[lab,anchor=east,align=right] at (-1.5,-1.3)
    {memoria $1/2$};
  \draw[mcsky!80,dashed] (48.3,-1.3) -- (28.6,-1.3);
  \fill[mcblue] (48.3,-1.3) circle (2pt) node[below,labt]{Prange $178{,}3$};
  \fill[mcteal] (32.8,-1.3) circle (2pt) node[below right,labt]{Stern $162{,}8$};
  \fill[mcorange] (32.2,-1.3) circle (1.4pt);
  \fill[mcgreen] (28.6,-1.3) circle (2.6pt) node[below left,labt,xshift=2pt]{MO $158{,}6$};
  % anotaciones de salto
  \draw[{Stealth[length=1.8mm]}-{Stealth[length=1.8mm]},mcred,thick]
    (43.4,-2.3) -- (10.8,-2.3)
    node[midway,below,labt,color=mcred]
      {$32{,}6$ bits de mejora total};
  \draw[{Stealth[length=1.8mm]}-{Stealth[length=1.8mm]},mcgreen,thick]
    (21.4,-2.85) -- (10.8,-2.85)
    node[midway,below,labt,color=mcgreen]
      {solo $10{,}6$ bits desde Stern};
\end{tikzpicture}
\caption{La escalera de los ataques contra \texttt{mceliece348864}: de
Prange (1962) al mejor ataque de 2022, con memoria libre (arriba) y
pagando la memoria con raíz cuadrada (abajo). El primer salto ---el
cruce de listas de Stern--- es el grande; el resto es afinar. BJMM
aparece en ambos carriles para mostrar el aplanamiento: $141{,}9$ con
memoria libre se convierte en $162{,}2$ pagándola.}
\label{fig:isd-escalera}
\end{figure}

\section{Lo que las mejoras no cambian}

Hay una manera engañosa y una manera honesta de contar estas seis
décadas de ataques, y la guía de seguridad las separa con cuidado.
La manera asintótica: si $k= Rn$ con $0<R<1$ y $t$ crece despacio
(como en McEliece, $t\approx n/\log_2 n$), el coste de Prange es
\[
  \Bigl(\frac{1}{1-R}+o(1)\Bigr)^{t},
\]
y ---esto es lo notable--- \emph{ninguna} de las mejoras desde 1962
ha cambiado el término principal: solo el $o(1)$~\cite{cmsecurity}.
En el sentido de los exponentes asintóticos, Stern, BJMM y
May--Ozerov no han hecho mejor que Prange; han reescrito las
constantes. Cuando se leen artículos que reportan ganancias
asintóticas, casi siempre están midiendo otro régimen ---valores de
$t$ del orden de $n$, como media distancia de Gilbert--Varshamov; por
ejemplo, allí May--Ozerov reportan $O(2^{0{,}0473n})$ frente a
$O(2^{0{,}0576n})$ de Prange~\cite{mo2015}--- que no es el régimen de
McEliece.

\begin{rigor}[$o(1)$ no es $0$]
De $\bigl(1/(1-R)\bigr)^{t}$ no se sigue que «Prange cuesta
$5^{t}$»: para $R=0{,}8$, Prange cuesta $(5+o(1))^{t}$ y las mejoras
también, pero reemplazar $o(1)$ por $0$ produce dos afirmaciones
contradictorias ---que Prange cuesta exactamente $5^{t}$ y que las
mejoras cuestan exactamente $5^{t}$---, mientras que en los números
concretos las mejoras reducen el coste de verdad (de $173{,}4$ a
$140{,}8$ con memoria libre). La resolución es que el $o(1)$ esconde
factores que a tamaño criptográfico no son pequeños, y que difieren
entre algoritmos: el análisis concreto \emph{es} el análisis del
$o(1)$. Por eso la tabla~\ref{tab:isd-guia} se calcula con
estimadores, no con teoremas: los exponentes asintóticos dicen qué
problema es difícil; los estimadores dicen cuánto cuesta hoy.
\end{rigor}

\begin{observacion}{¿Por qué $n=4608$ no es sospechosamente corto?}{isd-acortamiento}
El conjunto \texttt{mceliece460896} usa $n=4608$ mientras que su
cuerpo tiene $8192$ elementos: $n<q$, por debajo del $n=q$ original
de 1978. La duda ---¿y si el caso corto fuera más fácil?--- la
resuelve una reducción de acortamiento: cualquier ataque contra
$(q,n,t)$, con $n\le q$, se transporta a un ataque contra
$(q,q,t)$ ---el sistema original, con cuarenta años de estudio--- y
su probabilidad de éxito se multiplica exactamente por
$\binom{n}{t}/\binom{q}{t}$. Para $(4608,96)$ y $(8192,96)$, la
pérdida de seguridad de este traslado es de $0{,}8$ bits en el
modelo $1/2$: analizar el sistema original de McEliece es, a efectos
prácticos, tan exigente como analizar el acortado. La reducción
completa está en la sección~4.1 de~\cite{cmsecurity}.
\end{observacion}

\section{El debate metodológico: qué cuenta un $2^b$}

La tabla~\ref{tab:isd-mem} introduce el matiz más importante de toda
la parte: \emph{el modelo de memoria cambia las conclusiones}. Los
modelos oficiales del estimador son tres~\cite{esserbellini2022}:

\begin{itemize}[leftmargin=1.5em,itemsep=3pt]
  \item \textbf{mem $=0$: memoria libre.} Los accesos a memoria
        arbitrariamente grande no cuestan. Es el modelo de la
        tabla~\ref{tab:isd-guia}; la propia guía lo llama
        «físicamente implausible»~\cite{cmsecurity}.
  \item \textbf{mem $=1/3$ y $1/2$: costes de raíz cúbica y raíz
        cuadrada.} Cada modelo asigna al acceso a memoria un coste
        plausible en función de su tamaño. En estos modelos, las
        mejoras modernas valen mucho menos: «la mayor parte de la
        aceleración que reivindican estos algoritmos a tamaño
        criptográfico depende de asignar costes irrealmente bajos a
        los accesos a memoria»~\cite{cmsecurity}.
\end{itemize}

El segundo matiz es la \emph{unidad}. El estimador cuenta cada
operación vectorial ---un XOR entre vectores de $n$ bits, digamos---
como una operación: encontrar $C$ colisiones entre dos listas de $L$
vectores cuesta $2L+C$, «sin importar la longitud del
vector»~\cite{cmsecurity}. Es una unidad cómoda para comparar
algoritmos entre sí, pero no es una operación de bits. De ahí la
paradoja aparente que ya señaló el capítulo~\ref{ch:parametros}: el
$140{,}8$ de \texttt{mceliece348864} «cae por debajo» del $143$ de
AES-128 solo si se comparan unidades distintas; en operaciones de
bits, ese ataque cuenta bastante más de $2^{143}$. Y hay una tercera
capa de realismo, la única que valida implementaciones: el tiempo
medido en hardware. La guía lo resume con las dos formas de comparar
\texttt{mceliece6960119} con AES-256: con memoria libre, «las
variantes modernas de ISD han reducido el número de operaciones de
bits considerablemente por debajo de $2^{256}$»; con costes
realistas, «se espera que este conjunto sea más caro de romper que
AES-256 pre-cuántico y mucho más caro que AES-256
post-cuántico»~\cite{cmsecurity}. Por eso las categorías del equipo
son $1,3,5,5,5$: están calculadas con la memoria pagada.

\begin{trampa}[Comparar el número equivocado con la referencia equivocada]
Tres errores de lectura, frecuentes incluso en publicaciones. (1)
Comparar una cifra de memoria libre con las referencias AES: en
$1/2$, para \texttt{mceliece6960119}, May--Ozerov cuesta $279{,}2$
bits, por encima de AES-256 ($272$): la «tensión» con la categoría~5
que algunos artículos señalan nace de comparar $245{,}7$ (libre) con
$272$. (2) Olvidar que los estimadores subestiman las operaciones de
bits: la guía espera que los números «sean sustituidos por números
mayores procedentes de estimadores que cuenten operaciones de
bits»~\cite{cmsecurity}. (3) Creer que un número bajo implica un
ataque ejecutable: $2^{140{,}8}$ operaciones vectoriales con memoria
$2^{86{,}6}$ describen un atacante hipotético con acceso perfecto a
más memoria de la que ningún computador real puede ofrecer
(\texttt{mceliece348864} no está amenazado en la práctica; el propio
estimador la asume gratuita y perfectamente direccionable).
\end{trampa}

\section{Criptoanálisis cuántico}\label{sec:isd-cuantico}

Los ataques de esta sección admiten aceleración cuántica, y el
resultado es más modesto de lo que la palabra «cuántico» sugiere.
Grover ---y los paseos cuánticos, su generalización--- ofrece una
raíz cuadrada: si un ataque requiere $2^{\lambda}$ operaciones
clásicas, la versión cuántica ronda $2^{\lambda/2}$, pero solo del
bucle externo, y con sobrecostes grandes en el bucle
interno~\cite{grover1996,bernstein2010,kachigar2017,kirshanova2018}.
La propia guía lo formula así: «está claro estructuralmente que los
ataques cuánticos conocidos producen un efecto menor que en, por
ejemplo, la búsqueda de claves AES», porque Grover aplicado a la
decodificación por conjuntos de información sufre un sobrecoste
mucho mayor en el bucle interior~\cite{cmsecurity}.

Con esa advertencia, el modelo estándar para asignar seguridad
post-cuántica (el de ISO, capítulo~\ref{ch:parametros}) es simple:
tomar el mejor ataque clásico sin memoria y partir su exponente.
Para Stern sin memoria~\cite{cmsecurity}, y con la mitad:

\begin{center}
\small
\begin{tabular}{lccccc}
\toprule
& \texttt{348864} & \texttt{460896} & \texttt{6688128} & \texttt{6960119} & \texttt{8192128}\\
\midrule
clásico (Stern, mem.\ $0$) & $151{,}4$ & $193{,}3$ & $268{,}0$ & $268{,}4$ & $303{,}2$\\
mitad (modelo cuántico) & $75{,}7$ & $96{,}7$ & $134{,}0$ & $134{,}2$ & $151{,}6$\\
\bottomrule
\end{tabular}
\end{center}

Dos consecuencias. La primera: en ese modelo pesimista, los dos
conjuntos pequeños caen por debajo de $2^{128}$, y por eso la guía
solo recomienda los tres mayores para horizontes largos (la
historia completa, incluida la decisión de ISO de estandarizar la
familia entera, está en el capítulo~\ref{ch:parametros}). La
segunda: el mecanismo de rechazo implícito añade una defensa
independiente. La semilla maestra $\delta$ (que genera la clave de rechazo
$s$) tiene $32$ bytes ---$256$ bits---, y el término correspondiente del
teorema QROM
(teorema~\ref{thm:kem-qrom}) es la raíz de su espacio: un atacante
cuántico necesitaría del orden de $2^{128}$ consultas al hash solo
para explotar esa vía~\cite{cmsecurity}. El cuello de botella sigue
siendo SD, pero el diseño no deja flancos más baratos.

\section{Ejercicios}

\begin{enumerate}[leftmargin=1.6em,itemsep=4pt]
  \item \textbf{Prange a mano.} Para \texttt{mceliece348864}
        ($n=3488$, $k=2720$, $t=64$): (a)~calcule el exponente de
        $\bigl((n-k)/n\bigr)^{t}$ con dos decimales; (b)~estime el
        coste de una iteración como $k^{3}$ y componga el total;
        (c)~compare con el $173{,}4$ de la tabla~\ref{tab:isd-guia}
        y explique el medio bit de diferencia (pista: invertibilidad
        y constantes de la eliminación).
  \item \textbf{El primer salto.} La escalera de
        \texttt{mceliece348864} va de $173{,}4$ (Prange) a $151{,}4$
        (Stern) y de ahí a $140{,}8$ (May--Ozerov). Calcule cuántos
        bits y cuántos factores de $2^{20}$ aporta cada salto.
        ¿Por qué decimos que «el primer salto es el grande»?
  \item \textbf{La paradoja del $5^{t}$.} Explique con el recuadro
        de rigor: ¿cómo pueden Prange y May--Ozerov tener el mismo
        exponente asintótico $(1/(1-R))^{t}$ y, a la vez, costes
        concretos que difieren en $32{,}6$ bits para
        \texttt{mceliece348864}?
  \item \textbf{Las dos referencias.} Con las
        tablas~\ref{tab:isd-guia} y~\ref{tab:isd-mem}: (a)~¿cuántos
        bits cuesta pagar la memoria en el mejor ataque de cada
        modelo para \texttt{mceliece6960119}? (b)~¿Cómo se compara
        cada cifra con AES-256 ($272$ operaciones de bits) y por qué
        la categoría del equipo es la~5?
  \item \textbf{El acortamiento.} Con la
        observación~\ref{obs:isd-acortamiento}: escriba en cinco
        líneas por qué saber que $(8192,96)$ es duro «basta» para
        creer que $(4608,96)$ lo es, y en qué se pierde el $0{,}8$
        de bit.
  \item \textbf{Optimismo cuántico.} La tabla cuántica de esta
        sección parte el exponente de Stern, pero la guía avisa de
        sobrecostes del bucle interno. Explique por qué la mitad es
        una \emph{cota optimista} para el atacante, y qué papel
        juegan en el balance los $256$ bits de la semilla maestra
        $\delta$ (que genera la clave de rechazo $s$).
\end{enumerate}


\chapter{Ataques estructurales: recuperar la clave privada}\label{ch:estructura}

\begin{roadmap}
Los ataques de decodificación del capítulo~\ref{ch:isd} tratan la
clave pública como una matriz aleatoria y van a por el error $e$.
Esta otra familia hace lo contrario: explota que detrás de la matriz
hay \emph{estructura} ---un polinomio de Goppa $g$ y un soporte
ordenado $\alpha_0,\dots,\alpha_{n-1}$--- y ataca directamente la
clave privada. Es la familia que durante casi medio siglo pareció
una promesa rota: soportes que se dividen, distinguidores que
reconocen a Goppa, y aun así un muro de $2^{9116{,}8}$ donde el
mejor ataque genérico cuesta $2^{140{,}8}$. Veremos por qué el diseño
se blindó contra ella ($n<q$), mediremos el abismo con una figura, y
dejaremos el terreno preparado para el terremoto de 2025--2026 del
capítulo~\ref{ch:2026}. Este es el capítulo que responde a la
pregunta «¿y si atacan la clave en vez del mensaje?».
\end{roadmap}

\section{Dos cajas fuertes: la anatomía de la clave privada}

Recordemos qué hay dentro de una clave privada (capítulo~\ref{ch:goppa}):
un polinomio de Goppa
\[
  g(x) \;=\; g_t x^{t} + g_{t-1}x^{t-1} + \cdots + g_1 x + g_0
  \;\in\; \GF{q}[x],
\]
sin raíces en $\GF{q}$, y un soporte ordenado
$\alpha_0,\dots,\alpha_{n-1}$ de elementos distintos de $\GF{q}$
---las raíces del polinomio $\prod_i (x-\alpha_i)$---. Con la clave
privada $(g,\alpha_0,\dots,\alpha_{n-1})$ uno descifra exactamente
igual que el receptor: cualquier atacante que la recupere rompe la
tarea OW-CPA sin tocar el problema SD. La pregunta estructural es,
por tanto: \emph{¿cuánto cuesta enumerar las claves privadas
posibles?}

Dos hechos de conteo elemental, que la guía y la rationale del
equipo subrayan~\cite{cmsecurity,cmrationale}:

\begin{itemize}[leftmargin=1.5em,itemsep=3pt]
  \item Los polinomios $g$ candidatos (mónicos, de grado $t$ sobre
        $\GF{q}$, sin raíces en $\GF{q}$) son del orden de $q^{t}/t$.
        Para el conjunto clásico $m=12$, $t=64$: $q^{t}/t =
        4096^{64}/64 = 2^{768-6} = 2^{762}$. Para los parámetros
        originales de 1978, $2^{494{,}4}$.
  \item Las tuplas ordenadas $(\alpha_0,\dots,\alpha_{n-1})$ de
        elementos distintos de $\GF{q}$ son $q(q-1)\cdots(q-n+1)$:
        para $\texttt{mceliece348864}$, unos $38\,499$ bits; para los
        parámetros originales ($n=q$), $1024!\approx 2^{8\,769}$.
\end{itemize}

A primera vista ambos números son abrumadores, y ahí está el primer
malentendido: \emph{los ataques estructurales no enumeran a lo
bruto}. Dividen el trabajo. La figura~\ref{fig:est-rutas} muestra
las dos rutas conocidas hacia la clave, y la caja negra en la que se
apoyan: el algoritmo de división del soporte de Sendrier
(2000)~\cite{sendrier2000}.

\begin{figure}[htbp]
\centering
\begin{tikzpicture}[
    nd/.style={draw=mcgray,fill=white,rounded corners=3pt,align=flush center,
      font=\scriptsize,inner sep=4pt,minimum height=0.95cm,text width=2.5cm},
    pk/.style={nd,draw=mcblue,fill=mcsky!45,text width=2.2cm,font=\footnotesize},
    ok/.style={nd,draw=mcgreen,fill=mcgreen2!50,text width=2.9cm},
    sup/.style={nd,draw=mcorange,fill=mcorange2!50,text width=3.1cm},
    gen/.style={nd,draw=mcteal,fill=mcteal2!55,text width=2.9cm},
    arr/.style={-{Stealth[length=2.2mm]},thick,mcgray},
    ann/.style={font=\tiny,color=mcgray,align=center}
  ]
  % fila estructural
  \node[pk] (pk1) at (0,2.3) {clave pública\\ ($g$ y $\alpha$'s ocultos)};
  \node[nd,draw=mcorange] (g1) at (3.5,2.3) {adivinar $g$:\\ $2^{762}$ opciones\\ (caso $m=12$, $t=64$)};
  \node[sup] (ss) at (7.3,2.3) {support splitting\\ (caja negra: usa $g$ y el\\ conjunto $\{\alpha_i\}$)};
  \node[ok] (priv) at (11.4,2.3) {$(\alpha_0,\dots,\alpha_{n-1})$ ordenados:\\ descifrar como el receptor};
  \draw[arr] (pk1) -- (g1);
  \draw[arr] (g1) -- (ss);
  \draw[arr] (ss) -- (priv);
  % puerta de abajo
  \node[nd,draw=mcred,fill=mcred2!45] (a1) at (3.5,-1.3)
    {adivinar $(\alpha_0,\dots,\alpha_{n-k})$:\\ $q(q-1)\cdots(q-n+k)$ opciones,\\ $\approx 2^{9116{,}8}$ para \texttt{348864}};
  \draw[arr,mcred] (a1) -- (ss) ;
  \draw[arr,mcred,dashed] (pk1.south) |- (a1.west);
  % fila ISD
  \node[pk] (pk2) at (0,-3.7) {clave pública\\ ($g$ y $\alpha$'s ocultos)};
  \node[gen] (isd) at (3.5,-3.7) {decodificar (ISD,\\ capítulo~\ref{ch:isd}):\\ $2^{140{,}8}$};
  \node[ok] (e2) at (7.3,-3.7) {recuperar $e$:\\ descifrar};
  \draw[arr] (pk2) -- (isd);
  \draw[arr] (isd) -- (e2);
  \draw[arr,dashed,mcgray] (pk2.north) -- (pk1.south)
    node[midway,right,ann,text width=1.6cm]{el atacante elige ruta};
\end{tikzpicture}
\caption{Las rutas conocidas hacia la clave privada, a fecha de 2022.
Arriba, la ruta estructural: adivinar $g$ y dividir el soporte con el
algoritmo de Sendrier~\cite{sendrier2000}; con $n=q$ basta adivinar
$g$, y con $n<q$ la puerta exige además las posiciones
$(\alpha_0,\dots,\alpha_{n-k})$ (ataque de Kirshanova--May, 2022, con
dato parcial~\cite{kirshanova2022}). Abajo, la ruta genérica de
decodificación, que no usa estructura alguna. Los números son los de
las fuentes: $2^{762}$ y $2^{9116{,}8}$ para \texttt{mceliece348864},
$2^{140{,}8}$ de la tabla~\ref{tab:isd-guia}~\cite{cmsecurity}.}
\label{fig:est-rutas}
\end{figure}

\section{La división del soporte}\label{sec:est-support}

El algoritmo de Sendrier merece presentación pausada, porque es la
pieza que hace pensables los ataques estructurales. Su tesis:
\emph{dados $g$ y el desordenado conjunto $\{\alpha_0,\dots,\alpha_{n-1}\}$
---no las posiciones, solo el conjunto---, el polinomio de Goppa
permite reconstruir el orden $\alpha_0,\dots,\alpha_{n-1}$ en tiempo
rápido}. El nombre viene de ahí: el algoritmo «divide» el soporte en
las partes que $g$ distingue, y de esas divisiones deduce qué
columna del código corresponde a qué elemento del
cuerpo~\cite{sendrier2000}. La guía lo enuncia así: recupera
$(\alpha_0,\dots,\alpha_{n-1})$ «a partir de la clave pública, $g$
y $\{\alpha_0,\dots,\alpha_{n-1}\}$, bajo supuestos
menores»~\cite{cmsecurity}. Los detalles finos viven en el artículo
original; para contar el ataque completo basta tratarlo como caja
negra.

La consecuencia para el McEliece original es inmediata. Su diseño
tomaba $n=q$: el soporte es \emph{todo} el cuerpo $\GF{q}$, un
conjunto público, y por tanto el ataque estructural se reduce a
enumerar el polinomio $g$ ---$q^t/t\approx 2^{494{,}4}$ posibilidades
para los parámetros de 1978--- y dejar que la división del soporte
haga el resto. En 2000 esto no rompió a McEliece, pero convirtió lo
que parecía un espacio de claves de $1024!\approx 2^{8\,769}$ (las
tuplas ordenadas completas) en una enumeración de $q^{t}/t\approx
2^{494{,}4}$ polinomios $g$ ---el orden del soporte lo aporta el
propio algoritmo---, y dejó claro cuál era la pieza a defender.

La defensa, ya lo habrá adivinado el lector, es $n<q$: si el cuerpo
tiene $q$ elementos y solo se usan $n$, el conjunto
$\{\alpha_0,\dots,\alpha_{n-1}\}$ \emph{deja de ser público}.
Aparecen $\binom{q}{n}$ conjuntos posibles, y la caja negra ya no
tiene una de sus entradas. La rationale del equipo lo dice
explícitamente: mantener $n$ por debajo de $q$ es «una defensa extra
contra la división del soporte»; para \texttt{mceliece6960119} hay
$\binom{8192}{6960}\approx 2^{4997}$ conjuntos posibles ---«aunque
podría haber maneras de fusionar el trabajo entre
posibilidades»~\cite{cmrationale}. (Merece la pena notar el doble
filo: contra los ataques de decodificación, $n<q$ no ayuda ni
perjudica; el capítulo~\ref{ch:isd} cuantificó ese efecto en menos
de un bit. La asimetría es real: el mismo truco que es neutro contra
ISD es el que sostiene el muro contra lo estructural.)

En 2022, Kirshanova y May añadieron una pieza más al
rompecabezas~\cite{kirshanova2022}: dado un \emph{dato parcial} ---las
posiciones $(\alpha_0,\dots,\alpha_{n-k})$ de la clave ordenada---,
se recupera el resto de la clave privada con un algoritmo rápido.
Combinado con una búsqueda a lo bruto sobre ese dato parcial, esto
da la cota superior estándar para la recuperación de clave: hay
$q(q-1)\cdots(q-n+k)$ tuplas posibles, «aproximadamente
$2^{9116{,}797}$ posibilidades para
$(q,n,k,t)=(4096,3488,2720,64)$», los parámetros de
\texttt{mceliece348864}~\cite{cmsecurity}. La tabla~\ref{tab:est-espacios}
resume todos los espacios de búsqueda en juego.

\begin{table}[htbp]
\centering
\footnotesize
\setlength{\tabcolsep}{5pt}
\begin{tabular}{llll}
\toprule
ruta & espacio a enumerar & coste (ejemplo) & fuente\\
\midrule
adivinar $g$; soporte conocido ($n=q$) & $q^{t}/t$ &
  $2^{762}$ ($q=4096$, $t=64$) & \cite{cmsecurity}\\
adivinar $g$; con $n<q$ & $\times\,\binom{q}{n}$ &
  $2^{4997}$ ($\binom{8192}{6960}$) & \cite{cmrationale}\\
dato parcial $(\alpha_0,\dots,\alpha_{n-k})$ & $q\cdots(q-n+k)$ &
  $2^{9116{,}8}$ (\texttt{348864}) & \cite{cmsecurity,kirshanova2022}\\
parámetros originales (1978) & $q^{t}/t$ \ y \ $q\cdots(q-n+k)$ &
  $2^{494{,}4}$ \ y \ $2^{4794{,}7}$ & \cite{cmrationale}\\
\midrule
\emph{referencia: decodificación genérica} & --- &
  $2^{140{,}8}$ (\texttt{348864}) & \cite{cmsecurity}\\
\bottomrule
\end{tabular}
\caption{Los espacios de búsqueda de los ataques estructurales, con
los ejemplos que citan las fuentes. La última fila no es un ataque
estructural: es el rival genérico del capítulo~\ref{ch:isd}, incluido
para que la comparación de la figura~\ref{fig:est-abismo} tenga a
mano sus dos extremos.}
\label{tab:est-espacios}
\end{table}

\begin{historia}[Sendrier y la tradición de la equivalencia de códigos]
Los ataques estructurales no nacieron contra McEliece: nacieron de
una pregunta más vieja de la teoría de códigos, la
\emph{equivalencia}: ¿cuándo dos matrices generadoras distintas
describen el mismo código con las posiciones permutadas? Sendrier
llevaba años construyendo herramientas para responderla cuando, en
2000, publicó la división del soporte~\cite{sendrier2000}. Desde
entonces la criptografía basada en códigos la adoptó como parte
estándar del oficio; el capítulo de Overbeck--Sendrier del libro
clásico de criptografía post-cuántica la presenta junto a los ataques
de decodificación como la segunda gran familia~\cite{overbeck2009}.
La lista de trabajos de criptoanálisis que la rationale del equipo
recopila ---treinta publicaciones entre 1981 y 2022, con hitos como
Canteaut--Sendrier (1998) o la reciente
Carrier--Debris-Alazard--Meyer-Hilfiger--Tillich (2022)--- muestra
las dos tradiciones entrelazadas durante cuatro
décadas~\cite{cmrationale,canteautsendrier1998,carrier2022}.
\end{historia}

\begin{observacion}{El doble filo de $n<q$, en una línea}{est-nq}
Para ISD, elegir $n<q$ cuesta a lo sumo fracciones de bit
(observación~\ref{obs:isd-acortamiento}). Para lo estructural,
elegir $n<q$ multiplica el trabajo de la división del soporte por el
número de conjuntos posibles ($\approx 2^{4997}$ para
\texttt{6960119}, según~\cite{cmrationale}). Un mismo parámetro,
dos efectos de signo opuesto: por eso los conjuntos se eligen con el
estimador delante, no a ojo.
\end{observacion}

\section{Distinguidores: saben que es Goppa, ¿y qué?}\label{sec:est-distinguidores}

Hay una familia de resultados estructurales que no recupera la
clave ni decodifica: \emph{distingue}. Un distinguidor es un
algoritmo que, mirando la clave pública, decide si proviene de una
Goppa binaria o de una matriz aleatoria con probabilidad de acierto
apreciable, y con un coste que en los parámetros reales sería
ridículo ($2^{40}$, $2^{60}$\dots). Como en el capítulo~\ref{ch:seguridad}
adelantamos en forma de trampa, un distinguidor viola el supuesto
«pk parece aleatoria», pero no implica por sí solo ni recuperar $e$
ni recuperar la clave. La historia de estos ataques es la historia de
sus \emph{márgenes}: todos funcionan en regímenes de parámetros que
McEliece no usa, y todos mueren mucho antes de llegar a los
conjuntos reales. La tabla~\ref{tab:est-margenes} los recoge.

\begin{table}[htbp]
\centering
\footnotesize
\setlength{\tabcolsep}{6pt}
\begin{tabular}{llll}
\toprule
ataque & criterio & funciona en & McEliece usa\\
\midrule
Faugère et al.\ 2013~\cite{faugere2013} & pocos errores &
  $\le 8$ ($n=1024$);\ $\le 20$ ($n=8192$) & $50$;\ $96$--$128$\\
Carrier et al.\ 2022~\cite{carrier2022} & tasa de código baja &
  $<0{,}3$ & $0{,}7$--$0{,}8$\\
Couvreur et al.\ 2017~\cite{couvreur2017} & grado pequeño &
  $m\le 2$ & $m=12,13$\\
\bottomrule
\end{tabular}
\caption{Los distinguidores conocidos y sus márgenes, según la
rationale del equipo~\cite{cmrationale}. La segunda columna dice qué
región del espacio de parámetros explota cada ataque; la tercera,
hasta dónde llega; la cuarta, dónde viven los conjuntos de Classic
McEliece. El ataque de 2013 es el llamado «distinguidor de alta
tasa» de Faugère--Gauthier-Umaña--Otmani--Perret--Tillich; el de
Couvreur--Otmani--Tillich ataca las variantes «wild».}
\label{tab:est-margenes}
\end{table}

\begin{rigor}[Distinguir no es romper]
La lógica, para que no queden cabos sueltos. (1)~El objetivo
criptográfico del capítulo~\ref{ch:seguridad} es OW-CPA, y el
lema~\ref{lem:seg-ow-sd} lo redujo a SD: recuperar $e$. (2)~Un
distinguidor devuelve un bit ---«¿es Goppa?»--- y nada más; no
produce ni $e$ ni $(g,\alpha_0,\dots)$. (3)~No existe reducción
conocida del primero al segundo: que la clave pública sea
reconocible como Goppa no da, hoy por hoy, ninguna vía directa para
invertir SD más rápido que los ataques genéricos. (4)~El «caso
peor lógico» está identificado desde el propio distinguidor de
2013: que alguien encadene el distinguidor con un ataque de
inversión. Los autores del distinguidor y el equipo de Classic
McEliece han vigilado esa posibilidad durante una década ---y los
márgenes de la tabla~\ref{tab:est-margenes} son precisamente la
distancia que la cadena tendría que cubrir. (5)~Nada de esto es un
teorema de imposibilidad: es un historial. El capítulo~\ref{ch:2026}
contará cómo en 2025 un distinguidor directo ---no una cadena---
cambió el tono de la conversación.
\end{rigor}

\begin{trampa}[«Si su familia cayó, la mía también»]
Cuando una propuesta basada en otra familia de códigos se rompe
---y ha ocurrido muchas veces, como recuerda la
rationale~\cite{cmrationale}---, los autores de los ataques suelen
probarlos contra Goppa binaria. La rationale recoge el resultado
invariable: «estos estudios muestran consistentemente que el sistema
de McEliece está mucho más allá de todos los ataques conocidos»
(cada golpe contra Goppa queda en la situación de la
tabla~\ref{tab:est-margenes}: funcionan en regiones que los
parámetros evitan). La lección no es que Goppa sea invulnerable: es
que cada familia tiene su anatomía, y el criptoanálisis se trasplanta
entre familias con mucha menos facilidad de la que sugieren los
titulares. El capítulo~\ref{ch:2026} mostrará la otra cara del mismo
argumento: cuando el trasplante \emph{sí} prendió.
\end{trampa}

\section{El balance de 2022}\label{sec:est-balance}

Con todo lo anterior, la foto a fecha de la guía de seguridad es
nítida: el mejor ataque de recuperación de clave contra
\texttt{mceliece348864} cuesta $2^{9116{,}8}$ (tabla~\ref{tab:est-espacios})
y el mejor ataque genérico, $2^{140{,}8}$ (tabla~\ref{tab:isd-guia}).
La figura~\ref{fig:est-abismo} pone las tres barras relevantes en una
misma escala logarítmica: decodificar, adivinar $g$, adivinar los
$\alpha$'s.

\begin{figure}[htbp]
\centering
\begin{tikzpicture}[
    bar/.style={draw=mcgray,thick,rounded corners=1.5pt},
    lab/.style={font=\scriptsize,align=left},
    labt/.style={font=\tiny,color=mcgray}
  ]
  \def\sc{0.001383}
  % eje
  \draw[mcgray,thick,-{Stealth[length=2.2mm]}] (0,-0.5) -- (13.2,-0.5);
  \foreach \b in {2000,4000,6000,8000}
    \draw[mcgray] (\sc*\b,-0.62) -- (\sc*\b,-0.38)
      node[below,labt] {$\b$};
  % barras
  \fill[mcsky!70] (0,-0.05) rectangle (\sc*140.8,0.55);
  \draw[bar,mcblue] (0,-0.05) rectangle (\sc*140.8,0.55);
  \node[lab,anchor=west] at (0.30,0.25)
    {decodificar (ISD): $2^{140{,}8}$};
  \fill[mcorange2!80] (0,-1.35) rectangle (\sc*762,-0.75);
  \draw[bar,mcorange] (0,-1.35) rectangle (\sc*762,-0.75);
  \node[lab,anchor=west] at (1.15,-1.05)
    {adivinar $g$ con el soporte conocido: $2^{762}$};
  \fill[mcorange2!80] (0,-2.65) rectangle (\sc*9116.8,-2.05);
  \draw[bar,mcorange] (0,-2.65) rectangle (\sc*9116.8,-2.05);
  \node[lab,anchor=west] at (0.4,-2.35)
    {adivinar $(\alpha_0,\dots,\alpha_{n-k})$: $2^{9116{,}8}$};
  % llave entre barras
  \draw[decorate,decoration={brace,amplitude=3pt},
    mcgray] (\sc*140.8,0.75) -- (\sc*9116.8,0.75);
  \node[lab,text=mcgray] at (6.4,1.05)
    {el abismo: un factor $2^{8976}$};
  % nota
  \node[labt,text width=13cm,align=center] at (6.6,-3.3)
    {cada barra mide $\log_2(\text{coste})$ en la misma escala
     lineal: la primera ocupa aquí $\approx 0{,}2$ cm, y la
     última, casi la figura entera};
\end{tikzpicture}
\caption{El balance de 2022 para \texttt{mceliece348864}, en escala
de exponentes (bits). Decodificar cuesta $2^{140{,}8}$; recuperar la
clave por la vía de $g$ cuesta $2^{762}$ (y esta vía requiere haber
resuelto antes el conjunto, o $n=q$); recuperarla por la vía del
dato parcial cuesta $2^{9116{,}8}$. Datos:
tablas~\ref{tab:est-espacios} y~\ref{tab:isd-guia}.}
\label{fig:est-abismo}
\end{figure}

La lectura del abismo es obligada en este punto del libro. Durante
cuarenta y cinco años, todos los intentos de explotar la estructura
de Goppa se estrellaron contra distancias como esa: la cita de la
rationale es tajante ---«los ataques de recuperación de clave del
estado del arte son enormemente más lentos que la decodificación por
conjuntos de información»~\cite{cmrationale}--- y la conclusión
estratégica del equipo fue doble. Primera: conservar los códigos de
Goppa binarios, «la elección tradicional, conservadora y bien
estudiada»~\cite{cmrationale}, pese a que existen alternativas con
claves más pequeñas. Segunda: documentar cada margen, porque los
márgenes \emph{son} la seguridad en esta familia de ataques: no
existe un teorema que diga «Goppa es rígida»; existe la cuenta de
todos los ataques que se han quedado cortos.

Ese es el muro tal como se veía en 2022. El capítulo~\ref{ch:2026}
cuenta cómo, entre 2025 y 2026, aparecieron las primeras grietas
---primero un distinguidor directo, después ataques de recuperación
de clave--- y cómo reaccionaron el equipo, ISO, BSI y NIST.

\section{Ejercicios}

\begin{enumerate}[leftmargin=1.6em,itemsep=4pt]
  \item \textbf{Aritmética del muro.} Verifique con una calculadora:
        (a)~$4096^{64}/64 = 2^{762}$; (b)~$1024^{50}/50 \approx
        2^{494{,}4}$; (c)~$\log_2\bigl(q(q-1)\cdots(q-769+1)\bigr)
        \approx 9116{,}8$ para $q=4096$ (puede usar la identidad
        $\log_2 P(q,r) = \sum_{j=0}^{r-1}\log_2(q-j)$, una hoja de
        cálculo, o la aproximación de Stirling).
  \item \textbf{Conjuntos y tuplas.} Estime $\log_2\binom{4096}{3488}$
        con la fórmula de Stirling para entropías
        ($\approx 2^{2476}$) y compárela con $2^{9116{,}8}$: ¿por qué
        una va con $n$ y la otra con $n-k$? (pista: la segunda pide
        posiciones \emph{concretas} del soporte ordenado, no solo el
        conjunto).
  \item \textbf{El abismo en potencias.} (a)~Calcule
        $9116{,}8-140{,}8$ y exprese el cociente de costes como
        potencia de $2$. (b)~Compare esa potencia con $2^{256}$
        (el margen de AES-256) y comente en dos líneas por qué el
        titular «recuperar claves de McEliece cuesta $2^{9116}$»
        es, además de correcto en 2022, casi cómico como amenaza
        práctica.
  \item \textbf{Distinguir no es romper.} Use el
        recuadro de rigor: escriba la cadena lógica que convertiría
        un distinguidor en un ataque de inversión de SD, señale el
        eslabón que falta y explique qué papel juegan los márgenes
        de la tabla~\ref{tab:est-margenes}.
  \item \textbf{El doble filo de $n<q$.} Con la
        observación~\ref{obs:isd-acortamiento} (capítulo~\ref{ch:isd})
        y la observación~\ref{obs:est-nq}: explique en un párrafo por
        qué los conjuntos de Classic McEliece pueden tomar $n<q$ sin
        pagar seguridad contra ISD y cobrar, en cambio, ventaja
        contra la división del soporte.
  \item \textbf{Tabla de supervivencia.} Ordene los tres
        distinguidores de la tabla~\ref{tab:est-margenes} por
        «holgura» (distancia relativa entre su límite y los
        parámetros reales) y discuta cuál de los tres le parece el
        candidato más peligroso a escalar. (No hay solución única:
        se trata de razonar sobre los regímenes.)
\end{enumerate}


\chapter{El terremoto: 2025--2026}\label{ch:2026}

\begin{roadmap}
El capítulo~\ref{ch:estructura} cerró con un muro: recuperar la clave
privada costaba, en 2022, unos $2^{9116{,}8}$ frente a los
$2^{140{,}8}$ de la decodificación genérica. Este capítulo cuenta cómo
ese muro se resquebrajó en tres actos. Primero la grieta
(2024--2025): un distinguidor subexponencial y su réplica formal.
Después la primavera de 2026: ataques algebraicos que rompen retos
públicos y acercan las estimaciones, aún por encima de los niveles
declarados. Y por fin el verano: el método \emph{holdout}, un
distinguidor cuasipolinomial de origen insospechado, y una cascada de
estimaciones de recuperación de clave que caen a $2^{89}$--$2^{149}$
\emph{por debajo} de los niveles declarados de todos los conjuntos
---con supuestos heurísticos y sin ataque práctico alguno---. Cerraremos
con el debate técnico (¿cuánto cuesta mover los bits?), la respuesta
institucional (el equipo, el BSI, ISO, NIST) y una lectura sobria de
lo que un lector de este libro debe llevarse. La regla del capítulo:
separar a cada paso lo demostrado de lo heurístico, y lo heurístico de
lo titulado.
\end{roadmap}

\section{La primera grieta (2024--2025)}\label{sec:2026-grieta}

La tradición del capítulo~\ref{ch:estructura} tenía una moraleja
estadística: cada intento estructural se quedaba corto por miles de
bits. Esa estadística empezó a romperse en 2024, cuando Hugues
Randriambololona publicó el \emph{distinguidor de sizigias}
(«syzygy distinguisher»)\footnote{En álgebra conmutativa, una \emph{sizigia}
es una relación lineal entre un conjunto de polinomios o generadores
($\sum a_i p_i = 0$). En el criptoanálisis de códigos, permite buscar
relaciones algebraicas no visibles a simple vista calculando el núcleo de
matrices de productos, es decir, mediante álgebra lineal pura.}~\cite{randriam2024},
publicado después en EUROCRYPT~2025. Su hallazgo: las claves públicas de Goppa binaria se
pueden distinguir de matrices aleatorias con complejidad
\emph{subexponencial} en la capacidad de corrección del código. Para
\texttt{mceliece348864}, la estimación rondaba los $2^{529}$; para
los conjuntos recomendados por el BSI, al menos $2^{683}$ (véase su
evaluación de 2026~\cite{bsi2026}). Frente a los $2^{9000}$ de la
tabla~\ref{tab:est-espacios}, la grieta era visible; frente a los
$2^{140{,}8}$ del mejor ataque genérico, seguía siendo abisal.

La respuesta del equipo llegó el 17 de abril de 2025, en una nota
fechada titulada precisamente \emph{Notes on a recent claim that a
\texttt{mceliece348864} distinguisher uses only $2^{529}$
operations}~\cite{cm529}. Es una pieza ejemplar de metodología, y
conviene resumir sus tres argumentos porque enmarcan todo lo que
viene después.

\begin{enumerate}[leftmargin=1.6em,itemsep=3pt]
  \item \textbf{La seguridad no descansa en la indistinguibilidad.}
        La cadena del capítulo~\ref{ch:seguridad} va de arriba
        abajo: IND-CCA2 en el QROM se apoya en OW-CPA del cifrado,
        y OW-CPA en la dureza de la decodificación. Que la clave
        pública sea reconocible como Goppa no toca ningún eslabón:
        el equipo lo ilustra con RSA, cuyas claves se distinguen de
        enteros aleatorios de forma trivial sin que ello amenace la
        factorización~\cite{cm529}.
  \item \textbf{Hay un suelo trivial de $2^{256}$.} La clave privada
        se genera a partir de una semilla de 256 bits; un atacante
        puede, en el peor caso, enumerar semillas. Cualquier ataque
        de recuperación de clave por debajo de $2^{256}$ sería una
        noticia de primera magnitud; los $2^{529}$ no lo son. Esta
        observación, casi trivial, se convertirá en la vara de medir
        del capítulo.
  \item \textbf{El historial manda prudentemente.} La nota repasa la
        larga lista de «avances» algebraicos previos ---desde los
        ataques de 2002 que se anunciaban como amenaza para AES
        hasta un trabajo de 2010 que se describía como
        «breakthrough»--- todos ellos digeridos sin consecuencias:
        el método muere al comparar constantes, o no escala, o
        confunde distinguir con romper.
\end{enumerate}

Añadamos dos piezas de contexto. En marzo de 2025, NIST cerró su
cuarta ronda eligiendo HQC como respaldo basado en códigos y dejó a
Classic McEliece fuera de la estandarización, anotando en el
informe IR~8545 que el criptoanálisis reciente «debilita en cierta
medida el argumento de tratarlo como una opción especialmente
conservadora» y señalando que \texttt{mceliece460896} ya quedaba
corto frente a su categoría~\cite{nist8545}. Y en noviembre de 2025,
el equipo publicó una nota sobre \emph{You Only Decapsulate Once},
un ataque de canal lateral que explota variaciones de tiempo
introducidas por compiladores modernos en software criptográfico:
una advertencia de que la superficie de ataque de una implementación
no termina en su algoritmo~\cite{cm2025yodo}.

\begin{historia}[De 2002 a 2024: el largo cementerio algebraico]
La nota del equipo de abril de 2025~\cite{cm529} y la de junio de
2026~\cite{cm610} enumeran, con paciencia de archivero, las
promesas algebraicas caídas. La de 2026 abre con un ejemplo de 2002
que anunciaba coste «subexponencial» contra cifrados de bloque y
resultó inofensivo; sigue con la lista de trabajos contra McEliece
---2010: «breakthrough in code-based cryptography»; 2014:
reducción de Wild--Goppa; 2016: el método de claves compactas de
Faugère--Otmani--Perret--de~Portzamparc--Tillich; 2023: el
modelo Pfaffiano de Couvreur--Mora--Tillich; 2024: las sizigias de
Randriambololona---. La lección del cementerio no es que los papers
fueran malos: es que durante dos décadas cada funeral corporativo
se celebraba antes de la autopsia. En 2026, por primera vez, el
ritual se invertirá: primero las réplicas, después los títulos.
\end{historia}

\section{La primavera de 2026: de CFS a los retos TII}\label{sec:2026-primavera}

El 3 de marzo de 2026 apareció el primer golpe con \emph{víctima
real} fuera de los parámetros de Classic McEliece: Bardet, Lemoine y
Tillich rompieron el esquema de firmas CFS, de 25 años de
antigüedad, construido sobre códigos de Goppa binarios de tasa muy
alta, aprovechando el modelo Pfaffiano de 2023\footnote{Un \emph{Pfaffiano}
es un polinomio que calcula la raíz cuadrada del determinante de una matriz
antisimétrica. El modelo Pfaffiano expresa las relaciones algebraicas del
código alternante subyacente como un sistema de ecuaciones polinómicas (un
ideal algebraico) sobre los menores antisimétricos de una matriz asociada.}
y una búsqueda de cuadrados en el ideal correspondiente. El mismo documento recuperó
las claves de retos de TII «con seguridad declarada de hasta 210
bits»~\cite{bardet2026}. La letra pequeña salvó a Classic McEliece:
el ataque vive de tasas de código extremas, y los conjuntos del
esquema tienen tasas de $0{,}7$--$0{,}8$. Pero el cementerio acababa
de enterrar un esquema entero.

En junio llegaron dos piezas más. El 10 de junio, Briaud, Lemoine,
Randriambololona y Tillich (BLRT) publicaron un ataque algebraico
\emph{subexponencial heurístico} basado en encontrar matrices de
rango~2 en el código de relaciones cuadráticas asociado a la Goppa;
con ellas se recupera la estructura secreta y, de ahí, una clave
equivalente. Rompieron retos TII de 83, 89, 119, 166, 210 y 248
bits declarados, y una clave CFS en 14 horas y 24\,GB de
RAM~\cite{briaud2026}. Para los parámetros de Classic McEliece
nivel~1, los propios autores cifran el ataque en «a lo sumo del
orden de $2^{454}$». El 23 de junio, el equipo respondió con una
nota fechada~\cite{cm610}: auditando el cómputo, el coste varía con
las suposiciones algebraicas entre $2^{454}$ (Wiedemann), $2^{513}$,
$2^{515}$ (con exponente de multiplicación $2{,}373$), $2^{607}$ y
$2^{610}$; ninguna variante baja del suelo trivial $2^{256}$
(ninguno de esos trabajos lo compara, observa el equipo con
retintín); los retos TII rotos son juguetes ($n\le 1409$) que
«omiten una de las defensas centrales» al reducir $t$ y con él las
opciones de $g$; y la narrativa del «ataque que rompe el esquema»
repite, palabra por palabra, la de 1991.

El 29 de junio, Tobias Hemmert ---del propio BSI--- convirtió en
recuperación de clave el distinguidor de «anulación de orden
superior» (\emph{higher-order vanishing}) que él y Andreas Wiemers
habían presentado en CRYPTO~2026~\cite{hemmert2026,hemmertwiemers2026}:
los polinomios homogéneos que se anulan con multiplicidad alta en
las columnas de la matriz de paridad no solo distinguen Goppa de
aleatorio; su estructura permite reconstruir la clave. Conjetura
del autor: el coste sería «aproximadamente el mismo que el del
distinguidor», que el BSI evalúa en al menos $2^{550}$ para los
conjuntos que recomienda~\cite{bsi2026}. Markku-Juhani Saarinen
extendió el método (HOVER) sin mejora significativa para estos
parámetros~\cite{saarinen2026hover}. Balance de la primavera: las
estimaciones estructurales habían bajado de miles de bits a
$2^{454}$--$2^{742}$ ---inmensas, pero ya del orden de la decena
de miles de millones de veces «solo» por encima de los objetivos,
y con la dirección de la flecha en contra.

\begin{historia}[1991, 1997, 2026: la misma obra, el mismo papel]
La nota de junio de 2026~\cite{cm610} recuerda el precedente: un
«criptoanálisis del criptosistema de McEliece» de 1991, que resultó
ser un ataque sobre un juguete con $n=63$; la historia la contó
Berson en 1997, y este libro la citó en el capítulo~\ref{ch:kem}
como anécdota de época. El patrón se repite: una demostración en
tamaño de juguete, una extrapolación en el titular, una espera de
quince años hasta saber si el ataque es un embrión o un espejismo.
La diferencia de 2026 es la velocidad: lo que en 1991 tardó un año
en replicarse, aquí tarda semanas ---y la propia respuesta del
equipo está fechada con precisión notarial.
\end{historia}

\section{El método \emph{holdout} (agosto de 2026)}\label{sec:2026-holdout}

El 7 de agosto de 2026, Ashrujit Ghoshal, Yuval Ishai, Aayush Jain y
Nuozhou Sun (GIJS) publicaron un distinguidor \emph{cuasipolinomial}
para Goppa--McEliece en el régimen asintótico del esquema
($m=\Theta(\log n)$, $t=\Theta(n/\log n)$), con tiempo
$n^{\mathcal{O}(\log n)}$ y ventaja $1-o(1)$. El origen del método
es delicioso: nació de un intento \emph{fallido} de construir
protocolos de recuperación privada de información (PIR) a partir de
códigos localmente decodificables. El título del
artículo~\cite{ghoshal2026} lo confiesa: «PIR Meets McEliece».

La idea, en la figura~\ref{fig:2026-holdout}. Las columnas de la
clave pública son evaluaciones de polinomios secretos: yacen
---salvo reescalado--- sobre una curva algebraica oculta. El test
\emph{retira} (hold out) una columna y busca polinomios de bajo
grado que se anulen con multiplicidad alta en todas las demás. En
una clave real, esos polinomios también se anulan en la columna
retirada: es la estructura de curva forzándolo. En una matriz
aleatoria, solo ocurriría por azar. Ese sesgo, amplificado sobre
muchas retiradas, es el distinguidor. Así descrito parece poco,
pero trae la propiedad que ningún distinguidor previo tenía: la
complejidad crece \emph{cuasipolinomialmente} en el régimen de los
parámetros, y el método aplica a los cinco conjuntos, con
estimaciones de $2^{114}$--$2^{124}$ operaciones de bits para el
test (frente a los $2^{151}$--$2^{287}$ de la decodificación
genérica), aunque con almacenamiento de exabytes ($2^{66}$--$2^{71}$
bits) que los autores no intentaron optimizar~\cite{ghoshal2026}.

\begin{figure}[htbp]
\centering
\begin{tikzpicture}[
    col/.style={draw=mcblue,fill=mcsky!35,line width=0.4pt},
    colx/.style={draw=mcorange,fill=mcorange2!70,line width=0.4pt},
    dot/.style={circle,fill=mcred,inner sep=1.1pt},
    lbl/.style={font=\scriptsize,align=center},
    lblt/.style={font=\tiny,color=mcgray,align=center}
  ]
  % columnas de la clave publica
  \foreach \i in {0,1,2,3,4,5,6} {
    \draw[col] (0.95*\i,0) rectangle +(0.7,1.6);
  }
  \draw[colx] (7.1,0) rectangle +(0.7,1.6);
  \node[lblt] at (2.85,-0.25) {columnas de $G$ (clave pública)};
  \node[lblt,mcorange] at (7.45,-0.25) {retirada};
  \draw[-{Stealth[length=2mm]},mcorange,thick]
    (9.0,2.75) -- (7.85,2.45);
  \node[lbl,mcorange,text width=3.4cm] at (11.0,2.95)
    {se retira una columna y se busca un polinomio $P$ que se anule con
     multiplicidad alta en las demás};
  % puntos: P se anula
  \foreach \i in {0,1,2,3,4,5,6} {
    \node[dot] at (0.95*\i+0.35,1.85) {};
  }
  \node[lbl,mcred] at (3.0,2.15) {$P$ se anula en todas ellas};
  % la retirada: signo de pregunta
  \node[font=\Large,color=mcorange] at (7.45,1.95) {?};
  % paneles de respuesta
  \node[draw=mcgreen,fill=mcgreen2!45,rounded corners=3pt,lbl,
        text width=3.4cm,inner sep=5pt] at (11.9,1.15)
    {\textbf{clave Goppa real:}\\[2pt] $P$ también se anula en la
     retirada (lo fuerza la curva oculta)};
  \node[draw=mcred,fill=mcred2!45,rounded corners=3pt,lbl,
        text width=3.4cm,inner sep=5pt] at (11.9,-0.7)
    {\textbf{matriz aleatoria:}\\[2pt] se anula solo por azar; el
     test lo delata};
\end{tikzpicture}
\caption{El test \emph{holdout} de GIJS~\cite{ghoshal2026}: retirar
una columna y buscar polinomios de bajo grado que se anulen con
multiplicidad alta en todas las demás. Sobre una clave real el
polinomio se anula también en la retirada (estructura de curva);
sobre una matriz aleatoria, no. La extensión a recuperación de
clave usa la geometría local de la curva en la columna
retirada~\cite{vedenev2026,weis2026}.}
\label{fig:2026-holdout}
\end{figure}

El 27 de agosto, GIJS añadieron a su artículo los detalles de un
ataque de \emph{recuperación de clave} heurístico basado en la
misma maquinaria, con estimaciones de $2^{130{,}4}$ a $2^{149{,}0}$
según el conjunto, y de un ataque de recuperación de texto plano
que el equipo evaluó en unos $2^{266}$ para
\texttt{mceliece6960119}~\cite{ghoshal2026,cm610}. La revisión
también declaraba, con honestidad ejemplar, que el ataque «no es
todavía concretamente eficiente» pero que «puede ser relevante para
los niveles de seguridad de NIST». En paralelo, la réplica pública
del equipo (21 de agosto) recordaba que un distinguidor no ataca
los objetivos de seguridad del esquema ---el argumento de
abril de 2025--- y que la maquinaria no produce, por ahora, ningún
ataque ejecutable~\cite{bernstein2026,cm529}.

\section{La cascada de septiembre}\label{sec:2026-cascada}

Lo que ocurrió después fue una carrera de tres semanas entre grupos
---con préstamos cruzados, porque todos parten del test de GIJS---
para llevar el \emph{holdout} de distinguidor a recuperación de
clave. Conviene registrar los hitos con fechas, porque la velocidad
es parte de la historia.

\begin{itemize}[leftmargin=1.5em,itemsep=3pt]
  \item \textbf{Vedenev (20 de agosto; revisión de octubre).} Extiende
        el método a la recuperación de clave para subcódigos de
        GRS (y por tanto alternantes y Goppa)\footnote{Los códigos
        Reed--Solomon generalizados (GRS) y los códigos alternantes son
        familias más amplias que engloban a los códigos de Goppa: un código
        de Goppa binario es el subcódigo subcampo sobre $\F_2$ de un código
        GRS.}: un \emph{único} tipo
        de ejecución del test recupera la clave, con coste estimado
        de $2^{89{,}9}$ a $2^{98{,}6}$ operaciones de bits para los
        parámetros de Classic McEliece. Demostración: claves de los
        retos TII-173, TII-249 y TII-253, «los dos primeros no
        resueltos hasta donde sabemos»~\cite{vedenev2026}.
  \item \textbf{Saarinen (24 de agosto; revisiones).} Publica una
        contabilidad de costes del método («Holdout»), de $2^{126{,}8}$
        a $2^{145{,}2}$ según la revisión de septiembre, y resuelve
        el reto TII-254 ---el más difícil resuelto hasta la fecha
        en su métrica original, $2^{254}$--- con una variante de dos
        anclas: dos núcleos de relaciones de dimensión 121, un
        núcleo común de 80, una geometría proyectiva sobre
        $\GF{2^8}$, y una secuencia de Krylov final\footnote{Un
        subespacio de Krylov $\mathcal{K}_r(M, v) = \mathrm{span}\{v, Mv,
        \dots, M^{r-1}v\}$ permite calcular vectores del núcleo o resolver
        sistemas lineales dispersos gigantescos (como en los algoritmos de
        Lanczos o Wiedemann) mediante aplicaciones sucesivas de la matriz,
        sin requerir su inversión explícita.} de 27{,}2
        horas-GPU en aceleradoras NVIDIA~\cite{saarinen2026a,saarinen2026b}.
  \item \textbf{Weis (11 de septiembre; versión 0.5 de octubre).}
        Stephen Weis, de Anthropic, mejora el ataque en dos frentes.
        Primero, el coste de una ejecución baja a $2^{89}$--$2^{98}$
        operaciones de bits (sin cargo de memoria); pagando cada
        operación por la memoria que direcciona, a coste de raíz
        cuadrada, $2^{107}$--$2^{117}$; bajo la regla contable de la
        guía de seguridad ---que carga el coste de toda la memoria
        en cada operación---, $2^{117}$--$2^{128}$, contra
        $2^{159}$--$2^{316}$ de la decodificación genérica en ese
        mismo modelo. Segundo, dos rutas de recuperación de clave:
        una apoyada en un teorema que demuestra y que exige entre
        100 y 1400 ejecuciones; otra, de una sola ejecución, que
        construye sobre el key recovery de GIJS. «Nada de esto está
        cerca de lo práctico ---escribe---, y una ejecución a tamaño
        completo no puede probarse»~\cite{weis2026}.
  \item \textbf{Apon (26--28 de agosto) y la réplica.} Daniel Apon
        publica una cota inferior de la ruta de Vedenev: por razones
        algebraico-geométricas, el número de columnas retiradas que
        su paso crítico necesita es mayor que $2t+3$; para
        \texttt{mceliece8192128}, eso implica un coste superior a
        $2^{1500}$ operaciones de bits. Vedenev respondió con una
        revisión que añade una segunda ruta
        («código menor») y sostiene que las objeciones ya no aplican
        a la versión nueva~\cite{apon2026,vedenev2026}.
\end{itemize}

La tabla~\ref{tab:2026-costes} reúne las estimaciones publicadas con
el mismo protocolo: operaciones de bits, sin cargo de memoria, para
los cinco conjuntos. Es la tabla que en 2022 habría sido ciencia
ficción: la columna estructural ya no está en los miles de bits.

\begin{table}[htbp]
\centering
\footnotesize
\setlength{\tabcolsep}{6pt}
\begin{tabular}{lccccc}
\toprule
conjunto & decodificación & GIJS & Saarinen & Weis &
AES ref.\\
 & (ISD, cap.~\ref{ch:parametros}) & (27-ago) & (15-sep) & (v0.5) & (NIST)\\
\midrule
\texttt{mceliece348864} & $2^{151{,}4}$ & $2^{130{,}4}$ & $2^{126{,}8}$ & $2^{94{,}3}$ & $2^{143}$\\
\texttt{mceliece460896} & $2^{193{,}3}$ & $2^{149{,}0}$ & $2^{145{,}2}$ & $2^{99{,}3}$ & $2^{207}$\\
\texttt{mceliece6688128} & $2^{268{,}0}$ & $2^{141{,}3}$ & $2^{137{,}5}$ & $2^{102{,}4}$ & $2^{272}$\\
\texttt{mceliece6960119} & $2^{268{,}4}$ & $2^{140{,}5}$ & $2^{136{,}7}$ & $2^{101{,}6}$ & $2^{272}$\\
\texttt{mceliece8192128} & $2^{303{,}2}$ & $2^{141{,}6}$ & $2^{137{,}5}$ & $2^{102{,}4}$ & $2^{272}$\\
\bottomrule
\end{tabular}
\caption{La fotografía de octubre de 2026, en operaciones de bits y
sin cargo de memoria. La columna ISD es el mejor ataque genérico sin
trucos de memoria (guía de seguridad; capítulos~\ref{ch:parametros}
y~\ref{ch:isd}). Las columnas de 2026 son estimaciones
\emph{condicionales}, apoyadas en supuestos heurísticos, de ataques
de recuperación de clave: GIJS (revisión del 27 de agosto, su
tabla~3), Saarinen (revisión del 15 de septiembre) y Weis (versión
0.5 de octubre, tabla~1)~\cite{ghoshal2026,saarinen2026a,weis2026}.
La última columna es la referencia de NIST para las categorías
1/3/5 (búsqueda de clave sobre el AES correspondiente). Ninguna de
las cifras de 2026 corresponde a un ataque ejecutado ni ejecutable
en hardware real.}
\label{tab:2026-costes}
\end{table}

Un matiz imprescindible, del propio equipo: los retos TII resueltos
---$n=214$, $n=223$--- tienen $t$ diminutos y por tanto pocas
opciones de polinomio $g$; «omiten una de las defensas centrales»
del esquema~\cite{cm610}. Que la maquinaria digiera juguetes no
prueba que digiera los conjuntos reales: eso lo dirá el tiempo, o
una ejecución a escala que nadie ha hecho.

\section{El debate de los costes}\label{sec:2026-costes}

Los números de la tabla~\ref{tab:2026-costes} comparten un supuesto
silencioso: que una operación de bits cuesta lo mismo que cualquier
otra. Daniel Bernstein lo atacó de frente en el foro de NIST. Un
\emph{run} del ataque es una cadena de $2^{32}$ a $2^{36}$ pasadas
mutuamente dependientes sobre un array de 28 a 540~TiB
($2^{47{,}8}$--$2^{52{,}1}$ bits); el ataque contra
\texttt{mceliece348864} movería del orden de $2^{93}$ bits entre
memoria y unidades aritméticas. Enrutar esos bits a la velocidad de
interconexión de los superordenadores actuales, calculaba Bernstein,
llevaría unos 90\,000 años, solo para esa componente; y las cotas
de área-tiempo de Brent--Kung sugieren que ningún trazado de
circuito elimina el peaje. La misma asimetría se esconde en la
comparación con AES: la búsqueda de claves no toca memoria, se
amortiza entre muchos objetivos y gana con Grover; el
\emph{holdout} no hace nada de eso~\cite{bernstein2026}.

Weis, lejos de esquivar el golpe, lo incorpora: su artículo admite
que no conoce distribución de datos que haga local el acceso, y
ofrece la escalera completa de modelos contables. Con cargo de
raíz cuadrada ---un modelo con precedente en la propia guía---, una
ejecución cuesta $2^{107}$--$2^{117}$: aún por debajo de la
decodificación genérica con memoria libre. Con la regla más severa
de la guía, $2^{117}$--$2^{128}$. La discusión no es bizantina: es
\emph{el mismo} debate metodológico del capítulo~\ref{ch:isd},
ahora aplicado a la otra orilla. Y tiene precedente ilustre: en
2020, el propio Bernstein argumentó que NIST aplicó a Kyber-512 la
vara de las compuertas clásicas ($2^{143}$), y NIST respondió que
tiende a dar el beneficio de la duda a un conjunto cuyos ataques
quedan a pocos bits de la categoría. Aquí la distancia no es de
pocos bits: bajo el modelo más severo, Weis estima 25 bits bajo la
categoría~1 para \texttt{mceliece348864} y 82--145 bits bajo las
categorías del resto de conjuntos; y Apon declaró el 1 de octubre
tener un análisis no publicado que sitúa el nivel~1 en
$2^{89}$--$2^{92}$ \emph{incluyendo} costes de memoria, pero
decidió no publicarlo «para no perjudicar a los usuarios
actuales»~\cite{weis2026,bernstein2026}.

\begin{rigor}[Demostrado, heurístico, estimado: el mapa de 2026]
La contabilidad honesta de esta historia tiene tres columnas.
(1)~\textbf{Demostrado}: el distinguidor de GIJS está probado (en
su régimen asintótico), y Weis prueba varios enunciados sobre
Goppa binaria que sostienen su ruta. (2)~\textbf{Heurístico}: la
recuperación de clave a tamaño real depende de supuestos de rango
y genericidad ---Weis enumera cinco, los prueba en claves pequeñas
y en retos TII ($n=214$ frente a $n\ge3488$) y argumenta por qué
las desviaciones observadas no ocurrirían a escala---; el grado
del ataque es 7; el paso caro (el cómputo de núcleos) no se ha
ejecutado contra ningún conjunto real. (3)~\textbf{Estimado}: todas
las cifras de la tabla~\ref{tab:2026-costes} dependen de un modelo
de costes ---el de las operaciones de bits desnudas--- que el
propio equipo disputa, y que mueve las cuentas 20--30 bits según
el cargo. Nada de la columna (2)~(3) convierte a 2026 en «McEliece
roto»; nada de la columna (1) permite ya decir «McEliece inmune».
El estado del arte es exactamente eso: un arte, con supuestos
declarados y réplicas en curso.
\end{rigor}

\section{La respuesta institucional}\label{sec:2026-institucional}

El equipo de Classic McEliece sostuvo durante todo el otoño una
posición pública en dos capas. Como equipo: la réplica del 21 de
agosto al distinguidor ---no ataca los objetivos, no hay ataque
ejecutable--- y la evaluación del ataque de texto plano
($\approx 2^{266}$ para \texttt{mceliece6960119}). Y la respuesta
fechada de junio~\cite{cm610}. Como individuos: Bernstein mantuvo,
el 30 de septiembre, que ninguna nación-estado tiene recursos para
ejecutar los ataques, que estos no ganan nada con múltiples
objetivos ni con computación cuántica, y que su evaluación relativa
---ML-KEM con más riesgo que McEliece--- no había cambiado. La
pregunta directa de Filippo Valsorda ---¿sigue el equipo
recomendando sus conjuntos estandarizados?--- no había recibido
respuesta etiquetada como del equipo a fecha de cierre de este
libro~\cite{bernstein2026}.

El movimiento decisivo fue del BSI, el 1 de octubre de 2026, en una
nota de cuatro páginas fechada el 29 de septiembre~\cite{bsi2026}.
Su evaluación, en orden: los distinguidores previos no amenazaban
los supuestos (el de sizigias, $\ge 2^{683}$; el de
Hemmert--Wiemers, $\ge 2^{550}$; BLRT, $\ge 2^{742}$); los métodos
de agosto --octubre sí merecen acción, con cifras de $2^{120}$--$2^{124}$
(distinguidor GIJS) y $2^{140}$--$2^{148}$ (recuperación GIJS) para
los conjuntos de su guía, y $2^{99}$--$2^{102}$ en la contabilidad
de Weis; los resultados «no permiten todavía un ataque práctico»,
pero la mejora estructural «es significativa» y se esperan más.
Conclusión operativa: Classic McEliece «no debería seguir usándose
en nuevos desarrollos»; para intercambio de claves, preferir
ML-KEM o FrodoKEM, y HQC cuando su estándar aparezca; los
despliegues híbridos existentes conservan al menos la seguridad de
la mitad clásica; y la recomendación formal se revisará en
TR-02102-1 a principios de 2027. El BSI añadió una lección de
migración: los resultados «subrayan la importancia de la
\emph{crypto-agility}», y ninguno de los ataques toca a los otros
esquemas recomendados ---ni a HQC, que no usa códigos de Goppa ni la estructura
oculta que estos ataques explotan~\cite{bsi2026}.

En paralelo, el ecosistema: NIST dejó la puerta abierta a
estandarizar McEliece más adelante, pero su IR~8545 ya había
elegido HQC y anotado los recortes de \texttt{mceliece460896}; ISO
publicó en junio de 2026 su enmienda con los cuatro conjuntos
grandes ---sin \texttt{mceliece348864}, el más expuesto a los
nuevos métodos--- y cada uno de esos cuatro tiene ahora
estimaciones por debajo de su nivel declarado; las agencias
europeas divergen (ANSSI ya lo situaba por debajo de ML-KEM,
FrodoKEM y HQC; el manual holandés lo tenía por conservador; la
NCSC sueca solo bendice estándares NIST); y los despliegues reales
---túneles VPN de Mullvad, Rosenpass, PQConnect---, todos en modo
híbrido, activaron sus protocolos de revisión~\cite{bernstein2026,bsi2026,nist8545}.

\begin{trampa}[Los dos titulares fáciles]
Desde agosto de 2026 circulan dos frases gemelas, y ambas son
falsas. «\emph{McEliece ha caído}»: nadie ha recuperado una clave
ni descifrado un texto de ningún conjunto real; los propios
autores de los ataques lo repiten; el ataque más caro de ejecutar
---el cómputo de núcleos--- no se ha corrido a escala; los retos
rotos son juguetes. «\emph{McEliece es inmune}»: las estimaciones
de recuperación de clave de tres grupos independientes están por
debajo de los niveles que el propio esquema declaró, el BSI ---que
lo recomendaba desde 2020--- ha pedido detener nuevos despliegues,
y la incógnita no es si los supuestos heurísticos fallarán, sino
cuándo y cuánto. La versión utilizable de la verdad tiene tres
cláusulas y ninguna es un titular: nadie ha roto nada, muchas
estimaciones condicionales han bajado de los niveles declarados, y
la comunidad tarda años en digerir cuánto pesa cada supuesto. Quien
lea una sola de las dos frases, que relea este recuadro.
\end{trampa}

\section{Lectura final}\label{sec:2026-lectura}

¿Qué debe llevarse el lector de este libro de todo el terremoto?
Cuatro cosas, ninguna grandilocuente.

\textbf{Primera: la criptografía finita se juega en los márgenes.}
Los conjuntos de Classic McEliece llevaban cuarenta y ocho años
fuera de rango de todos los ataques estructurales conocidos, y
estaban elegidos contra \emph{un modelo} de ataque (la
decodificación genérica) con márgenes que la propia guía declaraba
para el otro modelo en miles de bits. El método \emph{holdout}
apareció fuera de ese modelo, traído de la teoría de PIR, y la
lección no es «los márgenes eran falsos» sino la del
capítulo~\ref{ch:seguridad}: \emph{seguridad es la mejor estimación
disponible más los supuestos, y los supuestos envejecen}.

\textbf{Segunda: el eje manda.} Weis midió que el coste del
\emph{holdout} crece unos ocho bits por cada unidad del grado $m$
del cuerpo, y apenas depende de $n$: por eso sus conjuntos
recomendados (todos con $m=12$ o $13$) se comportan en bloque. Sus
ilustraciones ---$m=18$, $26$, $34$, con claves de $0{,}6$, $2{,}1$ y
$7{,}1$ megabytes--- y la instancia de Bernstein ($m=18$,
$n=6176$, $t=112$, clave de~1\,MB, texto cifrado de 252 bytes,
ejecución estimada en $2^{142}$) dibujan la dirección de una
futura re-parametrización: cuerpos más grandes, $n$ bien por debajo
de $2^m$, claves más grandes. Nadie ha propuesto tal conjunto para
estandarización, y Weis lo etiqueta «ilustraciones»~\cite{weis2026}.

\textbf{Tercera: la asimetría de las pruebas.} El distinguidor está
probado; la recuperación de clave, no. Los experimentos existen
solo en tamaño de juguete. Quien compare las cotas del capítulo
con las de este debe recordar la dirección del rigor: una
estimación heurística baja no convierte a un esquema en roto ---ni
una superioridad histórica lo convierte en intacto. La única
moneda válida es la misma que usa este libro desde la
página~1: enunciar el modelo, citar la fuente, y separar lo
demostrado de lo conjeturado.

\textbf{Cuarta: esto es historia viva.} A la fecha de este libro
(octubre de 2026), las réplicas siguen circulando (la última
revisión de Vedenev es del 7 de octubre), el BSI revisará su guía a
principios de 2027, ISO mantiene publicado su estándar, y las
estimaciones de los distintos grupos ---dentro del mismo marco---
difieren en decenas de bits. Cualquier lector que llegue a estas
páginas años después debería buscar la actualización de la
guía~\cite{cmsecurity}, las respuestas fechadas del equipo
---la serie de notas continuará--- y las réplicas de las revistas
que aún están por llegar. El resto del libro no depende de la
resolución del terremoto: el mecanismo, los conjuntos, la ISO, el
laboratorio, todo lo anterior sigue siendo exactamente lo que las
páginas describen. Lo que cambia es la conversación sobre si
conviene usarlo, y esa conversación, por primera vez en la historia
del esquema, está abierta.

\section{Ejercicios}

\begin{enumerate}[leftmargin=1.6em,itemsep=4pt]
  \item \textbf{El muro de 2022 contra el suelo de 2026.}
        (a)~Compare la mejor estimación estructural de la
        tabla~\ref{tab:est-espacios} para \texttt{348864}
        ($2^{9116{,}8}$) con la de Weis ($2^{94{,}3}$): calcule la
        diferencia de exponentes y comente en dos líneas por qué
        \emph{no} es el «mismo ataque mejorado» (pista: ¿qué
        familia de métodos estaba agotada y cuál es nueva?).
        (b)~Compare ahora $2^{94{,}3}$ con el suelo trivial $2^{256}$
        de la respuesta del equipo~\cite{cm529}: ¿qué implica que
        la estimación esté por debajo?
  \item \textbf{Una ejecución, cien ejecuciones.} Con la ruta de
        Weis que exige entre 100 y 1400 ejecuciones, calcule el
        coste total si una ejecución cuesta $2^{89}$ (sin memoria)
        y si cuesta $2^{117}$ (con cargo severo). Encaje ambos
        resultados con la afirmación «esta es nuestra mejor
        estimación; con el cargo de la guía queda por debajo de la
        decodificación genérica en ese modelo».
  \item \textbf{El eje m.} Partiendo de $2^{94{,}3}$ ($m=12$) y de
        la regla de Weis (unos 8 bits por unidad de $m$), estime el
        coste para $m=18$ y compárelo con la instancia de Bernstein
        ($n=6176$, $t=112$), cifrada en $\approx 2^{142}$.
        ¿Cuánto sube la clave pública (en MB) frente a las actuales?
  \item \textbf{Distinguir, recuperar, romper.} Clasifique cada
        afirmación como demostrada, heurística o estimada, usando
        el recuadro de rigor: (a)~«las claves de Goppa binaria se
        distinguen de matrices aleatorias con ventaja
        $1-o(1)$»; (b)~«recuperar la clave de \texttt{8192128}
        cuesta $2^{102{,}4}$»; (c)~«el ataque de Vedenev cuesta
        esencialmente una ejecución»; (d)~«nadie ha ejecutado el
        cómputo de núcleos contra un conjunto real».
  \item \textbf{El debate de la memoria, aplicado.} Use la
        analogía de AES del capítulo~\ref{ch:isd} (tabla de
        equivalencias) y explique, en cinco líneas, por qué
        comparar $2^{94{,}3}$ con $2^{143}$ (AES-128) favorece al
        ataque, mientras que comparar $2^{117}$--$2^{128}$ con
        $2^{143}$ favorece al esquema. ¿Qué información falta para
        cerrar la comparación en un sentido o en otro?
  \item \textbf{Los titulares.} Escriba su propio titular de
        longitud máquina (menos de 80 caracteres) para la noticia
        del 1 de octubre de 2026, que sea \emph{verdadero} según el
        recuadro de trampa, y otro deliberadamente falso en cada
        dirección. Intercámbielos con un colega y discutan qué
        cláusula falta en cada uno.
\end{enumerate}

\begin{center}\rule{0.55\textwidth}{0.4pt}\end{center}
\medskip
\noindent\textit{Cierre de la parte~\ref{part:seguridad}.} Cuatro
capítulos para responder a la pregunta «¿qué significa romper
McEliece?»: los objetivos y la cadena de reducciones
(capítulo~\ref{ch:seguridad}); la trinchera de la decodificación
genérica y sus pliegos de costes (capítulo~\ref{ch:isd}); el muro
estructural y su anatomía (capítulo~\ref{ch:estructura}); y el
terremoto de 2025--2026 con su balance aún abierto
(capítulo~\ref{ch:2026}). Nadie ha roto Classic McEliece ---no
existe un ataque práctico, y cada estimación baja de los niveles
declarados viaja con supuestos explícitos--- pero las cifras de esa
orilla han cambiado de siglo en cuatro años, y con ellas la
conversación pública. La parte final del libro vuelve al suelo: las
implementaciones, el ecosistema que ya existe, y el balance de lo
que este esquema significa para quien lee.


% ======================================================================
%  CAPÍTULO 17 — PREGUNTAS FRECUENTES
% ======================================================================

\part{Cierre}\label{part:cierre}

\chapter{Preguntas frecuentes}\label{ch:faq}

\begin{roadmap}
Catorce capítulos de teoría, dos laboratorios y dos años de
literatura convulsa dejan, al final, las preguntas de siempre: ¿qué
es esto?, ¿por qué así?, ¿está roto?, ¿lo uso? Este capítulo las
responde en formato breve y con punteros exactos a las partes donde
cada respuesta está desarrollada. Si solo se va a leer una sección
del libro después de la introducción, que sea esta; si algo de aquí
resulta sorprendente o insuficiente, los punteros llevan al capítulo
correspondiente.
\end{roadmap}

\section{El esquema, en preguntas cortas}

\begin{enumerate}[leftmargin=1.9em,itemsep=5pt]
  \item \textbf{¿Qué es Classic McEliece, en una frase?} Un KEM
        post-cuántico: el criptosistema de McEliece (1978) en la
        forma dual de Niederreiter, con códigos de Goppa binarios,
        estandarizado por ISO en 2026. Es famoso por dos extremos:
        las claves públicas más grandes de la criptografía
        moderna---de $261\,120$ a $1\,357\,824$ bytes---y algunos de
        los criptogramas más pequeños ($96$--$208$ bytes). Su
        seguridad se ancla en la dureza conjetural de la
        decodificación por síndromes, un problema bajo escrutinio
        desde 1962~\cite{cmspec,cmsecurity}.

  \item \textbf{¿Es «el McEliece de 1978» o es otra cosa?} El mismo
        esqueleto (códigos de Goppa, trampa de Patterson), con tres
        piezas reingenierizadas: la forma dual, donde el mensaje es
        un vector de error de peso fijo y no un texto; la forma
        sistemática con semilla de 32 bytes, que reduce el
        almacenamiento y hace reproducible la generación
        (capítulo~\ref{ch:keygen}); y el transformado a KEM con
        rechazo implícito (capítulos~\ref{ch:kem}
        y~\ref{ch:encapdecap}). La trampa---decodificar es fácil con
        $g$ y $\alpha$ y es el problema duro sin
        ellos---es exactamente la de~\cite{mceliece1978}.

  \item \textbf{¿Por qué las claves públicas son gigantescas?}
        Porque la clave pública es, precisamente, una matriz binaria
        aleatoria de $mt\times k$ bits (la $T$ de
        $(\mathrm{I}_{mt}|T)$). No hay anillo ni estructura que la
        comprima, y esa ausencia de estructura es parte del argumento
        de seguridad: menos álgebra visible, menos superficie para el
        criptoanálisis estructural (capítulo~\ref{ch:estructura}).
        El contraste está en la tabla~\ref{tab:param-contexto}:
        frente a ML-KEM-768, entre $400$ y $1\,200$ veces las claves
        (según el conjunto); a cambio, criptogramas entre $5$ y $8$
        veces más pequeños~\cite{cmsecurity}.

  \item \textbf{¿Cifra mensajes o claves?} Es un KEM: acuerda una
        clave de sesión de $32$ bytes, que después cifra el mensaje
        en un esquema simétrico (AES-GCM, ChaCha20). En un protocolo
        real se combina en híbrido con un intercambio clásico
        (X25519): si un lado cae, el otro sostiene; los despliegues
        conocidos usan todos el modo
        híbrido~\cite{cmimpl,bernstein2026}.

  \item \textbf{¿Qué tan rápido es?} En los términos de conteo del
        capítulo~\ref{ch:encapdecap}: la encapsulación es una
        multiplicación matriz-vector sobre la clave pública
        (XORs; $\sim2^{23{,}4}$ operaciones de bit en el conjunto
        mayor) más los hashes; la desencapsulación añade Patterson.
        Lo caro es la \emph{generación} (reintentos de
        $\textsf{MatGen}$: $\sim3{,}4$ de media en la forma base,
        una sola en las variantes \texttt{f}), pero ocurre una vez
        por clave~\cite{cmrationale}.

  \item \textbf{¿Por qué el «mensaje» es un vector de error?} Porque
        el problema duro de un código es el síndrome: dada
        $C=He$, recuperar $e$ es exactamente
        decodificar (capítulo~\ref{ch:isd}). La especificación no
        disimula el corazón del problema: se lo pide al atacante en
        su forma desnuda, y restaura la uniformidad con el
        transformado a KEM (capítulos~\ref{ch:kem}
        y~\ref{ch:encapdecap}).

  \item \textbf{¿Qué hay dentro de una clave secreta?} En el
        conjunto mayor, $14\,120$ bytes: los $12\,800$ que codifican
        el orden de los elementos $\alpha$ (red de Benes), los
        $1\,024$ del vector de rechazo $s$, los $256$ del polinomio
        $g$ y las semillas ($40$). El orden de $\alpha$, no el
        código---que es público---, es el secreto
        (capítulo~\ref{ch:keygen}).
\end{enumerate}

\section{Ecosistema y estándares}

\begin{enumerate}[leftmargin=1.9em,itemsep=5pt,start=8]
  \item \textbf{¿Está estandarizado?} Sí, por ISO/IEC: la
        Enmienda~2 a la 18033-2 (junio de 2026) incluye Classic
        McEliece con cuatro tamaños y cuatro variantes---$16$
        perfiles---junto a ML-KEM y FrodoKEM~\cite{iso18033,cmiso}.
        El NIST \emph{no} lo estandarizó: su proceso de KEMs cerró
        en marzo de 2025 con HQC, y en 2025 aplazó la decisión sobre
        McEliece citando la estandarización ISO en
        curso~\cite{nist8545,cmnist}. Norma ISO hay; FIPS, no.

  \item \textbf{¿Dónde se usa hoy?} En despliegues ligados a la
        transición post-cuántica: túneles VPN (Mullvad), Rosenpass,
        PQConnect, y en el ecosistema de implementaciones
        liboqs/PQClean y SUPERCOP~\cite{bernstein2026,cmimpl}.
        Todos en modo híbrido.

  \item \textbf{¿Por qué el NIST lo dejó fuera?} Por ingeniería, no
        por criptoanálisis: cientos de kilobytes de clave pública no
        encajan en sus casos de uso objetivo (TLS de alto
        volumen), y en 2025 pesó la coexistencia con ISO. Es el
        patrón de toda la historia del esquema: seguridad cómoda,
        ergonomía incómoda (capítulo~\ref{ch:parametros}).

  \item \textbf{¿Qué fue «You Only Decapsulate Once»?} Un trabajo de
        2025 que afirmaba recuperar la clave privada mediante un ataque de
        canal lateral de tiempo con una única traza de desencapsulación,
        explotando que optimizaciones de compiladores modernos (como clang)
        convertían código de enmascaramiento en tiempo constante en
        bifurcaciones condicionales. El equipo respondió el 20 de noviembre de
        2025~\cite{cm2025yodo} señalando que el análisis apuntaba a una
        versión obsoleta de PQClean, que ellos mismos habían alertado
        públicamente de ese problema en 2024 y que el software oficial ya
        utilizaba la biblioteca \texttt{cryptoint} y las pruebas automáticas
        de TIMECOP para prevenirlo (capítulo~\ref{ch:2026}).
\end{enumerate}

\section{Seguridad: qué se sabe y qué no}

\begin{enumerate}[leftmargin=1.9em,itemsep=5pt,start=12]
  \item \textbf{¿Está roto?} No. A octubre de 2026 no hay ataque
        práctico contra ningún conjunto, ni recuperación de clave
        ejecutada contra una clave real, ni rotura de los retos de
        la especificación. Pero el tono cambió: aparecieron
        distinguidores que alcanzan el régimen real (GIJS, BLRT,
        Hemmert--Wiemers) y estimaciones de recuperación muy por
        debajo de ciertas cifras declaradas---de $2^{99}$--$2^{102}$
        en la contabilidad de una sola ejecución de Weis a
        $2^{140}$--$2^{148}$ en las estimaciones del BSI---frente a
        los $2^{140{,}8}$--$2^{275{,}6}$ del modelo clásico de
        memoria libre según el conjunto~\cite{ghoshal2026,weis2026,bsi2026}.

  \item \textbf{¿Qué significa «distinguir», exactamente?} Que un
        algoritmo decide si una matriz pública es de Goppa o
        aleatoria. No recupera $e$ ni la clave: es una propiedad más
        débil (figura~\ref{fig:seg-cadena}). Durante décadas los
        distinguidores vivieron en regímenes que McEliece no usa;
        los de 2026 alcanzan los parámetros reales, aunque con
        costes aún fuera de alcance y sin reducción conocida de
        distinguir a romper~\cite{ghoshal2026}.
        Capítulos~\ref{ch:estructura} y~\ref{ch:2026}.

  \item \textbf{¿Y la amenaza cuántica?} No hay rotura cuántica
        conocida: Shor no aplica (no hay grupo abeliano oculto), y
        los mejores ataques cuánticos ganan a lo sumo la mitad del
        exponente~\cite{kachigar2017}. ISO partió esos exponentes y
        exigió $\ge2^{128}$: excluyó a \texttt{mceliece348864} y
        admitió los demás con matices~\cite{iso18033,cmiso}
        (capítulo~\ref{ch:parametros}).

  \item \textbf{¿Qué dice el BSI?} Su nota del 29 de septiembre de
        2026 es la advertencia institucional más dura recibida: los
        métodos de agosto--octubre merecen tomarse en serio; «no
        debería seguir usándose en nuevos desarrollos»; los
        despliegues híbridos existentes conservan la mitad clásica;
        recomienda ML-KEM, FrodoKEM y HQC, y revisará su
        TR-02102-1~\cite{bsi2026}.

  \item \textbf{¿Y qué responde el equipo?} Que ningún ataque
        completo está verificado; que las cifras agresivas dependen
        de contabilidades discutibles (operaciones de bits, memoria
        gratis, pasos heurísticos no ejecutados); y que el suelo
        trivial de recuperación de clave es $2^{256}$, la búsqueda
        exhaustiva de la semilla de 32 bytes~\cite{cm529,cm610,bernstein2026}.
        El desacuerdo es técnico, público y está abierto
        (capítulo~\ref{ch:2026}).
\end{enumerate}

\section{Práctica: elegir, desplegar, verificar}

\begin{enumerate}[leftmargin=1.9em,itemsep=5pt,start=17]
  \item \textbf{¿Qué conjunto y qué variante elegir?} La
        recomendación del equipo es simétrica y simple: los seis
        dígitos para uso a largo plazo (\texttt{mceliece6688128},
        \texttt{mceliece6960119}); tamaños intermedios solo con
        presupuesto ajustado; \texttt{f} si el coste de generación
        importa ($<2^{-30}$ de reintento, sin efecto apreciable en
        seguridad); \texttt{pc}/\texttt{pcf} si se quiere
        confirmación de conocimiento del error ($+32$
        bytes)~\cite{cmsecurity,cmiso}. La guía completa está en el
        capítulo~\ref{ch:parametros}.

  \item \textbf{¿Cómo verifico una implementación?} Con los
        vectores de prueba de la especificación: la clave entera es
        reproducible bit a bit desde la semilla de 32 bytes, así
        que cualquiera regenera clave y criptograma y compara. Para
        C portable están PQClean y liboqs; para medidas
        comparables, SUPERCOP~\cite{cmimpl}.

  \item \textbf{¿Y si tengo que decidir hoy?} Tres reglas
        prudentes, coherentes con la nota del
        BSI~\cite{bsi2026}: (1)~en entornos con guía BSI, no
        iniciar despliegues nuevos y preferir ML-KEM, FrodoKEM o
        HQC; (2)~donde ya haya Classic McEliece, mantener el modo
        híbrido y monitorizar la literatura (la mitad clásica cubre
        la eventualidad de una degradación del componente
        post-cuántico); (3)~si el diseño exige diversidad de
        familias, tratar los tamaños \texttt{6*} como el mínimo
        razonable y revisar los resultados de 2025--2026 antes de
        congelar parámetros~\cite{cm610,cmsecurity}.
\end{enumerate}

\begin{observacion}{La respuesta honesta a «¿lo uso?»}{faq-uso}
Ningún documento de 2026---ni los ataques, ni la nota del BSI, ni
las réplicas del equipo---afirma que Classic McEliece esté roto
(octubre de 2026). Y ninguno autoriza a ignorar la tendencia. La
cifra que mejor resume el momento es la del capítulo~\ref{ch:2026}:
las estimaciones agresivas bajaron de $2^{9116}$ a $2^{94}$--$2^{102}$
en \emph{una parte} de la contabilidad, mientras la otra sigue en
$2^{142}$--$2^{150}$; la distancia entre ambas contabilidades es,
literalmente, un problema abierto. Un usuario racional hace lo que
hacen los usuarios reales: híbrido, tamaños grandes, seguimiento
activo, y cero apuestas a que el desacuerdo se resuelva a su favor.
\end{observacion}

% ======================================================================
%  CAPÍTULO 18 — REPASO Y SOLUCIONES (parte 1: repaso)
% ======================================================================

\chapter{Repaso y soluciones}\label{ch:soluciones}

\begin{roadmap}
El libro termina como empezó: con lápiz y papel. La
sección~\ref{sec:sol-repaso} propone doce preguntas de repaso que
recorren el argumento completo de las partes~I--IV: quien las
responda de memoria puede considerar que entendió el esquema. La
sección~\ref{sec:sol-soluciones} da las respuestas concisas. Y la
sección~\ref{sec:sol-ejercicios} resuelve, capítulo a capítulo, los
$94$ ejercicios del libro, con las cifras ya calculadas y
verificadas. En los laboratorios, el programa del
apéndice~\ref{app:codigo} reproduce todos los números.
\end{roadmap}

\section{Repaso: doce preguntas}\label{sec:sol-repaso}

\begin{enumerate}[leftmargin=1.9em,itemsep=5pt]
  \item Defina el problema SD en su versión de búsqueda y el juego
        OW-CPA del cifrado de Niederreiter, y enuncie su equivalencia
        (capítulo~\ref{ch:seguridad}).
  \item ¿Por qué los conjuntos de Classic McEliece pueden tomar
        $n<q$? Dé las dos caras: el acortamiento y la división del
        soporte.
  \item Escriba las fórmulas de tamaño del criptograma, la clave
        pública y la clave secreta de un conjunto $(m,t,n,k)$, y
        diga qué componente domina la clave secreta.
  \item Enumere los cuatro pasos de $\textsf{Decode}$ y explique por
        qué los dos chequeos finales (peso y síndrome) son
        imprescindibles para el KEM.
  \item ¿Con qué probabilidad tiene éxito $\textsf{MatGen}$ con
        $(\mu,\nu)=(0,0)$? ¿Qué compra y qué cuesta la forma
        semi-sistemática $(32,64)$?
  \item Describa la tubería $\textsf{Encap}/\textsf{Decap}$ con
        rechazo implícito y diga qué es exactamente $s$.
  \item Enuncie los tres términos de la cota ROM del
        teorema~\ref{thm:kem-rom} y evalúelos para
        \texttt{mceliece348864} con $q_D=q_H=2^{64}$.
  \item Distinga las tres familias de ataque que aparecen en el
        libro (ISD, estructurales, distinguidores) con un método y
        un capítulo para cada una.
  \item Escriba la probabilidad de éxito de una iteración de
        Prange, evalúela para $(n,k,t)=(3488,2720,64)$ y componga
        el coste total. ¿De dónde sale el medio bit de diferencia
        con la guía?
  \item Explique «$o(1)$ no es $0$» y su consecuencia para comparar
        Prange con May--Ozerov.
  \item ¿Qué cuenta exactamente el «muro» $2^{9116{,}8}$ y qué
        \emph{no} es? ¿A qué distancia quedó de él la familia de
        2025--2026?
  \item Resuma el año 2025--2026 en cuatro cifras y una moraleja
        (capítulo~\ref{ch:2026}).
\end{enumerate}

\section{Soluciones del repaso}\label{sec:sol-soluciones}

\begin{enumerate}[leftmargin=1.9em,itemsep=5pt]
  \item \textbf{SD y OW-CPA.} SD (búsqueda): dadas $H\in\Ftwo^{mt
        \times n}$, $C\in\Ftwo^{mt}$ y $t$, hallar $e$ de peso
        $\le t$ con $He=C$; encontrar el criptograma que genera
        una clave pública. OW-CPA: dados $(T,\,C=\textsf{Encode}(e,T))$
        con $e$ uniforme de peso $t$, hallar $e$. La
        equivalencia (lema del capítulo~\ref{ch:seguridad}): el
        cifrado es una biyección lineal en $e$, así que invertirlo
        es exactamente decodificar el síndrome; sin decodificador
        no hay inversión y viceversa.

  \item \textbf{$n<q$.} (a)~Acortamiento: cualquier ataque contra
        $(q,n,t)$ se transporta a un ataque contra $(q,q,t)$ con el
        éxito multiplicado por $\binom{n}{t}/\binom{q}{t}$; la pérdida
        de seguridad del traslado es de fracciones de bit (unos
        $0{,}8$ bits para $460896$ en el modelo $1/2$): bajar $n$ no
        regala seguridad contra ISD
        (capítulo~\ref{ch:isd}). (b)~División del soporte: los
        ataques que adivinan los elementos $\alpha$ deben elegir
        $n$ de $q$ columnas, y pagan $\binom{q}{n}$: tomar
        $n<q$ los encarece (capítulo~\ref{ch:estructura}). El
        diseño cobra en el segundo frente lo que no pierde en el
        primero.

  \item \textbf{Tamaños.} $ct=\lceil mt/8\rceil$ bytes;
        $ek=mt\lceil k/8\rceil$ (la $T$ de
        $(\mathrm{I}_{mt}|T)$, una fila por byte redondeado);
        $sk=32+8+t\lceil m/8\rceil+\lceil(2m-1)2^{m-4}\rceil
        +\lceil n/8\rceil$ (semilla, contador, $g$, orden de
        $\alpha$ en red de Benes, $s$). En \texttt{8192128}:
        $14120$ bytes, dominados por los $12800$ de $\alpha$;
        la pública es toda $T$ (capítulo~\ref{ch:keygen}).

  \item \textbf{Decode.} (1)~$v=(C,\mathbf{0})$; (2)~el único
        $c\in\Gamma_2$ con $Hc=0$ y $d(v,c)\le t$ (Patterson con
        $g$ y el orden de $\alpha$); si no existe, $\bot$;
        (3)~$e=v+c$; (4)~aceptar si y solo si $\wt(e)=t$ y
        $He=C$. Los chequeos caracterizan la imagen de
        $\textsf{FixedWeight}$: sin ellos, entrarían errores de
        peso menor y el «mensaje» aceptado dejaría de ser uniforme en
        el conjunto de peso $t$ que exige la reducción.

  \item \textbf{MatGen.} Con $(0,0)$ el éxito es
        $\prod_{i=0}^{mt-1}(1-2^{i-mt})\approx0{,}289$: unos
        $3{,}4$ intentos de media. Con $(32,64)$, el fallo por
        intento cae bajo $2^{-30}$ ($\approx1$ intento), al
        precio de $\le2$ bits de tensión en la reducción y de
        cómputo en tiempo constante dentro de la ventana
        (capítulos~\ref{ch:keygen} y~\ref{ch:parametros}).

  \item \textbf{Encap/Decap.}
        $\textsf{Encap}$: $e\leftarrow\textsf{FixedWeight}$,
        $C=\textsf{Encode}(e,T)$, $K=H(1,e,C)$.
        $\textsf{Decap}$: $e'\leftarrow\textsf{Decode}(C)$; si
        $e'=\bot$ entonces $e'\leftarrow s$, $b\leftarrow0$, si no
        $b\leftarrow1$; $K=H(b,e',C)$. $s$ es un vector uniforme
        de $n$ bits guardado en la clave secreta: la clave de la
        rama de rechazo. El rechazo implícito nunca devuelve
        error: dos criptogramas (válido e inválido) producen
        siempre $32$ bytes pseudoaleatorios (capítulo~\ref{ch:kem}).

  \item \textbf{Cota ROM.}
        $\mathrm{Adv}\le\Pr[\text{OW-CPA}]+\dfrac{q_D}{|P|}
        +\dfrac{2q_H}{|S|}$. Para \texttt{348864}:
        $2^{64}/2^{456}=2^{-392}$ y
        $2^{65}/2^{256}=2^{-191}$: el tercer término domina (acotado
        por los $256$ bits de entropía de la semilla $\delta$ que genera
        a $s\in\Ftwo^{3488}$) y aun así es despreciable; el verdadero
        supuesto es la dureza OW-CPA, la conjetura de SD.

  \item \textbf{Tres familias.} (1)~Decodificación genérica: Prange,
        Stern, May--Ozerov\dots (capítulo~\ref{ch:isd}): atacan el
        problema SD con álgebra lineal y listas. (2)~Estructurales:
        división del soporte, álgebra sobre el código, distinguidores
        con inversión (Faug\`ere, Carrier, Kirshanova, el
        \emph{holdout} de GIJS\dots) (capítulos~\ref{ch:estructura}
        y~\ref{ch:2026}): buscan $g$ y el orden de $\alpha$, o
        recuperan $e$ estructuralmente. (3)~Distinguidores
        (capítulo~\ref{ch:estructura}): deciden si la clave es de
        Goppa; no rompen por sí solos, pero erosionan supuestos y,
        desde 2026, operan en el régimen real.

  \item \textbf{Prange.} Una iteración elige $k$ columnas y espera
        que las $t$ posiciones del error caigan fuera:
        $\bigl((n-k)/n\bigr)^{t}=(768/3488)^{64}=2^{-139{,}73}$;
        con coste por iteración $\approx k^3=2^{34{,}23}$, el
        total es $2^{173{,}96}\approx2^{174{,}0}$. La guía cotiza
        $173{,}4$: el medio bit restante son las constantes de la
        eliminación (pivotes, invertibilidad) que el estimador
        fino sí cuenta (capítulo~\ref{ch:isd}).

  \item \textbf{$o(1)$ no es $0$.} Asintóticamente, Prange y las
        mejoras comparten exponente $\bigl(1/(1-R)\bigr)^{t}$:
        difieren solo en el $o(1)$. En concreto, ese $o(1)$ no
        es pequeño: para $348864$, Prange cuesta $2^{173{,}4}$ y
        May--Ozerov $2^{140{,}8}$ ($32{,}6$ bits, un factor
        $6{,}5\times10^{9}$). Las dos afirmaciones son ciertas a la vez:
        misma pendiente asintótica, constantes muy distintas; de
        ahí que las tablas se construyan con estimadores concretos
        (capítulo~\ref{ch:isd}).

  \item \textbf{El muro.} Es $\log_2$ del número de datos
        parciales $q(q-1)\cdots(q-769+1)\approx2^{9116{,}8}$ que
        la familia de ataques de 2022 necesitaría inspeccionar
        para \texttt{348864}: un muro \emph{de su modelo}, no de
        la criptografía. La familia de 2025--2026 (el
        \emph{holdout} y sus derivados) no lo escala: lo evita,
        cambiando el objeto de la búsqueda, y sus estimaciones
        bajan a $2^{94}$--$2^{150}$ según la contabilidad
        (capítulos~\ref{ch:estructura} y~\ref{ch:2026}).

  \item \textbf{Cuatro cifras y una moraleja.} (i)~$2^{9116{,}8}\to
        2^{94{,}3}$: la caída del modelo estructural de 2022;
        (ii)~$2^{120}$--$2^{124}$: el coste cotizado de los
        distinguidores de 2026 en el régimen real; (iii)~$2^{99}$
        --$2^{102}$ frente a $2^{140}$--$2^{148}$: las dos
        contabilidades que no acuerdan; (iv)~$2^{256}$: el suelo
        trivial (buscar la semilla). Moraleja: a octubre de 2026
        nadie ha roto Classic McEliece ---ni un ataque práctico,
        ni una clave real recuperada--- y nadie puede ya decir
        «inmune»: hay conjeturas con supuestos, contabilidades en
        disputa y vigías institucionales (capítulo~\ref{ch:2026}).
\end{enumerate}

\section{Soluciones de los ejercicios de los capítulos 1--16}
\label{sec:sol-ejercicios}

Las páginas que siguen resuelven, capítulo a capítulo, los $94$
ejercicios del libro. Las cifras están calculadas con la aritmética
exacta de los apéndices (y de \texttt{toy.py} cuando el ejercicio es
de laboratorio); los argumentos completos están en el texto del
capítulo correspondiente.

% ======================================================================
%  CAPÍTULO 18 — SOLUCIONES (parte 2: capítulos 1-8)
% ======================================================================

\subsection*{Capítulo 1 --- La amenaza cuántica y el giro
post-cuántico (capítulo~\ref{ch:pqc})}

\begin{enumerate}[leftmargin=1.9em,itemsep=4pt]
  \item Grover busca la clave en $2^{k/2}$ evaluaciones: AES-256
        cae a $2^{128}$ (dentro del listón) y AES-128 a $2^{64}$
        (muy por debajo). La aceleración cuántica es cuadrática en el
        exponente, no total: basta doblar el tamaño de clave.
  \item \emph{Criptogramas}, de menor a mayor: \texttt{mceliece460896}
        ($156$~B) $<$ ML-KEM-768 ($1088$~B) $<$ HQC-128
        ($\approx4497$~B). \emph{Claves públicas}: ML-KEM-768
        ($1184$~B) $<$ HQC-128 ($\approx2249$~B) $<$
        \texttt{mceliece460896} ($\num{524160}$~B). McEliece
        intercambia los extremos: el criptograma más pequeño y la
        clave más grande.
  \item Porque el tráfico cifrado \emph{hoy} ya se graba: cuando
        exista el ordenador cuántico, se descifrará retroactivamente.
        Los datos de vida útil larga están expuestos desde el momento
        en que se capturan; migrar después no repara lo ya grabado, de
        ahí actuar ahora (híbridos incluidos).
  \item \textnormal{(}Investigación.\textnormal{)} La nota del BSI
        (1 de octubre de 2026, \cite{bsi2026}) sostiene, en esencia:
        (a)~los métodos de agosto--octubre (distinguidores de
        $2^{120}$--$2^{124}$; recuperaciones estimadas de
        $2^{140}$--$2^{148}$; el cálculo de una sola ejecución de
        Weis, $2^{99}$--$2^{102}$) estrechan el margen lo bastante
        para tomarse en serio, aunque nada esté roto; (b)~principio de
        precaución: no iniciar despliegues nuevos, preferir ML-KEM,
        FrodoKEM y HQC, y conservar la mitad clásica de los híbridos
        existentes, revisando la TR-02102-1. (Respuesta abierta: lo
        evaluable es distinguir lo demostrado de lo estimado y leer
        la recomendación como decisión de política.)
\end{enumerate}

\subsection*{Capítulo 2 --- Códigos lineales
(capítulo~\ref{ch:lineales})}

\begin{enumerate}[leftmargin=1.9em,itemsep=4pt]
  \item Si $\vx$ y $\vz$ difieren en la posición $j$, entonces
        $\vx_j\neq\mathbf{y}_j$ o $\mathbf{y}_j\neq\vz_j$ (si ambas igualdades se
        dieran, $\vx_j=\vz_j$). Así
        $\supp(\vx+\vz)\subseteq\supp(\vx+\mathbf{y})\cup\supp(\mathbf{y}+\vz)$,
        y tomando cardinales se obtiene la desigualdad triangular.
  \item Una base: las filas $1100$ y $0011$ sirven como $G$ y también
        como $\Hmat$ (las comprobaciones son $x_1+x_2=0$ y
        $x_3+x_4=0$). Las palabras no nulas son $1100$, $0011$ y
        $1111$, de pesos $2$, $2$ y $4$: $d=2$.
  \item Para $y=0101010$: síndrome $=\Hmat_2+\Hmat_4+\Hmat_6
        =011+101+010=100=\Hmat_5$: se corrige la posición $5$
        ($0$-indexada $4$): $\hat y=0101110$, mensaje $0101$. Para
        $y'=1100110$: síndrome $=\Hmat_1+\Hmat_2+\Hmat_5+\Hmat_6
        =110+011+100+010=011=\Hmat_2$: se corrige la posición $2$:
        $\hat y'=1000110$, mensaje $1000$.
  \item Hay $\binom{n}{t}$ palabras de peso $t$: $\binom71=7$ para
        el Hamming; $\binom{3488}{64}\approx2^{456{,}3}$ para
        \texttt{348864}: unos $10^{137}$ candidatos.
  \item Sea $e''$ de peso mínimo con $\Syn(e'')=\Syn(y)$
        (existe: al menos $e$). Entonces $y+e''$ es un codeword (su
        síndrome es $0$) y todo codeword $c$ cumple
        $\wt(y+c)\ge\wt(e'')$, porque $y+c$ tiene el mismo síndrome
        que $e''$ y $e''$ es el mínimo. Así el codeword más cercano
        dista de $y$ exactamente $\wt(e'')$, que es a lo sumo
        $\wt(e)$.
  \item Un código sirve para «publicar $G'=SGP$» si tiene (a)
        decodificación eficiente hasta $t$ (sin la descripción
        secreta no; con ella, Patterson) y (b) ocultación: el disfraz
        hace que $G'$ parezca una matriz aleatoria. Los dos recursos
        de quien descifra son el computacional (Patterson) y el
        algebraico (la descripción $g$ y el orden de $\alpha$, que
        convierten $G'$ en un código con decodificador).
\end{enumerate}

\subsection*{Capítulo 3 --- Cuerpos finitos
(capítulo~\ref{ch:cuerpos})}

\begin{enumerate}[leftmargin=1.9em,itemsep=4pt]
  \item $z^5=z\cdot z^4=z(z+1)=z^2+z$; y
        $z^{-5}=z^{15-5}=z^{10}=z^2+z+1=\hx7$.
  \item $\Tr(1)=1+1+1+1=0$; $\Tr(z)=z+z^2+z^4+z^8
        =z+z^2+(z+1)+(z^2+1)=0$. Tienen traza $1$ exactamente la
        mitad: $8$ elementos (el núcleo de $\Tr$ es un subespacio
        de dimensión $m-1=3$).
  \item $\sqrt{z^3+z^2}$: como $z^3+z^2=z^6$ (pues
        $z^6=z^4z^2=(z+1)z^2=z^3+z^2$), la raíz es $z^3$.
        $\sqrt{z^2+z+1}$: como $z^2+z+1=z^{10}$, la raíz es
        $z^5=z^2+z=\hx6$. En general, $\sqrt{\beta}=\beta^{2^{m-1}}$
        (para $m=4$, elevar al exponente $8$).
  \item $\Tr(\beta^2)=\beta^2+\beta^4+\beta^8+\beta^{16}
        =\beta^2+\beta^4+\beta^8+\beta=\Tr(\beta)$ (en $\FF4$,
        $\beta^{16}=\beta$). La linealidad sobre $\Ftwo$: Frobenius da
        $(\beta+\gamma)^{2^i}=\beta^{2^i}+\gamma^{2^i}$ para cada
        $i$, y sumar conserva la aditividad; $\Tr(a\beta)=a\Tr(\beta)$
        para $a\in\Ftwo$ es inmediato.
  \item $x^2+x+1$ no tiene raíces en $\Ftwo$ ($P(0)=P(1)=1$), luego
        es irreducible; $x^2+x=x(x+1)$ es reducible. El criterio de la
        traza ($x^2+x+a$ irreducible $\iff\Tr(a)=1$) coincide:
        $\Tr(1)=1$ y $\Tr(0)=0$. (Sobre $\FF4$ ambos se parten:
        $x^2+x+1$ tiene las raíces de orden $3$.)
  \item Suma: un XOR por bit. Producto: red de AND+XOR de ancho
        $\approx m^2$. Raíz cuadrada: casi gratis, es $\Ftwo$-lineal
        (una matriz de XORs). Inversa: la más cara (cadena de
        elevar-al-cuadrado y multiplicar, $x\mapsto x^{q-2}$; o
        Itoh--Tsujii; o tablas), tanto por profundidad como por
        dependencia de datos. Patterson la usa $O(t)$ veces.
\end{enumerate}

\subsection*{Capítulo 4 --- Códigos de Goppa
(capítulo~\ref{ch:goppa})}

\begin{enumerate}[leftmargin=1.9em,itemsep=4pt]
  \item Si $g(\alpha_i)=0$, entonces $(x+\alpha_i)$ divide a $g$ y
        $1/(x+\alpha_i)$ no existe módulo $g$ (es un divisor de
        cero); la columna de síndrome
        $\sigma_i=(x+\alpha_i+1)/g(\alpha_i)$ no se puede calcular.
        Equivalente: todas las condiciones se anulan en la posición
        $i$ (el término $c_i/(x+\alpha_i)$ queda libre): aparecerían
        palabras de peso $1$ y la distancia se derrumba.
  \item $A(x)/(x-\alpha_i)=\prod_{j\neq i}(x-\alpha_j)$: evaluado
        en $\alpha_i$ da $\prod_{j\neq i}(\alpha_i-\alpha_j)$. Y
        $A'(x)$ es la suma de los productos que omiten un factor;
        al evaluar en $\alpha_i$, los términos que conservan el
        factor $(x-\alpha_i)$ se anulan, y queda exactamente el
        mismo producto. Es la identidad que traduce la suma de
        fracciones a la fórmula polinomial de la especificación.
  \item Con el ordenamiento natural del juguete ($\alpha_0=0$,
        $\alpha_1=1$): $g(0)=g(1)=z^3$, luego
        $1/g(0)=z^{-3}=z^{12}=\hx f$; la columna $0$ es
        $(\hx f,0)$ y la $1$ es $(\hx f,\hx f)$. (La
        tabla~\ref{tab:lab-columnas} del laboratorio muestra el
        resto.)
  \item $d\ge2t+1$ dice exactamente que la bola de radio $2t$
        alrededor de $0$ solo contiene al codeword nulo: ningún peso
        $\le2t$. No contradice $d>2t+1$ porque es una cota
        \emph{inferior}: el lema garantiza al menos $2t+1$; la
        distancia real puede ser mayor (en el juguete es $5=2t+1$,
        pero la cota no prohíbe $6$ o más en otros códigos).
  \item Con $g=h^2$: la hipótesis del
        teorema~\ref{thm:goppa-2t} (libre de cuadrados) falla, pero la
        respuesta de fondo es la identidad de van Lint
        $\Gamma_2(L,g)=\Gamma_2(L,g^2)$: tomando $g=h$ queda
        $\Gamma_2(L,h^2)=\Gamma_2(L,h)$. El código con $g=h^2$ es
        \emph{el mismo} que con $h$: el «$t$ efectivo» sigue siendo
        $\deg h$, no se gana distancia ni dimensión. McEliece nunca
        lo intentó porque duplicar el grado nominal no cambia el
        código: promesa vacía.
  \item Semillas $\to$ $g$ (con $\textsf{Irreducible}$) y orden de
        $\alpha$ (con $\textsf{FieldOrdering}$) $\to$ el código
        $\Gamma_2(L,g)=\{c:\sum_j c_j/(x+\alpha_j)\equiv0\pmod g\}$
        $\to$ la tabla de síndromes $\tilde{\Hmat}$ ($2\times n$
        sobre $\FF4$, columnas $\alpha_i^j/g(\alpha_i)$) $\to$ la
        matriz binaria $\hat{\Hmat}$ ($8\times16$) $\to$
        eliminación gaussiana $\to$ $(\mathrm{I}_8\,|\,T)$. Se
        publica solo $T$ (junto con $t$); se mantienen secretos $g$,
        el orden $\alpha$ (y las semillas, y $s$ en el KEM). El
        código $\ker(\mathrm{I}_{mt}|T)$ es público.
\end{enumerate}

\subsection*{Capítulo 5 --- Laboratorio 1
(capítulo~\ref{ch:labgoppa})}

\begin{enumerate}[leftmargin=1.9em,itemsep=4pt]
  \item $g(z)=z^2+z+z^3=z^3+z^2+z=\hx e$; $1/\hx e=z^{-11}=z^4
        =z+1=\hx3$; $z\cdot\hx3=z^2+z=\hx6$. La columna $j=2$ de
        la tabla~\ref{tab:lab-columnas} es exactamente $(\hx3,\hx6)$:
        coinciden.
  \item $\tilde{\Hmat}e_5=\tilde{\Hmat}_5=(\hx8,\hx e)$ (pues
        $g(z^2+1)=\hx f$, $1/\hx f=z^3=\hx8$ y
        $(z^2+1)z^3=\hx e$); $\tilde{\Hmat}e_9=(\hx a,\hx5)$ (pues
        $g(z^3+1)=z^6$, $1/z^6=z^9=\hx a$ y
        $(z^3+1)z^9=\hx5$). La suma es
        $(\hx8+\hx a,\hx e+\hx5)=(\hx2,\hx b)$. $e_5+e_9$ se decodifica
        de forma única porque dos palabras de peso $\le2$ con el
        mismo síndrome diferirían en un codeword de peso
        $\le4<d=5$: imposible.
  \item Multiplicando y usando $z^4=z+1$ en cada coeficiente:
        constante $4\cdot\hx a=z^2(z^3+z)=z^3+z^2+z=\hx e$; lineal
        $4\cdot7+9\cdot\hx a=\hx f+\hx5=\hx a$; cuadrático
        $9\cdot7=z^3+z=\hx a$. En el lado derecho,
        $1+(z+z^3)g(x)$ tiene esa misma forma: $z+z^3=\hx a$,
        $1+\hx a z^3=\hx e$. Es el certificado de Patterson.
  \item Con $E=\{j\}$: $\sigma(x)=\prod_{i\in E}(x+\alpha_i)
        =x+\alpha_j$ (grado $1\le t$), y $\sigma'=1$: la ecuación
        clave se satisface con $\sigma=x+\alpha_j$. Patterson
        devuelve el localizador; su raíz $\alpha_j$ señala la
        posición del error.
  \item $(\vx P^{-1})_1=\vx_{\sigma^{-1}(1)}=\vx_{15}=1$ y
        $(\vx P^{-1})_3=\vx_{\sigma^{-1}(3)}=\vx_{13}=0$: la misma
        permutación envía los dos errores de $\vz$ a las posiciones
        $0$ y $10$ ($\supp(\vz P^{-1})=\{\sigma(3),\sigma(12)\}
        =\{10,0\}$). Es exactamente lo que el laboratorio decodifica
        con Patterson.
  \item Dos ataques: (a)~enumerar las $137$ palabras de peso $\le2$
        y quedarse con la que comparte síndrome con $C$
        ($120$ de peso $2$, $16$ de peso $1$, la nula);
        (b)~tabla completa síndrome$\to$palabra mínima ($256$
        entradas, $137$ no triviales): un oráculo de descifrado
        instantáneo. Con $n\approx8000$ el espacio de síndromes pasa
        a $\ge2^{64}$ y $\binom{n}{t}$ es astronómico: mueren la
        tabla y la enumeración, y quedan los algoritmos de la
        parte~\ref{part:seguridad} (ISD, métodos estructurales).
\end{enumerate}

\subsection*{Capítulo 6 --- El McEliece de 1978
(capítulo~\ref{ch:1978})}

\begin{enumerate}[leftmargin=1.9em,itemsep=4pt]
  \item La cadena es
        $\vx P^{-1}=(\vu S)G+\vz P^{-1}$: $\wt(\vz P^{-1})=\wt(\vz)$
        (permutar no cambia el peso), y $\wt(\vz)\le t$ es lo que
        permite aplicar el corolario de corrección al codeword
        $(\vu S)G$. El paso fino: para otro codeword $c'$,
        $\dist(\vx P^{-1},c')\ge d-\wt(\vz)\ge 2t+1-t=t+1>t$:
        el codeword verdadero está \emph{estrictamente} más cerca.
        Con errores de peso $t+1$ «a veces», esa desigualdad se
        rompe (podría haber dos codewords a distancia $t+1$):
        Patterson podría devolver otro codeword y el descifrado
        fallaría o mentiría.
  \item $G'$ para $(1024,524)$: $524\cdot1024=\num{536576}$ bits
        ($\approx65{,}5$~KiB). Para \texttt{348864}:
        $2720\cdot3488=\num{9487360}$ bits
        ($\approx1{,}13$~MB). La parte sistemática
        $(\mathrm{I}_{2720}|T)$ ocupa $mt\lceil k/8\rceil
        =768\cdot340=\num{261120}$~B (la clave pública): el ahorro
        frente a $G'$ completa es $1185920/261120=4{,}54$ veces.
  \item Si $(\vu',\vz')\ne(\vu,\vz)$ diera el mismo $\vx$, entonces
        $(\vu+\vu')G'=\vz+\vz'$: un codeword de peso
        $\le2t<2t+1\le d$, imposible salvo que sea $0$; y si es
        $0$, $\vu=\vu'$ y $\vz=\vz'$.
  \item $\binom{974}{524}/\binom{1024}{524}=2^{-53{,}61}$: es la
        probabilidad de que un conjunto de $k$ coordenadas elegido al
        azar evite las $t=50$ posiciones del error; con
        $2^{53{,}61}$ intentos de coste $k^3\approx2^{27{,}1}$ sale
        $2^{80{,}7}\approx2\cdot10^{24}$ (el «ataque de las $k$
        coordenadas», cotizado en $10^{19}$ en el artículo: el
        recálculo favorece al defensor). $\sum_{i\le50}\binom{1024}{i}
        =2^{284{,}11}$, así que un síndrome al azar se decodifica con
        probabilidad $2^{284{,}11-500}=2^{-215{,}89}$ (el artículo
        decía $2^{-215{,}59}$: coincide en lo esencial).
  \item Porque la NP-dureza es un enunciado \emph{del peor caso},
        \emph{asintótico} y sobre \emph{instancias arbitrarias}:
        (a)~no dice nada de las instancias aleatorias de una familia
        estructurada (los códigos de Goppa); (b)~para $n=3488$ fijo
        no acota nada concreto; (c)~y la seguridad del esquema es una
        conjetura de dureza \emph{en media} sobre su distribución
        (capítulo~\ref{ch:seguridad}).
  \item Con el oráculo que avisa: se envía $\vx\oplus\ve_j$ para cada
        $j$. Si $j$ está en el soporte del error, el peso baja a
        $t-1$ y el descifrado «funciona»; si no, sube a $t+1$ y
        «falla». Bastan $n=1024$ consultas para aprender el soporte
        de $\vz$ y, con él, $\vu$. En el KEM real cada consulta
        devuelve una clave pseudoaleatoria (rechazo implícito): no
        hay aviso que aprender (capítulo~\ref{ch:encapdecap}).
\end{enumerate}

\subsection*{Capítulo 7 --- La reformulación de Niederreiter
(capítulo~\ref{ch:niederreiter})}

\begin{enumerate}[leftmargin=1.9em,itemsep=4pt]
  \item Con $H_0=(\mathrm{I}_{n-k}\,|\,T_0)$, la palabra
        $v=(\gamma,\mathbf{0})$ tiene
        $vH_0^{\top}=\mathrm{I}\cdot\gamma+0=\gamma$. El paso (2) del
        descifrado pide el único $\vf$ de peso $\le t$ con
        $\vf H_0^{\top}=\gamma$: exactamente el problema de
        decodificación del capítulo~\ref{ch:goppa} (encontrar el
        vector de menor peso con un síndrome dado), y Patterson lo
        resuelve operando sobre el polinomio síndrome del vector
        rellenado.
  \item Si $\vf H_0^{\top}=\vf'H_0^{\top}$, entonces
        $(\vf+\vf')H_0^{\top}=0$: $\vf+\vf'$ es un codeword de
        $\Gamma$ de peso $\le2t<d$, luego $\vf+\vf'=0$ y
        $\vf=\vf'$.
  \item La demostración: primero $G'(H')^{\top}
        =S_1GP\,P^{\top}H_0^{\top}S_2^{\top}$: aquí se usa
        $PP^{\top}=\mathrm{I}_n$; después
        $S_1(GH_0^{\top})S_2^{\top}=0$: aquí, la ortogonalidad
        $GH_0^{\top}=0$ (los codewords tienen síndrome nulo).
        Con permutaciones distintas quedaría $P_1P_2^{\top}\ne
        \mathrm{I}_n$ en medio y la ortogonalidad no se seguiría:
        sin ella no hay reducción (el síndrome de $\vx$ ya no es el
        de $\vz$).
  \item Factores $k/(n-k)$: juguete $8/8=1$;
        \texttt{348864}, $2720/768=3{,}54$;
        \texttt{460896}, $3360/1248=2{,}69$;
        \texttt{6688128}, $5024/1664=3{,}02$;
        \texttt{6960119}, $5413/1547=3{,}50$;
        \texttt{8192128}, $6528/1664=3{,}92$.
        El original $(1024,524)$: $k/(n-k)=524/500=1{,}05$: con tasa
        $\approx1/2$, la matriz dual era casi tan grande como la
        generadora; el ahorro, mísero (del orden del $5\%$).
  \item Tamaño del espacio de mensajes: el descifrado debe
        funcionar sobre la imagen de
        $\textsf{FixedWeight}$; admitir pesos menores agrandaría la
        imagen con preimágenes «fáciles» (pesos pequeños, búsqueda
        más ligera) sin ganar entropía. Uniformidad: los teoremas de
        las transformaciones a KEM (capítulo~\ref{ch:kem}) y la
        proposición~\ref{prop:fixedweight-uniform} exigen $\ve$
        uniforme en \emph{el} conjunto de vectores de peso $t$
        exacto; mezclar pesos rompería el modelo y la simetría del
        reto.
  \item Los dos problemas de inversión:
        $\vx=\vu G'+\vz\mapsto\vu$ (McEliece) y
        $\vc=\ve(H')^{\top}\mapsto\ve$ (Niederreiter). Son «el mismo
        problema» por la proposición~\ref{prop:nied-reduccion}: el
        síndrome público de $\vx$ es el de $\vz$ (ortogonalidad
        $G'(H')^{\top}=0$), y recuperar $\ve$ equivale a recuperar
        $\vz$; en ambos casos el fondo es el problema
        de decodificación por síndromes: dada la imagen de un error
        de peso $\le t$, encontrarlo.
\end{enumerate}

\subsection*{Capítulo 8 --- Del cifrado a un KEM
(capítulo~\ref{ch:kem})}

\begin{enumerate}[leftmargin=1.9em,itemsep=4pt]
  \item En la transformación, $\textsf{Encap}$ genera $e$, publica
        $C=\textsf{Enc}(e)$ (aquí $C=(\mathrm{I}_{mt}|T)e$) y deriva
        $K=H(1,e,C)$; $\textsf{Decap}$ calcula $e'=\textsf{Decode}(C)$
        y $K'=H(1,e',C)$. Las dos partes coinciden si y solo si
        $\textsf{Decode}(\textsf{Enc}(e))=e$ con probabilidad $1$: el
        único punto donde entra la corrección de la PKE (el
        $(\varepsilon_{\mathrm{corr}})$ de los teoremas).
  \item Con $K=H(e)$, la consulta al oráculo con la segunda
        codificación $\tilde C$ de $C^*$ (distinta como bytes,
        idéntica como elemento, por el relleno del
        capítulo~\ref{ch:parametros}) devuelve $H(e^*)=K^*$: el reto
        cae con una consulta que no es $C^*$. Con $K=H(1,e,C)$, el
        criptograma entra en el hash: la respuesta del oráculo es
        $H(1,e^*,\tilde C)$, una clave distinta (ligada a la
        codificación exacta) e inútil para el ataque; la etiqueta
        $1$ además separa el dominio de la rama de rechazo
        $H(0,s,C)$.
  \item Devolver $\bot$ obliga a los protocolos a reaccionar
        (alertas, reintentos) y esa reacción es observable: es un
        oráculo de validez. Devolver siempre una clave
        pseudoaleatoria elimina la señal: el atacante no distingue
        «no decodificable» de «decodificable con valor improbable».
        La clave $s$ aparece como secreto de la rama de rechazo: en
        el teorema~\ref{thm:kem-rom} en el término $2q_H/|S|$ (hay
        que adivinar $s$ para simular) y en el
        teorema~\ref{thm:kem-qrom} en los términos de PRF sobre $s$.
  \item Segundo término: $q_D/|P|=2^{64}/2^{456}=2^{-392}$. Tercero:
        $2q_H/|S|\le 2\cdot2^{64}/2^{256}=2^{-191}$ (acotando la clave
        $s\in\Ftwo^{3488}$ por la entropía de los 256 bits de la semilla
        $\delta$). Para que el tercero dejara de ser despreciable haría
        falta $q_H\approx2^{255}$: un número de consultas imposible
        (y entonces fallarían antes los demás términos del modelo).
  \item (a)~Con ventaja $<1/1000$: $2\sqrt{\varepsilon_{\mathrm{OW}}}
        <1/1000$ exige
        $\varepsilon_{\mathrm{OW}}<(1/2000)^2\approx2^{-21{,}9}$.
        (b)~Con ventaja $<2^{-128}$:
        $\varepsilon_{\mathrm{OW}}<2^{-258}$. Es decir: el teorema
        convierte el objetivo CCA de $128$ bits en un requisito
        OW-CPA de $\approx2^{-258}$: los parámetros deben tener
        margen doble, no nulo (es la raíz cuadrada del QROM, que el
        modelo de costes realista suaviza pero no elimina).
  \item La tabla~\ref{tab:kem-cm} lo dice: (1)~sin reencriptar: la
        verificación $\textsf{Enc}(\mathrm{pk},e')=C$ no hace
        falta porque $\textsf{Decode}$ ya exige $\wt(e')=t$ y
        $(\mathrm{I}_{mt}|T)e'=C$; (2)~sin $\mathrm{pk}'$ dentro de
        la clave secreta (no hay reencriptación que lo use);
        (3)~$s$ de $n$ bits en vez de $256$; (4)~$\ve$ de
        $\textsf{FixedWeight}$ en vez de uniforme-genérico. La
        verificación de $\textsf{Decode}$ basta porque, para un
        criptograma, «ser imagen de un $e'$ de peso $t$» y «cumplir
        $\wt(e')=t\ \wedge\ (\mathrm{I}_{mt}|T)e'=C$» son la misma
        condición: la primera se comprueba sin reencriptar.
\end{enumerate}

% ======================================================================
%  CAPÍTULO 18 — SOLUCIONES (parte 3: capítulos 9-16)
% ======================================================================

\subsection*{Capítulo 9 --- La generación de claves
(capítulo~\ref{ch:keygen})}

\begin{enumerate}[leftmargin=1.9em,itemsep=4pt]
  \item Todo el proceso, incluidos los reintentos, es una función
        determinista de la semilla $\delta$: al fallar
        $\textsf{FieldOrdering}$, se toma $\delta'\leftarrow
        \textnormal{(últimos $256$ bits de }E)$ y se repite sin pedir
        azar nuevo. Las pruebas de vectores fijan $\delta$ y el
        resultado esperado: cada implementación consume el mismo flujo
        $E$ y reproduce los mismos reintentos, comparables byte a
        byte.
  \item Con $(7,1,9,3)$: al ordenar los pares $(a_i,i)$ queda
        $1<3<7<9$, es decir $\pi=(1,3,0,2)$; leyendo los bits de
        $\pi(i)$ como $\alpha_i=\sum_{j=0}^{m-1}\pi(i)_jz^{m-1-j}$
        resulta $(\alpha_0,\dots,\alpha_3)=(z,\,z+1,\,0,\,1)$. Con
        $(0,1,2,3)$: $\pi$ es la identidad
        y $(\alpha_0,\dots,\alpha_3)=(0,\,z,\,1,\,z+1)$. Con
        $(0,0,0,0)$ los enteros no son distintos:
        $\textsf{FieldOrdering}$ devuelve $\bot$ y
        $\textsf{MatGen}$ reintenta con $\delta\leftarrow\delta'$
        (determinista).
  \item $\beta$ recorre uniformemente $\GF{q^t}$ (los bloques de
        $\sigma_1=16$ bits por coeficiente; solo se usan los primeros
        $m$ de cada bloque). El polinomio mínimo de $\beta$ sobre
        $\GF q$ es mónico, irreducible, y de grado $t$ exactamente
        cuando $\beta$ no cae en un subcuerpo propio
        $\GF{q^s}$, $s\mid t$, $s<t$ (si cae, el grado es $s$ y
        $\textsf{Irreducible}$ devuelve $\bot$). Cada mónico
        irreducible de grado $t$ divide a $x^{q^t}-x$ y tiene
        exactamente $t$ raíces, todas simples, en $\GF{q^t}$; sus
        $t$ raíces son exactamente las $\beta$ cuyo mínimo es él. Así
        cada $g$ admite $t$ de las $q^t$ elecciones: $\beta$
        uniforme $\Rightarrow$ $g$ uniforme. La masa rechazada
        ---elementos que caen en subcuerpos propios--- está dominada
        por $\GF{q^{t/2}}$: $\approx q^{t/2}/q^t=q^{-t/2}$,
        despreciable para $t\ge16$.
  \item La probabilidad de que las primeras $mt$ columnas de una
        matriz binaria aleatoria $mt\times n$ sean independientes es
        $\prod_{i=0}^{mt-1}(1-2^{i-mt})
        =\prod_{j=1}^{mt}(1-2^{-j})\approx0{,}289$:
        $3{,}46$ intentos de media, y $P(\ge3)\approx(1-0{,}289)^2
        \approx0{,}5055$. Con $(\mu,\nu)=(32,64)$ el fallo por
        intento cae bajo $2^{-30}$: un solo intento en la práctica,
        a cambio de $\le2$ bits de tensión en la reducción y una
        ventana de intercambios de tamaño fijo (tiempo constante).
  \item $pk=mt\lceil k/8\rceil=1664\cdot816=\num{1357824}$~B;
        $sk=32+8+t\lceil m/8\rceil+\lceil(2m-1)2^{m-4}\rceil
        +\lceil n/8\rceil=40+256+12\,800+1024=\num{14120}$~B.
        En la clave secreta domina $\alpha$ (los $12\,800$ bytes de
        la red de Benes, el $91\%$); en la pública, la matriz $T$
        entera.
  \item El código $\ker(\mathrm{I}_{mt}|T)$ es público: dimensión
        $k=2720$ y $2^{2720}$ codewords en \texttt{348864}
        (cualquiera con la clave pública comprueba membresías). Quien
        descubra $g$ y el orden $\alpha$ obtiene la descripción de
        Goppa del código público: un decodificador eficiente, es
        decir, una clave equivalente: descifrar todo
        (observación~\ref{obs:code-public} y
        proposición~\ref{prop:nied-reduccion}: atacar el esquema es
        decodificar el código público; el secreto es su
        descripción).
\end{enumerate}

\subsection*{Capítulo 10 --- Cifrado y descifrado
(capítulo~\ref{ch:encapdecap})}

\begin{enumerate}[leftmargin=1.9em,itemsep=4pt]
  \item $v=(C,\mathbf{0})$; como $C=\hat{\Hmat}e$, para $c=v+e$
        vale $\hat{\Hmat}c=\hat{\Hmat}v+\hat{\Hmat}e=C+C=0$: es un
        codeword, y $d(c,v)=\wt(e)=2=t$. Con $C'=\hat{\Hmat}e'$ y
        $\wt(e')=1$: el paso~2 de $\textsf{Decode}$ encuentra el
        único codeword $c'$ a distancia $\le2$ (que es $v' + e'$);
        el paso~3 devuelve $e''=e'$ (peso $1$); el paso~4
        \emph{rechaza} ($\wt(e')\ne t$) y $\textsf{Decap}$ toma la
        rama de rechazo.
  \item (a)~Los trozos de $16$ bits bastan para nombrar cualquier
        posición porque todo índice es $<n\le q=2^m\le2^{16}$ (en
        los conjuntos reales $m\le13$; el juguete usa $m=4$);
        (b)~cada
        elección usa un trozo fresco: las elecciones son
        independientes; (c)~descartar los trozos fuera de rango y
        repetir no sesga: condicionado a ser aceptado, cada residuo
        es equiprobable. Productos: juguete:
        $1-1/16=0{,}9375$; \texttt{348864}:
        $e^{-t(t-1)/(2n)}=e^{-64\cdot63/6976}\approx0{,}561$.
  \item Si $c_1\ne c_2$ estuvieran ambos a distancia $\le t$ de
        $v$, entonces $\wt(c_1+c_2)\le2t<2t+1\le d$: la diferencia
        sería un codeword no nulo de peso bajo la distancia mínima.
        Imposible: el candidato es único.
  \item Con $C'=\hat{\Hmat}e'$, $\wt(e')=1$: el paso~2 devuelve
        $e'$ (la corrección funciona: $1\le t$); el paso~4 lo
        rechaza por el chequeo de peso. Sin ese chequeo,
        $\textsf{Decap}$ aceptaría todo síndrome con preimagen de
        peso $\le t$, y los «mensajes» aceptados serían uniformes en
        $\{$peso $\le t\}$ en vez de en $\{$peso $=t\}$: se
        romperían la distribución y la rigidez que las pruebas del
        capítulo~\ref{ch:kem} usan (la aceptación debe caracterizar
        exactamente los criptogramas de $\textsf{FixedWeight}$).
  \item \textsf{Encode}: $mt\cdot k=1664\cdot6528=\num{10862592}$
        bits $\approx2^{23{,}4}$; síndrome de \textsf{Decode}:
        $tn=128\cdot8192=\num{1048576}$ $\approx2^{20}$ operaciones
        de cuerpo. \textsf{Encode} toca unos $10$ veces más bits,
        pero son XORs sobre datos públicos (vectorizables, sin
        secretos); las de \textsf{Decode} son multiplicaciones en
        $\GF{2^{13}}$, más Patterson ($t^2$ operaciones) y la
        verificación: en coste efectivo domina \textsf{Decode}. Y
        esa asimetría es justo la que quiere un KEM: el lado sin
        clave (cifrar) por el camino barato y paralelo, y el lado
        con la clave (descifrar) pagando el álgebra fuerte.
  \item Para \texttt{6960119}: $H(b,v,C)$ recibe $1+\lceil
        n/8\rceil+\lceil mt/8\rceil=1+870+194=1065$ bytes. Piezas:
        la etiqueta ($0/1/2$) separa dominios (rechazo, éxito,
        confirmación \texttt{pc}); $v=e$ fija el error; $C$ ata la
        clave a la codificación exacta del criptograma (el escenario
        del capítulo~\ref{ch:kem}). No se usa la lista de posiciones
        del soporte porque el vector ya es una codificación de
        longitud fija y canónica ($\lceil n/8\rceil$ bytes): una
        lista exigiría convenios de orden, tendría longitud variable
        y no aportaría nada.
\end{enumerate}

\subsection*{Capítulo 11 --- Laboratorio 2
(capítulo~\ref{ch:labkem})}

\begin{enumerate}[leftmargin=1.9em,itemsep=4pt]
  \item Ejecutando el programa del apéndice~\ref{app:codigo}:
        $\textsf{SHAKE256}(\hx{01}\,\|\,\hx{00\,24}\,\|\,\hx{98})$
        da exactamente la $K$ de la
        figura~\ref{fig:labkem-pipeline}
        (\texttt{597f\ldots6278}), y $s=\hx{CBAE}$ se serializa
        como los bytes $\hx{AE}\,\hx{CB}$ (little-endian por grupos
        de $8$).
  \item $e=9216=2^{10}+2^{13}$: las posiciones $10$ y $13$; cada
        byte empaqueta $8$ posiciones con la posición $j$ en el bit
        $j\bmod8$ del byte $j/8$: el primer byte
        (posiciones $0$--$7$) es $\hx{00}$ y el segundo
        (posiciones $8$--$15$) es $2^2+2^5=\hx{24}$: los bytes
        $\hx{00}\,\hx{24}$. La convención «posición $0$ a la
        izquierda» describe cómo se \emph{imprimen} los vectores;
        la serialización empaqueta en little-endian. Igual para
        $s$: $\hx{CBAE}$ escribe $\hx{AE}\,\hx{CB}$.
  \item Si $c=v+e$ es un codeword y $S_c=\sum_jc_j\sigma_j\equiv0
        \pmod g$, entonces
        $S_v=\sum_jv_j\sigma_j\equiv S_c+S_e\equiv S_e\pmod g$
        (en $\Ftwo$, $v_j+e_j=c_j$). Para $\tilde C$ el vector
        $\tilde v+e$ ya \emph{no} es un codeword: su síndrome es la
        columna del bit volteado ($\ne0$), la cancelación falla y el
        síndrome cambia (en el juguete, $S=\hx4+\hx9x$ pasa a
        $\hx b+\hx6x$).
  \item Si $\textsf{Decap}$ dijera «error», el atacante voltearía
        bits en $C$ y aprendería de la respuesta si $\tilde C$
        decodifica: un oráculo de validez sobre los $256$ bytes
        (los $119$ inalcanzables marcan los que siempre fallan; cada
        consulta revelaría clases). En el KEM real no hay señal: la
        rama de rechazo devuelve $H(0,s,\tilde C)$, indistinguible
        de una $K$ válida; los $119$ síndromes inalcanzables se
        tratan como cualquier otro fracaso.
  \item $1+16+120=137$ síndromes alcanzables con peso $\le2$; los de
        peso $3$ ya no están garantizados como únicos porque la
        diferencia de dos palabras de peso $3$ puede pesar hasta
        $6>d=5$ (solo los pesos con sumas $<d$ son únicos).
        $137/256\approx53{,}5\%$ es la fracción de bytes alcanzables
        por algún error de peso $\le2$; un byte manipulado al azar
        falla tras el chequeo de peso con probabilidad
        $(256-120)/256=136/256\approx53{,}1\%$: casi la misma cifra
        (los $17$ síndromes de peso $0/1$, rechazados por el
        chequeo, reparten la diferencia).
  \item Clave pública: $\num{261120}/8=\num{32640}$ veces la del
        juguete; criptograma: $96/1=96$ veces. Mecanismos omitidos:
        $\sigma_1=16$ (bits sobrantes por coeficiente en
        $\textsf{Irreducible}$), $\sigma_2=32$ (bits por elemento en
        $\textsf{FieldOrdering}$), la forma semi-sistemática
        $(\mu,\nu)$ y la red de Benes (compresión del orden
        $\alpha$): los cuatro, en el capítulo~\ref{ch:keygen} (y el
        tercero también en el~\ref{ch:parametros}).
\end{enumerate}

\subsection*{Capítulo 12 --- Los conjuntos y las variantes
(capítulo~\ref{ch:parametros})}

\begin{enumerate}[leftmargin=1.9em,itemsep=4pt]
  \item Con $ct=\lceil mt/8\rceil$ y $ek=mt\lceil k/8\rceil$:
        \texttt{348864}: $96$ y $768\cdot340=\num{261120}$;
        \texttt{460896}: $156$ y $1248\cdot420=\num{524160}$;
        \texttt{6688128}: $208$ y $1664\cdot628=\num{1044992}$;
        \texttt{6960119}: $194$ y $1547\cdot677=\num{1047319}$;
        \texttt{8192128}: $208$ y $1664\cdot816=\num{1357824}$.
        En todos menos \texttt{6960119}, $k$ es múltiplo de $8$ y
        $mt\lceil k/8\rceil=\lceil kmt/8\rceil$. En
        \texttt{6960119}, $k=5413$: $\lceil kmt/8\rceil
        =\lceil\num{8373911}/8\rceil=\num{1046739}$, frente a
        $\num{1047319}$: la fila se redondea a byte y se pierden
        $580$~B (el relleno de la última fila, $3$ bits por fila).
  \item Con $H(2,e)$: $1+\lceil n/8\rceil$ bytes;
        $H(1,e,C)$ y $H(0,s,C)$: $1+\lceil n/8\rceil+\lceil mt/8\rceil
        +32$ en las variantes \texttt{pc} (el $32$ es la
        confirmación). Para \texttt{460896pc}: $577$ y $765$~B; la
        serie: $437/577/837/871/1025$ (para $H(2,e)$) y, en
        \texttt{pc}, $565/765/1077/1097/1265$ (para los otros dos).
  \item Mitades de Stern: $75{,}7$, $96{,}7$, $134{,}0$,
        $134{,}2$, $151{,}6$. Superan $2^{128}$ los tres conjuntos
        $\texttt{6*}$: exactamente el borrador ISO de 2023. ISO
        estandarizó \emph{además} \texttt{460896}, sabiendo que no
        alcanza el listón ingenuo, porque la raíz cuadrada es
        optimista para el atacante (Grover limitada por latencia):
        el borrador lo dejó fuera y la versión final lo incluyó con
        esa salvedad (la historia del capítulo~\ref{ch:parametros}).
  \item $p\approx0{,}289$: $3{,}46$ intentos de media y
        $P(>2)=(1-p)^2\approx0{,}5055$. Con $(\mu,\nu)=(32,64)$ el
        fallo es $<2^{-30}$: $\approx1$ intento. La variante
        \texttt{f} se ahorra $\approx3{,}5$ generaciones por clave
        (y elimina la cola y la varianza).
  \item Dos cadenas de $194$ bytes que difieran solo en los $5$
        bits de relleno representan el mismo elemento de
        $\Ftwo^{1547}$. En modo Simply el relleno se ignora: ambas
        producen la \emph{misma} $K$ (un solo valor distinto). En
        modo Narrowly, la cadena con relleno no nulo se rechaza:
        una produce la $K$ legítima y la otra la clave de la rama de
        rechazo. Debe fijarlo el protocolo porque una implementación
        Simply y una Narrowly, mezcladas, discreparían en qué clave
        corresponde a los mismos bytes: es la maleabilidad benigna
        de la figura~\ref{fig:param-narrowly}.
  \item Mensajería: \texttt{mceliece6960119} (el criptograma más
        pequeño entre los \texttt{6*}: $194$~B; \texttt{pc} si se
        quieren las defensas de fallo, $+32$~B): su clave pública de
        $\approx1$~MB se cachea una vez y se amortiza entre millones
        de mensajes. IoT: \texttt{mceliece460896f} (o
        \texttt{6960119f} si se exige largo plazo): la variante
        \texttt{f} porque la generación con reintentos es el punto
        débil de un dispositivo lento, y el tamaño de la clave
        pública allí importa menos que poder generarla de una vez.
\end{enumerate}

\subsection*{Capítulo 13 --- ¿Qué significa romper McEliece?
(capítulo~\ref{ch:seguridad})}

\begin{enumerate}[leftmargin=1.9em,itemsep=4pt]
  \item (a)~Con una clave equivalente ($g$ y el orden $\alpha$) se
        decodifica el código público: todo criptograma se abre, y
        con él todo mensaje. (b)~Un rompemensajes «de juego» triunfa
        sobre la \emph{distribución} del reto (por ejemplo con
        conocimiento del reparto uniforme de errores) y no tiene por
        qué resolver el problema de decodificación sobre síndromes
        arbitrarios: no hay reducción conocida de juego a
        decodificador general, y hay instancias estructuradas
        débiles. (c)~Un distinguidor solo decide «¿Goppa o
        aleatoria?»: no descifra; su papel es erosionar el supuesto
        de que la clave pública es pseudoaleatoria (el modelo que usan
        las reducciones) y, a veces, ser la primera etapa de la
        recuperación estructural; la reducción
        distinguidor$\to$rotura es justamente la que no se conoce.
  \item Las $137$ palabras de peso $\le2$ y el síndrome
        $(\hx2,\hx b)$: el único candidato es $e_5+e_9$, porque
        $\tilde{\Hmat}e_5+\tilde{\Hmat}e_9
        =(\hx8,\hx e)+(\hx a,\hx5)=(\hx2,\hx b)$ (la unicidad estaba
        garantizada por $d=5$: dos candidatos diferirían en un
        codeword de peso $\le4<d$). Enumerar las palabras de peso
        $64$ de \texttt{348864} cuesta $\approx2^{456}$; Prange se
        cotiza en $2^{173{,}4}$: la diferencia ($\approx283$ bits)
        vive en el algoritmo: la estructura lineal convierte el
        «gran catálogo» en una probabilidad por iteración con
        iteraciones baratas (el acortamiento del
        capítulo~\ref{ch:isd}).
  \item La conjetura, con cuantificadores: \emph{para todo}
        algoritmo $A$ de tiempo subexponencial, \emph{para} la
        instancia $(H,\,C=He)$ generada por la distribución del
        esquema ($e$ uniforme de peso $t$), \emph{con éxito}
        $=\Pr[A(H,C)=e]$ despreciable (con los mejores ataques
        conocidos $\ge2^{140}$--$2^{150}$ en los modelos serios).
        No equivalen: (a)~«NP-duro» habla del peor caso,
        asintótico y de instancias arbitrarias: ni media, ni
        distribución, ni parámetros; (b)~«no existe algoritmo
        polinómico» ni fija la distribución ni cuantifica el éxito
        (y no implica la dureza concreta de las instancias del
        esquema); (c)~«el mejor ataque cuesta $2^{140{,}8}$» es una
        cota \emph{superior} empírica: evidencia, no teorema, y
        puede mejorar (2025--2026 lo mostró). En la práctica, el
        que sostiene las decisiones es (c): las categorías y los
        márgenes se calculan con el mejor ataque conocido; la
        conjetura es el marco de la promesa.
  \item $2^{140{,}8}$ son operaciones \emph{vectoriales} en el
        modelo de memoria libre, donde una operación de vector
        grande cuenta como una: no es directamente comparable con
        las $2^{143}$ operaciones de bits de buscar una clave
        AES-128; con el cargo realista de la memoria ($2^{86{,}6}$),
        el mismo ataque sube a $2^{158{,}6}$: factor
        $2^{158{,}6-140{,}8}=2^{17{,}8}\approx228\,210$. La
        categoría~$1$ se certifica con el modelo que carga la
        memoria ($158{,}6\ge143$): de ahí que no haya
        contradicción.
  \item Mitad en el exponente: $151{,}4/2=75{,}7$, muy por debajo
        del listón $2^{128}$: \texttt{348864} no pasa el test
        cuántico ingenuo. ISO estandarizó igualmente
        \texttt{460896} ($96{,}7$, también por debajo) porque la
        raíz cuadrada sobreestima al atacante: Grover está limitada
        por la latencia del bucle interno, y el coste cuántico real
        es mayor; con esa salvedad documentada, ISO prefirió
        incluir el conjunto de claves pequeñas antes que dejar el
        perfil sin cobertura.
  \item (a)~Distinguidor (familia estructural, capítulos
        \ref{ch:seguridad} y~\ref{ch:estructura}): rompe supuestos
        de indistinguibilidad; no descifra. (b)~Mejora de ISD
        (capítulo~\ref{ch:isd}): mueve la frontera de la
        decodificación genérica; sola no rompe el esquema: cambia
        la estimación. (c)~Canal lateral (capítulos
        \ref{ch:kem}--\ref{ch:parametros} y~\ref{ch:2026}): rompe
        un despliegue real (recupera la clave de una instalación
        concreta midiendo tiempos), no el esquema abstracto.
\end{enumerate}

\subsection*{Capítulo 14 --- Decodificación por conjuntos de
información (capítulo~\ref{ch:isd})}

\begin{enumerate}[leftmargin=1.9em,itemsep=4pt]
  \item (a)~$64\log_2(768/3488)=-139{,}73$;
        (b)~$k^3=2720^3$: $3\log_2 2720=34{,}23$; total
        $173{,}96$; (c)~la guía cotiza $173{,}4$: el medio bit son
        las constantes finas del estimador (trabajo de la
        eliminación, pivotaje, normalización del éxito) que
        nuestro producto grueso sobrecuenta. Diferencia de
        constantes, no de modelo.
  \item Primer salto: $22{,}0$ bits $=2^{22}\approx4{,}19$ millones
        $\approx4$ factores de $2^{20}$; segundo salto: $10{,}6$
        bits $\approx1\,552=1{,}5\cdot2^{10}$. El primero es
        el grande: el cruce de listas de Stern aporta
        $\approx4\cdot2^{20}$ sobre Prange; la búsqueda de vecinos
        de May--Ozerov añade «solo» $1{,}5\cdot2^{10}$.
  \item Porque ambos son $\bigl(1/(1-R)\bigr)^{t}$ por
        $2^{o(t)}$: con $R\approx0{,}78$, la base es
        $1/(1-R)\approx4{,}55$ ($2^{2{,}18t}$; para $t=64$,
        $2^{139{,}8}$). Prange arrastra un coste polinómico
        ($k^3=2^{34{,}2}$): de $173{,}96$; May--Ozerov mete el
        trabajo en las listas y los vecinos y se queda cerca del
        término desnudo ($140{,}8$). El «$5^t$» es taquigrafía del
        término dominante ($5\approx1/(1-R)$ redondeado): mismo
        exponente, $o(1)$ distinto, y a $t=64$ ese $o(1)$ vale
        $33$ bits.
  \item (a)~Memoria gratis vs realista para \texttt{6960119}:
        $279{,}2-245{,}7=33{,}5$ bits; (b)~$245{,}7<272$ (el
        ataque con memoria gratis sale «más barato» que agotar
        AES-256), pero $279{,}2>272$ (con el cargo de la memoria
        ya no): la categoría del equipo es la $5$ porque certifica
        con el presupuesto realista.
  \item (8192,96) cubre a (4608,96) por acortamiento: cualquier
        ataque contra el acortado se transporta al original $(q,q,t)$
        ---el sistema de McEliece, con su historia de estudio--- con
        el éxito multiplicado por $\binom{4608}{96}/\binom{8192}{96}$;
        la pérdida de seguridad del traslado es de solo
        $\approx0{,}8$ bits en el modelo 1/2. Por eso basta analizar
        el original: bajar $n$ no regala seguridad contra ISD.
  \item La mitad del exponente es una cota optimista porque
        supone que cada iteración de Grover cuesta lo mismo que su
        versión clásica y que el paralelismo no tiene penalización:
        en la práctica el bucle interno (accesos a memoria, la
        búsqueda en listas) domina y la ventaja cuántica real es
        menor que la cuadrática. Los $256$ bits de $\delta$ fijan
        el suelo complementario: la búsqueda cuántica de la semilla
        cuesta $2^{128}$ (Grover sobre $256$ bits, con una
        generación por comprobación), de modo que la seguridad
        cuántica de un conjunto es $\min($mitad-ISD$,2^{128})$: por
        eso el listón de ISO es exactamente $2^{128}$ y por eso los
        conjuntos por debajo ($75{,}7$, $96{,}7$) son el punto
        débil, aunque el modelo de la mitad sea optimista.
\end{enumerate}

\subsection*{Capítulo 15 --- Ataques estructurales: recuperar la
clave privada (capítulo~\ref{ch:estructura})}

\begin{enumerate}[leftmargin=1.9em,itemsep=4pt]
  \item (a)~$4096^{64}/64=2^{768-6}=2^{762}$;
        (b)~$1024^{50}/50\approx2^{500-5{,}6}=2^{494{,}4}$;
        (c)~$\log_2P(4096,769)=\sum_{j=0}^{768}\log_2(4096-j)
        \approx9116{,}8$.
  \item $\binom{4096}{3488}$ cuenta \emph{conjuntos} de $n$
        posiciones: $\approx2^{2476}$ (entropía de «qué elementos
        forman el soporte»); $P(4096,769)$ cuenta \emph{tuplas
        ordenadas} de $769$ posiciones concretas: el dato parcial
        que el modelo de 2022 necesitaba reconstruir. La diferencia
        ($\approx6640$ bits) es el precio del orden y de las
        posiciones concretas: una va con el tamaño del soporte
        ($n$) y la otra con los elementos libres ($n-k+1$).
  \item (a)~$9116{,}8-140{,}8=8976$ bits; el cociente de costes es
        $2^{8976}$ (¡$2^{8720}$ veces el margen entero de
        AES-256!). (b)~Comparada con $2^{256}$, la cifra de 2022 es
        inalcanzable por cualquier medio físico imaginable: el
        titular era «correcto» dentro de un modelo de datos
        parciales que estaba a punto de agotarse, y como amenaza
        práctica era casi cómico (nadie ejecutaría jamás tal
        cómputo).
  \item La cadena sería: distinguidor con ventaja $1-o(1)$
        $\to$ estadístico que expone estructura en las claves
        reales $\to$ (eslabón que falta) inversión eficiente en
        recuperación de clave $\to$ descifrado. Los márgenes de
        la tabla~\ref{tab:est-margenes} miden la distancia del régimen de
        cada método a los parámetros reales: Faugère ($t\le20$
        frente a $64$--$128$) tiene holgura enorme; Carrier
        (tasa $<0{,}3$ frente a $0{,}7$) moderada; el
        \emph{holdout}/anulación de orden superior ya alcanza el
        $t$ real: es el candidato natural a escalar, y el eslabón
        que falta es la barrera que hoy separa el estadístico de
        la clave.
  \item $n<q$ es doblemente bueno: contra ISD no se paga nada
        (el acortamiento transporta la dureza con $0{,}8$ bits de
        pérdida: observación~\ref{obs:isd-acortamiento}), y contra
        la división del soporte el atacante debe adivinar qué $n$
        de las $q$ columnas forman $L$ y en qué orden: paga
        $\binom{q}{n}$ (y tuplas ordenadas: $2^{9116{,}8}$); con
        $n=q$ ese frente se desvanece
        (observación~\ref{obs:est-nq}). El diseño cobra en un
        frente lo que no pierde en el otro.
  \item Holguras (menor = más peligroso): (1)~anulación de orden
        superior / \emph{holdout} (su régimen alcanza el $t$ real:
        la frontera está pegada a los parámetros, aunque sus
        costes aún no escalen); (2)~Carrier (necesita tasa baja,
        $<0{,}3$, y los conjuntos reales van a $0{,}7$--$0{,}8$:
        margen moderado en el otro lado de la frontera);
        (3)~Faugère ($t\le20$ frente a $\ge64$: holgura enorme).
        El más peligroso a escalar es el primero: es el único
        cuyo \emph{régimen} ya contiene a los conjuntos reales.
        (Ejercicio de razonamiento: la respuesta admite matices,
        pero no intercambiar los tres regímenes.)
\end{enumerate}

\subsection*{Capítulo 16 --- El terremoto: 2025--2026
(capítulo~\ref{ch:2026})}

\begin{enumerate}[leftmargin=1.9em,itemsep=4pt]
  \item (a)~$9116{,}8-94{,}3=9022{,}5$ bits de diferencia. No es el
        «mismo ataque mejorado»: la familia de 2022 (datos
        parciales y reconstrucción del soporte) estaba agotada
        contra su muro; el \emph{holdout} (GIJS/Weis) es una
        familia nueva, importada de la teoría de la información
        privada: no busca el soporte, evalúa el código con familias
        de funciones correlacionadas. (b)~$94{,}3<256$: por debajo
        del suelo trivial (enumerar la semilla de $256$ bits, con
        una generación por candidato). Una estimación así solo es
        compatible con (i)~una contabilidad que no cobra el trabajo
        real por candidato/paso (memoria gratis, «operaciones de
        bits», heurísticas no ejecutadas) o (ii)~una rotura
        histórica. De ahí la exigencia de verificación pública: a
        octubre de 2026, ninguna de las dos cosas está demostrada.
  \item Total con $100$--$1400$ ejecuciones: sin memoria,
        $2^{89}\cdot[100,1400]=2^{95{,}6}$--$2^{99{,}5}$; con cargo
        severo, $2^{117}\cdot[100,1400]=2^{123{,}6}$--$2^{127{,}5}$.
        La decodificación genérica en ese modelo para
        \texttt{348864} está en $2^{158{,}6}$: ambos rangos quedan
        por debajo: «nuestra mejor estimación; con el cargo de la
        guía queda por debajo de la decodificación genérica en ese
        modelo». Es exactamente la afirmación del equipo, y la
        razón del sobresalto: incluso el rango con cargo severo
        ($\le2^{127{,}5}$) roza el listón ISO de $2^{128}$.
  \item $94{,}3+8\cdot6=142{,}3\approx2^{142}$: la regla de Weis
        (unos $8$ bits por unidad de $m$) reproduce la instancia de
        Bernstein ($n=6176$, $t=112$, $m=18$), cifrada en
        $\approx2^{142}$, con error de $0{,}3$ bits. Clave pública
        allí: $mt\lceil k/8\rceil=2016\cdot520=\num{1048320}$~B
        $\approx1{,}0$~MB: cuatro veces la de \texttt{348864}
        ($0{,}25$~MB) y del orden de la mayor actual
        (\texttt{8192128}: $1{,}36$~MB). El eje $m$ compra
        seguridad exponencial a precio lineal en la clave.
  \item (a)~Demostrada (en su régimen asintótico: es un teorema
        con ventaja $1-o(1)$). (b)~Estimada: el método de
        recuperación es heurístico (cinco supuestos enumerados por
        Weis, probados solo en claves pequeñas y retos TII) y la
        cifra concreta $2^{102{,}4}$ depende de la contabilidad.
        (c)~Heurística y disputada: la afirmación de Vedenev
        «esencialmente una ejecución» no está verificada de forma
        independiente (hay réplicas y contra-réplicas). (d)~Hecho
        (negativo y verificable): nadie ha ejecutado el cómputo de
        núcleos contra un conjunto real; los retos TII tienen
        $n=214$.
  \item $2^{94{,}3}$ frente a $2^{143}$ (AES-128): en el mismo
        registro de «operaciones de bits», el ataque saldría
        $48{,}7$ bits más barato: alarmante. Pero ese registro no
        cobra la memoria ni la latencia; con el cargo severo
        ($2^{117}$--$2^{128}$), el ataque cuesta tanto o más que
        agotar AES-128: favorable al esquema. Falta: (i)~un modelo
        de costes consensuado (cómo cobrar memoria y accesos) y
        (ii)~una ejecución de extremo a extremo (nadie ha corrido
        el método contra un conjunto real): sin ambas, la
        comparación está abierta (BSI elige la precaución; el
        equipo, su contabilidad clásica).
  \item Ejemplos (límite de $80$ caracteres):
        \emph{verdadero} ---~«McEliece: ataques de 2026 estimados,
        ninguno ejecutado ni verificado»; \emph{falso por
        alarma}~--- «McEliece roto: un ataque de $2^{94}$ desarma
        el estándar ISO» (le falta: estimación, modelo de costes
        disputado, no ejecutado); \emph{falso por
        complacencia}~--- «McEliece intacto: nada ha cambiado»
        (le falta: las estimaciones bajaron $9000$ bits en cuatro
        años y el BSI desaconseja despliegues nuevos). (Ejercicio
        abierto: intercambiar con un colega y discutir la cláusula
        ausente.)
\end{enumerate}

% ======================================================================
%  CAPÍTULO 19 — GLOSARIO (parte final de la parte V)
% ======================================================================

\chapter{Glosario}\label{ch:glosario}

\begin{roadmap}
Este glosario reúne, en orden alfabético, los términos que el libro
usa con un significado técnico preciso. Cada entrada recuerda en una
o dos líneas la definición y remite al capítulo donde se desarrolla.
Hace además las veces de índice: quien busque «¿dónde se habló de
X?» puede venir aquí primero.
\end{roadmap}

\medskip
\noindent Los nombres de algoritmos y los parámetros aparecen con su
notación del texto; los capítulos se citan por su número de
referencia.

\begin{description}[leftmargin=2.6em,style=nextline,itemsep=2.5pt,parsep=0pt]

  \item[\textbf{Acortamiento} (\emph{shortening})]
        Técnica que fija a cero $s$ coordenadas de un código
        $(n,k)$; en ISD transporta la dureza entre instancias con
        pérdida $\le0{,}8$ bits, y es la cara amable de tomar
        $n<q$ (caps.~\ref{ch:isd} y~\ref{ch:estructura}).

  \item[\textbf{AES-128 / AES-256}]
        Cifrado simétrico de referencia. Grover reduce su búsqueda
        de clave a $2^{64}$ y $2^{128}$: el «listón de 128 bits»
        que traducen las categorías (caps.~\ref{ch:pqc}
        y~\ref{ch:isd}).

  \item[\textbf{Alfa} ($\alpha$; orden del soporte)]
        Lista ordenada $\alpha_0,\dots,\alpha_{q-1}$ de todos los
        elementos de $\GF q$; el soporte del código es
        $L=\{\alpha_0,\dots,\alpha_{n-1}\}$. Su \emph{orden} es
        secreto y se sortea con $\textsf{FieldOrdering}$
        (cap.~\ref{ch:keygen}).

  \item[\textbf{Anulación de orden superior} (\emph{higher-order
        vanishing})]
        Familia de distinguidores y ataques de 2025--2026 basada en hallar
        polinomios homogéneos multivariados de grado creciente que se
        anulan con alta multiplicidad en las palabras del código; de ella
        surgen el \emph{holdout} y las estimaciones de GIJS/Weis
        (caps.~\ref{ch:estructura} y~\ref{ch:2026}).

  \item[\textbf{Benes, red de}]
        Representación de la permutación $\pi$ (el orden $\alpha$)
        con $2m-1$ etapas de $2^{m-1}$ conmutadores:
        $\lceil(2m-1)2^{m-4}\rceil$ bytes (5\,888 para $m=12$);
        es como la clave secreta guarda $\alpha$
        (cap.~\ref{ch:keygen}).

  \item[\textbf{Berson, ataque de reacción de}]
        Con un descifrado que avisa del fallo, $n$ consultas
        $x\oplus e_j$ revelan el soporte del error. El rechazo
        implícito del KEM lo desactiva (cap.~\ref{ch:1978}).

  \item[\textbf{BSI} (\emph{Bundesamt für Sicherheit in der
        Informationstechnik})]
        Oficina federal alemana de seguridad. Su nota del
        1-oct-2026 desaconseja Classic McEliece en nuevos
        despliegues tras las estimaciones de 2026 y recomienda
        ML-KEM, FrodoKEM o HQC (cap.~\ref{ch:2026}).

  \item[\textbf{Canal lateral}]
        Ataque a la implementación (tiempos, cachés, accesos)
        que no rompe el esquema abstracto; el ataque de tiempos de
        Shoup contra la decodificación es el ejemplo clásico
        (caps.~\ref{ch:kem} y~\ref{ch:2026}).

  \item[\textbf{Categorías (NIST)}]
        Escala que traduce «romper el esquema» a costes tipo AES:
        categoría~1 $\approx2^{143}$ operaciones de puerta
        (AES-128), categoría~5 $\approx2^{272}$ (AES-256). Classic
        McEliece se certifica con el modelo de costes que carga la
        memoria (caps.~\ref{ch:parametros} y~\ref{ch:isd}).

  \item[\textbf{Clase lateral}]
        Conjunto $y+C$ de vectores con el mismo síndrome que $y$;
        decodificar es elegir en la clase el representante de peso
        mínimo (cap.~\ref{ch:lineales}).

  \item[\textbf{Clave equivalente}]
        Objeto (por ejemplo $g$ y el orden $\alpha$, o un
        decodificador del código público) con el que se descifra
        todo sin ser la clave secreta literal; su obtención es una
        rotura total (caps.~\ref{ch:keygen}, \ref{ch:seguridad}
        y~\ref{ch:estructura}).

  \item[\textbf{Código de Goppa} ($\Gamma_2(L,g)$)]
        Código $\{c:\sum_j c_j/(x+\alpha_j)\equiv0\pmod g\}$, con
        $g$ de grado $t$ sin raíces en $L$; la versión binaria
        (subcuerpo) con $g$ irreducible tiene distancia
        $\ge2t+1$: es el código que usa Classic McEliece
        (cap.~\ref{ch:goppa}).

  \item[\textbf{Código lineal}]
        Subespacio de $\Ftwo^n$ descrito por una matriz generadora
        $G$ o una de paridad $\Hmat$; su distancia mínima $d$ es
        el peso del codeword no nulo más ligero
        (cap.~\ref{ch:lineales}).

  \item[\textbf{Conjetura de dureza (SD en media)}]
        «Para todo algoritmo subexponencial, el éxito sobre
        instancias de la distribución del esquema es
        despreciable»: el supuesto en el que descansa la seguridad
        (caps.~\ref{ch:seguridad} y~\ref{ch:isd}).

  \item[\textbf{Cuerpo finito} ($\GF q$ con $q=2^m$; $\Ftwo$)]
        El terreno donde viven $g$, $\alpha$ y los síndromes;
        $\Ftwo$ es el cuerpo base y $\GF q$ el de extensión
        (cap.~\ref{ch:cuerpos}).

  \item[\textbf{Delta} ($\delta$; semilla)]
        Los $32$ bytes que siembran toda la generación: de ella se
        derivan deterministamente el código, el orden y el secreto
        de rechazo $s$; la clave secreta completa es regenerable
        desde $\delta$ (cap.~\ref{ch:keygen}).

  \item[\textbf{Decode / Encode}]
        Primitivas del KEM: $\textsf{Encode}(e)=\hat{\Hmat}e$ y
        $\textsf{Decode}(C)$, que corrige hasta $t$ errores y solo
        acepta vectores de peso exactamente $t$
        (caps.~\ref{ch:encapdecap} y~\ref{ch:labkem}).

  \item[\textbf{Despliegues}]
        Implementaciones reales y usos del esquema: SUPERCOP y
        \texttt{cmimpl}, liboqs y PQClean, y los despliegues
        híbridos (Mullvad, Rosenpass, PQConnect), todos anteriores
        al aviso del BSI de 2026 (cap.~\ref{ch:2026}).

  \item[\textbf{Distancia de Hamming; distancia mínima} ($d$)]
        El peso de la diferencia entre dos vectores; $d$ es el
        mínimo peso no nulo de un codeword. El teorema de Goppa da
        $d\ge2t+1$ y de ahí la decodificación única
        (caps.~\ref{ch:lineales} y~\ref{ch:goppa}).

  \item[\textbf{Distinguidor}]
        Algoritmo que decide «¿Goppa de verdad o matriz aleatoria?»
        con ventaja no despreciable. Erosiona el supuesto de
        indistinguibilidad de la clave pública, pero no descifra
        por sí solo (caps.~\ref{ch:seguridad}--\ref{ch:2026}).

  \item[\textbf{Dumer, algoritmo de}]
        Mejora de Prange (1991) que precomputa columnas y
        amplifica la probabilidad de éxito por iteración; primer
        escalón relevante tras Prange (cap.~\ref{ch:isd}).

  \item[\textbf{$e$ (vector de error)}]
        Vector de peso exactamente $t$ que cifra:
        $C=\hat{\Hmat}e$. Su soporte es el «mensaje-error» del
        esquema (caps.~\ref{ch:1978}--\ref{ch:encapdecap}).

  \item[\textbf{Escalera de ataques}]
        La sucesión que baja la cotización de ISD para cada
        conjunto: Prange (1962) $\to$ Lee--Brickell/Leon (1988)
        $\to$ Stern/Dumer (1989/91) $\to$ MMT/BJMM (2011/12)
        $\to$ May--Ozerov (2015) $\to$ Both--May (2018)
        $\to$ Esser--Bellini y Carrier (2022) $\to$ las revisiones
        de 2026 (cap.~\ref{ch:isd}).

  \item[\textbf{Espacio de criptogramas} ($P$)]
        La imagen de los errores de peso $t$ bajo $\textsf{Encode}$
        ($\binom{n}{t}$ elementos; $|P|>2^{456}$ en
        \texttt{348864}); aparece en la cota ROM
        (cap.~\ref{ch:kem}).

  \item[\textbf{Faugère, ataque de}]
        Recuperación algebraica de la clave por sistemas
        polinómicos con interrupción temprana (XLS de Faugère et
        al.); hoy limitado a $t\le20$: holgura enorme frente a los
        $t$ reales (caps.~\ref{ch:estructura} y~\ref{ch:2026}).

  \item[\textbf{FieldOrdering}]
        Algoritmo de la especificación que convierte $32q$ bits en
        el orden uniforme $\alpha$; el choque de enteros provoca
        reintento determinista (cap.~\ref{ch:keygen};
        proposición~\ref{prop:field-ordering}).

  \item[\textbf{FIPS 203/204/205}]
        Los estándares post-cuánticos del NIST (ML-KEM, ML-DSA,
        SLH-DSA); el ecosistema del que Classic McEliece es
        alternativa y complemento (caps.~\ref{ch:pqc}
        y~\ref{ch:2026}).

  \item[\textbf{FixedWeight}]
        Algoritmo que muestrea errores uniformes de peso exacto
        $t$ con trozos de $16$ bits y rechazo; su uniformidad es
        una proposición (cap.~\ref{ch:encapdecap};
        proposición~\ref{prop:fixedweight-uniform}).

  \item[\textbf{FO (Fujisaki--Okamoto)}]
        Transformación genérica de PKE a KEM en el ROM; Classic
        McEliece usa una versión simplificada (sin reencriptación,
        con rechazo implícito y secretos más pequeños)
        (cap.~\ref{ch:kem}).

  \item[\textbf{FrodoKEM}]
        KEM de retículos (LWE) considerado conservador; una de las
        alternativas que el BSI recomienda mientras se revisa la
        TR-02102-1 (cap.~\ref{ch:2026}).

  \item[\textbf{$G'$ (matriz pública de McEliece)}]
        La generadora disfrazada $G'=SGP$; su versión sistemática
        se reduce a la clave pública $T$ (caps.~\ref{ch:1978}
        y~\ref{ch:niederreiter}).

  \item[\textbf{$g$ (polinomio de Goppa)}]
        Polinomio mónico de grado $t$ sobre $\GF q$ sin raíces en
        $L$; el secreto central del código, sorteado con
        $\textsf{Irreducible}$ (caps.~\ref{ch:goppa}
        y~\ref{ch:keygen}).

  \item[\textbf{Gaussiana (eliminación)}]
        Procedimiento para llevar $\hat{\Hmat}$ a la forma
        $(\mathrm{I}_{mt}\,|\,T)$; también el corazón de las
        iteraciones de ISD (caps.~\ref{ch:labgoppa}
        y~\ref{ch:isd}).

  \item[\textbf{Grover, algoritmo de}]
        Búsqueda cuántica con aceleración cuadrática en el
        exponente: AES-256 $\to2^{128}$; en ISD, la «mitad» del
        exponente, con la latencia como límite realista
        (caps.~\ref{ch:pqc} y~\ref{ch:parametros}).

  \item[\textbf{Hamming, código de} {($[7,4,3]$)}]
        El ejemplo recurrente para síndromes y corrección: matriz
        de paridad con las siete columnas no nulas y decodificación
        de un error (cap.~\ref{ch:lineales}).

  \item[\textbf{Harvest now, decrypt later}]
        Grabar hoy el tráfico cifrado para descifrarlo cuando
        exista el ordenador cuántico; la razón de migrar antes de
        que sea urgente (cap.~\ref{ch:pqc}).

  \item[\textbf{Híbrido}]
        Combinación de un KEM post-cuántico con uno clásico (por
        ejemplo X25519); los despliegues reales de Classic
        McEliece son híbridos, lo que conserva la mitad clásica
        ante roturas (caps.~\ref{ch:pqc} y~\ref{ch:2026}).

  \item[\textbf{Holgura} (\emph{tightness})]
        Pérdida entre una cota de seguridad y el juego real: el
        QROM pierde la raíz cuadrada del éxito OW-CPA
        (cap.~\ref{ch:kem}).

  \item[\textbf{HQC}]
        KEM de códigos cuasi-cíclicos seleccionado por el NIST en
        2025; la otra alternativa que el BSI pone por delante de
        McEliece (caps.~\ref{ch:parametros} y~\ref{ch:2026}).

  \item[\textbf{Irreducible}]
        Algoritmo de la especificación que convierte $16t$ bits en
        un $g$ mónico irreducible uniforme; su demostración es la
        proposición de uniformidad (cap.~\ref{ch:keygen}).

  \item[\textbf{ISD (decodificación por conjuntos de
        información)}]
        Familia de algoritmos de decodificación genérica (Prange,
        Stern, MMT, BJMM, May--Ozerov, \dots) que fija la
        seguridad clásica de los conjuntos (cap.~\ref{ch:isd}).

  \item[\textbf{ISO/IEC 18033-2 (Amd.2)}]
        Norma con los $4$ conjuntos $\times$ $4$ variantes $=16$
        perfiles de Classic McEliece (jun-2026); descarta
        \texttt{348864} (caps.~\ref{ch:parametros} y~\ref{ch:2026}).

  \item[\textbf{Juguete (el)}]
        La instancia diminuta usada en los laboratorios:
        $m=4$ ($q=16$), $n=16$, $t=2$, elaborada paso a paso en el
        capítulo~\ref{ch:labgoppa} y completa en
        el~\ref{ch:labkem}; \texttt{toy.py}
        (apéndice~\ref{app:codigo}) reproduce sus números.

  \item[\textbf{$K$ (clave del KEM)}]
        Los $32$ bytes que encapsula y desencapsula el KEM:
        $H(1,e,C)$ en el éxito y $H(0,s,C)$ en el rechazo
        (caps.~\ref{ch:kem} y~\ref{ch:labkem}).

  \item[\textbf{KAT (vectores de prueba)}]
        Archivos de la especificación con claves, criptogramas y
        secretos de referencia; comparables byte a byte porque todo
        el proceso (incluidos los reintentos) es determinista
        (cap.~\ref{ch:keygen}).

  \item[\textbf{KEM (encapsulamiento de claves)}]
        Primitiva \textsf{Encap}/\textsf{Decap} que produce una
        clave compartida a partir de una clave pública; la que
        Classic McEliece deriva de su PKE (cap.~\ref{ch:kem}).

  \item[\textbf{Latencia (cuántica)}]
        Límite práctico de Grover: las iteraciones se serializan y
        la raíz cuadrada sobreestima al atacante real; de ahí que
        la «mitad» del exponente se compare con cautela
        (caps.~\ref{ch:parametros} y~\ref{ch:2026}).

  \item[\textbf{Lee--Brickell}]
        Primer algoritmo de ISD con lista (1988): fija $p$
        posiciones «con error» y busca en ellas; semilla de la
        línea Stern (cap.~\ref{ch:isd}).

  \item[\textbf{Leon}]
        Variante de lista por información canónica (1988),
        hermana de Lee--Brickell (cap.~\ref{ch:isd}).

  \item[\textbf{Localizador (polinomio $\sigma$)}]
        Polinomio cuyas raíces señalan las posiciones del error:
        $\sigma(x)=\prod_{j\in\supp(e)}(x+\alpha_j)$; Patterson lo
        recupera del síndrome (caps.~\ref{ch:goppa}
        y~\ref{ch:labgoppa}).

  \item[\textbf{$m$}]
        Exponente del cuerpo: $q=2^m$. Classic McEliece usa
        $m=12$ ($q=4096$) o $m=13$ ($q=8192$); la seguridad
        estructural crece $\approx8$ bits por unidad de $m$
        (caps.~\ref{ch:parametros} y~\ref{ch:estructura}).

  \item[\textbf{MatGen}]
        Subalgoritmo de generación de claves: forma
        semi-sistemática de las $mt$ primeras columnas, con
        reintentos deterministas ($\delta\leftarrow\delta'$)
        (cap.~\ref{ch:keygen}).

  \item[\textbf{May--Ozerov / Both--May}]
        Ataques que llevan MMT/BJMM y la búsqueda de vecinos a
        ISD; May--Ozerov (2015) cotiza \texttt{348864} en
        $2^{140{,}8}$ (memoria libre) y Both--May (2018) afinan el
        régimen de tasas de error altas (cap.~\ref{ch:isd}).

  \item[\textbf{McEliece (criptosistema)}]
        El cifrado de 1978: $x=uG'+z$ con la generadora disfrazada
        y error de peso $t$; quien conoce el código lo decodifica
        con Patterson (cap.~\ref{ch:1978}).

  \item[\textbf{Memoria (coste de la)}]
        En ISD, las tablas grandes dominan el coste real: cargarla
        añade decenas de bits (hasta $33{,}5$ en \texttt{6960119})
        y decide las categorías (caps.~\ref{ch:isd}
        y~\ref{ch:2026}).

  \item[\textbf{ML-KEM}]
        El KEM post-cuántico estandarizado por el NIST (FIPS 203);
        referencia de comparación para la criptografía de códigos
        (caps.~\ref{ch:pqc} y~\ref{ch:2026}).

  \item[\textbf{Modelo de costes}]
        La convención con la que se cuenta el trabajo de un ataque
        (puertas, bits, operaciones vectoriales, memoria,
        latencia); las columnas «libre» y «con cargo» de las
        tablas (caps.~\ref{ch:isd}, \ref{ch:estructura}
        y~\ref{ch:2026}).

  \item[\textbf{NCP (decodificación / palabra más cercana)}]
        El problema hermano de SD: dado $y$, hallar el codeword de
        peso mínimo a distancia de $y$; el descifrado de McEliece
        es una instancia (caps.~\ref{ch:lineales}
        y~\ref{ch:seguridad}).

  \item[\textbf{Niederreiter (reformulación)}]
        Versión dual basada en síndromes: $C=e(H')^{\top}$;
        equivalente a McEliece por la reducción de la
        proposición~\ref{prop:nied-reduccion}
        (cap.~\ref{ch:niederreiter}).

  \item[\textbf{NIST}]
        Organizador de la competición post-cuántica (2016--2024);
        estandarizó ML-KEM/ML-DSA/SLH-DSA, aplazó su decisión
        sobre McEliece y seleccionó HQC como KEM extra
        (caps.~\ref{ch:pqc} y~\ref{ch:2026}).

  \item[\textbf{OW-CPA}]
        Juego de seguridad del PKE: recuperar $e$ dado $C$; la
        raíz que el QROM divide por dos en su cota
        (caps.~\ref{ch:kem} y~\ref{ch:seguridad}).

  \item[\textbf{Parámetros (conjuntos de)}]
        Los cinco del estándar: \texttt{348864}, \texttt{460896},
        \texttt{6688128}, \texttt{6960119}, \texttt{8192128}
        ($m=12$ o $13$, $t=64$--$128$); sus tablas de tamaños y
        seguridad (caps.~\ref{ch:parametros}--\ref{ch:2026}).

  \item[\textbf{Paridad (matriz de, $\Hmat$)}]
        Matriz cuyas filas ortogonalizan el código
        ($GH^{\top}=0$); en forma sistemática
        $\hat{\Hmat}=(\mathrm{I}_{mt}\,|\,T)$
        (caps.~\ref{ch:lineales} y~\ref{ch:keygen}).

  \item[\textbf{Patterson (algoritmo de)}]
        Decodificador de los códigos de Goppa: forma el síndrome
        $S$, resuelve la ecuación clave (Euclides/Berlekamp) y
        halla las raíces del localizador; $O(t^2)$ operaciones en
        $\GF q$ más inversiones (caps.~\ref{ch:goppa},
        \ref{ch:labgoppa} y~\ref{ch:encapdecap}).

  \item[\textbf{pcf / pc / f (variantes)}]
        Sufijos de los conjuntos: \texttt{f} = generación sin
        reintentos visibles (forma $(32,64)$); \texttt{pc} =
        confirmación de texto claro ($C=(C_0,C_1)$ con
        $C_1=H(2,e)$); \texttt{pcf} = ambas (cap.~\ref{ch:parametros}).

  \item[\textbf{Permutación ($P$, $\pi$)}]
        El disfraz de coordenadas, $(\vv P)_j=\vv_{\sigma(j)}$; su
        orden es el secreto que FieldOrdering sortea
        (caps.~\ref{ch:labgoppa} y~\ref{ch:keygen}).

  \item[\textbf{Peso} ($\wt$)]
        Número de posiciones no nulas. Los parámetros del esquema
        son pesos: $\wt(e)=t$ y $\wt(z)\le t$
        (caps.~\ref{ch:lineales} y~\ref{ch:1978}).

  \item[\textbf{Pi} ($\pi$)]
        La permutación que resulta de ordenar los enteros de
        FieldOrdering; se guarda comprimida en la red de Benes
        (cap.~\ref{ch:keygen}).

  \item[\textbf{PKE (cifrado de clave pública)}]
        El cifrado subyacente (McEliece/Niederreiter) del que el
        KEM se deriva mediante la transformación del
        capítulo~\ref{ch:kem}.

  \item[\textbf{PQC (criptografía post-cuántica)}]
        El campo entero: algoritmos diseñados para resistir
        atacantes con ordenador cuántico
        (caps.~\ref{ch:pqc} y~\ref{ch:2026}).

  \item[\textbf{Prange}]
        ISD original (1962): elegir $k$ columnas y comprobar si el
        error cae fuera de ellas; iteración barata y probabilidad
        pequeña. La referencia contra la que se mide todo lo demás
        (cap.~\ref{ch:isd}).

  \item[\textbf{QROM}]
        Modelo del oráculo aleatorio cuántico: la cota de FO
        pierde la raíz cuadrada del éxito OW-CPA
        (caps.~\ref{ch:kem} y~\ref{ch:estructura}).

  \item[\textbf{Recuperación de clave} (\emph{key recovery})]
        El objetivo más fuerte de la taxonomía estructural:
        obtener $g$, el orden $\alpha$ o un decodificador; a él
        apuntan las cifras de 2026 ($2^{140}$--$2^{148}$,
        $2^{99}$--$2^{102}$ según contabilidad)
        (caps.~\ref{ch:estructura} y~\ref{ch:2026}).

  \item[\textbf{Rechazo implícito}]
        \textsf{Decap} nunca devuelve $\bot$: entrega una clave
        pseudoaleatoria $H(0,s,C)$; elimina el oráculo
        observacional de los protocolos (caps.~\ref{ch:kem}
        y~\ref{ch:encapdecap}).

  \item[\textbf{ROM}]
        Modelo del oráculo aleatorio clásico; la cota aditiva de
        la transformación de Fujisaki--Okamoto (cap.~\ref{ch:kem}).

  \item[\textbf{$s$ (secreto de rechazo)}]
        Cadena uniforme de $n$ bits dentro de la clave secreta;
        la clave de la rama de rechazo es $H(0,s,C)$
        (caps.~\ref{ch:keygen} y~\ref{ch:encapdecap}).

  \item[\textbf{SD (Syndrome Decoding)}]
        El problema: dado $(H,y)$ con $y=He$, hallar $e$ de peso
        $\le t$. NP-duro en el peor caso; duro en media (conjetura)
        sobre la distribución del esquema (caps.~\ref{ch:seguridad}
        y~\ref{ch:isd}).

  \item[\textbf{Semi-sistemática (forma, $(\mu,\nu)$)}]
        Truco de \textsf{MatGen}: permutar hasta que las $mt$
        primeras columnas sean independientes;
        $(\mu,\nu)=(32,64)$ en las variantes \texttt{f}
        (cap.~\ref{ch:keygen}).

  \item[\textbf{SHAKE256}]
        El XOF de la especificación: expande semillas, genera
        $g$ y $\alpha$, y deriva las claves $K$ y de rechazo
        (caps.~\ref{ch:keygen} y~\ref{ch:labkem}).

  \item[\textbf{Shor, algoritmo de}]
        Rompe RSA y ECC en tiempo polinómico: la razón de ser de
        toda la migración post-cuántica (cap.~\ref{ch:pqc}).

  \item[\textbf{Shoup, ataque de tiempos de}]
        Ataque de canal lateral clásico contra implementaciones de
        decodificación; la motivación histórica de decodificar en
        tiempo constante (cap.~\ref{ch:kem}).

  \item[\textbf{Síndrome}]
        $S(y)=\Hmat y^{\top}$; en el código de Goppa también
        $S(x)=\sum_j c_j/(x+\alpha_j)$ con columnas
        $\sigma_j=(x+\alpha_j+1)/g(\alpha_j)$; su espacio tiene
        dimensión $n-k$ (caps.~\ref{ch:lineales}
        y~\ref{ch:labgoppa}).

  \item[\textbf{Soporte} ($\supp$)]
        El conjunto de posiciones no nulas de un vector. El error
        del KEM es uniforme sobre los $\binom{n}{t}$ soportes de
        tamaño $t$ (caps.~\ref{ch:lineales}
        y~\ref{ch:encapdecap}).

  \item[\textbf{Stern (algoritmo de)}]
        ISD con cruce de listas (1989): el primer gran salto de la
        escalera ($22$ bits sobre Prange en \texttt{348864})
        (cap.~\ref{ch:isd}).

  \item[\textbf{$t$}]
        Grado del polinomio de Goppa y peso exacto del error
        ($t=64$--$128$ en los conjuntos); fija la distancia
        $\ge2t+1$ y el trabajo de Patterson
        (caps.~\ref{ch:goppa} y~\ref{ch:parametros}).

  \item[\textbf{$\Tmat$ (clave pública)}]
        La mitad derecha de $\hat{\Hmat}=(\mathrm{I}_{mt}|T)$: la
        clave pública, de $mt\lceil k/8\rceil$ bytes
        ($261\,120$ a $1\,357\,824$) (caps.~\ref{ch:keygen}
        y~\ref{ch:parametros}).

  \item[\textbf{Terremoto (2025--2026)}]
        El periodo de distinguidores, estimaciones y contabilidades
        en disputa que movió la conversación pública sin romper el
        esquema (caps.~\ref{ch:estructura} y~\ref{ch:2026}).

  \item[\textbf{TII}]
        Retos públicos de instancias Goppa pequeñas ($n=214$)
        usados para probar heurísticas estructurales; no existe
        (a octubre de 2026) un cómputo de núcleos contra un
        conjunto real (cap.~\ref{ch:2026}).

  \item[\textbf{Traza} ($\Tr$)]
        El funcional $\Tr(\beta)=\sum_{i=0}^{m-1}\beta^{2^i}$;
        criterios de irreducibilidad, raíces cuadradas y
        linealidad sobre $\Ftwo$ (cap.~\ref{ch:cuerpos}).

  \item[\textbf{Vectores (notación)}]
        Las cadenas binarias se imprimen con la posición $0$ a la
        izquierda; al serializar se empaqueta en little-endian
        (posición $j$ en el bit $j\bmod8$ del byte $j/8$)
        (caps.~\ref{ch:labkem} y~\ref{app:codigo}).

  \item[\textbf{Weis (ataques de 2026)}]
        Recuperación heurística de clave vía \emph{holdout},
        con cifras entre $2^{89}$ y $2^{128}$ según la
        contabilidad, y la regla «$\approx8$ bits por unidad de
        $m$» (caps.~\ref{ch:estructura} y~\ref{ch:2026}).

\end{description}

\medskip
\noindent\textit{Fin del glosario. Los apéndices que siguen cierran el
libro: el código comentado del juguete (apéndice~\ref{app:codigo}),
la cronología de medio siglo de historia y las tablas del cuerpo
$\FF4$.}

% ======================================================================
%  APÉNDICE A — EL PROGRAMA DEL JUGUETE
% ======================================================================

\appendix
\chapter{El programa del juguete: \texttt{toy.py}}\label{app:codigo}

\begin{roadmap}
Este apéndice reproduce, íntegro y ejecutable, el programa que generó
las cifras de los laboratorios de los
capítulos~\ref{ch:labgoppa} y~\ref{ch:labkem}: el McEliece de juguete
sobre $\FF4$ con $n=16$ y $t=2$. El listado no es decorativo: cada
número impreso en esos capítulos está aquí calculado, y cada
afirmación delicada está protegida por una aserción que detiene el
programa si algo no cuadra.
\end{roadmap}

\section{Qué es y cómo se ejecuta}

\texttt{toy.py} es un único archivo de Python~3, sin más dependencia
que la biblioteca estándar (\texttt{hashlib.shake\_256}); ocupa unas
$450$ líneas. Se ejecuta con \texttt{python3 toy.py} y termina
imprimiendo \texttt{TODAS LAS COMPROBACIONES OK} si todo cuadra; si
una aserción falla, se detiene en la primera discrepancia. Reproduce,
en orden:

\begin{enumerate}[leftmargin=1.8em,itemsep=2pt]
  \item el cuerpo $\FF4$ y las $16$ columnas de la matriz $\tilde H$
        de la tabla~\ref{tab:lab-columnas} (capítulo~\ref{ch:labgoppa});
  \item la forma sistemática $(\mathrm{I}_8\,|\,T)$, la matriz
        $G=(T^{\top}\,|\,\mathrm{I}_8)$ y la permutación
        $\sigma$ del laboratorio (figura~\ref{fig:lab-SPG});
  \item el síndrome del error $e$ con soporte $\{10,13\}$ y la
        traza completa de Patterson, incluidos sus certificados;
  \item la comprobación de tipo van Lint--MacWilliams: las $2^{16}$
        palabras del código, contadas dos veces y comparadas;
  \item los $120$ errores de peso $\le 2$, los $137$ síndromes
        alcanzables y la lista completa de síndromes;
  \item la cadena del capítulo~\ref{ch:labkem}: la semilla $s$, el
        error $e$, el criptograma $C$, la clave $K$ y el criptograma
        manipulado $\tilde C$ con su rechazo implícito $K''$.
\end{enumerate}

\section{El listado completo}

\begin{lstlisting}
#!/usr/bin/env python3
# ----------------------------------------------------------------------
# toy.py - el juguete de bolsillo de los capitulos 4, 5 y 11
# Reproduce, paso a paso, todas las cifras impresas del libro.
# Uso:  python3 toy.py
# ----------------------------------------------------------------------

from hashlib import shake_256

# ======================================================================
# 1. GF(16) = F2[z]/(z^4+z+1)
# Cada elemento es un entero 0..15 (un digito hexadecimal): el bit l
# es el coeficiente de z^l.  z = 2, y z^4+1 = z.
# ======================================================================

def _tablas():
    exp, log = [0] * 15, [0] * 16
    v = 1
    for e in range(15):
        exp[e] = v
        log[v] = e
        v <<= 1
        if v & 0x10:
            v ^= 0x13                      # z^4 = z+1
    return exp, log

EXP, LOG = _tablas()

def mul(a, b):
    return 0 if not a or not b else EXP[(LOG[a] + LOG[b]) % 15]

def inv(a):
    return EXP[(-LOG[a]) % 15]

def pw(a, k):
    return 0 if not a else EXP[(LOG[a] * k) % 15]

# ======================================================================
# 2. Polinomios sobre GF(16): listas de coeficientes, indice = grado
# ======================================================================

def trim(p):
    p = p[:]
    while len(p) > 1 and p[-1] == 0:
        p.pop()
    return p

def p_add(a, b):
    r = [0] * max(len(a), len(b))
    for i, c in enumerate(a):
        r[i] ^= c
    for i, c in enumerate(b):
        r[i] ^= c
    return trim(r)

def p_esc(c, p):                           # multiplica por el escalar c
    return trim([mul(c, x) for x in p])

def p_mul(a, b):
    r = [0] * (len(a) + len(b) - 1)
    for i, x in enumerate(a):
        if x:
            for j, y in enumerate(b):
                if y:
                    r[i + j] ^= mul(x, y)
    return trim(r)

def p_divmod(a, b):                        # a = q*b + r,  deg r < deg b
    a = trim(a)
    q = [0] * (max(0, len(a) - len(b)) + 1)
    while a != [0] and len(a) >= len(b):
        c = mul(a[-1], inv(b[-1]))
        k = len(a) - len(b)
        q[k] ^= c
        for i, y in enumerate(b):
            a[i + k] ^= mul(c, y)
        a = trim(a)
    return trim(q), a

def p_mod(a, m):
    return p_divmod(a, m)[1]

def p_eval(p, x):
    r = 0
    for c in reversed(p):
        r = mul(r, x) ^ c
    return r

def p_inv_mod(a, m):                       # a^{-1} mod m  (gcd(a,m)=1)
    r0, r1 = trim(m), p_mod(a, m)
    u0, u1 = [0], [1]                      # invariante: r_i = u_i*a mod m
    while len(r1) > 1:                     # hasta llegar a una constante
        assert r1 != [0], "sin inverso: gcd != 1"
        q, r2 = p_divmod(r0, r1)
        r0, r1 = r1, r2
        u0, u1 = u1, p_add(u0, p_mul(q, u1))
    assert r1 != [0], "sin inverso: gcd != 1"
    lam = r1[0]                            # r1 = lam == u1*a  ->  1 = (u1/lam)*a
    return p_mod(p_esc(inv(lam), u1), m)

def p_sqrt_mod(v, g):                      # deg g = 2: anillo GF(256)
    w = p_mod(v, g)                        # raiz = v^(2^7)
    for _ in range(7):
        w = p_mod(p_mul(w, w), g)
    return w

def p_hex(p):
    def t(i):
        return "" if i == 0 else ("x" if i == 1 else f"x^{i}")
    s = '+'.join(f"{c:x}{t(i)}" for i, c in enumerate(p) if c)
    return s or '0'

# ======================================================================
# 3. El codigo de Goppa del juguete (m = 4, n = 16, t = 2)
# ======================================================================

N, M4 = 16, 4
ALPHA = list(range(N))                     # L = (0, 1, z, 1+z, ...)
G_POLY = [0x8, 0x1, 0x1]                   # g(x) = x^2 + x + z^3
INV_G = [inv(p_eval(G_POLY, a)) for a in ALPHA]
COLS = [(INV_G[j], mul(ALPHA[j], INV_G[j])) for j in range(N)]

def sbits(v, n):
    return ''.join(str((v >> i) & 1) for i in range(n))

def vbits(bs):
    return sum(b << i for i, b in enumerate(bs))

def syndrome(word):
    # S(x) = suma_j palabra_j * (x + a_j + 1)/g(a_j)   (mod g)
    S = [0]
    for j, b in enumerate(word):
        if b:
            S = p_add(S, p_mod(p_mul([ALPHA[j] ^ 1, 1], [INV_G[j]]), G_POLY))
    return S

def patterson(S, traza=None):
    R = p_inv_mod(S, G_POLY)                    # R = S^{-1}
    W = p_sqrt_mod(p_add(R, [0, 1]), G_POLY)    # W^2 = x + R
    tau = p_inv_mod(W, G_POLY)                  # tau = W^{-1}
    A, B = p_divmod(G_POLY, tau)                # g = tau*A + B
    sigma = p_add(p_mul(A, A), p_mul([0, 1], p_mul(B, B)))
    if traza is not None:
        traza.update(R=R, W=W, tau=tau, A=A, B=B)
    return sigma

def raices(sigma):
    return [j for j in range(N) if p_eval(sigma, ALPHA[j]) == 0]

# ---- expansion a bits, forma sistematica (I8|T) y G = (T^t|I8) -------

def hhat_variante(variante):
    # variante: (intercambio, invertir_bits)
    swap, msb = variante
    u = lambda j: COLS[j][1] if swap else COLS[j][0]   # filas 0-3
    v = lambda j: COLS[j][0] if swap else COLS[j][1]   # filas 4-7
    H = []
    for r in range(4):
        H.append([(u(j) >> (3 - r if msb else r)) & 1 for j in range(N)])
    for r in range(4):
        H.append([(v(j) >> (3 - r if msb else r)) & 1 for j in range(N)])
    return H

def sistematica(H):
    A = [row[:] for row in H]
    nr = len(A)
    for c in range(nr):
        piv = next((r for r in range(c, nr) if A[r][c]), None)
        if piv is None:
            return None
        A[c], A[piv] = A[piv], A[c]
        for r in range(nr):
            if r != c and A[r][c]:
                A[r] = [x ^ y for x, y in zip(A[r], A[c])]
    return [row[nr:] for row in A]

EXPECT_G = ["1100010110000000", "1100101001000000", "0110110000100000",
            "1001110000010000", "1011101000001000", "0111010100000100",
            "0110101100000010", "1001011100000001"]

def elige_convencion():
    for variante in ((0, 0), (0, 1), (1, 0), (1, 1)):
        H = hhat_variante(variante)
        T = sistematica(H)
        if T is None:
            continue
        G = [sbits(sum(T[a][i] << a for a in range(8)), 8)
             + sbits(1 << i, 8) for i in range(8)]
        if G == EXPECT_G:
            return H, T, G, variante
    raise SystemExit("ninguna convencion reproduce la G del libro")

HHAT, T, G, VARIANTE = elige_convencion()

print("=" * 70)
print("0. GF(16) = F2[z]/(z^4+z+1);  1/g y a/g se calculan en GF(16).")
print("   Potencias de z:", ' '.join(f"z^{e}={EXP[e]:x}" for e in range(15)))

# ----------------------------------------------------------------------
print()
print("=" * 70)
print("1. Las 16 columnas de H~ (tabla del libro):  (1/g(a_j), a_j/g(a_j))")
EXPECT_COLS = {0: (0xf, 0x0), 1: (0xf, 0xf), 2: (0x3, 0x6), 3: (0x3, 0x5),
               4: (0x8, 0x6), 5: (0x8, 0xe), 6: (0x2, 0xc), 7: (0x2, 0xe),
               8: (0xa, 0xf), 9: (0xa, 0x5), 10: (0xc, 0x1), 11: (0xc, 0xd),
               12: (0x5, 0x9), 13: (0x5, 0xc), 14: (0x4, 0xd), 15: (0x4, 0x9)}
for j in range(N):
    a, b = COLS[j]
    assert (a, b) == EXPECT_COLS[j], f"columna {j} discrepa"
    print(f"   j={j:2d}  a_j={ALPHA[j]:x}  g(a_j)={p_eval(G_POLY, ALPHA[j]):x}"
          f"  (1/g, a/g)=({a:x},{b:x})")

# ----------------------------------------------------------------------
print()
print("=" * 70)
print("2. Forma sistematica: filas de G = (T^t | I8)  [T unica]")
for i, row in enumerate(G):
    assert row == EXPECT_G[i]
    print(f"   G_{i} = {row}")
print("   columnas de T (hex):",
      ' '.join(f"{sum(T[a][i] << a for a in range(8)):02x}" for i in range(8)))

# ----------------------------------------------------------------------
print()
print("=" * 70)
print("3. Sindrome del error e = posiciones {10, 13}")
e2 = [0] * N
e2[10] = e2[13] = 1
s0 = 0
s1 = 0
for j in (10, 13):
    s0 ^= COLS[j][0]
    s1 ^= COLS[j][1]
assert (s0, s1) == (0x9, 0xd)              # H~ e = (9, d)
S_e = syndrome(e2)
print(f"   H~ e = (s0, s1) = ({s0:x}, {s1:x})")
print(f"   S(x) = s0*x + (s0+s1) = {p_hex(S_e)}   con S_e = [{S_e[0]:x}, {S_e[1]:x}]")
assert S_e == [0x4, 0x9]

# ----------------------------------------------------------------------
print()
print("=" * 70)
print("4. Patterson paso a paso para S(x) = 4 + 9x (con certificados)")
tr = {}
sig49 = patterson([0x4, 0x9], traza=tr)
R, W, tau, A, B = tr["R"], tr["W"], tr["tau"], tr["A"], tr["B"]
assert R == [0xa, 0x7]
assert p_mul([0x4, 0x9], R) == p_add([1], p_esc(0xa, G_POLY))
assert p_add(R, [0, 1]) == [0xa, 0x6]
assert W == [0xc, 0x7]
assert p_mul(W, W) == p_add([0xa, 0x6], p_esc(0x6, G_POLY))
assert tau == [0x3, 0xb]
assert p_mul(W, tau) == p_add([1], p_esc(0x4, G_POLY))
assert A == [0x3, 0x5] and B == [0xd]
assert p_add(p_mul(tau, A), B) == G_POLY
assert sig49 == [0x5, 0xe, 0x2]
assert raices(sig49) == [10, 13]
q1, r1 = p_divmod(sig49, [0xa, 1])
q2, r2 = p_divmod(sig49, [0xd, 1])
assert r1 == [0] and r2 == [0] and p_add(q1, q2) == [0xe]
assert p_mul(sig49, [0x4, 0x9]) == p_add([0xe], p_mul([0xe, 1], G_POLY))
print(f"   R = S^-1 = {p_hex(R)}        [certificado: S*R = 1 + (z+z^3)*g]")
print(f"   x + R = {p_hex(p_add(R, [0,1]))}")
print(f"   W = raiz(x+R) = {p_hex(W)}   [certificado: W^2 = (x+R) + (z+z^2)*g]")
print(f"   tau = W^-1 = {p_hex(tau)}    [certificado: W*tau = 1 + z^2*g]")
print(f"   g = tau*A + B con A = {p_hex(A)}, B = {p_hex(B)}  (division exacta)")
print(f"   sigma = A^2 + x*B^2 = {p_hex(sig49)}")
print(f"   raices de sigma: {raices(sig49)} = posiciones del error")
print(f"   sigma' = e (0xe) y sigma*S = sigma' + (x+e)*g  [OK]")

# ----------------------------------------------------------------------
print()
print("=" * 70)
print("5. Identidad de van Lint: Gamma2(L,g) = Gamma2(L,g^2)")

def nucleo(H):
    nr = len(H)
    colmask = [sum(H[r][j] << r for r in range(nr)) for j in range(N)]
    tabL, tabH = [0] * 256, [0] * 256
    for b in range(256):
        sL = sH = 0
        for j in range(8):
            if (b >> j) & 1:
                sL ^= colmask[j]
                sH ^= colmask[8 + j]
        tabL[b], tabH[b] = sL, sH
    return frozenset(w for w in range(1 << N)
                     if tabL[w & 255] ^ tabH[w >> 8] == 0)

H4 = []
g2inv = [inv(mul(p_eval(G_POLY, a), p_eval(G_POLY, a))) for a in ALPHA]
for i in range(4):
    for l in range(4):
        H4.append([(mul(1 if i == 0 else pw(ALPHA[j], i), g2inv[j]) >> l) & 1
                   for j in range(N)])
C1, C2 = nucleo(HHAT), nucleo(H4)
assert len(C1) == len(C2) == 256 and C1 == C2
print(f"   |nucleo(H~)| = {len(C1)},  |nucleo(4 filas a^i/g^2)| = {len(C2)}")
print("   las 65536 palabras: coinciden palabra a palabra  [OK]  (y son 256)")

# ----------------------------------------------------------------------
print()
print("=" * 70)
print("6. Correccion: los 120 errores de peso 2 y los 137 sindromes")
ok = 0
for i in range(N):
    for j in range(i + 1, N):
        w = [0] * N
        w[i] = w[j] = 1
        if raices(patterson(syndrome(w))) == [i, j]:
            ok += 1
assert ok == 120
print(f"   Patterson corrige {ok}/120 errores de peso 2  [OK]")
sind = {0}
for j in range(N):
    sind.add((COLS[j][0] << 4) | COLS[j][1])
for i in range(N):
    for j in range(i + 1, N):
        sind.add((COLS[i][0] ^ COLS[j][0]) << 4 | (COLS[i][1] ^ COLS[j][1]))
assert len(sind) == 1 + 16 + 120 == 137
print(f"   de los 2^8 = 256 sindromes, alcanzables con peso <= 2: {len(sind)}"
      "  (1 + 16 + 120)  [OK]")

# ----------------------------------------------------------------------
print()
print("=" * 70)
print("7. McEliece 1978 en miniatura (capitulo 5): G' = S*G*P")

def mmul(A, B):
    return [[sum(A[i][k] * B[k][j] for k in range(len(B))) % 2
             for j in range(len(B[0]))] for i in range(len(A))]

def perm_pull(M, perm):                    # (M P)[i][j] = M[i][perm[j]]
    return [[row[perm[j]] for j in range(len(row))] for row in M]

def inv_f2(M):
    n = len(M)
    A = [row[:] + [1 if i == j else 0 for j in range(n)]
         for i, row in enumerate(M)]
    for c in range(n):
        piv = next(r for r in range(c, n) if A[r][c])
        A[c], A[piv] = A[piv], A[c]
        for r in range(n):
            if r != c and A[r][c]:
                A[r] = [x ^ y for x, y in zip(A[r], A[c])]
    return [row[n:] for row in A]

S8 = [[int(b) for b in row] for row in
      ("01101010", "10010001", "01010001", "01010110",
       "01110110", "00111110", "10011100", "00010101")]
SIG = [14, 2, 7, 10, 4, 8, 12, 6, 9, 15, 11, 5, 0, 3, 13, 1]
SIGINV = [SIG.index(j) for j in range(N)]
Gm = [[int(b) for b in row] for row in EXPECT_G]
GP = mmul(S8, perm_pull(Gm, SIG))
GP_fig = ["1111001110010101", "0000110101111001", "0010000011101001",
          "1000100010100011", "1101000010110010", "1001001000110111",
          "0000011100111110", "0100100101110111"]
for i, row in enumerate(GP):
    txt = sbits(vbits(row), 16)
    assert txt == GP_fig[i], f"G' fila {i} != figura"
print("   G' = S*G*P reproduce la figura del libro  [OK]")

u = [int(c) for c in "01001101"]
z = [int(c) for c in "0001000000001000"]
uGP = [sum(u[i] * GP[i][j] for i in range(8)) % 2 for j in range(N)]
assert sbits(vbits(uGP), 16) == "0000011010001011"
x = [a ^ b for a, b in zip(uGP, z)]
assert sbits(vbits(x), 16) == "0001011010000011"
xPinv = [x[SIGINV[j]] for j in range(N)]
zPinv = [z[SIGINV[j]] for j in range(N)]
assert sbits(vbits(xPinv), 16) == "0100000011101100"
assert sbits(vbits(zPinv), 16) == "1000000000100000"
c = [a ^ b for a, b in zip(xPinv, zPinv)]
assert sbits(vbits(c), 16) == "1100000011001100"
print(f"   u = {sbits(vbits(u), 8)}  (4 unos)   z = {sbits(vbits(z), 16)}")
print(f"   x = u*G' + z = {sbits(vbits(x), 16)}")
print(f"   x*P^-1 = {sbits(vbits(xPinv), 16)}   z*P^-1 = {sbits(vbits(zPinv), 16)}")
print(f"   c = x*P^-1 + z*P^-1 = {sbits(vbits(c), 16)}")
S_x = syndrome(xPinv)
assert S_x == [0x2, 0x3]                   # S(x) = 2 + 3x
sig = patterson(S_x)
assert sig == [0, 0xc, 0xf]                # c x + f x^2
r = raices(sig)
assert r == [0, 10]
print(f"   sindrome: {p_hex(S_x)};  sigma = {p_hex(sig)} (multiplo escalar del localizador)")
print(f"   raices: {r} -> z' recuperado = {sbits(vbits(zPinv), 16)}  [OK]")
uS = c[8:]
u_rec = [sum(uS[i] * inv_f2(S8)[i][k] for i in range(8)) % 2 for k in range(8)]
assert sbits(vbits(uS), 8) == "11001100"
assert u_rec == u
print(f"   u*S = {sbits(vbits(uS), 8)};  u = (u*S)*S^-1 = {sbits(vbits(u_rec), 8)}"
      "  -> mensaje recuperado  [OK]")

# ----------------------------------------------------------------------
print()
print("=" * 70)
print("8. KEM del juguete (capitulo 11)")

DELTA = bytes([0x5a]) * 32
flujo = shake_256(bytes([0x40]) + DELTA).digest(8)   # etiqueta 64 = 0x40
s_val = flujo[0] | (flujo[1] << 8)
assert s_val == 0xCBAE
s_str = sbits(s_val, 16)
print(f"   delta = 5a repetido 32 veces")
print(f"   s = {s_str[:4]} {s_str[4:8]} {s_str[8:12]} {s_str[12:]}"
      f" = {s_val:04x}  (bytes AE CB, little-endian)")

e_alto, e_bajo = 0x00, 0x24                # e = (e_alto, e_bajo), posiciones 10 y 13
C_bits = [(e_alto >> r) & 1 ^
          sum(T[r][j] * ((e_bajo >> j) & 1) for j in range(8)) % 2
          for r in range(8)]
C_val = vbits(C_bits)
assert C_val == 0x98
print(f"   e = 0000000000100100 (posiciones 10 y 13)")
print(f"   C = e_alto + T*(e_bajo) = col2(T) + col5(T) = 36 ^ AE = {C_val:02x}"
      f" = {sbits(C_val, 8)}")

K = shake_256(bytes([0x01]) + bytes([e_alto, e_bajo]) + bytes([C_val])).digest(32)
assert K.hex() == \
    "597f46fbba325c68f0b0b45c6ea984961461eb220d1f24024ed888912b336278"
print(f"   K  = SHAKE256(01|00 24|98) = {K.hex()}")

v = [(C_val >> r) & 1 for r in range(8)] + [0] * 8
S_v = syndrome(v)
assert S_v == [0x4, 0x9]
sig_v = patterson(S_v)
assert sig_v == [0x5, 0xe, 0x2]
r_v = raices(sig_v)
assert r_v == [10, 13]
e_dec = [0] * N
for p in r_v:
    e_dec[p] = 1
assert e_dec == e2
K2 = shake_256(bytes([0x01]) + bytes([e_alto, e_bajo]) + bytes([C_val])).digest(32)
assert K2 == K
print(f"   decap: v = (C, 0); S(v) = {p_hex(S_v)}; sigma = {p_hex(sig_v)}")
print(f"   raices {r_v} -> e correcto;  K' = K  [OK]")

C_t = 0x99
v_t = [(C_t >> r) & 1 for r in range(8)] + [0] * 8
S_t = syndrome(v_t)
sig_t = patterson(S_t)
assert raices(sig_t) == []
K2t = shake_256(bytes([0x00]) + bytes([flujo[0], flujo[1]]) + bytes([C_t])).digest(32)
assert K2t.hex() == \
    "ba1a72ac26b0fcf20c55c108bb681e27aa9c270c60802a821601130b7f1e1494"
print(f"   C~ = C+1 = {sbits(C_t, 8)}: S~ = {p_hex(S_t)}, sigma~ = {p_hex(sig_t)},"
      " sin raices")
print(f"   K'' = SHAKE256(00|AE CB|99) = {K2t.hex()}  (rechazo)")

print()
print("=" * 70)
print("TODAS LAS COMPROBACIONES OK")
\end{lstlisting}

\section{La salida}

La salida siguiente es la de una ejecución real (octubre de 2026),
sin más recorte que los espacios en blanco finales de línea:

\begin{lstlisting}[numbers=none]
======================================================================
0. GF(16) = F2[z]/(z^4+z+1);  1/g y a/g se calculan en GF(16).
   Potencias de z: z^0=1 z^1=2 z^2=4 z^3=8 z^4=3 z^5=6 z^6=c z^7=b z^8=5 z^9=a z^10=7 z^11=e z^12=f z^13=d z^14=9

======================================================================
1. Las 16 columnas de H~ (tabla del libro):  (1/g(a_j), a_j/g(a_j))
   j= 0  a_j=0  g(a_j)=8  (1/g, a/g)=(f,0)
   j= 1  a_j=1  g(a_j)=8  (1/g, a/g)=(f,f)
   j= 2  a_j=2  g(a_j)=e  (1/g, a/g)=(3,6)
   j= 3  a_j=3  g(a_j)=e  (1/g, a/g)=(3,5)
   j= 4  a_j=4  g(a_j)=f  (1/g, a/g)=(8,6)
   j= 5  a_j=5  g(a_j)=f  (1/g, a/g)=(8,e)
   j= 6  a_j=6  g(a_j)=9  (1/g, a/g)=(2,c)
   j= 7  a_j=7  g(a_j)=9  (1/g, a/g)=(2,e)
   j= 8  a_j=8  g(a_j)=c  (1/g, a/g)=(a,f)
   j= 9  a_j=9  g(a_j)=c  (1/g, a/g)=(a,5)
   j=10  a_j=a  g(a_j)=a  (1/g, a/g)=(c,1)
   j=11  a_j=b  g(a_j)=a  (1/g, a/g)=(c,d)
   j=12  a_j=c  g(a_j)=b  (1/g, a/g)=(5,9)
   j=13  a_j=d  g(a_j)=b  (1/g, a/g)=(5,c)
   j=14  a_j=e  g(a_j)=d  (1/g, a/g)=(4,d)
   j=15  a_j=f  g(a_j)=d  (1/g, a/g)=(4,9)

======================================================================
2. Forma sistematica: filas de G = (T^t | I8)  [T unica]
   G_0 = 1100010110000000
   G_1 = 1100101001000000
   G_2 = 0110110000100000
   G_3 = 1001110000010000
   G_4 = 1011101000001000
   G_5 = 0111010100000100
   G_6 = 0110101100000010
   G_7 = 1001011100000001
   columnas de T (hex): a3 53 36 39 5d ae d6 e9

======================================================================
3. Sindrome del error e = posiciones {10, 13}
   H~ e = (s0, s1) = (9, d)
   S(x) = s0*x + (s0+s1) = 4+9x   con S_e = [4, 9]

======================================================================
4. Patterson paso a paso para S(x) = 4 + 9x (con certificados)
   R = S^-1 = a+7x        [certificado: S*R = 1 + (z+z^3)*g]
   x + R = a+6x
   W = raiz(x+R) = c+7x   [certificado: W^2 = (x+R) + (z+z^2)*g]
   tau = W^-1 = 3+bx    [certificado: W*tau = 1 + z^2*g]
   g = tau*A + B con A = 3+5x, B = d  (division exacta)
   sigma = A^2 + x*B^2 = 5+ex+2x^2
   raices de sigma: [10, 13] = posiciones del error
   sigma' = e (0xe) y sigma*S = sigma' + (x+e)*g  [OK]

======================================================================
5. Identidad de van Lint: Gamma2(L,g) = Gamma2(L,g^2)
   |nucleo(H~)| = 256,  |nucleo(4 filas a^i/g^2)| = 256
   las 65536 palabras: coinciden palabra a palabra  [OK]  (y son 256)

======================================================================
6. Correccion: los 120 errores de peso 2 y los 137 sindromes
   Patterson corrige 120/120 errores de peso 2  [OK]
   de los 2^8 = 256 sindromes, alcanzables con peso <= 2: 137  (1 + 16 + 120)  [OK]

======================================================================
7. McEliece 1978 en miniatura (capitulo 5): G' = S*G*P
   G' = S*G*P reproduce la figura del libro  [OK]
   u = 01001101  (4 unos)   z = 0001000000001000
   x = u*G' + z = 0001011010000011
   x*P^-1 = 0100000011101100   z*P^-1 = 1000000000100000
   c = x*P^-1 + z*P^-1 = 1100000011001100
   sindrome: 2+3x;  sigma = cx+fx^2 (multiplo escalar del localizador)
   raices: [0, 10] -> z' recuperado = 1000000000100000  [OK]
   u*S = 11001100;  u = (u*S)*S^-1 = 01001101  -> mensaje recuperado  [OK]

======================================================================
8. KEM del juguete (capitulo 11)
   delta = 5a repetido 32 veces
   s = 0111 0101 1101 0011 = cbae  (bytes AE CB, little-endian)
   e = 0000000000100100 (posiciones 10 y 13)
   C = e_alto + T*(e_bajo) = col2(T) + col5(T) = 36 ^ AE = 98 = 00011001
   K  = SHAKE256(01|00 24|98) = 597f46fbba325c68f0b0b45c6ea984961461eb220d1f24024ed888912b336278
   decap: v = (C, 0); S(v) = 4+9x; sigma = 5+ex+2x^2
   raices [10, 13] -> e correcto;  K' = K  [OK]
   C~ = C+1 = 10011001: S~ = b+6x, sigma~ = 5+8x+ax^2, sin raices
   K'' = SHAKE256(00|AE CB|99) = ba1a72ac26b0fcf20c55c108bb681e27aa9c270c60802a821601130b7f1e1494  (rechazo)

======================================================================
TODAS LAS COMPROBACIONES OK
\end{lstlisting}

\section{Simplificaciones del juguete}

El juguete no implementa la especificación real: la ilustra. Sus
diferencias con Classic McEliece, todas deliberadas:

\begin{itemize}[leftmargin=1.5em,itemsep=3pt]
  \item \textbf{El error es fijo.} \textsf{FixedWeight} devuelve
        siempre el error de soporte $\{10,13\}$: alcanza para
        verificar la tubería, pero la especificación lo sortea
        uniformemente entre las $\binom{n}{t}$ posibilidades
        (proposición~\ref{prop:fixedweight-uniform}).
  \item \textbf{No hay semillas, ni reintentos, ni variantes.} El
        soporte, $g$ y $\sigma$ están fijados a mano: no aparecen
        $\sigma_1$ y $\sigma_2$, ni la forma semi-sistemática
        $(\mu,\nu)=(32,64)$, ni la red de Benes que comprime
        $\alpha$ en la clave secreta (todo el
        capítulo~\ref{ch:keygen}); tampoco las variantes
        \texttt{f}, \texttt{pc} y \texttt{pcf}
        (capítulo~\ref{ch:parametros}).
  \item \textbf{La seguridad es nula.} Con $n=16$, enumerar las
        $2^{16}$ palabras es un parpadeo: el juguete verifica la
        corrección, nunca la dificultad. El problema «duro» empieza
        en los capítulos~\ref{ch:isd} y~\ref{ch:estructura}.
  \item \textbf{Las etiquetas del hash sí son las de la
        especificación}: \texttt{0x40} para el generador
        $G(\delta)$, \texttt{0x01} y \texttt{0x00} para las dos
        ramas de la clave de sesión (capítulo~\ref{ch:labkem}), con
        la convención de bits LSB-primero y bytes little-endian de
        la sección~6.2 de~\cite{cmspec}.
\end{itemize}

\medskip
\noindent Para las cuentas a mano, el apéndice~\ref{app:gf16} tiene
las tablas de $\FF4$ en formato de bolsillo.

% ======================================================================
%  APÉNDICE B — CRONOLOGÍA
% ======================================================================

\chapter{Cronología: de Goppa (1970) a octubre de 2026}\label{app:cronologia}

\begin{roadmap}
Medio siglo de fechas, en orden: los cimientos (códigos de Goppa,
decodificación, el criptosistema de 1978), las seis décadas de
criptoanálisis clásico y cuántico, los procesos de estandarización y,
con más detalle, el año que cambió el tono (2025--2026). Las
referencias apuntan a la bibliografía; los hitos de 2026 citan los
trabajos originales, no las crónicas.
\end{roadmap}

\vspace{-0.4\baselineskip}
\begin{longtable}{@{}p{2.6cm}p{12.96cm}@{}}
\toprule
fecha & hito\\
\midrule
\endfirsthead
\toprule
fecha & hito\\
\midrule
\endhead
\bottomrule
\endlastfoot
\multicolumn{2}{@{}l}{\textit{Fundaciones (1970--1991)}}\\
\midrule
1970 & Valeri Goppa introduce los códigos algebraico-geométricos
  sobre cuerpos finitos que hoy llevan su
  nombre~\cite{goppa1970}.\\
1975 & Patterson publica el algoritmo de decodificación que hace
  utilizables los códigos de Goppa~\cite{patterson1975}.\\
1978 & Robert J. McEliece propone el primer criptosistema de clave
  pública basado en códigos~\cite{mceliece1978}; Berlekamp, McEliece
  y van Tilborg demuestran que la decodificación por síndromes es
  NP-completa~\cite{bcmv1978}.\\
1986 & Niederreiter reformula el esquema en forma dual, con vectores
  de error de peso fijo como
  mensajes~\cite{niederreiter1986}.\\
1988 & Primeros ataques por conjuntos de información: Lee--Brickell
  y Leon toleran unos pocos errores en las $k$ columnas
  elegidas~\cite{lee1988,leon1988}.\\
1989 & Stern convierte el ataque en un cruce de listas
  (\emph{meet-in-the-middle})~\cite{stern1989}.\\
1991 & Dumer reutiliza y optimiza las listas de
  Stern~\cite{dumer1991}.\\
\midrule
\multicolumn{2}{@{}l}{\textit{La amenaza cuántica y el
criptoanálisis clásico (1994--2015)}}\\
\midrule
1994 & Shor: factorización y logaritmo discreto en tiempo
  polinómico en un ordenador cuántico~\cite{shor1994}; nace la
  agenda post-cuántica.\\
1996 & Grover: búsqueda cuántica con aceleración
  cuadrática~\cite{grover1996}.\\
1997 & Berson: ataques de reacción contra
  McEliece~\cite{berson1997}.\\
1998 & Canteaut y Sendrier: el ataque por división del
  soporte~\cite{canteautsendrier1998}.\\
2000 & Sendrier consolida la división del soporte como la segunda
  gran familia de criptoanálisis estructural~\cite{sendrier2000}.\\
2001 & Shoup: ataques de temporización contra la
  decodificación~\cite{shoup2001}.\\
2009 & Overbeck y Sendrier publican el panorama del criptoanálisis
  basado en códigos~\cite{overbeck2009}.\\
2010 & Bernstein: «Grover vs.\ McEliece»; los costes de puerta
  limitan la ventaja cuántica incluso en el peor
  caso~\cite{bernstein2010}.\\
2011 & May--Meurer--Thomae~\cite{mmt2011} y
  Bernstein--Lange--Peters (colisión de bolas)~\cite{blp2011}:
  representaciones y radios óptimos.\\
2012 & Becker--Joux--May--Meurer refinan las
  representaciones~\cite{bjmm2012}.\\
2013 & Faug\`ere et al.: primer distinguidor de claves de Goppa, en
  el régimen de tasa alta~\cite{faugere2013}.\\
2015 & May--Ozerov: la búsqueda de vecinos cercanos que fija los
  mejores números del estimador~\cite{mo2015}.\\
\midrule
\multicolumn{2}{@{}l}{\textit{Estandarización (2016--2024)}}\\
\midrule
2016 & El NIST convoca su proceso de estandarización
  post-cuántica.\\
2017 & Classic McEliece se presenta al concurso (llegará a la ronda
  final de KEMs)~\cite{cmspec}; Kachigar--Tillich: la decodificación
  por síndromes en el modelo cuántico~\cite{kachigar2017}.\\
2018 & Both--May mejoran las representaciones~\cite{bothmay2018};
  Bernstein--Persichetti prueban la conversión a KEM en el modelo de
  oráculo aleatorio~\cite{bp2018}.\\
2019 & Bindel et al.: cotas más finas de seguridad CCA en el
  modelo de oráculo aleatorio cuántico~\cite{bindel2019}.\\
2020 & Classic McEliece pasa a la cuarta ronda del NIST, junto a
  BIKE, HQC y SIKE.\\
2022 & El NIST selecciona Kyber (con Dilithium, Falcon y
  SPHINCS+) y cierra la ronda 3; SIKE cae ante el ataque de
  Castryck--Decru; Esser--Bellini publican el estimador de ISD que
  sustentará las tablas del equipo~\cite{esserbellini2022}; Carrier
  et al.\ y el ataque de fallos contra
  NTRU~\cite{carrier2022,ntrufault}; el equipo documenta los
  cambios de la ronda 4~\cite{cmmods3}.\\
2023 & ISO/IEC abre la enmienda con Classic McEliece: borrador del
  19 de abril~\cite{cmiso}.\\
2024 & El NIST finaliza FIPS~203 (ML-KEM), 204 (ML-DSA) y 205
  (SLH-DSA)~\cite{fips203}, y publica el plan de transición
  IR~8547~\cite{nist8547}.\\
\midrule
\multicolumn{2}{@{}l}{\textit{2025: la primera sacudida}}\\
\midrule
17 abr 2025 & El equipo responde a la afirmación de un distinguidor
  de \texttt{mceliece348864} «de solo $2^{529}$
  operaciones»~\cite{cm529}.\\
mar 2025 & El NIST cierra su proceso de KEMs con HQC; Classic
  McEliece queda fuera de la selección~\cite{nist8545}.\\
20 nov 2025 & El equipo se pronuncia sobre «You Only Decapsulate
  Once»~\cite{cm2025yodo}.\\
\midrule
\multicolumn{2}{@{}l}{\textit{2026: el año del terremoto}}\\
\midrule
3 mar 2026 & Bardet et al.: ataques estructurales sobre CFS y los
  retos TII~\cite{bardet2026}.\\
jun 2026 & ISO/IEC publica la Enmienda~2 a la 18033-2 con Classic
  McEliece: cuatro tamaños, cuatro variantes~\cite{iso18033}.\\
10 jun 2026 & Briaud et al.\ (BLRT): un nuevo distinguidor
  algebraico~\cite{briaud2026}.\\
23 jun 2026 & El equipo responde a la especulación de una recuperación
  de clave «de $2^{610}$ operaciones»~\cite{cm610}.\\
29 jun 2026 & Hemmert--Wiemers: distinguidor por anulación de orden
  superior (CRYPTO~2026)~\cite{hemmertwiemers2026,hemmert2026}.\\
7 ago 2026 & Ghoshal--Ishai--Jain--Sun (GIJS): el test
  \emph{holdout}, con préstamos de la teoría de la información
  privada~\cite{ghoshal2026}.\\
20 ago 2026 & Vedenev extiende el método a la recuperación de
  clave; resuelve retos TII~\cite{vedenev2026}.\\
21 ago 2026 & El equipo publica su réplica a los nuevos
  métodos~\cite{bernstein2026}.\\
24 ago 2026 & Saarinen: contabilidad de costes del \emph{holdout} y
  desafíos resueltos (tres anclas; TII-254 con GPU dedicada durante
  27{,}2 horas)~\cite{saarinen2026a,saarinen2026b,saarinen2026hover}.\\
26--28 ago 2026 & Apon et al.\ publican una cota inferior ($2^{1500}$
  operaciones de bits) para el paso crítico del método de GIJS;
  Vedenev responde con una ruta
  alternativa~\cite{apon2026,vedenev2026}.\\
27 ago 2026 & GIJS añaden a su artículo el método de recuperación de
  clave aplicado a los conjuntos reales~\cite{ghoshal2026}.\\
11 sep 2026 & Weis: el \emph{holdout} aplicado directamente; el
  equipo estimará después que las cuentas de una sola ejecución
  quedan en $2^{99}$--$2^{102}$~\cite{weis2026}.\\
15 sep 2026 & Saarinen revisa sus cifras a
  $2^{126{,}8}$--$2^{145{,}2}$~\cite{saarinen2026b}.\\
30 sep 2026 & Bernstein defiende la contabilidad clásica del
  equipo~\cite{bernstein2026}.\\
1 oct 2026 & El BSI publica su postura, fechada el 29 de septiembre:
  los métodos de agosto--octubre merecen tomarse en serio; «no debería
  seguir usándose en nuevos desarrollos»~\cite{bsi2026}.\\
7 oct 2026 & Vedenev publica la revisión de su
  ataque~\cite{vedenev2026}.\\
16 oct 2026 & Cierre de estos apuntes.\\
\end{longtable}

\begin{observacion}{Cómo leer una cronología en movimiento}{cronologia-viva}
Tres convenciones, para no confundir categorías:
(1)~las fechas de 2026 son las de las \emph{versiones públicas} de
los trabajos en el ePrint del IACR, no necesariamente las de su
concepción;
(2)~«resuelve el reto TII-$x$» describe desafíos de instancias de
laboratorio, no la rotura de ningún conjunto de
Classic McEliece~\cite{bardet2026,saarinen2026b};
(3)~las cifras de coste de esta cronología son las que cada autor
reclama, con su contabilidad: el debate del
capítulo~\ref{ch:2026} precisamente es sobre la contabilidad.
\end{observacion}

% ======================================================================
%  APÉNDICE C — GF(16) DE BOLSILLO
% ======================================================================

\chapter{GF(16) de bolsillo}\label{app:gf16}

El cuerpo $\FF4=\F_2[z]/(z^4+z+1)$ es el banco de pruebas de todo el
libro (ejemplo~\ref{ej:gf16} y los laboratorios de los
capítulos~\ref{ch:labgoppa} y~\ref{ch:labkem}). Este apéndice reúne,
para consulta rápida, las tablas que los ejercicios y las cuentas a
mano necesitan: potencias, logaritmos, inversos, raíces cuadradas,
traza y la tabla de multiplicar completa. Convenio: cada elemento se
escribe como dígito hexadecimal con el bit $i$-ésimo igual al
coeficiente de $z^i$ (en particular $z=\mathtt{2}$ y $z^4=z+1$); los
dígitos hex van en minúsculas, como en el programa del
apéndice~\ref{app:codigo}.

\section{Potencias y logaritmos}

\begin{table}[htbp]
\centering
\begin{tabular}{@{}clc@{\hspace{3.5em}}clc@{}}
\toprule
$i$ & hex & polinomio & $i$ & hex & polinomio\\
\midrule
0 & \texttt{1} & $1$ & 8 & \texttt{5} & $1+z^2$\\
1 & \texttt{2} & $z$ & 9 & \texttt{a} & $z+z^3$\\
2 & \texttt{4} & $z^2$ & 10 & \texttt{7} & $1+z+z^2$\\
3 & \texttt{8} & $z^3$ & 11 & \texttt{e} & $z+z^2+z^3$\\
4 & \texttt{3} & $1+z$ & 12 & \texttt{f} & $1+z+z^2+z^3$\\
5 & \texttt{6} & $z+z^2$ & 13 & \texttt{d} & $1+z^2+z^3$\\
6 & \texttt{c} & $z^2+z^3$ & 14 & \texttt{9} & $1+z^3$\\
7 & \texttt{b} & $1+z+z^3$ &  & &\\
\bottomrule
\end{tabular}
\caption{Las potencias $z^i$, $i=0,\dots,14$: el valor de $z^i$ es el
dígito hex y su expresión polinómica, la de al lado. El cuerpo tiene
$15$ elementos no nulos: $z$ es primitivo y su ciclo los agota todos
(figura~\ref{fig:ciclo-z}).}
\end{table}

\begin{table}[htbp]
\centering
\begin{tabular}{@{}cc@{\hspace{4em}}cc@{}}
\toprule
$a$ & $\log_z a$ & $a$ & $\log_z a$\\
\midrule
\texttt{1} & \texttt{0} & \texttt{9} & \texttt{e}\\
\texttt{2} & \texttt{1} & \texttt{a} & \texttt{9}\\
\texttt{3} & \texttt{4} & \texttt{b} & \texttt{7}\\
\texttt{4} & \texttt{2} & \texttt{c} & \texttt{6}\\
\texttt{5} & \texttt{8} & \texttt{d} & \texttt{d}\\
\texttt{6} & \texttt{5} & \texttt{e} & \texttt{b}\\
\texttt{7} & \texttt{a} & \texttt{f} & \texttt{c}\\
\texttt{8} & \texttt{3} & \texttt{0} & ---\\
\bottomrule
\end{tabular}
\caption{Logaritmos en base $z$ (los exponentes, también en hex: por
ejemplo $\log_z\mathtt{6}=\mathtt{5}$, pues $z^5=z+z^2$, y el $0$ no
tiene logaritmo). Para multiplicar: $ab=z^{\log a+\log b}$ con el
exponente módulo $15$.}
\end{table}

\section{Inversos, cuadrados, raíces y traza}

\begin{table}[htbp]
\centering
\begin{tabular}{@{}cc@{\hspace{4em}}cc@{}}
\toprule
$a$ & $a^{-1}$ & $a$ & $a^{-1}$\\
\midrule
\texttt{1} & \texttt{1} & \texttt{9} & \texttt{2}\\
\texttt{2} & \texttt{9} & \texttt{a} & \texttt{c}\\
\texttt{3} & \texttt{e} & \texttt{b} & \texttt{5}\\
\texttt{4} & \texttt{d} & \texttt{c} & \texttt{a}\\
\texttt{5} & \texttt{b} & \texttt{d} & \texttt{4}\\
\texttt{6} & \texttt{7} & \texttt{e} & \texttt{3}\\
\texttt{7} & \texttt{6} & \texttt{f} & \texttt{8}\\
\texttt{8} & \texttt{f} & \texttt{0} & ---\\
\bottomrule
\end{tabular}
\caption{Inversos multiplicativos. Comprobación rápida:
$\mathtt{3}\cdot\mathtt{e}=z^{4}\cdot z^{11}=z^{15}=1$.}
\end{table}

\begin{table}[htbp]
\centering
\begin{tabular}{@{}cccc@{\hspace{3em}}cccc@{}}
\toprule
$x$ & $x^2$ & $\sqrt{x}$ & $\Tr(x)$ & $x$ & $x^2$ & $\sqrt{x}$ & $\Tr(x)$\\
\midrule
\texttt{0} & \texttt{0} & \texttt{0} & 0 & \texttt{8} & \texttt{c} & \texttt{a} & 1\\
\texttt{1} & \texttt{1} & \texttt{1} & 0 & \texttt{9} & \texttt{d} & \texttt{b} & 1\\
\texttt{2} & \texttt{4} & \texttt{5} & 0 & \texttt{a} & \texttt{8} & \texttt{f} & 1\\
\texttt{3} & \texttt{5} & \texttt{4} & 0 & \texttt{b} & \texttt{9} & \texttt{e} & 1\\
\texttt{4} & \texttt{3} & \texttt{2} & 0 & \texttt{c} & \texttt{f} & \texttt{8} & 1\\
\texttt{5} & \texttt{2} & \texttt{3} & 0 & \texttt{d} & \texttt{e} & \texttt{9} & 1\\
\texttt{6} & \texttt{7} & \texttt{7} & 0 & \texttt{e} & \texttt{b} & \texttt{d} & 1\\
\texttt{7} & \texttt{6} & \texttt{6} & 0 & \texttt{f} & \texttt{a} & \texttt{c} & 1\\
\bottomrule
\end{tabular}
\caption{Cuadrados, raíces cuadradas ($\sqrt{x}=x^{8}$, la mitad de
$2^{m-1}$ porque la característica es $2$) y traza absoluta
($\Tr(x)=x+x^2+x^4+x^8\in\F_2$). Exactamente $8$ elementos tienen
traza $1$: el núcleo de $\Tr$ es un subespacio de dimensión $3$.}
\end{table}

\section{La tabla de multiplicar}

\begin{table}[htbp]
\centering
\setlength{\tabcolsep}{4.2pt}
\begin{tabular}{@{}c@{}c@{}c@{}c@{}c@{}c@{}c@{}c@{}c@{}c@{}c@{}c@{}c@{}c@{}c@{}c@{}c@{}}
\toprule
$\cdot$ & \texttt{0} & \texttt{1} & \texttt{2} & \texttt{3} & \texttt{4} & \texttt{5} & \texttt{6} & \texttt{7} & \texttt{8} & \texttt{9} & \texttt{a} & \texttt{b} & \texttt{c} & \texttt{d} & \texttt{e} & \texttt{f}\\
\midrule
\texttt{0} & \texttt{0} & \texttt{0} & \texttt{0} & \texttt{0} & \texttt{0} & \texttt{0} & \texttt{0} & \texttt{0} & \texttt{0} & \texttt{0} & \texttt{0} & \texttt{0} & \texttt{0} & \texttt{0} & \texttt{0} & \texttt{0}\\
\texttt{1} & \texttt{0} & \texttt{1} & \texttt{2} & \texttt{3} & \texttt{4} & \texttt{5} & \texttt{6} & \texttt{7} & \texttt{8} & \texttt{9} & \texttt{a} & \texttt{b} & \texttt{c} & \texttt{d} & \texttt{e} & \texttt{f}\\
\texttt{2} & \texttt{0} & \texttt{2} & \texttt{4} & \texttt{6} & \texttt{8} & \texttt{a} & \texttt{c} & \texttt{e} & \texttt{3} & \texttt{1} & \texttt{7} & \texttt{5} & \texttt{b} & \texttt{9} & \texttt{f} & \texttt{d}\\
\texttt{3} & \texttt{0} & \texttt{3} & \texttt{6} & \texttt{5} & \texttt{c} & \texttt{f} & \texttt{a} & \texttt{9} & \texttt{b} & \texttt{8} & \texttt{d} & \texttt{e} & \texttt{7} & \texttt{4} & \texttt{1} & \texttt{2}\\
\texttt{4} & \texttt{0} & \texttt{4} & \texttt{8} & \texttt{c} & \texttt{3} & \texttt{7} & \texttt{b} & \texttt{f} & \texttt{6} & \texttt{2} & \texttt{e} & \texttt{a} & \texttt{5} & \texttt{1} & \texttt{d} & \texttt{9}\\
\texttt{5} & \texttt{0} & \texttt{5} & \texttt{a} & \texttt{f} & \texttt{7} & \texttt{2} & \texttt{d} & \texttt{8} & \texttt{e} & \texttt{b} & \texttt{4} & \texttt{1} & \texttt{9} & \texttt{c} & \texttt{3} & \texttt{6}\\
\texttt{6} & \texttt{0} & \texttt{6} & \texttt{c} & \texttt{a} & \texttt{b} & \texttt{d} & \texttt{7} & \texttt{1} & \texttt{5} & \texttt{3} & \texttt{9} & \texttt{f} & \texttt{e} & \texttt{8} & \texttt{2} & \texttt{4}\\
\texttt{7} & \texttt{0} & \texttt{7} & \texttt{e} & \texttt{9} & \texttt{f} & \texttt{8} & \texttt{1} & \texttt{6} & \texttt{d} & \texttt{a} & \texttt{3} & \texttt{4} & \texttt{2} & \texttt{5} & \texttt{c} & \texttt{b}\\
\texttt{8} & \texttt{0} & \texttt{8} & \texttt{3} & \texttt{b} & \texttt{6} & \texttt{e} & \texttt{5} & \texttt{d} & \texttt{c} & \texttt{4} & \texttt{f} & \texttt{7} & \texttt{a} & \texttt{2} & \texttt{9} & \texttt{1}\\
\texttt{9} & \texttt{0} & \texttt{9} & \texttt{1} & \texttt{8} & \texttt{2} & \texttt{b} & \texttt{3} & \texttt{a} & \texttt{4} & \texttt{d} & \texttt{5} & \texttt{c} & \texttt{6} & \texttt{f} & \texttt{7} & \texttt{e}\\
\texttt{a} & \texttt{0} & \texttt{a} & \texttt{7} & \texttt{d} & \texttt{e} & \texttt{4} & \texttt{9} & \texttt{3} & \texttt{f} & \texttt{5} & \texttt{8} & \texttt{2} & \texttt{1} & \texttt{b} & \texttt{6} & \texttt{c}\\
\texttt{b} & \texttt{0} & \texttt{b} & \texttt{5} & \texttt{e} & \texttt{a} & \texttt{1} & \texttt{f} & \texttt{4} & \texttt{7} & \texttt{c} & \texttt{2} & \texttt{9} & \texttt{d} & \texttt{6} & \texttt{8} & \texttt{3}\\
\texttt{c} & \texttt{0} & \texttt{c} & \texttt{b} & \texttt{7} & \texttt{5} & \texttt{9} & \texttt{e} & \texttt{2} & \texttt{a} & \texttt{6} & \texttt{1} & \texttt{d} & \texttt{f} & \texttt{3} & \texttt{4} & \texttt{8}\\
\texttt{d} & \texttt{0} & \texttt{d} & \texttt{9} & \texttt{4} & \texttt{1} & \texttt{c} & \texttt{8} & \texttt{5} & \texttt{2} & \texttt{f} & \texttt{b} & \texttt{6} & \texttt{3} & \texttt{e} & \texttt{a} & \texttt{7}\\
\texttt{e} & \texttt{0} & \texttt{e} & \texttt{f} & \texttt{1} & \texttt{d} & \texttt{3} & \texttt{2} & \texttt{c} & \texttt{9} & \texttt{7} & \texttt{6} & \texttt{8} & \texttt{4} & \texttt{a} & \texttt{b} & \texttt{5}\\
\texttt{f} & \texttt{0} & \texttt{f} & \texttt{d} & \texttt{2} & \texttt{9} & \texttt{6} & \texttt{4} & \texttt{b} & \texttt{1} & \texttt{e} & \texttt{c} & \texttt{3} & \texttt{8} & \texttt{7} & \texttt{5} & \texttt{a}\\
\bottomrule
\end{tabular}
\caption{Producto $a\cdot b$ en $\FF4$, con $a$ en la fila y $b$ en la
columna (como en todas las tablas de este apéndice). Cada fila no
nula es una permutación de $\{1,\dots,\mathtt{f}\}$: la tabla lo
muestra y la estructura de cuerpo lo garantiza.}
\end{table}

\section{Cómo usar estas tablas}

\begin{itemize}[leftmargin=1.5em,itemsep=3pt]
  \item \textbf{Sumar} es XOR sin acarreo, bit a bit:
        $\mathtt{6}\oplus\mathtt{5}=\mathtt{3}$, es decir
        $(z+z^2)+(1+z^2)=1+z$.
  \item \textbf{Multiplicar} pasa por logaritmos:
        $a\,b=z^{\log_z a+\log_z b \pmod{15}}$. Ejemplo:
        $\mathtt{6}\cdot\mathtt{5}=z^{5}z^{8}=z^{13}=\mathtt{d}$;
        la versión polinómica es
        $(z+z^2)(1+z^2)=1+z^2+z^3$.
  \item \textbf{Invertir} es cambiar de signo el logaritmo. Con eso,
        la división es una resta de exponentes, y los
        $1/g(\alpha_i)$ del capítulo~\ref{ch:labgoppa} se calculan
        sin escribir una sola división polinómica.
  \item \textbf{La raíz cuadrada} no necesita algoritmos iterativos:
        en característica $2$, $\sqrt{x}=x^{2^{m-1}}=x^8$; Patterson
        la usa constantemente (capítulo~\ref{ch:goppa}).
  \item \textbf{La traza} decide irreducibilidad y pertenencia a
        subcuerpos; en el juguete, $\Tr(\beta)=\Tr(\beta^2)$ para
        todo $\beta$ (ejercicio~\ref{ej:traza-ej}).
  \item \textbf{Verificación cruzada.} El programa del
        apéndice~\ref{app:codigo} construye estas mismas tablas en
        una función de dos líneas y las usa para comprobar, número a
        número, todas las cuentas de los
        capítulos~\ref{ch:labgoppa} y~\ref{ch:labkem}.
\end{itemize}

\addcontentsline{toc}{chapter}{Bibliografía}

\begin{thebibliography}{99}
\raggedright

\bibitem{hqcspec}
C.~Aguilar-Melchor, N.~Aragon, S.~Bettaieb, L.~Bidoux, O.~Blazy,
et al. \emph{Hamming Quasi-Cyclic (HQC): specification}, submission
to the NIST post-quantum cryptography standardization process,
round~4, 2023.

\bibitem{nist8545}
G.~Alagic et al. \emph{Status report on the fourth round of the NIST
post-quantum cryptography standardization process}, NIST Internal
Report 8545, March 2025.

\bibitem{apon2026}
D.~Apon. An algebraic-geometry lower bound against the ePrint
2026/1747 McEliece key-recovery attack, Cryptology ePrint Archive,
Report 2026/1810, 2026.

\bibitem{bardet2026}
M.~Bardet, A.~Lemoine, J.-P.~Tillich. An attack on the CFS scheme
and on TII McEliece challenges. In \emph{Advances in Cryptology ---
CRYPTO 2026}, Springer, 2026.

\bibitem{bjmm2012}
A.~Becker, A.~Joux, A.~May, A.~Meurer. Decoding random binary linear
codes in $2^{n/20}$: How $1+1=0$ improves information set decoding.
In \emph{Advances in Cryptology --- EUROCRYPT 2012}, LNCS 7237,
pp.~520--536, Springer, 2012.

\bibitem{bcmv1978}
E.~R.~Berlekamp, R.~J.~McEliece, H.~C.~A.~van Tilborg. On the
inherent intractability of certain coding problems. \emph{IEEE
Transactions on Information Theory} 24(3):384--386, 1978.

\bibitem{bernstein2010}
D.~J.~Bernstein. Grover vs.\ McEliece. In \emph{Post-Quantum
Cryptography --- PQCrypto 2010}, LNCS 6061, pp.~73--80, Springer,
2010.

\bibitem{ntrufault}
D.~J.~Bernstein. A one-time single-bit fault leaks all previous
NTRU-HRSS session keys to a chosen-ciphertext attack, Cryptology
ePrint Archive, Report 2022/1125, 2022.

\bibitem{bernstein2026}
D.~J.~Bernstein et al. Posturas y respuestas del equipo Classic
McEliece, lista de correo del \emph{pqc-forum} del NIST, agosto--octubre
de 2026, \url{https://groups.google.com/a/list.nist.gov/g/pqc-forum}.

\bibitem{cmspec}
D.~J.~Bernstein, T.~Chou, T.~Lange, et al. \emph{Classic McEliece:
conservative code-based cryptography} (specification), 23 de octubre
de 2022, \url{https://classic.mceliece.org/mceliece-spec-20221023.pdf}.

\bibitem{blp2011}
D.~J.~Bernstein, T.~Lange, C.~Peters. Smaller decoding exponents:
ball-collision decoding. In \emph{Advances in Cryptology ---
CRYPTO 2011}, LNCS 6841, pp.~743--760, Springer, 2011.

\bibitem{bp2018}
D.~J.~Bernstein, E.~Persichetti. Towards KEM unification, Cryptology
ePrint Archive, Report 2018/526, 2018.

\bibitem{berson1997}
T.~A.~Berson. Failure of the McEliece public-key cryptosystem under
message-resend and related-message attack. In \emph{Advances in
Cryptology --- CRYPTO~'97}, LNCS 1294, pp.~213--220, Springer, 1997.

\bibitem{bindel2019}
N.~Bindel, M.~Hamburg, K.~H\"ovelmanns, A.~H\"ulsing, E.~Persichetti.
Tighter proofs of CCA security in the quantum random oracle model.
In \emph{Theory of Cryptography --- TCC 2019}, LNCS 11892,
pp.~61--90, Springer, 2019.

\bibitem{bothmay2018}
L.~Both, A.~May. Decoding linear codes with high error rate and its
impact for LPN security. In \emph{Post-Quantum Cryptography ---
PQCrypto 2018}, LNCS 10786, pp.~25--46, Springer, 2018.

\bibitem{briaud2026}
P.~Briaud, A.~Lemoine, H.~Randriambololona, J.-P.~Tillich. A
heuristic subexponential attack on the McEliece cryptosystem. In
\emph{Advances in Cryptology --- ASIACRYPT 2026}, Springer, 2026.

\bibitem{bsi2026}
Bundesamt f\"ur Sicherheit in der Informationstechnik. \emph{Notes on
recent developments concerning Classic McEliece}, 29 de septiembre
de 2026, \url{https://www.bsi.bund.de}.

\bibitem{canteautsendrier1998}
A.~Canteaut, N.~Sendrier. Cryptanalysis of the original McEliece
cryptosystem. In \emph{Advances in Cryptology --- ASIACRYPT~'98},
LNCS 1514, pp.~187--199, Springer, 1998.

\bibitem{carrier2022}
K.~Carrier, T.~Debris-Alazard, C.~Meyer-Hilfiger, J.-P.~Tillich.
Statistical decoding 2.0: reducing decoding to LPN,
arXiv:2208.02201, 2022.

\bibitem{cm529}
Classic McEliece Team. \emph{Notes on a recent claim that a
\texttt{mceliece348864} distinguisher uses only $2^{529}$
operations}, 17 de abril de 2025,
\url{https://classic.mceliece.org/mceliece-529-20250417.pdf}.

\bibitem{cm610}
Classic McEliece Team. \emph{Notes on recent speculation that a
\texttt{mceliece348864} key-recovery attack uses only $2^{610}$
operations}, 23 de junio de 2026,
\url{https://classic.mceliece.org/mceliece-610-20260623.pdf}.

\bibitem{cm2025yodo}
Classic McEliece Team. \emph{Statement regarding the paper ``You Only
Decapsulate Once''}, 20 de noviembre de 2025,
\url{https://classic.mceliece.org/mceliece-20251120.txt}.

\bibitem{cmimpl}
Classic McEliece Team. \emph{Classic McEliece: conservative
code-based cryptography: guide for implementors}, 23 de octubre de
2022,
\url{https://classic.mceliece.org/mceliece-impl-20221023.pdf}.

\bibitem{cmiso}
Classic McEliece Team. P\'agina de estandarizaci\'on ISO,
\url{https://classic.mceliece.org/iso.html}, consultada en octubre de
2026.

\bibitem{cmmods3}
Classic McEliece Team. \emph{Classic McEliece: conservative
code-based cryptography: description of modifications from round 3
to round 4}, 23 de octubre de 2022,
\url{https://classic.mceliece.org/mceliece-mods3-20221023.pdf}.

\bibitem{cmnist}
Classic McEliece Team. P\'agina de estandarizaci\'on NIST,
\url{https://classic.mceliece.org/nist.html}, consultada en octubre de
2026.

\bibitem{cmpc}
Classic McEliece Team. \emph{What plaintext confirmation means},
23 de octubre de 2022,
\url{https://classic.mceliece.org/mceliece-pc-20221023.pdf}.

\bibitem{cmrationale}
Classic McEliece Team. \emph{Classic McEliece: conservative
code-based cryptography: design rationale}, 23 de octubre de 2022,
\url{https://classic.mceliece.org/mceliece-rationale-20221023.pdf}.

\bibitem{cmsecurity}
Classic McEliece Team. \emph{Classic McEliece: conservative
code-based cryptography: guide for security reviewers}, 23 de
octubre de 2022,
\url{https://classic.mceliece.org/mceliece-security-20221023.pdf}.

\bibitem{couvreur2017}
A.~Couvreur, A.~Otmani, J.-P.~Tillich. Polynomial time attack on
wild McEliece over quadratic extensions. \emph{IEEE Transactions on
Information Theory} 63(1):404--427, 2017.

\bibitem{dumer1991}
I.~Dumer. On minimum distance decoding of linear codes. En
\emph{Proc.\ 5th Joint Soviet-Swedish International Workshop on
Information Theory}, pp.~50--52, 1991.

\bibitem{esserbellini2022}
A.~Esser, E.~Bellini. Syndrome decoding estimator. In
\emph{Public-Key Cryptography --- PKC 2022}, LNCS 13177,
pp.~112--141, Springer, 2022.

\bibitem{faugere2013}
J.-C.~Faug\`ere, V.~Gauthier-Uma\~na, A.~Otmani, L.~Perret,
J.-P.~Tillich. A distinguisher for high-rate McEliece cryptosystems.
\emph{IEEE Transactions on Information Theory} 59(10):6830--6844,
2013.

\bibitem{fips203}
National Institute of Standards and Technology. \emph{Module-lattice-based
key-encapsulation mechanism standard}, FIPS 203, agosto de 2024.

\bibitem{fips204}
National Institute of Standards and Technology. \emph{Module-lattice-based
digital signature standard}, FIPS 204, agosto de 2024.

\bibitem{fips205}
National Institute of Standards and Technology. \emph{Stateless
hash-based digital signature standard}, FIPS 205, agosto de 2024.

\bibitem{foj1999}
E.~Fujisaki, T.~Okamoto. Secure integration of asymmetric and
symmetric encryption schemes. In \emph{Advances in Cryptology ---
CRYPTO~'99}, LNCS 1666, pp.~537--554, Springer, 1999.

\bibitem{ghoshal2026}
A.~Ghoshal, Y.~Ishai, A.~Jain, N.~Sun. Quasipolynomial cryptanalysis
of the McEliece cryptosystem (or: PIR meets McEliece), Cryptology
ePrint Archive, Report 2026/1630, 2026.

\bibitem{gidney2021}
C.~Gidney, M.~Eker\r{a}. How to factor 2048 bit RSA integers in
8 hours using 20 million noisy qubits. \emph{Quantum} 5:433, 2021.

\bibitem{goppa1970}
V.~D.~Goppa. A new class of linear correcting codes.
\emph{Problemy Peredachi Informatsii} 6(3):24--30, 1970.

\bibitem{grover1996}
L.~K.~Grover. A fast quantum mechanical algorithm for database
search. En \emph{Proc.\ 28th ACM Symposium on Theory of Computing
(STOC)}, pp.~212--219, 1996.

\bibitem{hemmert2026}
T.~Hemmert. Key recovery for the McEliece cryptosystem using
higher-order vanishing, Cryptology ePrint Archive, Report 2026/1339,
2026.

\bibitem{hemmertwiemers2026}
T.~Hemmert, A.~Wiemers. Distinguishing Goppa codes using
higher-order vanishing. In \emph{Advances in Cryptology ---
CRYPTO 2026}, pp.~191--223, Springer, 2026.

\bibitem{hhk2017}
D.~Hofheinz, K.~H\"ovelmanns, E.~Kiltz. A modular analysis of the
Fujisaki--Okamoto transformation. In \emph{Theory of Cryptography ---
TCC 2017}, LNCS 10677, pp.~341--371, Springer, 2017.

\bibitem{iso18033}
ISO/IEC. \emph{Information technology --- Security techniques ---
Encryption algorithms --- Part~2: Asymmetric ciphers ---
Amendment~2: Classic McEliece}, ISO/IEC 18033-2:2006/Amd~2:2026,
junio de 2026.

\bibitem{kachigar2017}
G.~Kachigar, J.-P.~Tillich. Quantum information set decoding
algorithms. In \emph{Post-Quantum Cryptography --- PQCrypto 2017},
LNCS 10346, pp.~69--89, Springer, 2017.

\bibitem{kirshanova2018}
E.~Kirshanova. Improved quantum information set decoding. In
\emph{Post-Quantum Cryptography --- PQCrypto 2018}, LNCS 10786,
pp.~507--527, Springer, 2018.

\bibitem{kirshanova2022}
E.~Kirshanova, A.~May. Decoding McEliece with a hint --- secret
Goppa key parts reveal everything. In \emph{Security and Cryptography
for Networks --- SCN 2022}, LNCS 13409, Springer, 2022.

\bibitem{lee1988}
P.~J.~Lee, E.~F.~Brickell. An observation on the security of
McEliece's public-key cryptosystem. In \emph{Advances in Cryptology
--- EUROCRYPT~'88}, LNCS 330, pp.~275--280, Springer, 1988.

\bibitem{leon1988}
J.~S.~Leon. A probabilistic algorithm for computing minimum weights
of large error-correcting codes. \emph{IEEE Transactions on
Information Theory} 34(5):1354--1359, 1988.

\bibitem{lidw1994}
Y.~X.~Li, R.~H.~Deng, X.~M.~Wang. On the equivalence of McEliece's
and Niederreiter's public-key cryptosystems. \emph{IEEE Transactions
on Information Theory} 40(1):271--273, 1994.

\bibitem{lidl}
R.~Lidl, H.~Niederreiter. \emph{Finite Fields}, Encyclopedia of
Mathematics and its Applications 20, Cambridge University Press,
2.ª edici\'on, 1997.

\bibitem{msbook}
F.~J.~MacWilliams, N.~J.~A.~Sloane. \emph{The Theory of
Error-Correcting Codes}, North-Holland, 1977.

\bibitem{mmt2011}
A.~May, A.~Meurer, E.~Thomae. Decoding random linear codes in
$\tilde{\mathcal{O}}(2^{0.054n})$. In \emph{Advances in Cryptology ---
ASIACRYPT 2011}, LNCS 7073, pp.~107--124, Springer, 2011.

\bibitem{mo2015}
A.~May, I.~Ozerov. On computing nearest neighbors with applications
to decoding of binary linear codes. In \emph{Advances in Cryptology
--- EUROCRYPT 2015}, LNCS 9056, pp.~203--228, Springer, 2015.

\bibitem{mceliece1978}
R.~J.~McEliece. A public-key cryptosystem based on algebraic coding
theory. \emph{DSN Progress Report 42-44}, Jet Propulsion Laboratory,
pp.~114--116, 1978.

\bibitem{nist8547}
D.~Moody et al. \emph{Transition to post-quantum cryptography
standards}, NIST Internal Report 8547 (borrador p\'ublico inicial),
noviembre de 2024.

\bibitem{niederreiter1986}
H.~Niederreiter. Knapsack-type cryptosystems and algebraic coding
theory. \emph{Problems of Control and Information Theory}
15(2):159--166, 1986.

\bibitem{overbeck2009}
R.~Overbeck, N.~Sendrier. Code-based cryptography. En
\emph{Post-Quantum Cryptography} (D.~J.~Bernstein, J.~Buchmann,
E.~Dahmen, eds.), pp.~95--145, Springer, 2009.

\bibitem{patterson1975}
N.~J.~Patterson. The algebraic decoding of Goppa codes. \emph{IEEE
Transactions on Information Theory} 21(2):203--207, 1975.

\bibitem{persichetti2012}
E.~Persichetti. \emph{Improving the efficiency of code-based
cryptography}, tesis doctoral, University of Auckland, 2012.

\bibitem{prange1962}
E.~Prange. The use of information sets in decoding cyclic codes.
\emph{IRE Transactions on Information Theory} 8(5):5--9, 1962.

\bibitem{randriam2024}
H.~Randriambololona. The syzygy distinguisher. In \emph{Advances in
Cryptology --- EUROCRYPT 2025}, LNCS 15606, Springer, 2025
(revisi\'on ampliada: Cryptology ePrint Archive, Report 2024/1193).

\bibitem{saarinen2026a}
M.-J.~O.~Saarinen. Bit operation cost of ``holdout'' key-recovery
attacks against Classic McEliece, Cryptology ePrint Archive,
Report 2026/1786, 2026.

\bibitem{saarinen2026b}
M.-J.~O.~Saarinen. Two-anchor holdout/Hermite: solving the TII-254
McEliece key recovery challenge, Cryptology ePrint Archive,
Report 2026/1986, 2026.

\bibitem{saarinen2026hover}
M.-J.~O.~Saarinen. HOVER: Higher-order vanishing endomorphism
recovery, Cryptology ePrint Archive, Report 2026/1778, 2026
(NORDSEC 2026, LNCS, en prensa).

\bibitem{sendrier2000}
N.~Sendrier. Finding the permutation between equivalent linear
codes: the support splitting algorithm. \emph{IEEE Transactions on
Information Theory} 46(4):1193--1203, 2000.

\bibitem{shor1994}
P.~W.~Shor. Algorithms for quantum computation: discrete logarithms
and factoring. En \emph{Proc.\ 35th IEEE Symposium on Foundations of
Computer Science (FOCS)}, pp.~124--134, 1994.

\bibitem{shoup2001}
V.~Shoup. A proposal for an ISO standard for public key encryption,
Cryptology ePrint Archive, Report 2001/112, 2001.

\bibitem{stern1989}
J.~Stern. A method for finding codewords of small weight. En
\emph{Coding Theory and Applications}, LNCS 388, pp.~106--113,
Springer, 1989.

\bibitem{vlint}
J.~H.~van Lint. \emph{Introduction to Coding Theory}, Graduate
Texts in Mathematics 86, Springer, 3.ª edici\'on, 1999.

\bibitem{vedenev2026}
K.~Vedenev. Extending distinguishing to key recovery for subfield
subcodes of GRS codes, Cryptology ePrint Archive, Report 2026/1747,
2026.

\bibitem{weis2026}
S.~A.~Weis. Improving GIJS key recovery for Classic McEliece,
Cryptology ePrint Archive, Report 2026/1984, 2026.

\end{thebibliography}

\clearpage

\end{document}
